% ======================================================================
%  Time Series -- Final Exam Preparation  (MATH-342, EPFL, S. Olhede)
%  Master file.  Compile with:  pdflatex main  (twice for TOC).
% ======================================================================
\documentclass[11pt]{article}
% ======================================================================
%  preamble.tex  --  Time Series Final Exam Preparation (MATH-342, EPFL)
%  Robust preamble: amsmath/amssymb/amsthm, tcolorbox theorem boxes,
%  math macros, hyperref, geometry.  Compiles with pdflatex.
% ======================================================================
\usepackage[utf8]{inputenc}
\usepackage[T1]{fontenc}
\usepackage[a4paper,margin=2.3cm]{geometry}
\usepackage{amsmath,amssymb,amsthm,mathtools}
\usepackage{bm}
\usepackage{enumitem}
\usepackage{xcolor}
\usepackage[most]{tcolorbox}
\usepackage{booktabs}
\usepackage{array}
\usepackage{microtype}
\usepackage{hyperref}

\hypersetup{
  colorlinks=true,
  linkcolor=blue!55!black,
  urlcolor=teal!70!black,
  citecolor=teal!70!black,
  pdftitle={Time Series -- Final Exam Preparation (MATH-342, EPFL)},
  pdfauthor={Revision document}
}

% --- robustness against overfull boxes -------------------------------
\emergencystretch=3em
\allowdisplaybreaks[1]
\setlength{\parindent}{0pt}
\setlength{\parskip}{0.4em}

% ======================================================================
%  MATH MACROS
% ======================================================================
\newcommand{\E}{\mathbb{E}}
\DeclareMathOperator{\Var}{Var}
\DeclareMathOperator{\Cov}{Cov}
\DeclareMathOperator{\Corr}{Corr}
\DeclareMathOperator*{\argmin}{arg\,min}
\DeclareMathOperator*{\argmax}{arg\,max}
\DeclareMathOperator{\MSE}{MSE}
\DeclareMathOperator{\trace}{tr}
\DeclareMathOperator{\diag}{diag}
\DeclareMathOperator{\sgn}{sgn}
\newcommand{\Reals}{\mathbb{R}}
\newcommand{\Ints}{\mathbb{Z}}
\newcommand{\Comp}{\mathbb{C}}
\newcommand{\Nat}{\mathbb{N}}
\newcommand{\Prob}{\mathbb{P}}
\newcommand{\Tr}{^{\mathsf{T}}}            % transpose
\newcommand{\Herm}{^{\mathsf{H}}}          % conjugate transpose
\newcommand{\conj}[1]{\overline{#1}}
\newcommand{\abs}[1]{\left\lvert #1 \right\rvert}
\newcommand{\norm}[1]{\left\lVert #1 \right\rVert}
\newcommand{\ind}{\mathbf{1}}              % indicator
\newcommand{\eps}{\varepsilon}
\newcommand{\bx}{\mathbf{x}}
\newcommand{\bX}{\mathbf{X}}
\newcommand{\diff}{\nabla}                  % difference operator
\newcommand{\acvs}{\gamma}
\newcommand{\acf}{\rho}
\newcommand{\pacf}{\alpha}
\newcommand{\sdf}{\mathcal{S}}              % spectral density function
\newcommand{\Fej}{\mathcal{F}}             % Fejer kernel
\newcommand{\Dir}{\mathcal{D}}             % Dirichlet kernel
\newcommand{\filt}{\mathcal{L}}            % LTI filter operator
\newcommand{\Bsh}{B}                        % backshift operator
\newcommand{\ms}{\;\overset{\mathrm{ms}}{=}\;}
\newcommand{\eqdist}{\overset{d}{=}}
\newcommand{\todist}{\xrightarrow{d}}
\newcommand{\toprob}{\xrightarrow{P}}
\newcommand{\iid}{\overset{\text{iid}}{\sim}}
\newcommand{\indep}{\overset{\text{ind}}{\sim}}

% ======================================================================
%  TCOLORBOX THEOREM ENVIRONMENTS
%  Shared counter (thmcounter) for theory objects -> sequential per
%  section: Def 1.1, Thm 1.2, Prop 1.3, ...
%  Separate counters for exercises and key techniques.
% ======================================================================
\newcounter{thmcounter}
\numberwithin{thmcounter}{section}
\newcounter{exocounter}
\numberwithin{exocounter}{section}
\newcounter{ktcounter}
\numberwithin{ktcounter}{section}

\tcbset{
  thmbox/.style={
    enhanced, breakable,
    fonttitle=\bfseries, coltitle=black,
    boxrule=0.7pt, arc=1.4mm,
    left=2.4mm, right=2.4mm, top=3.0mm, bottom=1.4mm,
    before skip=9pt, after skip=7pt,
    attach boxed title to top left={yshift=-3.6mm, xshift=2.6mm},
    boxed title style={boxrule=0pt, arc=0.6mm, colframe=black!0},
    fonttitle=\bfseries
  }
}

\newtcbtheorem[use counter=thmcounter]{definition}{Definition}{%
  thmbox, colback=blue!4!white, colframe=blue!58!black,
  colbacktitle=blue!58!black, coltitle=white}{def}

\newtcbtheorem[use counter=thmcounter]{theorem}{Theorem}{%
  thmbox, colback=red!4!white, colframe=red!62!black,
  colbacktitle=red!62!black, coltitle=white}{thm}

\newtcbtheorem[use counter=thmcounter]{proposition}{Proposition}{%
  thmbox, colback=teal!5!white, colframe=teal!62!black,
  colbacktitle=teal!62!black, coltitle=white}{prop}

\newtcbtheorem[use counter=thmcounter]{lemma}{Lemma}{%
  thmbox, colback=violet!5!white, colframe=violet!60!black,
  colbacktitle=violet!60!black, coltitle=white}{lem}

\newtcbtheorem[use counter=thmcounter]{corollary}{Corollary}{%
  thmbox, colback=violet!4!white, colframe=violet!45!black,
  colbacktitle=violet!45!black, coltitle=white}{cor}

\newtcbtheorem[use counter=thmcounter]{example}{Example}{%
  thmbox, colback=black!3!white, colframe=black!55!white,
  colbacktitle=black!55!white, coltitle=white}{exmp}

\newtcbtheorem[use counter=ktcounter]{keytech}{Key Technique}{%
  thmbox, colback=orange!7!white, colframe=orange!75!black,
  colbacktitle=orange!75!black, coltitle=white}{kt}

\newtcbtheorem[use counter=exocounter]{exercise}{Exercise}{%
  thmbox, colback=cyan!4!white, colframe=cyan!55!black,
  colbacktitle=cyan!55!black, coltitle=white}{exo}

% --- plain (unnumbered) boxes ----------------------------------------
\newtcolorbox{solution}{%
  enhanced, breakable, colback=green!3!white, colframe=green!48!black,
  title=\textbf{Solution}, coltitle=white, colbacktitle=green!48!black,
  fonttitle=\bfseries, boxrule=0.6pt, arc=1.2mm,
  left=2.4mm, right=2.4mm, top=1.2mm, bottom=1.2mm,
  before skip=4pt, after skip=8pt}

\newtcolorbox{intuition}[1][Intuition / Why it matters]{%
  enhanced, breakable, colback=yellow!10!white, colframe=orange!70!black,
  title={\textbf{#1}}, coltitle=black, colbacktitle=yellow!35!white,
  fonttitle=\bfseries, boxrule=0.5pt, arc=1mm, sharp corners=south,
  left=2.4mm, right=2.4mm, top=1.2mm, bottom=1.2mm,
  before skip=6pt, after skip=6pt}

\newtcolorbox{pitfall}{%
  enhanced, breakable, colback=red!6!white, colframe=red!70!black,
  title={\textbf{\ding{43}\,Classic Pitfall}}, coltitle=white,
  colbacktitle=red!70!black, fonttitle=\bfseries,
  boxrule=0.6pt, arc=1mm,
  left=2.4mm, right=2.4mm, top=1.2mm, bottom=1.2mm,
  before skip=6pt, after skip=6pt}

\newtcolorbox{recallbox}[1][Recall]{%
  enhanced, breakable, colback=blue!3!white, colframe=blue!40!black,
  title={\textbf{#1}}, coltitle=white, colbacktitle=blue!40!black,
  fonttitle=\bfseries, boxrule=0.5pt, arc=1mm,
  left=2.4mm, right=2.4mm, top=1.2mm, bottom=1.2mm,
  before skip=6pt, after skip=6pt}

% pifont for the dingbat in pitfall title
\usepackage{pifont}

% nicer enumerate defaults
\setlist[enumerate]{leftmargin=1.8em, topsep=2pt, itemsep=1pt}
\setlist[itemize]{leftmargin=1.6em, topsep=2pt, itemsep=1pt}


\begin{document}

% ----------------------------------------------------------------------
%  TITLE PAGE
% ----------------------------------------------------------------------
\begin{titlepage}
\centering
\vspace*{2.2cm}
{\Huge\bfseries Time Series Analysis\par}
\vspace{0.6cm}
{\LARGE Complete Final-Exam Preparation\par}
\vspace{0.4cm}
{\large MATH-342 -- \'Ecole Polytechnique F\'ed\'erale de Lausanne\par}
{\large Course by Prof.\ Sofia C.\ Olhede\par}
\vspace{1.6cm}

\begin{tcolorbox}[width=0.86\textwidth, colback=blue!4!white,
  colframe=blue!58!black, arc=2mm, boxrule=0.8pt,
  title=\textbf{What this document is}, coltitle=white,
  colbacktitle=blue!58!black, fonttitle=\large\bfseries]
A self-contained revision companion covering \emph{every} chapter of the
lecture notes: all \textbf{definitions}, \textbf{theorems / propositions /
lemmas} (statement $+$ proof idea $+$ why it matters), the \textbf{key
techniques} you reproduce in the exam, and \textbf{worked exercises}
(selected from the problem sheets and freshly created) with \textbf{full
solutions}. It ends with an \textbf{Exam Strategy Checklist} and a list of
\textbf{classic pitfalls}.
\end{tcolorbox}

\vfill
{\large Coverage: Stationarity \,$\cdot$\, ACVS/ACF/PACF \,$\cdot$\,
AR/MA/ARMA/ARIMA/SARIMA \,$\cdot$\, Causality \& Invertibility \,$\cdot$\,
Estimation (Yule--Walker, LS, MLE) \,$\cdot$\, Forecasting \,$\cdot$\,
Fourier \& Spectral Analysis \,$\cdot$\, Periodogram \& Tapering \,$\cdot$\,
Multivariate / VAR \,$\cdot$\, Diagnostics \& Model Selection \,$\cdot$\,
Long Memory \,$\cdot$\, ARCH/GARCH \,$\cdot$\, LTI Filters\par}
\vspace{1.2cm}
{\large Revision document compiled with \texttt{pdflatex}.\par}
\end{titlepage}

% ----------------------------------------------------------------------
%  HOW TO USE
% ----------------------------------------------------------------------
\section*{How to use this document}
\addcontentsline{toc}{section}{How to use this document}

\begin{tcolorbox}[colback=orange!7!white, colframe=orange!75!black,
  arc=1.5mm, boxrule=0.7pt, title=\textbf{Reading strategy},
  coltitle=white, colbacktitle=orange!75!black, fonttitle=\bfseries]
\begin{enumerate}[leftmargin=2.2em]
  \item \textbf{First pass.} Read each section's \emph{Definitions} and
        \emph{Theorems} boxes. For every theorem, learn the
        \emph{one-line statement}, the \emph{proof idea}, and the
        \emph{``why''} --- examiners reward knowing \emph{when} a result
        applies, not just \emph{that} it is true.
  \item \textbf{Second pass.} Work through the \emph{Key Technique} boxes.
        These are the recurring computational moves (compute an ACVS,
        check causality via roots, derive an MA$(\infty)$ representation,
        forecast recursively, get a spectral density). The exam is mostly
        these moves.
  \item \textbf{Third pass.} Do the \emph{Exercises} \textbf{with the
        solution hidden}. Exercises tagged \texttt{[Sheet $k$]} come from
        the problem sheets; those tagged \texttt{[New]} were written to
        probe understanding from a different angle. \emph{When you can do
        every exercise unaided, you are ready.}
  \item \textbf{Night before.} Read only the
        \hyperref[sec:strategy]{Exam Strategy Checklist} and the
        \hyperref[sec:pitfalls]{Classic Pitfalls}.
\end{enumerate}
\end{tcolorbox}

\begin{tcolorbox}[colback=blue!3!white, colframe=blue!45!black,
  arc=1.5mm, boxrule=0.6pt, title=\textbf{Colour legend},
  coltitle=white, colbacktitle=blue!45!black, fonttitle=\bfseries]
\textcolor{blue!58!black}{$\blacksquare$}~Definition \quad
\textcolor{red!62!black}{$\blacksquare$}~Theorem \quad
\textcolor{teal!62!black}{$\blacksquare$}~Proposition \quad
\textcolor{violet!60!black}{$\blacksquare$}~Lemma/Corollary \quad
\textcolor{orange!75!black}{$\blacksquare$}~Key Technique \quad
\textcolor{cyan!55!black}{$\blacksquare$}~Exercise \quad
\textcolor{green!48!black}{$\blacksquare$}~Solution
\end{tcolorbox}

\newpage
\tableofcontents
\newpage

% ----------------------------------------------------------------------
%  CONTENT PARTS
% ----------------------------------------------------------------------
% ======================================================================
%  PART 1 -- Foundations: Stationarity, Autocovariance, Ergodicity
% ======================================================================
\section{Foundations: Stationarity, Autocovariance \& Ergodicity}
\label{sec:foundations}

A \emph{time series} $\{X_t\}$ is both the data recorded in time order and the
stochastic process $\{X_t : t\in T\}$ on a probability space
$(\Omega,\mathcal F,\Prob)$ it is a realisation of. The index set $T$ is a set
of time points (here $T=\Ints$ unless stated otherwise; we take the sampling
step $\Delta=1$). Time series differ from classical statistics because the
observations are \textbf{dependent}: ``there are many more ways to be dependent
than to be independent''. Everything in this course is built on describing that
dependence through second-order structure.

% ----------------------------------------------------------------------
\subsection{Definitions}

\begin{definition}{(Weak / second-order) stationarity}{wkstat}
$\{X_t\}$ is \textbf{second-order} (equivalently \textbf{weakly} or
\textbf{covariance}) \textbf{stationary} if all first- and second-order joint
moments exist, are finite, and are invariant to time shifts. Concretely, for
all $t,s\in T$ and all admissible shifts $\tau$:
\[
\text{(i) } \E[X_t]=\mu;\qquad
\text{(ii) } \Var(X_t)=\sigma^2<\infty;\qquad
\text{(iii) } \E[X_tX_{t+\tau}]=\E[X_sX_{s+\tau}].
\]
Condition (iii) forces $\E[X_tX_{t+\tau}]$ to be a function of the lag $\tau$
only.
\end{definition}

\begin{definition}{Strong / strict stationarity}{ststat}
$\{X_t\}$ is \textbf{strictly} (completely) \textbf{stationary} if for every
$n\ge1$, every $t_1,\dots,t_n$ and every shift $\tau$, the joint distribution
of $(X_{t_1},\dots,X_{t_n})$ equals that of
$(X_{t_1+\tau},\dots,X_{t_n+\tau})$.
\end{definition}

\begin{intuition}
Neither notion implies the other in general.
\textbf{Strict $\not\Rightarrow$ weak}: i.i.d.\ Student-$t$ on $1$ d.o.f.\
(Cauchy) is strictly stationary but has no finite moments, so it is not weakly
stationary. \textbf{Weak $\not\Rightarrow$ strict}: weak stationarity only
constrains the first two moments and says nothing about higher ones. The two
coincide for \emph{Gaussian} processes (Definition~\ref{def:gauss:gp}).
\end{intuition}

\begin{definition}{Gaussian process}{gauss:gp}
$\{X_t\}$ is a \textbf{Gaussian process} if for all $n$ and $t_1,\dots,t_n$ the
vector $(X_{t_1},\dots,X_{t_n})$ is multivariate normal. Such a process is
completely characterised by its (finite) first two moments; hence for Gaussian
processes \textbf{weak $\Leftrightarrow$ strict stationarity}.
\end{definition}

\begin{definition}{Autocovariance sequence (ACVS)}{acvs}
For a discrete-time second-order stationary $\{X_t\}$, the \textbf{ACVS} is
\[
\acvs_\tau=\Cov(X_0,X_\tau)=\Cov(X_t,X_{t+\tau}),\qquad \tau\in\Ints .
\]
Immediate properties: $\acvs_0=\Var(X_t)\ge0$, symmetry $\acvs_{-\tau}=\acvs_\tau$,
and (Proposition~\ref{prop:psd}) positive semi-definiteness.
\end{definition}

\begin{definition}{Autocorrelation sequence (ACF)}{acf}
The \textbf{ACF} is $\acf_\tau=\Corr(X_t,X_{t+\tau})=\acvs_\tau/\acvs_0$, with
$\acf_0=1$ and $\abs{\acf_\tau}\le1$.
\end{definition}

\begin{definition}{Sample autocovariance estimators}{sampleacvs}
Given $X_1,\dots,X_n$ with sample mean $\bar X$, two estimators are used for
$\abs{\tau}\le n-1$ (both $0$ otherwise):
\[
\tilde\acvs_\tau=\frac{1}{n-\abs\tau}\sum_{t=1}^{n-\abs\tau}(X_t-\bar X)(X_{t+\abs\tau}-\bar X),
\qquad
\hat\acvs_\tau=\frac{1}{n}\sum_{t=1}^{n-\abs\tau}(X_t-\bar X)(X_{t+\abs\tau}-\bar X).
\]
With known mean, $\tilde\acvs_\tau$ is \emph{unbiased} but need not be
positive semi-definite; $\hat\acvs_\tau$ is \emph{biased}
($\E[\hat\acvs_\tau]=\tfrac{n-\abs\tau}{n}\acvs_\tau$) but \emph{is} positive
semi-definite, which is why it is the default. The plug-in ACF estimator is
$\hat\acf_\tau=\hat\acvs_\tau/\hat\acvs_0$.
\end{definition}

\begin{definition}{Mean ergodicity}{ergodic}
$\{X_t\}$ is \textbf{mean ergodic} if its first two moments are finite and
$\bar X\toprob \E[X_t]$ as $n\to\infty$ (``temporal averages
$=$ population averages''). Ergodicity and stationarity are \emph{not}
equivalent.
\end{definition}

% ----------------------------------------------------------------------
\subsection{Theorems \& Propositions}

\begin{theorem}{Kolmogorov existence theorem}{kolmo}
A family of finite-dimensional distribution functions $\{F_{\mathbf t}\}$ are
the f.d.d.'s of a stochastic process \emph{iff} they are
\textbf{consistent}: for every $n$, every $\mathbf t$ and every $1\le i\le n$,
\[
\lim_{x_i\to\infty}F_{\mathbf t}(\mathbf x)=F_{\mathbf t(i)}(\mathbf x(i)),
\]
where $\mathbf t(i),\mathbf x(i)$ delete the $i$-th coordinate.\\[2pt]
\textit{Why:} it guarantees an infinite-index stochastic process actually
exists given only its finite-dimensional ``marginals'', justifying modelling
$\{X_t\}_{t\in\Ints}$ even though we observe finitely many points.
\end{theorem}

\begin{proposition}{The ACVS is positive semi-definite}{psd}
For any $n$, any $t_1,\dots,t_n\in\Ints$ and any $a_1,\dots,a_n\in\Reals$,
\[
\sum_{j=1}^n\sum_{k=1}^n a_ja_k\,\acvs_{j-k}\ge0 .
\]
\textit{Proof idea:} set $W=\sum_j a_jX_j$; then
$0\le\Var(W)=\sum_{j,k}a_ja_k\Cov(X_j,X_k)=\sum_{j,k}a_ja_k\acvs_{j-k}$.\\
\textit{Why:} p.s.d.\ is the defining algebraic property of a valid ACVS; a
candidate sequence that is \emph{not} p.s.d.\ \textbf{cannot} be an
autocovariance (used constantly to reject ``fake'' ACVS, see
Exercise in Part~\ref{sec:spectral}).
\end{proposition}

\begin{proposition}{Variance of the sample mean \& consistency}{xbarvar}
If $\sum_{\tau}\abs{\acvs_\tau}<\infty$ then $\bar X$ is unbiased and
\[
\Var(\bar X)=\frac1{n^2}\sum_{i,j=1}^n\acvs_{j-i}
=\frac1{n}\sum_{\tau=-(n-1)}^{n-1}\Big(1-\tfrac{\abs\tau}{n}\Big)\acvs_\tau,
\qquad
n\Var(\bar X)\xrightarrow[n\to\infty]{}\sum_{\tau=-\infty}^{\infty}\acvs_\tau .
\]
Hence $\Var(\bar X)\to0$ and $\bar X\toprob\mu$ (consistency).\\
\textit{Proof idea:} expand the double sum, collect by lag $\tau$, then apply
the \textbf{Ces\`aro summability theorem} to replace
$\sum(1-\tfrac{\abs\tau}{n})\acvs_\tau$ by $\sum\acvs_\tau$.\\
\textit{Why:} this is the prototype ``when can a sample average replace a
population average?''. Strong positive correlation inflates $\Var(\bar X)$
(e.g.\ AR(1) with $\phi\to1$): dependence, not just sample size, controls
estimation accuracy.
\end{proposition}

% ----------------------------------------------------------------------
\subsection{Key Techniques}

\begin{keytech}{Proving (or disproving) weak stationarity}{checkstat}
Check the three conditions \emph{in order}, and stop at the first failure:
\begin{enumerate}
  \item \textbf{Mean constant in $t$:} compute $\E[X_t]$; it must not depend on $t$.
  \item \textbf{Variance finite \& constant:} compute $\Var(X_t)$; must be
        finite and $t$-free.
  \item \textbf{Covariance shift-invariant:} compute $\Cov(X_t,X_{t+\tau})$ and
        check it depends on $\tau$ only.
\end{enumerate}
Common disqualifiers: undefined moments (Cauchy), or a covariance that retains
a $t$-dependence such as $\cos(ct)\cos(c(t+\tau))$.
\end{keytech}

\begin{keytech}{ACVS of a finite linear filter of white noise}{acvsfilter}
For $X_t=\sum_k\theta_k\eps_{t-k}$ with $\eps$ white noise of variance
$\sigma^2$, use bilinearity and $\Cov(\eps_s,\eps_u)=\sigma^2\ind_{\{s=u\}}$:
\[
\acvs_\tau=\Cov\!\Big(\sum_k\theta_k\eps_{t+\tau-k},\sum_j\theta_j\eps_{t-j}\Big)
=\sigma^2\sum_{j}\theta_j\theta_{j+\abs\tau}.
\]
Only the terms whose noise indices coincide survive --- this single move powers
\emph{all} MA covariance computations.
\end{keytech}

% ----------------------------------------------------------------------
\subsection{Exercises}

\begin{exercise}{Student-$t$ noise \quad\textnormal{[Sheet 1]}}{e1}
Let $X_t$ be i.i.d.\ Student-$t$ on one degree of freedom. Is $\{X_t\}$ weakly
stationary?
\end{exercise}
\begin{solution}
A Student-$t$ on $1$ d.o.f.\ is a standard Cauchy: its mean and variance are
\emph{undefined}. Condition (ii) of weak stationarity (finite, constant
variance) fails, so $\{X_t\}$ is \textbf{not} weakly stationary --- even though
it is strictly stationary (i.i.d.). A clean illustration that strict
$\not\Rightarrow$ weak.
\end{solution}

\begin{exercise}{Four candidate processes \quad\textnormal{[Sheet 1]}}{e2}
Let $Y_t\iid \mathcal N(0,\sigma_Y^2)$ and $a,b,c$ constants. Classify each as
weakly / strongly stationary and give the mean and ACVS:
(1) $X_t=a+bY_t+cY_{t-1}$;\;
(2) $X_t=a+bY_0$;\;
(3) $X_t=Y_1\cos(ct)+Y_2\sin(ct)$;\;
(4) $X_t=Y_0\cos(ct)$.
\end{exercise}
\begin{solution}
\textbf{(1)} $\E[X_t]=a$ and, using independence,
\[
\acvs_\tau=\begin{cases}(b^2+c^2)\sigma_Y^2 & \tau=0,\\ bc\,\sigma_Y^2 & \tau=\pm1,\\ 0 &\text{else.}\end{cases}
\]
Depends on $\tau$ only, finite $\Rightarrow$ weakly stationary; a linear
combination of Gaussians $\Rightarrow$ Gaussian $\Rightarrow$ also strongly
stationary. (This is an MA(1).)\\
\textbf{(2)} $\E[X_t]=a$, $\acvs_\tau=\Cov(bY_0,bY_0)=b^2\sigma_Y^2$ for
\emph{all} $\tau$ (the process is constant in randomness across $t$). Weakly and
strongly stationary.\\
\textbf{(3)} $\E[X_t]=0$ and, using
$\cos\alpha\cos\beta+\sin\alpha\sin\beta=\cos(\alpha-\beta)$,
\[
\acvs_\tau=\sigma_Y^2\big[\cos(ct)\cos(c(t{+}\tau))+\sin(ct)\sin(c(t{+}\tau))\big]
=\sigma_Y^2\cos(c\tau),
\]
a function of $\tau$ only $\Rightarrow$ weakly (and, being Gaussian, strongly)
stationary.\\
\textbf{(4)} $\E[X_t]=0$ but
$\acvs_\tau=\sigma_Y^2\cos(ct)\cos(c(t+\tau))$ \emph{depends on $t$}
$\Rightarrow$ \textbf{not} stationary. The contrast (3) vs (4) is the classic
trap: two random coefficients (3) restore stationarity that a single one (4)
destroys.
\end{solution}

\begin{exercise}{Sum of uncorrelated stationary series \quad\textnormal{[Sheet 1]}}{e3}
Let $\{X_t\},\{Y_t\}$ be stationary and mutually uncorrelated ($X_t,Y_s$
uncorrelated for all $t,s$). Show $\{X_t+Y_t\}$ is stationary with
$\acvs^{X+Y}_\tau=\acvs^X_\tau+\acvs^Y_\tau$.
\end{exercise}
\begin{solution}
Mean: $\E[X_t+Y_t]=\E[X_t]+\E[Y_t]$ is constant. Variance:
$\Var(X_t+Y_t)=\Var(X_t)+\Var(Y_t)<\infty$ since $\Cov(X_t,Y_t)=0$.
Covariance, by bilinearity and uncorrelatedness of the cross terms,
\[
\Cov(X_t+Y_t,X_{t+\tau}+Y_{t+\tau})=\acvs^X_\tau+\underbrace{\Cov(X_t,Y_{t+\tau})}_{0}
+\underbrace{\Cov(Y_t,X_{t+\tau})}_{0}+\acvs^Y_\tau=\acvs^X_\tau+\acvs^Y_\tau,
\]
a function of $\tau$ only. Hence stationary with additive ACVS.
\end{solution}

\begin{exercise}{Superposition of $m$ orthogonal series \quad\textnormal{[Sheet 1]}}{e4}
Let $\{X_{t,1}\},\dots,\{X_{t,m}\}$ be zero-mean stationary with
$\acvs_j(\tau)=\Cov(X_{t,j},X_{t+\tau,j})$, and $\E[X_{t,j}X_{t+\tau,k}]=0$ for
$j\ne k$. Find the ACVS of $X_t=\sum_{j=1}^m X_{t,j}$.
\end{exercise}
\begin{solution}
By bilinearity, $\acvs_X(\tau)=\sum_{j=1}^m\sum_{k=1}^m\Cov(X_{t,j},X_{t+\tau,k})$.
The cross terms ($j\ne k$) vanish by orthogonality, leaving only the diagonal:
$\acvs_X(\tau)=\sum_{j=1}^m\acvs_j(\tau)$. (Geometrically: sum a matrix that is
zero off the diagonal.)
\end{solution}

\begin{exercise}{Signal plus independent noise \quad\textnormal{[Sheet 1]}}{e5}
Let $\{X_t\}$ be second-order stationary with ACVS $\acvs_\tau$ and let
$\eps_t$ be i.i.d.\ with variance $\sigma^2$, independent of $\{X_t\}$. Find the
ACVS of $Z_t=X_t+\eps_t$.
\end{exercise}
\begin{solution}
Independence kills all cross-covariances. For $\tau\ne0$,
$\acvs_Z(\tau)=\acvs_\tau$; at $\tau=0$, $\acvs_Z(0)=\acvs_0+\sigma^2$:
\[
\acvs_Z(\tau)=\acvs_\tau+\sigma^2\ind_{\{\tau=0\}} .
\]
Adding white noise inflates only the variance (the $\tau=0$ ``nugget''),
leaving every other lag unchanged.
\end{solution}

\begin{exercise}{Bias, variance and MSE of the two ACVS estimators \quad\textnormal{[Sheet 2]}}{e6}
For $X_t$ i.i.d.\ Gaussian white noise (variance $\sigma_X^2$, known mean
replaced for $\bar X$), find the mean, variance and MSE of $\tilde\acvs_\tau$
and $\hat\acvs_\tau$, and say which wins.
\end{exercise}
\begin{solution}
Both are unbiased for white noise: $\E[\tilde\acvs_\tau]=\E[\hat\acvs_\tau]=0$
for $\tau\ne0$ (and $\sigma_X^2$ at $\tau=0$). Using Isserlis' theorem for the
fourth moments one finds, for $\tau\ne0$,
\[
\Var(\hat\acvs_\tau)=\frac{n-\abs\tau}{n^2}\sigma_X^4,
\qquad
\Var(\tilde\acvs_\tau)=\frac{1}{n-\abs\tau}\sigma_X^4 .
\]
Since the estimators are unbiased here, $\MSE=\Var$, and
$\Var(\hat\acvs_\tau)=\big(\tfrac{n-\abs\tau}{n}\big)^2\Var(\tilde\acvs_\tau)
\le\Var(\tilde\acvs_\tau)$. The \textbf{biased} divide-by-$n$ estimator
$\hat\acvs_\tau$ has the \emph{smaller} MSE --- a recurring lesson that
unbiasedness is not the goal, low MSE is.
\end{solution}

\begin{exercise}{Sample autocovariances sum to zero \quad\textnormal{[Sheet 2 / New angle]}}{e7}
Show that for any data $\{x_1,\dots,x_n\}$, $\sum_{\abs\tau<n}\hat\acvs_\tau=0$.
\end{exercise}
\begin{solution}
\[
\sum_{\tau=-(n-1)}^{n-1}\hat\acvs_\tau
=\frac1n\sum_{t_1=1}^n\sum_{t_2=1}^n(x_{t_1}-\bar x)(x_{t_2}-\bar x)
=\frac1n\Big(\sum_{t=1}^n(x_t-\bar x)\Big)^2=0,
\]
because the double sum factorises and the centred data sum to zero. (The first
equality just re-indexes the triangular lag sum as a full $t_1,t_2$ square.)
\end{solution}

\begin{exercise}{ACVS of a two-tap filter \quad\textnormal{[New]}}{e8}
Let $\eps_t$ be mean-zero white noise with variance $\sigma^2$ and set
$X_t=\theta_1\eps_t+\theta_2\eps_{t-1}$. Compute the ACVS and the ACF, and state
for which $\theta_1,\theta_2$ the lag-1 correlation is maximised in magnitude.
\end{exercise}
\begin{solution}
By the white-noise filter technique,
\[
\acvs_\tau=\begin{cases}(\theta_1^2+\theta_2^2)\sigma^2 & \tau=0,\\
\theta_1\theta_2\sigma^2 & \tau=\pm1,\\ 0 &\text{else,}\end{cases}
\qquad
\acf_1=\frac{\theta_1\theta_2}{\theta_1^2+\theta_2^2}.
\]
Since $\abs{\theta_1\theta_2}\le\tfrac12(\theta_1^2+\theta_2^2)$ (AM--GM), we
have $\abs{\acf_1}\le\tfrac12$, with equality iff $\abs{\theta_1}=\abs{\theta_2}$.
So \textbf{no MA(1) can have $\abs{\acf_1}>\tfrac12$} --- a fact worth
remembering as a sanity check on any claimed MA(1) ACF.
\end{solution}

% ======================================================================
%  PART 2 -- ARMA Models: White Noise, MA, AR, ARMA and the PACF
% ======================================================================
\section{ARMA Models: White Noise, MA, AR, ARMA and the PACF}
\label{sec:arma}

The ARMA family is the workhorse of linear time series modelling. Starting from
\textbf{white noise} --- the indivisible ``atom'' of randomness --- we build
\textbf{moving average} (MA), \textbf{autoregressive} (AR) and combined
\textbf{ARMA} processes by feeding white noise through finite linear filters and
linear recursions. Two diagnostic curves run through the whole story: the
autocorrelation function (ACF, $\acf_\tau$) and its less-obvious partner, the
\textbf{partial autocorrelation function} (PACF, $\pacf_\tau$). The punchline,
encoded in one small table, is that the ACF identifies the MA order $q$ while
the PACF identifies the AR order $p$. We close with the
\textbf{Durbin--Levinson} recursions that compute the PACF directly from the
ACF.

% ----------------------------------------------------------------------
\subsection{Definitions}

\begin{definition}{White noise (purely random process)}{p2-wn}
$\{X_t\}$ is \textbf{white noise} (a purely random process) if it is a sequence
of \emph{uncorrelated}, mean-zero, finite-variance random variables: for all
$t\in\Ints$,
\[
\E[X_t]=0,\qquad \Var(X_t)=\sigma^2<\infty,
\]
and $\Cov(X_s,X_t)=0$ for $s\ne t$. Equivalently its second-order structure is
\[
\acvs_\tau=\begin{cases}\sigma^2 & \tau=0,\\ 0 & \tau\ne0,\end{cases}
\qquad
\acf_\tau=\begin{cases}1 & \tau=0,\\ 0 & \tau\ne0.\end{cases}
\]
White noise is automatically (weakly) stationary and is the \textbf{building
block} of every model below. We write $\eps_t$ for a mean-zero white noise of
variance $\sigma^2_\eps$.
\end{definition}

\begin{intuition}
``Uncorrelated'' is weaker than ``independent''. White noise only constrains the
first two moments, so a white-noise process need not be i.i.d.\ and need not be
Gaussian. Calling something \emph{white} refers to its flat spectral density
(equal power at all frequencies), which mirrors the flat ACF above.
\end{intuition}

\begin{definition}{Moving average process MA($q$)}{p2-maq}
Let $\{\eps_t\}$ be mean-zero white noise of variance $\sigma_\eps^2$. The
\textbf{$q$-th order moving average process} MA($q$) is
\[
X_t=\mu-\sum_{j=0}^{q}\theta_j\,\eps_{t-j},
\qquad \theta_0=-1,\ \ \theta_q\ne0 .
\]
(The convention $\theta_0=-1$ makes the leading term $-\theta_0\eps_t=\eps_t$.)
Simple examples use the equivalent zero-mean one-sign convention
$X_t=\eps_t-\theta\eps_{t-1}$ for MA(1). An MA($q$) is just a finite linear
filter of white noise, hence always stationary; the only requirement is
$\abs{\theta_j}<\infty$.
\end{definition}

\begin{definition}{Moments of an MA($q$)}{p2-mamom}
For the MA($q$) of Definition~\ref{def:p2-maq}, with $\theta_j:=0$ for $j>q$,
\[
\E[X_t]=\mu,
\qquad
\acvs_\tau=\Cov(X_{t+\tau},X_t)=
\begin{cases}
\sigma^2_\eps\displaystyle\sum_{j=0}^{q}\theta_j\theta_{j+\abs\tau} & \abs\tau\le q,\\[2mm]
0 & \abs\tau>q .
\end{cases}
\]
Hence the \textbf{ACVS of an MA($q$) cuts off (is exactly $0$) for lags
$\abs\tau>q$} --- this is the signature used to read $q$ off the ACF.
\end{definition}

\begin{definition}{Autoregressive process AR($p$)}{p2-arp}
Let $\{\eps_t\}$ be mean-zero white noise. The \textbf{$p$-th order
autoregressive process} AR($p$) is
\[
X_t=\sum_{j=1}^{p}\phi_j X_{t-j}+\eps_t,
\qquad \phi_p\ne0 .
\]
Unlike the MA case there \emph{are} constraints on $\{\phi_j\}$ for stationarity
(the roots of the AR polynomial must lie outside the unit circle), and there is
no simple closed form for the ACVS in general.
\end{definition}

\begin{definition}{ARMA($p,q$) process}{p2-arma}
$\{X_t\}$ is an \textbf{autoregressive moving average} process of orders $p$ and
$q$, ARMA($p,q$), if
\[
X_t=\sum_{j=1}^{p}\phi_j X_{t-j}-\sum_{k=0}^{q}\theta_k\eps_{t-k},
\qquad \theta_0=-1,
\]
with $\{\eps_t\}$ mean-zero white noise and $\phi_j,\theta_k$ as in the AR and MA
cases. Using the backshift operator $\Bsh$ (with $\Bsh^j X_t=X_{t-j}$) this is
written compactly as
\[
\Phi(\Bsh)X_t=\Theta(\Bsh)\eps_t,
\qquad
\Phi(z)=1-\sum_{j=1}^{p}\phi_j z^{j},
\quad
\Theta(z)=-\sum_{k=0}^{q}\theta_k z^{k}.
\]
As for AR, conditions on the $\phi_j$ are needed for the model to be valid, and
closed-form ACVS is generally unavailable (though computable).
\end{definition}

\begin{definition}{Causal ARMA($p,q$)}{p2-causal}
An ARMA($p,q$) process is \textbf{causal} (more precisely a causal function of
$\{\eps_t\}$) if there exists a sequence $\{\psi_j\}$ with
$\sum_{j=0}^{\infty}\abs{\psi_j}<\infty$ such that
\[
X_t=\sum_{j=0}^{\infty}\psi_j\,\eps_{t-j},
\qquad t=0,\pm1,\pm2,\dots
\]
i.e.\ $X_t$ depends only on \emph{present and past} noise. The $\{\psi_j\}$ are
the \textbf{$\psi$-weights} (the MA($\infty$) representation). A mean-zero
MA($q$) is \emph{always} causal (it already has this form with finitely many
nonzero $\psi_j$). For ARMA, causality holds iff all roots of $\Phi(z)$ lie
\emph{outside} the unit circle.
\end{definition}

\begin{recallbox}[Roots convention]
\textbf{Stationarity / causality:} all roots of the AR polynomial $\Phi(z)$ lie
\emph{outside} the unit circle ($\abs z>1$). \textbf{Invertibility:} all roots
of the MA polynomial $\Theta(z)$ lie \emph{outside} the unit circle. For AR(1),
$\Phi(z)=1-\phi z$ has root $z=1/\phi$, so $\abs z>1\iff\abs\phi<1$.
\end{recallbox}

\begin{definition}{Best linear predictor}{p2-blp}
Given mean-zero random variables $X_1,\dots,X_n,Y$, the \textbf{best linear
predictor} of $Y$ from $X_1,\dots,X_n$ is
\[
P_{X_1,\dots,X_n}(Y)=\sum_{j=1}^{n}\beta_j X_j,
\]
where $\beta_1,\dots,\beta_n$ minimise the mean-squared error
$\E\big[(Y-P_{X_1,\dots,X_n}(Y))^2\big]$. The minimiser is characterised by the
\textbf{prediction (orthogonality) equations}: the prediction error is
uncorrelated with every regressor,
\[
\Cov\big(Y-P_{X_1,\dots,X_n}(Y),\,X_k\big)=0,\qquad k=1,\dots,n .
\]
\end{definition}

\begin{definition}{Partial correlation}{p2-partcorr}
For random variables $X,Y$ and confounders $Z=(Z_1,\dots,Z_n)\Tr$, the
\textbf{partial correlation} of $X$ and $Y$ given $Z$ is
\[
\acf_{XY\bullet Z}=\Corr\big(X-P_Z(X),\,Y-P_Z(Y)\big),
\]
i.e.\ the correlation of the residuals after removing the best linear effect of
$Z$ from each. If $X,Y,Z$ are jointly Gaussian then
$\acf_{XY\bullet Z}=\E[\Corr(X,Y\mid Z)]$, and in fact
$\acf_{XY\bullet Z}\overset{\text{a.s.}}{=}\Corr(X,Y\mid Z)$.
\end{definition}

\begin{definition}{Partial autocorrelation function (PACF)}{p2-pacf}
For a mean-zero stationary $\{X_t\}$, the \textbf{PACF} at lag $\tau$ is the
partial correlation of $X_t$ and $X_{t+\tau}$ given the intervening values
$Z=(X_{t+1},\dots,X_{t+\tau-1})$. Writing
$\hat X_t=P_{X_{t+1},\dots,X_{t+\tau-1}}(X_t)$ and
$\hat X_{t+\tau}=P_{X_{t+1},\dots,X_{t+\tau-1}}(X_{t+\tau})$,
\[
\pacf_\tau=\Corr\big(X_{t+\tau}-\hat X_{t+\tau},\,X_t-\hat X_t\big).
\]
Equivalently (regression form): if
$P_{X_0,\dots,X_{\tau-1}}(X_\tau)=\sum_{j=1}^{\tau}\pacf_{\tau,j}X_{\tau-j}$ is
the best linear predictor of $X_\tau$ from its $\tau$ predecessors, then
$\pacf_\tau=\pacf_{\tau,\tau}$ (the \emph{last} regression coefficient).
Conventionally $\pacf_1=\acf_1$.
\end{definition}

\begin{intuition}
The PACF answers ``how much \emph{new} linear information about $X_{t+\tau}$ is
in $X_t$, once everything in between is accounted for?''. In an AR(1),
$\acf_\tau=\phi^{\abs\tau}\ne0$ for all $\tau$ even though $X_t$ only directly
depends on $X_{t-1}$: the correlation at lag $2$ is entirely \emph{transmitted}
through $X_{t-1}$. The PACF strips out that transmitted part, so it vanishes for
$\tau>1$ --- exactly identifying the order $p=1$.
\end{intuition}

% ----------------------------------------------------------------------
\subsection{Theorems \& Propositions}

\begin{proposition}{Causal AR(1) expansion and its ACVS}{p2-ar1}
For the AR(1) $X_t=\phi X_{t-1}+\eps_t$ with $\abs\phi<1$, iterating the
recursion gives the causal (MA($\infty$)) representation
\[
X_t=\sum_{k=0}^{\infty}\phi^{k}\eps_{t-k},
\qquad \psi_k=\phi^k,\ \ \sum_k\abs{\psi_k}=\tfrac1{1-\abs\phi}<\infty,
\]
with $\E[X_t]=0$ and autocovariance
\[
\acvs_\tau=\frac{\sigma_\eps^2}{1-\phi^2}\,\phi^{\abs\tau},
\qquad
\acf_\tau=\phi^{\abs\tau}.
\]
\textit{Proof idea:} substitute $X_{t-1}=\phi X_{t-2}+\eps_{t-1}$ repeatedly;
the remainder $\phi^{n}X_{t-n}\to0$ in mean square. Then
$\acvs_\tau=\sigma_\eps^2\sum_{k\ge0}\phi^k\phi^{k+\abs\tau}
=\sigma_\eps^2\phi^{\abs\tau}\sum_{k\ge0}\phi^{2k}
=\sigma_\eps^2\phi^{\abs\tau}/(1-\phi^2)$ by the geometric series.\\
\textit{Why:} the AR(1) is the canonical example of a causal process and the
template for ``write an AR as an infinite MA, then use the white-noise filter
trick to get the ACVS''. The ACF $\phi^{\abs\tau}$ \emph{tails off} (never hits
exactly $0$), motivating the PACF.
\end{proposition}

\begin{proposition}{ACF--PACF relation (prediction equations)}{p2-acfpacf}
For a mean-zero stationary $\{X_t\}$, fix $\tau>0$ and write
$\hat X_\tau=P_{X_0,\dots,X_{\tau-1}}(X_\tau)=\sum_{j=1}^{\tau}\pacf_{\tau,j}X_{\tau-j}$.
Then for each $k\in\{1,\dots,\tau\}$,
\[
\acf_k=\sum_{j=1}^{\tau}\pacf_{\tau,j}\,\acf_{k-j}.
\]
\textit{Proof idea:} decompose
$\acvs_k=\Cov(X_{\tau-k},X_\tau)
=\Cov(X_{\tau-k},\hat X_\tau)+\Cov(X_{\tau-k},X_\tau-\hat X_\tau)$;
the second term is $0$ by the orthogonality of the prediction error, and
expanding the first using $\hat X_\tau=\sum_j\pacf_{\tau,j}X_{\tau-j}$ gives
$\acvs_k=\sum_j\pacf_{\tau,j}\acvs_{k-j}$; divide by $\acvs_0$.\\
\textit{Why:} these are the \textbf{Yule--Walker-type} equations linking the ACF
to the prediction coefficients $\pacf_{\tau,j}$. They are exactly what the
Durbin--Levinson recursion solves to extract the PACF
$\pacf_\tau=\pacf_{\tau,\tau}$.
\end{proposition}

\begin{theorem}{ACF \& PACF signatures of ARMA models}{p2-table}
For \emph{causal and invertible} ARMA($p,q$) models the ACF and PACF behave as:
\begin{center}
\renewcommand{\arraystretch}{1.25}
\begin{tabular}{lccc}
\toprule
 & \textbf{AR($p$)} & \textbf{MA($q$)} & \textbf{ARMA($p,q$)}\\
\midrule
\textbf{ACF}  & Tails off            & Cuts off after lag $q$ & Tails off\\
\textbf{PACF} & Cuts off after lag $p$ & Tails off             & Tails off\\
\bottomrule
\end{tabular}
\end{center}
\textit{Proof idea:} the ACF cut-off for MA($q$) is the moment computation of
Definition~\ref{def:p2-mamom}; the PACF cut-off for AR($p$) is
Proposition~\ref{prop:p2-arpacf}; ``tails off'' means a geometric/damped-sine
decay with no exact zero (as for the AR(1) ACF $\phi^{\abs\tau}$).\\
\textit{Why:} this $2\times3$ table is \emph{the} model-identification tool:
read $q$ from where the ACF cuts, $p$ from where the PACF cuts; if neither cuts
(both tail off) the model is mixed ARMA.
\end{theorem}

\begin{proposition}{PACF of an AR($p$) vanishes beyond lag $p$}{p2-arpacf}
For a causal AR($p$),
\[
\pacf_\tau=0\qquad\text{for all }\tau>p .
\]
Moreover the \emph{estimated} PACF at lags $>p$ is, asymptotically, mean $0$
with variance $1/n$ ($n$ the sample size), giving the $\pm1.96/\sqrt n$
significance bands.\\
\textit{Proof idea:} for $\tau>p$ write the prediction error
$X_\tau-\sum_{j=1}^{\tau}\pacf_{\tau,j}X_{\tau-j}$. Because the process is
AR($p$), the true best predictor uses only $\phi_1,\dots,\phi_p$ and the noise
$\eps_\tau$; setting $\tilde\phi_j=\phi_j$ for $j\le p$ and $\tilde\phi_j=0$ for
$j>p$, the prediction equations
$\Cov(X_\tau-\sum_j\pacf_{\tau,j}X_{\tau-j},X_k)=0$ are solved by
$\pacf_{\tau,j}=\tilde\phi_j$ (using $\Cov(\eps_\tau,X_k)=0$ for $k<\tau$, by
causality). In particular $\pacf_\tau=\pacf_{\tau,\tau}=\tilde\phi_\tau=0$ for
$\tau>p$.\\
\textit{Why:} this is the theoretical backbone of the table: it is why a sharp
cut-off in the (sample) PACF at lag $p$ diagnoses an AR($p$), and the $1/n$
variance gives the threshold for ``cut off''.
\end{proposition}

\begin{proposition}{Non-identifiability of MA(1) from the ACVS}{p2-maident}
The MA(1) processes
\[
X_t=\eps_t-\theta\eps_{t-1}\ \ (\Var\eps_t=\sigma_\eps^2)
\qquad\text{and}\qquad
Y_t=\eta_t-\tfrac1\theta\eta_{t-1}\ \ (\Var\eta_t=\sigma_\eps^2\theta^2)
\]
have \emph{identical} autocovariance sequences, hence are indistinguishable from
second-order structure alone.\\
\textit{Proof idea:} both give
$\acvs_0=\sigma_\eps^2(1+\theta^2)$, $\acvs_{\pm1}=-\sigma_\eps^2\theta$, and $0$
elsewhere (direct computation, see Exercise~\ref{exo:p2-maident-ex}).\\
\textit{Why:} an MA(1) and its ``root-reciprocal'' twin share an ACVS; choosing
the \emph{invertible} representation ($\abs\theta<1$, root outside the unit
circle) pins down a unique model. This is the reason invertibility is imposed.
\end{proposition}

% ----------------------------------------------------------------------
\subsection{Key Techniques}

\begin{keytech}{Computing the ACVS of an MA($q$)}{p2-ktma}
Use bilinearity of covariance and
$\Cov(\eps_s,\eps_u)=\sigma_\eps^2\ind_{\{s=u\}}$. For
$X_t=\mu-\sum_{j=0}^q\theta_j\eps_{t-j}$ (constant $\mu$ drops out of
covariances),
\[
\acvs_\tau=\Cov\!\Big(\sum_{k}\theta_k\eps_{t+\tau-k},\sum_j\theta_j\eps_{t-j}\Big)
=\sigma_\eps^2\sum_{j}\theta_j\theta_{j+\abs\tau},
\]
the surviving terms being those with $k=j+\tau$. Steps: (1) write both sums; (2)
keep only matching noise indices; (3) read off that $\acvs_\tau=0$ for
$\abs\tau>q$.
\end{keytech}

\begin{keytech}{Checking causality / stationarity via roots}{p2-ktroots}
To test an AR or ARMA: (1) form $\Phi(z)=1-\sum_{j=1}^p\phi_j z^j$; (2) find its
roots; (3) the model is causal/stationary iff every root satisfies $\abs z>1$
(outside the unit circle). For AR(1) this reduces to $\abs\phi<1$; for AR(2)
$\Phi(z)=1-\phi_1 z-\phi_2 z^2$, the conditions are
$\phi_1+\phi_2<1,\ \phi_2-\phi_1<1,\ \abs{\phi_2}<1$. Invertibility uses the same
test on $\Theta(z)$.
\end{keytech}

\begin{keytech}{Deriving the MA($\infty$) ($\psi$-weight) representation}{p2-ktpsi}
Match coefficients in $\Phi(\Bsh)\,\psi(\Bsh)=\Theta(\Bsh)$ where
$\psi(z)=\sum_{j\ge0}\psi_j z^j$. Equivalently expand
$\psi(z)=\Theta(z)/\Phi(z)$ as a power series. For AR(1):
$\psi(z)=1/(1-\phi z)=\sum_{j\ge0}\phi^j z^j$, so $\psi_j=\phi^j$. For an
ARMA(1,1) written $X_t=\phi X_{t-1}+\eps_t+\vartheta\eps_{t-1}$
(so $\Phi(z)=1-\phi z$, $\Theta(z)=1+\vartheta z$):
$\psi_0=1$ and $\psi_j=(\phi+\vartheta)\phi^{\,j-1}$ for $j\ge1$. Once you have
$\{\psi_j\}$, the ACVS is
$\acvs_\tau=\sigma_\eps^2\sum_{j\ge0}\psi_j\psi_{j+\abs\tau}$.
\end{keytech}

\begin{keytech}{Durbin--Levinson recursion for the PACF}{p2-ktdl}
Given the ACF $\{\acf_\tau\}$, compute $\pacf_\tau=\pacf_{\tau,\tau}$ by:
\[
\pacf_{1,1}=\acf_1,\qquad
\pacf_{\tau,\tau}
=\frac{\acf_\tau-\sum_{j=1}^{\tau-1}\pacf_{\tau-1,j}\,\acf_{\tau-j}}
       {1-\sum_{j=1}^{\tau-1}\pacf_{\tau-1,j}\,\acf_{j}},
\]
then update the auxiliary coefficients
\[
\pacf_{\tau,j}=\pacf_{\tau-1,j}-\pacf_{\tau,\tau}\,\pacf_{\tau-1,\tau-j},
\qquad j=1,\dots,\tau-1 .
\]
Works with theoretical \emph{or} sample ACF. A clean cut to $0$ in
$\pacf_{\tau,\tau}$ at $\tau=p+1$ flags an AR($p$); the last diagonal
coefficients then equal the AR coefficients.
\end{keytech}

\begin{keytech}{Reading $(p,q)$ off ACF \& PACF plots}{p2-ktident}
\begin{enumerate}
  \item First confirm \textbf{stationarity}: stable mean/variance; an ACF that
        decays very slowly (almost linearly) signals non-stationarity ---
        difference first.
  \item \textbf{ACF cuts after lag $q$, PACF tails off} $\Rightarrow$ MA($q$).
  \item \textbf{PACF cuts after lag $p$, ACF tails off} $\Rightarrow$ AR($p$).
  \item \textbf{Both tail off} $\Rightarrow$ mixed ARMA($p,q$).
  \item \textbf{Neither significant at any nonzero lag} $\Rightarrow$ white noise.
\end{enumerate}
Use the $\pm1.96/\sqrt n$ band to decide ``significant''. Choices are rarely
unique --- justify with a sensible argument.
\end{keytech}

\begin{pitfall}
Two recurring sign/convention traps. (i) The notes' MA/ARMA definition carries a
\emph{minus} sign in front of the $\theta$-sum with $\theta_0=-1$, so
$X_t=\mu-\sum_{j=0}^q\theta_j\eps_{t-j}$ starts as $\eps_t+\cdots$; simple
examples instead write $X_t=\eps_t-\theta\eps_{t-1}$. Track which convention a
problem uses before plugging into the ACVS formula. (ii) ``Roots outside the
unit circle'' is for the polynomials $\Phi(z),\Theta(z)$ in $z$; in terms of the
coefficient it means $\abs\phi<1$ for AR(1). Do not confuse it with ``roots
inside''.
\end{pitfall}

% ----------------------------------------------------------------------
\subsection{Exercises}

\begin{exercise}{Why fix $\theta_0=-1$? (MA identifiability) \quad\textnormal{[Sheet 2]}}{p2-theta0}
When defining the MA($q$) we set $\theta_0=-1$. In terms of covariance
structure, why do we gain no generality by letting $\theta_0$ be arbitrary?
\end{exercise}
\begin{solution}
Rescaling the leading coefficient produces a \emph{different} parametrisation
with the \emph{same} covariance, so $\theta_0$ is not free --- we fix it for
identifiability. Take the zero-mean MA(1)
\[
X_t=-\theta_0\eps_t-\theta_1\eps_{t-1}=\eps_t-\theta_1\eps_{t-1},
\qquad \eps_t\sim\mathcal N(0,\sigma_\eps^2),\ \theta_0=-1 .
\]
Suppose a second parametrisation uses $\theta_0'=c\theta_0=c$ and $\theta_1'$
with noise $\eps_t'\sim\mathcal N(0,\sigma_{\eps'}^2)$:
$X_t=-\theta_0'\eps_t'-\theta_1'\eps_{t-1}'=c\eps_t'-\theta_1'\eps_{t-1}'$.
Rewriting the first form,
\[
X_t=\eps_t-\theta_1\eps_{t-1}
=c\,\tfrac{\eps_t}{c}-(c\theta_1)\tfrac{\eps_{t-1}}{c},
\]
so $\eps_t'=c^{-1}\eps_t$, $\theta_1'=c\theta_1$, and
$\sigma_{\eps'}^2=c^{-2}\sigma_\eps^2$. The two representations share moments:
\[
\Var(X_t)=(1+\theta_1^2)\sigma_\eps^2,
\qquad
\{(\theta_0')^2+(\theta_1')^2\}\sigma_{\eps'}^2
=c^2(1+\theta_1^2)c^{-2}\sigma_\eps^2=(1+\theta_1^2)\sigma_\eps^2,
\]
and $\Cov(X_t,X_{t-1})=\theta_0\theta_1\sigma_\eps^2$ is likewise unchanged.
Thus many MA(1) representations give the same $X_t$; fixing $\theta_0=-1$ makes
the representation unique. The same scaling argument extends verbatim to
MA($q$).
\end{solution}

\begin{exercise}{Rewriting an MA(1) recursively \quad\textnormal{[Sheet 2]}}{p2-marewrite}
Show that the MA(1) process $X_t=\eps_t-\theta\eps_{t-1}$ can be written, for any
positive integer $p$, as
\[
X_t=\eps_t-\sum_{j=1}^{p}\theta^{j}X_{t-j}-\theta^{p+1}\eps_{t-p-1}.
\]
\end{exercise}
\begin{solution}
For any $p\ge1$, $X_{t-p}=\eps_{t-p}-\theta\eps_{t-p-1}$, hence
$\eps_{t-p}=X_{t-p}+\theta\eps_{t-p-1}$. Substitute this recursively:
\[
\begin{aligned}
X_t&=\eps_t-\theta\eps_{t-1}
   =\eps_t-\theta\,(X_{t-1}+\theta\eps_{t-2})
   =\eps_t-\theta X_{t-1}-\theta^2\eps_{t-2}\\
   &=\eps_t-\theta X_{t-1}-\theta^2(X_{t-2}+\theta\eps_{t-3})
   =\eps_t-\theta X_{t-1}-\theta^2 X_{t-2}-\theta^3\eps_{t-3}\\
   &=\ \cdots\ =\eps_t-\sum_{j=1}^{p}\theta^{j}X_{t-j}-\theta^{p+1}\eps_{t-p-1}.
\end{aligned}
\]
If $\abs\theta<1$ (invertible MA), the remainder $\theta^{p+1}\eps_{t-p-1}\to0$
in mean square as $p\to\infty$, giving the AR($\infty$) representation
$X_t=\eps_t-\sum_{j\ge1}\theta^j X_{t-j}$.
\end{solution}

\begin{exercise}{Moments of a sum of two filtered noises \quad\textnormal{[Sheet 2]}}{p2-sumfilter}
Let $Y_t=\sum_{s=-p}^{p}g_{s-t}\eps_s$ and $Z_t=\sum_{s=-p}^{p}h_{s-t}\eps_s$ with
$\{\eps_t\}$ mean-zero white noise of variance $\sigma^2$, and define
$X_t=Y_t+Z_t$. Determine the first and second moments of $X_t$.
\end{exercise}
\begin{solution}
Combine the filters: $X_t=\sum_{s=-p}^{p}(g_{s-t}+h_{s-t})\eps_s$. Since
$\E[\eps_s]=0$,
\[
\E[X_t]=\sum_{s=-p}^{p}(g_{s-t}+h_{s-t})\,\E[\eps_s]=0 .
\]
For the (auto)covariance, use $\Cov(\eps_s,\eps_l)=\sigma^2\ind_{\{s=l\}}$:
\[
\begin{aligned}
\E[X_t X_{t+\tau}]=\Cov(X_t,X_{t+\tau})
&=\Cov\!\Big(\sum_{s=-p}^{p}(g_{s-t}+h_{s-t})\eps_s,\
            \sum_{l=-p}^{p}(g_{l-t-\tau}+h_{l-t-\tau})\eps_l\Big)\\
&=\sum_{s=-p}^{p}\sum_{l=-p}^{p}(g_{s-t}+h_{s-t})(g_{l-t-\tau}+h_{l-t-\tau})\,
   \Cov(\eps_s,\eps_l)\\
&=\sigma^2\sum_{s=-p}^{p}(g_{s-t}+h_{s-t})(g_{s-t-\tau}+h_{s-t-\tau}).
\end{aligned}
\]
Only the diagonal $s=l$ survives, collapsing the double sum to a single one.
(Note: because the filter is indexed by $s-t$ over the fixed window
$\{-p,\dots,p\}$, this is a covariance computation rather than a stationarity
claim --- the surviving sum can depend on $t$ near the window edges.)
\end{solution}

\begin{exercise}{PACF of an AR($p$) is zero beyond lag $p$ \quad\textnormal{[Sheet 2]}}{p2-arpacf-ex}
Prove that for a causal AR($p$), $\pacf_\tau=0$ for all $\tau>p$. Use the
prediction equations
$\Cov(Y-P_{X_1,\dots,X_n}(Y),X_k)=0$ and the causal fact
$\Cov(X_t,\eps_s)=0$ for $t<s$.
\end{exercise}
\begin{solution}
Recall $\pacf_\tau=\pacf_{\tau,\tau}$, the last coefficient in
$P_{X_0,\dots,X_{\tau-1}}(X_\tau)=\sum_{j=1}^{\tau}\pacf_{\tau,j}X_{\tau-j}$.
Take $\tau>p$. The prediction equations state, for all
$k\in\{0,\dots,\tau-1\}$,
\[
\Cov\!\Big(X_\tau-\sum_{j=1}^{\tau}\pacf_{\tau,j}X_{\tau-j},\,X_k\Big)=0 .
\]
Substitute the AR($p$) form $X_\tau=\sum_{j=1}^{p}\phi_j X_{\tau-j}+\eps_\tau$
and define $\tilde\phi_j=\phi_j$ for $j\le p$, $\tilde\phi_j=0$ for $j>p$:
\[
\Cov\!\Big(\sum_{j=1}^{\tau}(\tilde\phi_j-\pacf_{\tau,j})X_{\tau-j}+\eps_\tau,\,X_k\Big)
=\sum_{j=1}^{\tau}(\tilde\phi_j-\pacf_{\tau,j})\Cov(X_{\tau-j},X_k)
+\Cov(\eps_\tau,X_k).
\]
By causality $X_k$ depends only on $\eps_s$ with $s\le k<\tau$, so
$\Cov(\eps_\tau,X_k)=0$. Hence choosing $\pacf_{\tau,j}=\tilde\phi_j$ for all $j$
satisfies every prediction equation; by uniqueness of the best linear predictor
these are the coefficients. In particular
$\pacf_\tau=\pacf_{\tau,\tau}=\tilde\phi_\tau=0$ for $\tau>p$. $\qquad\blacksquare$
\end{solution}

\begin{exercise}{Model identification from ACF \& PACF \quad\textnormal{[Sheet 2]}}{p2-ident}
For six time series, each shown with its sample ACF and PACF, decide which of
white noise, MA($q$), AR($p$), ARMA($p,q$) is appropriate, with reasoning.
\end{exercise}
\begin{solution}
The diagnostic is the signature table (Theorem~\ref{thm:p2-table}). The (hidden)
truths and sample arguments were:
\begin{enumerate}
  \item Stationary-looking; \emph{both} ACF and PACF decay without a sharp
        cut-off $\Rightarrow$ \textbf{ARMA} (truth: ARMA(1,1); one could argue
        AR(2)).
  \item Stationary; ACF cuts sharply after lag $1$ while the PACF decays
        $\Rightarrow$ \textbf{MA(1)}.
  \item Stationary; ACF decays, PACF cuts after lag $1$ $\Rightarrow$
        \textbf{AR(1)} (an ARMA is also arguable).
  \item Stationary; neither ACF nor PACF significant at nonzero lags
        $\Rightarrow$ \textbf{white noise}.
  \item Very strong lag-1 PACF, ACF decaying very slowly, apparent changing
        variance $\Rightarrow$ likely \textbf{non-stationary} (could be a
        near-unit-root AR(1) with $\phi\approx1$).
  \item Stationary; ACF decays, PACF cuts after lag $1$ $\Rightarrow$
        \textbf{AR(1)}.
\end{enumerate}
The point is to make a defensible argument; model selection is not unique, and
ACF/PACF are sanity checks, not proofs.
\end{solution}

\begin{exercise}{Durbin--Levinson: $\acf_1=\tfrac25,\ \acf_2=-\tfrac1{20},\ \acf_3=-\tfrac18$ \quad\textnormal{[Lecture worked example]}}{p2-dl}
A stationary process has $\acf_1=\tfrac25,\ \acf_2=-\tfrac1{20},\
\acf_3=-\tfrac18$. Find $\pacf_1,\pacf_2,\pacf_3$, and identify a plausible model.
\end{exercise}
\begin{solution}
Apply the Durbin--Levinson recursion (Key Technique~\ref{kt:p2-ktdl}).

\textbf{Lag 1:} $\ \pacf_{1,1}=\acf_1=\dfrac25 .$

\textbf{Lag 2:}
\[
\pacf_{2,2}=\frac{\acf_2-\pacf_{1,1}\acf_1}{1-\pacf_{1,1}\acf_1}
=\frac{-\frac1{20}-\frac25\cdot\frac25}{1-\frac25\cdot\frac25}
=\frac{-\frac1{20}-\frac4{25}}{1-\frac4{25}}
=\frac{-\frac{21}{100}}{\frac{84}{100}}
=-\frac{21}{84}=-\frac14 .
\]
Update: $\ \pacf_{2,1}=\pacf_{1,1}-\pacf_{2,2}\pacf_{1,1}
=\tfrac25-(-\tfrac14)\tfrac25=\tfrac25+\tfrac1{10}=\tfrac12 .$

\textbf{Lag 3:}
\[
\pacf_{3,3}
=\frac{\acf_3-\pacf_{2,1}\acf_2-\pacf_{2,2}\acf_1}
       {1-\pacf_{2,1}\acf_1-\pacf_{2,2}\acf_2}.
\]
Numerator:
$-\tfrac18-\tfrac12\!\cdot\!(-\tfrac1{20})-(-\tfrac14)\!\cdot\!\tfrac25
=-\tfrac18+\tfrac1{40}+\tfrac1{10}
=-\tfrac{5}{40}+\tfrac{1}{40}+\tfrac{4}{40}=0$.
Hence $\pacf_{3,3}=0$.

So $\pacf_1=\tfrac25,\ \pacf_2=-\tfrac14,\ \pacf_3=0$. The PACF cutting to $0$ at
lag $3$ suggests an \textbf{AR(2)}. Indeed
\[
Y_t=\tfrac12 Y_{t-1}-\tfrac14 Y_{t-2}+\eps_t
\]
matches: its Yule--Walker ACF is
$\acf_1=\phi_1/(1-\phi_2)=\frac{1/2}{1+1/4}=\frac25$
and $\acf_2=\phi_1\acf_1+\phi_2=\frac12\cdot\frac25-\frac14=-\frac1{20}$, exactly
the given values; and the diagonal PACF coefficients
$\pacf_{2,1}=\tfrac12=\phi_1$, $\pacf_{2,2}=-\tfrac14=\phi_2$ recover the AR
coefficients.
\end{solution}

\begin{exercise}{MA(1) and its reciprocal twin share an ACVS \quad\textnormal{[New]}}{p2-maident-ex}
Let $X_t=\eps_t-\theta\eps_{t-1}$ with $\Var\eps_t=\sigma_\eps^2$ ($\theta\ne0$),
and let $Y_t=\eta_t-\tfrac1\theta\eta_{t-1}$ with $\Var\eta_t=\sigma_\eps^2\theta^2$.
Compute both ACVSs and show they coincide. Which representation is invertible?
\end{exercise}
\begin{solution}
For $X_t$, by the MA covariance technique
(Key Technique~\ref{kt:p2-ktma}) with taps $(1,-\theta)$,
\[
\acvs_0=\sigma_\eps^2(1+\theta^2),\qquad
\acvs_{\pm1}=-\sigma_\eps^2\theta,\qquad
\acvs_\tau=0\ (\abs\tau\ge2).
\]
For $Y_t$ the MA(1) coefficient is $1/\theta$ and the noise variance is
$\sigma_\eta^2=\sigma_\eps^2\theta^2$:
\[
\begin{aligned}
\tilde\acvs_0&=\sigma_\eta^2\Big(1+\tfrac1{\theta^2}\Big)
=\sigma_\eps^2\theta^2\cdot\frac{\theta^2+1}{\theta^2}
=\sigma_\eps^2(1+\theta^2),\\
\tilde\acvs_{\pm1}&=-\sigma_\eta^2\cdot\tfrac1\theta
=-\sigma_\eps^2\theta^2\cdot\tfrac1\theta=-\sigma_\eps^2\theta,
\qquad \tilde\acvs_\tau=0\ (\abs\tau\ge2).
\end{aligned}
\]
Thus $\tilde\acvs_\tau=\acvs_\tau$ for all $\tau$: the two MA(1)s are
\textbf{indistinguishable} from second-order structure. The MA polynomial roots
are $z=1/\theta$ (for $X_t$) and $z=\theta$ (for $Y_t$); the \textbf{invertible}
representation is the one with the root outside the unit circle, i.e.\ the one
with the smaller coefficient in modulus ($\abs\theta<1$ picks $X_t$,
$\abs\theta>1$ picks $Y_t$).
\end{solution}

\begin{exercise}{ACVS and ACF of an MA(2) \quad\textnormal{[New]}}{p2-ma2}
Let $X_t=\eps_t-\theta_1\eps_{t-1}-\theta_2\eps_{t-2}$ with white-noise variance
$\sigma_\eps^2$. Find $\acvs_\tau$ for all $\tau$ and write down $\acf_1,\acf_2$.
Verify $\acvs_\tau=0$ for $\abs\tau\ge3$.
\end{exercise}
\begin{solution}
The filter taps are $(1,-\theta_1,-\theta_2)$. By Key
Technique~\ref{kt:p2-ktma},
$\acvs_\tau=\sigma_\eps^2\sum_j(\text{tap}_j)(\text{tap}_{j+\abs\tau})$:
\[
\begin{aligned}
\acvs_0&=\sigma_\eps^2(1+\theta_1^2+\theta_2^2),\\
\acvs_{\pm1}&=\sigma_\eps^2\big[(1)(-\theta_1)+(-\theta_1)(-\theta_2)\big]
=\sigma_\eps^2(-\theta_1+\theta_1\theta_2)=\sigma_\eps^2\theta_1(\theta_2-1),\\
\acvs_{\pm2}&=\sigma_\eps^2\big[(1)(-\theta_2)\big]=-\sigma_\eps^2\theta_2,\\
\acvs_\tau&=0\quad(\abs\tau\ge3),
\end{aligned}
\]
the last because for $\abs\tau\ge3$ no pair of taps overlaps (the filter has
length $3$). The ACF cuts off after lag $q=2$, as the table predicts:
\[
\acf_1=\frac{\theta_1(\theta_2-1)}{1+\theta_1^2+\theta_2^2},
\qquad
\acf_2=\frac{-\theta_2}{1+\theta_1^2+\theta_2^2}.
\]
\end{solution}

\begin{exercise}{AR(2) causality region and ACF via Yule--Walker \quad\textnormal{[New]}}{p2-ar2yw}
Consider $X_t=\phi_1 X_{t-1}+\phi_2 X_{t-2}+\eps_t$. (a) State the causality
(stationarity) conditions on $(\phi_1,\phi_2)$. (b) For $\phi_1=\tfrac12$,
$\phi_2=-\tfrac14$, compute $\acf_1,\acf_2,\acf_3$ from the Yule--Walker
equations, and predict the PACF pattern.
\end{exercise}
\begin{solution}
\textbf{(a)} With $\Phi(z)=1-\phi_1 z-\phi_2 z^2$, causality requires both roots
outside the unit circle. Equivalently (the standard ``stationarity triangle''):
\[
\phi_1+\phi_2<1,\qquad \phi_2-\phi_1<1,\qquad \abs{\phi_2}<1 .
\]

\textbf{(b)} Here $\phi_1+\phi_2=\tfrac14<1$, $\phi_2-\phi_1=-\tfrac34<1$,
$\abs{\phi_2}=\tfrac14<1$: causal. The Yule--Walker recursion
$\acf_\tau=\phi_1\acf_{\tau-1}+\phi_2\acf_{\tau-2}$ ($\tau\ge1$, with
$\acf_0=1,\ \acf_{-1}=\acf_1$) gives
\[
\acf_1=\frac{\phi_1}{1-\phi_2}=\frac{1/2}{1+1/4}=\frac25,
\qquad
\acf_2=\phi_1\acf_1+\phi_2=\tfrac12\cdot\tfrac25-\tfrac14=-\tfrac1{20},
\]
\[
\acf_3=\phi_1\acf_2+\phi_2\acf_1
=\tfrac12\cdot(-\tfrac1{20})-\tfrac14\cdot\tfrac25
=-\tfrac1{40}-\tfrac1{10}=-\tfrac18 .
\]
These are exactly the values of Exercise~\ref{exo:p2-dl}: the PACF cuts off after
lag $2$ ($\pacf_3=0$), the AR(2) signature, while the ACF tails off.
\end{solution}

\begin{exercise}{One-step recursive forecast for an AR(1) \quad\textnormal{[New]}}{p2-forecast}
For a causal AR(1) $X_t=\phi X_{t-1}+\eps_t$ ($\abs\phi<1$), find the best linear
one-step-ahead predictor $P_{X_1,\dots,X_n}(X_{n+1})$ and its mean-squared error.
\end{exercise}
\begin{solution}
The prediction equations require the error $X_{n+1}-\hat X_{n+1}$ to be
uncorrelated with $X_1,\dots,X_n$. Try $\hat X_{n+1}=\phi X_n$: the error is
$X_{n+1}-\phi X_n=\eps_{n+1}$, and by causality $X_k$ depends only on
$\eps_s,\ s\le k\le n<n+1$, so $\Cov(\eps_{n+1},X_k)=0$ for all $k\le n$. The
orthogonality conditions hold, hence
\[
P_{X_1,\dots,X_n}(X_{n+1})=\phi X_n,
\qquad
\MSE=\E[\eps_{n+1}^2]=\sigma_\eps^2 .
\]
Only the most recent observation matters --- a direct consequence of the AR(1)
having PACF $0$ beyond lag $1$ (the older points carry no \emph{new} linear
information).
\end{solution}

% ======================================================================
%  PART 3 -- ARMA Revisited: Backshift Operator, Wold Decomposition,
%            Causality & Invertibility
% ======================================================================
\section{ARMA Revisited: Backshift Operator, Wold Decomposition, Causality \& Invertibility}
\label{sec:armarev}

Having met AR, MA and ARMA models as explicit sums of white noise, we now
repackage them with the \textbf{backshift operator} $\Bsh$, which turns the
difference equations into a clean algebra of \emph{polynomials} $\Phi(z)$ and
$\Theta(z)$. The payoff is conceptual: questions about a process
(``is it stationary? is it invertible? what is its $\mathrm{MA}(\infty)$
representation?'') become questions about the \emph{roots} of these
polynomials and the analyticity of a transfer function $G(z)$. Along the way
the \textbf{Wold decomposition theorem} tells us that essentially every
stationary process is a one-sided moving average of white noise, which is
exactly why the linear/ARMA world is so central. Throughout,
$\{\eps_t\}$ is mean-zero white noise with $\Var(\eps_t)=\sigma^2$, and we use
the lecture-note conventions: ARMA written $\Phi(\Bsh)X_t=\Theta(\Bsh)\eps_t$,
stationarity from roots of $\Phi$ \emph{outside} the unit circle, invertibility
from roots of $\Theta$ \emph{outside} the unit circle.

% ----------------------------------------------------------------------
\subsection{Definitions}

\begin{definition}{Backshift operator $\Bsh$}{p3-backshift}
The \textbf{backshift operator} maps a time series $\{X_t\}$ to the series
$\{Y_t\}=\Bsh[\{X_t\}]$ shifted back one step in time,
\[
Y_t = X_{t-1}\qquad\text{for all }t\in\Ints .
\]
We write informally $\Bsh X_t = X_{t-1}$. Iterating, $\Bsh^k X_t = X_{t-k}$
for $k\ge0$, with $\Bsh^0=I$ the identity. The operator is linear, and
because shifts commute, products of polynomials in $\Bsh$ may be multiplied
and reordered like ordinary polynomials.
\end{definition}

\begin{intuition}
$\Bsh$ is just ``$X_{t}\mapsto X_{t-1}$'', but treating it as an algebraic
symbol $z$ is what lets us \emph{factor}, \emph{invert} and \emph{expand}
model equations. Every result in this part is a statement about the polynomial
$\Phi(z)$ or $\Theta(z)$ obtained by replacing $\Bsh$ with the complex variable
$z$.
\end{intuition}

\begin{definition}{ARMA polynomial notation}{p3-armapoly}
An $\mathrm{ARMA}(p,q)$ process satisfies
\[
\sum_{j=0}^{p}\phi_j X_{t-j} \;=\; \sum_{k=0}^{q}\theta_k\eps_{t-k},
\qquad \phi_0=-1,\ \theta_0=-1,
\]
which, with the backshift operator, becomes the compact statement
\[
\Phi(\Bsh)X_t=\Theta(\Bsh)\eps_t,\qquad
\Phi(z)=\sum_{j=0}^{p}\phi_j z^{j},\quad
\Theta(z)=\sum_{k=0}^{q}\theta_k z^{k}.
\]
In the common ``signed'' form one writes
$X_t=\sum_{j=1}^{p}\phi_j X_{t-j}-\sum_{k=0}^{q}\theta_k\eps_{t-k}$;
simple examples instead use $X_t=\eps_t-\theta\eps_{t-1}$ etc. Pure cases:
$\mathrm{AR}(p)$ is $\Phi(\Bsh)X_t=\eps_t$ and $\mathrm{MA}(q)$ is
$X_t=\Theta(\Bsh)\eps_t$.
\end{definition}

\begin{definition}{Characteristic polynomials}{p3-charpoly}
For $\Phi(\Bsh)X_t=\Theta(\Bsh)\eps_t$ we call
\begin{itemize}
  \item $\Phi(z)$ the \textbf{characteristic polynomial of the autoregressive
        part}, and
  \item $\Theta(z)$ the \textbf{characteristic polynomial of the moving-average
        part}.
\end{itemize}
For a pure $\mathrm{AR}$, $\Phi(\Bsh)X_t=\eps_t$, ``the characteristic
polynomial'' means $\Phi(z)$; for a pure $\mathrm{MA}$,
$X_t=\Theta(\Bsh)\eps_t$, it means $\Theta(z)$. The \emph{roots} of these
polynomials govern stationarity and invertibility.
\end{definition}

\begin{definition}{Two-sided linear process}{p3-twosided}
A \textbf{(two-sided) linear process} is
\[
X_t=\sum_{k=-\infty}^{\infty}g_k\,\eps_{t-k},
\qquad \norm{g}_2^2=\sum_k g_k^2<\infty,
\]
with $\{g_k\}\subset\Reals$ and $\{\eps_t\}$ mean-zero white noise. Its first
two moments are, for all $t,\tau\in\Ints$,
\[
\E[X_t]=0,\qquad
\Var(X_t)=\sigma^2\norm{g}_2^2<\infty,\qquad
\Cov(X_t,X_{t+\tau})=\sigma^2\sum_{k=-\infty}^{\infty}g_k\,g_{k+\tau}.
\]
The summability $\norm{g}_2^2<\infty$ is what makes these moments finite, so the
process is (second-order) stationary.
\end{definition}

\begin{definition}{General linear process (one-sided)}{p3-genlin}
If in the two-sided process we set $g_{-1}=g_{-2}=\cdots=0$, we obtain a
\textbf{general linear process}
\[
X_t=\sum_{k=0}^{\infty}g_k\,\eps_{t-k},\qquad \norm{g}_2^2<\infty .
\]
Now $X_t$ depends only on the \emph{past and present} of $\{\eps_t\}$, making
the process \textbf{causal} (in the linear sense). The same mean/autocovariance
formulas apply, so it is stationary. Writing $G(z)=\sum_{k\ge0}g_kz^k$ gives the
operator form $X_t=G(\Bsh)\eps_t$. By convention $g_0=1$.
\end{definition}

\begin{definition}{Transfer function $G(z)$, zeros and poles}{p3-transfer}
For a linear process $X_t=G(\Bsh)\eps_t$, the function
\[
G(z)=\sum_{k=0}^{\infty}g_k z^{k}
\]
is the \textbf{transfer function}. When $G$ is a rational ratio
$G(z)=G_1(z)/G_2(z)$:
\begin{itemize}
  \item the roots of the numerator $G_1(z)$ are the \textbf{zeros} of $G$;
  \item the roots of the denominator $G_2(z)$ are the \textbf{poles} of $G$.
\end{itemize}
For an $\mathrm{ARMA}$ model $G(z)=\Theta(z)/\Phi(z)$: the zeros of $\Phi$
(roots of the AR polynomial) are the poles of $G$, and the zeros of $\Theta$
(roots of the MA polynomial) are the zeros of $G$. Whether the expansion of
$G$ involves only past/present (versus future) noise is decided by where these
roots sit relative to the unit circle $|z|=1$.
\end{definition}

\begin{definition}{Innovations outlier}{p3-innov}
If a single realised value of $\eps_t$ affects $X_t$ \emph{and all subsequent}
$X_{t+1},X_{t+2},\dots$, that value is called an \textbf{innovations outlier}.
This is exactly the behaviour of a one-sided ($\mathrm{MA}(\infty)$) process:
each innovation $\eps_t$ propagates forward through the weights $g_0,g_1,\dots$.
\end{definition}

% ----------------------------------------------------------------------
\subsection{Theorems \& Propositions}

\begin{theorem}{Wold decomposition theorem}{p3-wold}
Any (zero-mean) second-order stationary process $\{X_t\}$ can be written as
\[
X_t = U_t + V_t,
\qquad U_t\ \text{and}\ V_t\ \text{uncorrelated},
\]
where
\begin{itemize}
  \item $U_t$ has a \textbf{one-sided linear} ($\mathrm{MA}(\infty)$)
        representation
        \[
        U_t=\sum_{k=0}^{\infty}g_k\,\eps_{t-k},\qquad g_0=1,\ \norm{g}_2^2<\infty,
        \]
        with $\{\eps_t\}$ mean-zero white noise uncorrelated with $V$,
        i.e.\ $\E[\eps_s V_t]=0$ for all $s,t$; the sequences $\{g_k\}$ and
        $\{\eps_t\}$ are \emph{uniquely} determined;
  \item $V_t$ is \textbf{singular} (deterministic): it can be predicted from its
        own past with zero error.
\end{itemize}
\textit{Proof idea:} project $X_t$ onto the closed linear span of its own past;
the residual one-step innovations are the white noise $\eps_t$, the $g_k$ are the
projection coefficients, and $V_t$ is the part lying in the infinite-past
intersection of these spans (the purely deterministic component).\\
\textit{Why:} this is the structural justification for the whole linear-model
programme: up to a deterministic piece, \emph{every} stationary process is a
one-sided moving average of white noise, and ARMA models are the finite,
rational approximations to this $\mathrm{MA}(\infty)$.
\end{theorem}

\begin{intuition}
Wold says: the only ``new information'' arriving at time $t$ is the innovation
$\eps_t$; everything else is either an echo of past innovations ($U_t$) or
perfectly forecastable from the past ($V_t$, e.g.\ a fixed sinusoid). For the
processes we study $V_t=0$, leaving the pure $\mathrm{MA}(\infty)$ form
$X_t=G(\Bsh)\eps_t$.
\end{intuition}

\begin{proposition}{Laurent expansion of $1/G_2$ and the past/present criterion}{p3-laurent}
Let $G(z)=G_1(z)/G_2(z)$ and let $z_1,\dots,z_p$ be the zeros of $G_2$, ordered
so that $z_1,\dots,z_k$ lie \emph{inside} the unit circle and
$z_{k+1},\dots,z_p$ lie \emph{outside}. A partial-fraction plus geometric-series
(Laurent) expansion, convergent on $|z|=1$, gives after replacing $z\mapsto\Bsh$
\[
\frac{1}{G_2(\Bsh)}\eps_t
=\underbrace{\sum_{j=1}^{k}A_j\sum_{l=0}^{\infty}z_j^{\,l}\,\eps_{t+1+l}}_{\text{future}}
\;-\;\underbrace{\sum_{j=k+1}^{p}A_j\sum_{l=0}^{\infty}z_j^{-(l+1)}\,\eps_{t-l}}_{\text{past}+\text{present}} .
\]
Hence if \emph{all} roots of $G_2(z)$ are outside the unit circle ($k=0$), only
past and present values appear, and the one-sided general linear process exists.\\
\textit{Proof idea:} expand $A_j/(z-z_j)$ as a geometric series, choosing the
\emph{inside} expansion $\sum_l (z_j/z)^l$ for $|z_j|<1$ and the \emph{outside}
expansion for $|z_j|>1$, so that each series converges on $|z|=1$.\\
\textit{Why:} this is the precise mechanism behind ``roots outside the unit
circle $\Rightarrow$ causal''. It also gives the equivalent analytic statement:
$|G(z)|<\infty$ for $|z|\le1$, i.e.\ $G$ is \textbf{analytic inside and on the
unit circle}.
\end{proposition}

\begin{proposition}{Stationarity of the general linear process}{p3-glpstat}
A general linear process $X_t=\sum_{k\ge0}g_k\eps_{t-k}$ with
$\norm{g}_2^2<\infty$ is second-order stationary, with
\[
\E[X_t]=0,\qquad
\acvs_\tau=\sigma^2\sum_{k=0}^{\infty}g_k\,g_{k+\abs\tau}.
\]
\textit{Proof idea:} the mean is zero by linearity; bilinearity of covariance
together with $\Cov(\eps_s,\eps_u)=\sigma^2\ind_{\{s=u\}}$ collapses the double
sum to the single sum above, which is finite by Cauchy--Schwarz and
$\norm{g}_2^2<\infty$.\\
\textit{Why:} this is the engine for computing the ACVS of \emph{any} causal
ARMA once its $\mathrm{MA}(\infty)$ weights $\{g_k\}=\{\psi_k\}$ are known.
\end{proposition}

\begin{proposition}{Stationarity (causality) of $\mathrm{AR}(p)$}{p3-arstat}
For $\mathrm{AR}(p)$, $\Phi(\Bsh)X_t=\eps_t$, the (unique) stationary causal
solution is $X_t=\Phi^{-1}(\Bsh)\eps_t=G(\Bsh)\eps_t$ with
$G(z)=\Phi^{-1}(z)$. This power expansion is analytic on $|z|\le1$
\emph{iff} all roots of $\Phi(z)$ lie \emph{outside} the unit circle. Therefore
\[
\mathrm{AR}(p)\ \text{is stationary}\iff \text{all roots of }\Phi(z)\ \text{satisfy }\abs{z}>1 .
\]
\textit{Proof idea:} the poles of $G=1/\Phi$ are the roots of $\Phi$; by the
Laurent criterion (Prop.~\ref{prop:p3-laurent}) the one-sided expansion exists
exactly when those poles are outside $|z|=1$.\\
\textit{Why:} this is \emph{the} root test for AR stationarity; it reduces an
infinite-order question to checking the moduli of $p$ roots.
\end{proposition}

\begin{proposition}{$\mathrm{AR}(p)$ is always invertible; $\mathrm{MA}(q)$ is always stationary}{p3-alwaysinv}
\begin{itemize}
  \item For $\mathrm{AR}(p)$, $\Phi(\Bsh)X_t=\eps_t$ already expresses the white
        noise as a \emph{finite} polynomial in $X$, so $\mathrm{AR}(p)$ is
        \textbf{always invertible} (the AR polynomial $\Phi$ has finite order;
        there is no MA polynomial, $\Theta\equiv1$).
  \item For $\mathrm{MA}(q)$, $X_t=\Theta(\Bsh)\eps_t=G(\Bsh)\eps_t$ with
        $G(z)=\Theta(z)$ a \emph{finite} polynomial, so $|G(z)|<\infty$
        everywhere; hence $\mathrm{MA}(q)$ is \textbf{always stationary}.
\end{itemize}
\textit{Proof idea:} ``finite polynomial $\Rightarrow$ bounded on the unit
disk''; only the \emph{inverse} expansion can fail to converge, never the
finite direct one.\\
\textit{Why:} it tells you which check is \emph{free}: never test an AR for
invertibility or an MA for stationarity.
\end{proposition}

\begin{proposition}{Invertibility via $G^{-1}$ and the $\mathrm{MA}(q)$ root test}{p3-invert}
A linear process $X_t=G(\Bsh)\eps_t$ is \textbf{invertible} if it can be
rewritten as an $\mathrm{AR}(\infty)$,
\[
G^{-1}(\Bsh)X_t=\eps_t,\qquad G^{-1}(z)=\sum_{k=0}^{\infty}h_k z^{k},
\]
with $G^{-1}$ analytic on $|z|\le1$, equivalently $|G^{-1}(z)|<\infty$ for
$|z|\le1$, equivalently \emph{all poles of $G^{-1}$ are outside} the unit circle
(equivalently, all \emph{zeros} of $G$ are inside). For $\mathrm{MA}(q)$,
$G=\Theta$, so
\[
\mathrm{MA}(q)\ \text{is invertible}\iff \text{all roots of }\Theta(z)\ \text{satisfy }\abs{z}>1 .
\]
\textit{Proof idea:} inverting $X_t=\Theta(\Bsh)\eps_t$ needs the power series of
$1/\Theta$ to converge on $|z|\le1$; by the Laurent criterion this requires the
roots of $\Theta$ (the poles of $1/\Theta$) outside the unit circle.\\
\textit{Why:} invertibility is what makes the innovations $\eps_t$ recoverable
from the observed past $X_t,X_{t-1},\dots$, which is essential for forecasting
and for parameter identifiability.
\end{proposition}

\begin{intuition}
Stationarity (causality) and invertibility are the \emph{same} test applied to
\emph{two different} polynomials: ``all roots outside the unit circle'' for
$\Phi$ controls whether $X_t=G(\Bsh)\eps_t$ uses only past/present noise
(causal), while the identical test for $\Theta$ controls whether
$\eps_t=G^{-1}(\Bsh)X_t$ uses only past/present data (invertible). An AR has no
$\Theta$ to fail, so it is always invertible; an MA has no $\Phi$ to fail, so it
is always stationary.
\end{intuition}

\begin{corollary}{Stationarity \& invertibility summary for ARMA}{p3-summary}
For $\Phi(\Bsh)X_t=\Theta(\Bsh)\eps_t$:
\[
\begin{aligned}
\textbf{AR}(p):&\quad \text{stationary}\iff \text{roots of }\Phi\ \text{outside }\abs{z}=1;\quad \text{\emph{always} invertible.}\\
\textbf{MA}(q):&\quad \text{\emph{always} stationary};\quad \text{invertible}\iff \text{roots of }\Theta\ \text{outside }\abs{z}=1.\\
\textbf{ARMA}(p,q):&\quad \text{stationary}\iff \text{roots of }\Phi\ \text{outside};\ \ \text{invertible}\iff \text{roots of }\Theta\ \text{outside.}
\end{aligned}
\]
\textit{Why:} the single most-used reference fact of the chapter; combine the
two independent root tests.
\end{corollary}

\begin{pitfall}
``Roots outside the unit circle'' refers to the roots of the polynomial
$\Phi(z)=\sum_j\phi_j z^j$ in the variable $z$, i.e.\ $\abs{z}>1$. Equivalently,
in terms of the \emph{reciprocal} roots $g=1/z$, one needs $\abs{g}<1$ ---
exactly the AR(1) condition $\abs{\phi}<1$ (since the root of $1-\phi z$ is
$z=1/\phi$). Do not confuse the two conventions, and do not check
\emph{inside} the circle.
\end{pitfall}

\begin{proposition}{Causal $\mathrm{AR}(p)$ has $\pacf_\tau=0$ for $\tau>p$}{p3-pacfcut}
If $\{X_t\}$ is a causal $\mathrm{AR}(p)$, its partial autocorrelation function
satisfies $\pacf_\tau=0$ for all $\tau>p$.\\
\textit{Proof idea:} the order-$\tau$ best linear predictor is
$P_{X_0,\dots,X_{\tau-1}}(X_\tau)=\sum_{j=1}^{\tau}\pacf_{\tau,j}X_{\tau-j}$ and
$\pacf_\tau=\pacf_{\tau,\tau}$. Since
$X_\tau=\sum_{j=1}^p\phi_j X_{\tau-j}+\eps_\tau$, choosing
$\pacf_{\tau,j}=\tilde\phi_j$ (with $\tilde\phi_j=\phi_j$ for $j\le p$,
$\tilde\phi_j=0$ for $j>p$) makes the prediction residual equal $\eps_\tau$,
which by causality satisfies $\Cov(\eps_\tau,X_k)=0$ for all $k<\tau$ (here
$\Cov(X_t,\eps_s)=0$ when $t<s$). Thus the orthogonality (prediction) equations
are met, so these are the true coefficients; in particular
$\pacf_\tau=\pacf_{\tau,\tau}=\tilde\phi_\tau=0$ for $\tau>p$.\\
\textit{Why:} the PACF ``cuts off'' after lag $p$ for an $\mathrm{AR}(p)$ ---
the diagnostic mirror image of the ACF cutting off after lag $q$ for an
$\mathrm{MA}(q)$ (this proof is Sheet~2, Ex.~2.7).
\end{proposition}

% ----------------------------------------------------------------------
\subsection{Key Techniques}

\begin{keytech}{Root test for stationarity / invertibility}{p3-roottest}
Given $\Phi(\Bsh)X_t=\Theta(\Bsh)\eps_t$:
\begin{enumerate}
  \item Write the polynomials in $z$: $\Phi(z)=1-\phi_1 z-\dots-\phi_p z^p$ and
        $\Theta(z)=1-\theta_1 z-\dots-\theta_q z^q$ (matching your sign
        convention).
  \item \textbf{Stationarity:} solve $\Phi(z)=0$ and check \emph{every} root has
        $\abs{z}>1$. (AR/ARMA only; an MA is automatically stationary.)
  \item \textbf{Invertibility:} solve $\Theta(z)=0$ and check \emph{every} root
        has $\abs{z}>1$. (MA/ARMA only; an AR is automatically invertible.)
\end{enumerate}
For a quadratic $a z^2+bz+c$ use $z=\frac{-b\pm\sqrt{b^2-4ac}}{2a}$ and report
$\abs{z}=\sqrt{(\mathrm{Re}\,z)^2+(\mathrm{Im}\,z)^2}$. A handy AR(1)/AR(2)
shortcut: $1-\phi z$ has root $1/\phi$; $1-\phi_1 z-\phi_2 z^2$ is stationary iff
$\abs{\phi_2}<1$, $\phi_2+\phi_1<1$ and $\phi_2-\phi_1<1$ (the AR(2) triangle).
\end{keytech}

\begin{keytech}{$\psi$-weights by recurrence (matching coefficients)}{p3-recur}
To find the $\mathrm{MA}(\infty)$ weights in $X_t=\sum_{j\ge0}\psi_j\eps_{t-j}$
for a causal $\Phi(\Bsh)X_t=\Theta(\Bsh)\eps_t$, equate
$\Phi(z)\big(\sum_{j\ge0}\psi_j z^j\big)=\Theta(z)$ and match powers of $z$.
With $\Phi(z)=1-\phi_1 z-\dots-\phi_p z^p$ and
$\Theta(z)=1-\theta_1 z-\dots-\theta_q z^q$ this gives the recurrence
\[
\psi_j=\theta_j^{\star}+\sum_{i=1}^{p}\phi_i\,\psi_{j-i},
\qquad \psi_0=1,\ \psi_j=0\ (j<0),
\]
where $\theta_j^{\star}=-\theta_j$ for $1\le j\le q$ and $0$ otherwise (with the
$X_t=\eps_t-\theta\eps_{t-1}$ sign convention). Solve forward; for a constant-
coefficient linear recurrence the closed form is a sum
$\sum_i (\text{polynomial in }j)\,r_i^{-j}$ over the AR roots $r_i$ (with the
familiar $(c_0+c_1 j)r^{-j}$ shape for a repeated root).
\end{keytech}

\begin{keytech}{$\psi$-weights by partial fractions}{p3-partialfrac}
Alternatively, expand $G(z)=\Theta(z)/\Phi(z)$ directly:
\begin{enumerate}
  \item Factor $\Phi(z)=\phi_p\prod_{i}(z-r_i)$ with roots $r_i$ ($\abs{r_i}>1$).
  \item Partial-fraction $\dfrac{\Theta(z)}{\Phi(z)}
        =(\text{poly})+\sum_i\dfrac{A_i}{z-r_i}$, then write each
        $\dfrac{A_i}{z-r_i}=-\dfrac{A_i}{r_i}\dfrac{1}{1-z/r_i}
        =-\dfrac{A_i}{r_i}\sum_{j\ge0}(z/r_i)^{\,j}$ (valid for
        $\abs{z}<\abs{r_i}$, in particular on/inside the unit circle).
  \item Collect the coefficient of $z^j$ to read off $\psi_j$.
\end{enumerate}
Because $\abs{r_i}>1$ the geometric series converge inside the unit disk,
giving a decaying $\psi_j$. This is the method of Sheet~3, Ex.~3.1.
\end{keytech}

\begin{keytech}{ACVS of a causal ARMA via the difference equation}{p3-acvsdiff}
To get $\acvs_\tau$ for $\Phi(\Bsh)X_t=\Theta(\Bsh)\eps_t$ without summing
$\psi$-weights:
\begin{enumerate}
  \item Multiply the model equation by $X_{t-\tau}$ and take expectations. For
        $\tau$ larger than the MA order $q$ the right side vanishes, giving a
        \emph{homogeneous} recurrence in $\acvs_\tau$ whose characteristic
        equation is $\Phi$ (so $\acvs_\tau$ is a combination of $r_i^{-\tau}$).
  \item For the small lags $\tau=0,1,\dots,q$ the right side keeps cross-terms
        $\E[\eps_{t-k}X_{t-\tau}]=\sigma^2\psi_{k-\tau}$ (only $k\ge\tau$
        contribute, since $X_{t-\tau}$ is a one-sided MA in
        $\eps_{t-\tau},\eps_{t-\tau-1},\dots$). These give the initial
        conditions.
  \item Fit the general homogeneous solution to the initial conditions.
\end{enumerate}
This is the route used in the lecture-note worked example
(Exercise~\ref{exo:p3-arma22ex}).
\end{keytech}

\begin{keytech}{Spectral density of an ARMA from its polynomials}{p3-sdf}
For a causal, invertible ARMA $\Phi(\Bsh)X_t=\Theta(\Bsh)\eps_t$, the spectral
density (with frequency $f$, sampling step $1$) is
\[
\sdf(f)=\sigma^2\,\frac{\abs{\Theta(e^{-2\pi i f})}^2}{\abs{\Phi(e^{-2\pi i f})}^2},
\qquad -\tfrac12\le f\le\tfrac12 .
\]
For the pure $\mathrm{MA}(1)$ $X_t=\eps_t-\theta\eps_{t-1}$ this gives
$\sdf(f)=\sigma^2\abs{1-\theta e^{-2\pi i f}}^2
=\sigma^2\big(1+\theta^2-2\theta\cos(2\pi f)\big)$, and one checks
$\int_{-1/2}^{1/2}\sdf(f)\,df=\acvs_0=\sigma^2(1+\theta^2)$ and
$\int_{-1/2}^{1/2}\sdf(f)e^{2\pi i f}\,df=\acvs_1=-\theta\sigma^2$.
\end{keytech}

% ----------------------------------------------------------------------
\subsection{Exercises}

\begin{exercise}{$\psi$-weights via partial fractions \quad\textnormal{[Sheet 3]}}{p3-e31}
Find the coefficients $\psi_j$, $j=0,1,2,\dots$, in the representation
$X_t=\sum_{j\ge0}\psi_j\eps_{t-j}$ of the ARMA process
\[
(1-0.5\Bsh+0.04\Bsh^2)X_t=(1+0.25\Bsh)\eps_t,
\]
where $\eps_t$ is white noise.
\end{exercise}
\begin{solution}
Write $X_t=G(\Bsh)\eps_t$ with
\[
G(z)=\frac{1+0.25z}{1-0.5z+0.04z^2}
=\frac{1+0.25z}{0.04\,(z-2.5)(z-10)},
\]
since $0.04z^2-0.5z+1=0$ has roots $z=2.5$ and $z=10$ (both outside the unit
circle, so the process is causal). Partial-fraction the rational part
$\frac{1}{(z-2.5)(z-10)}=\frac{1}{7.5}\big(\frac{-1}{z-2.5}+\frac{1}{z-10}\big)$,
hence
\[
G(z)=\frac{1+0.25z}{0.04\cdot7.5}\left(\frac{-1}{z-2.5}+\frac{1}{z-10}\right)
=\frac{1+0.25z}{0.3}\left(\frac{1}{2.5}\sum_{j\ge0}0.4^j z^j-\frac{1}{10}\sum_{j\ge0}0.1^j z^j\right),
\]
using $\frac{-1}{z-2.5}=\frac{1}{2.5}\frac{1}{1-z/2.5}=\frac{1}{2.5}\sum_{j\ge0}(0.4)^j z^j$
and $\frac{1}{z-10}=-\frac{1}{10}\sum_{j\ge0}(0.1)^j z^j$. Distributing the
factor $(1+0.25z)/0.3$ and collecting powers of $z$,
\[
G(z)=\sum_{j\ge0}\left(\frac{0.4^j}{0.75}-\frac{0.1^j}{3}\right)z^j
+\sum_{l\ge1}\left(\frac{0.4^l}{1.2}-\frac{0.1^l}{1.2}\right)z^l,
\]
so for $j\ge1$,
\[
\psi_j=\frac{0.4^j}{0.75}-\frac{0.1^j}{3}+\frac{0.4^j}{1.2}-\frac{0.1^j}{1.2}
=\Big(\tfrac{1}{0.75}+\tfrac{1}{1.2}\Big)0.4^j-\Big(\tfrac{1}{3}+\tfrac{1}{1.2}\Big)0.1^j
=\frac{13}{6}\,0.4^j-\frac{7}{6}\,0.1^j,
\]
and $\psi_0=1$. Hence
\[
\boxed{\;\psi_j=\begin{cases}1 & j=0,\\[2pt]
\dfrac{13}{6}\,(0.4)^j-\dfrac{7}{6}\,(0.1)^j & j\ge1.\end{cases}}
\]
\textit{Check (coefficient-matching):} the recurrence
$\psi_0=1,\ \psi_1-0.5\psi_0=0.25\Rightarrow\psi_1=0.75,\
\psi_j=0.5\psi_{j-1}-0.04\psi_{j-2}\ (j\ge2)$ gives
$\psi_0,\psi_1,\psi_2,\psi_3=1,\,0.75,\,0.335,\,0.1375$, matching
$\tfrac{13}{6}0.4^j-\tfrac{7}{6}0.1^j$ exactly.
\end{solution}

\begin{exercise}{$\mathrm{MA}(\infty)$ representation of $\mathrm{AR}(1)$ and $\mathrm{AR}(2)$ \quad\textnormal{[Sheet 3]}}{p3-e32}
For a stationary causal $\mathrm{AR}(1)$ $(1-\phi\Bsh)Y_t=\eps_t$ and a
stationary $\mathrm{AR}(2)$ $(1-\phi_1\Bsh-\phi_2\Bsh^2)Y_t=\eps_t$, find the
representation of $Y_t$ in terms of $\eps_t$.
\end{exercise}
\begin{solution}
\textbf{AR(1).} By causality ($\abs\phi<1$),
\[
Y_t=\frac{1}{1-\phi\Bsh}\eps_t=\sum_{k=0}^{\infty}\phi^{k}\eps_{t-k}.
\]
\textbf{AR(2).} Factor $\Phi(z)=1-\phi_1 z-\phi_2 z^2=(1-g_1 z)(1-g_2 z)$, where
$g_1^{-1},g_2^{-1}$ are the roots (outside the unit circle, so
$\abs{g_1},\abs{g_2}<1$). Then
\[
Y_t=\frac{1}{(1-g_1\Bsh)(1-g_2\Bsh)}\eps_t
=\sum_{k_1=0}^{\infty}\sum_{k_2=0}^{\infty}g_1^{k_1}g_2^{k_2}\,\eps_{t-(k_1+k_2)}
=\sum_{n=0}^{\infty}\Big(\sum_{k=0}^{n}g_1^{k}g_2^{\,n-k}\Big)\eps_{t-n},
\]
collecting terms with $n=k_1+k_2$ (Cauchy product, justified by
$\abs{g_1},\abs{g_2}<1$). For $g_1\ne g_2$ the inner sum is a geometric series,
\[
\sum_{k=0}^{n}g_1^{k}g_2^{\,n-k}=\frac{g_1^{\,n+1}-g_2^{\,n+1}}{g_1-g_2},
\quad\text{so}\quad
Y_t=\frac{g_1}{g_1-g_2}\sum_{n\ge0}g_1^{n}\eps_{t-n}-\frac{g_2}{g_1-g_2}\sum_{n\ge0}g_2^{n}\eps_{t-n}.
\]
The reciprocal roots are
$g_{1,2}^{-1}=\dfrac{-\phi_1\mp\sqrt{\phi_1^2+4\phi_2}}{2\phi_2}$. When
$g_1=g_2=g$ the inner sum becomes $\sum_{k=0}^n g^n=(n+1)g^{n}$, so
$\psi_n=(n+1)g^{n}$.
\end{solution}

\begin{exercise}{$\mathrm{AR}(2)$ stationarity via roots \quad\textnormal{[Sheet 3]}}{p3-e33}
For the $\mathrm{AR}(2)$ process
$\frac{15}{16}Y_t=\frac14 Y_{t-1}-\frac{1}{16}Y_{t-2}+\eps_t$ (with $\eps_t$
Gaussian), write the characteristic AR polynomial. Is it stationary? Invertible?
\end{exercise}
\begin{solution}
Move all $Y$ terms to the left and use the backshift operator:
\[
\Big(\tfrac{15}{16}-\tfrac14\Bsh+\tfrac{1}{16}\Bsh^2\Big)Y_t=\eps_t,
\qquad
\Phi(z)=\tfrac{15}{16}-\tfrac14 z+\tfrac{1}{16}z^2 .
\]
Solve $\frac{1}{16}z^2-\frac14 z+\frac{15}{16}=0$, i.e.\ $z^2-4z+15=0$:
\[
z=\frac{4\pm\sqrt{16-60}}{2}=2\pm\sqrt{-11}=2\pm\sqrt{11}\,i
\approx 2\pm 3.3166\,i,
\qquad \abs{z}=\sqrt{4+11}=\sqrt{15}\approx3.873 .
\]
Both roots have modulus $\sqrt{15}>1$, i.e.\ they lie outside the unit circle,
so the process is \textbf{stationary}. An $\mathrm{AR}$ process is
\textbf{always invertible}.
\end{solution}

\begin{exercise}{$\mathrm{ARMA}(2,1)$ stationary \& invertible \quad\textnormal{[Sheet 3]}}{p3-e34}
For the $\mathrm{ARMA}(2,1)$ process
$Y_t-\frac14 Y_{t-1}+\frac{1}{16}Y_{t-2}=\eps_t-0.5\eps_{t-1}$, is it stationary?
Invertible?
\end{exercise}
\begin{solution}
In backshift form,
$\big(1-\tfrac14\Bsh+\tfrac{1}{16}\Bsh^2\big)Y_t=(1-0.5\Bsh)\eps_t$, so
\[
\Phi(z)=1-\tfrac14 z+\tfrac{1}{16}z^2,\qquad \Theta(z)=1-\tfrac12 z .
\]
\textbf{Stationarity.} $\Phi(z)=0\iff z^2-4z+16=0$ (multiply by $16$), giving
\[
z=\frac{4\pm\sqrt{16-64}}{2}=2\pm 2\sqrt3\,i,\qquad \abs{z}=\sqrt{4+12}=4>1,
\]
both outside the unit circle $\Rightarrow$ \textbf{stationary}.\\
\textbf{Invertibility.} $\Theta(z)=0\iff z=2$, and $\abs{2}=2>1$ is outside the
unit circle $\Rightarrow$ \textbf{invertible}. The process is both stationary
and invertible.
\end{solution}

\begin{exercise}{Same ACF; inverting a non-invertible MA \quad\textnormal{[Sheet 3]}}{p3-e35}
Let $\{U_t\}$ be stationary zero-mean. Define $X_t=(1-0.4\Bsh)U_t$ and
$W_t=(1-2.5\Bsh)U_t$.
\textbf{(a)} Express the ACFs of $\{X_t\}$ and $\{W_t\}$ via that of $\{U_t\}$.
\textbf{(b)} Show $\{X_t\}$ and $\{W_t\}$ have the \emph{same} ACF.
\textbf{(c)} Show $\eps_t=-\sum_{j\ge1}(0.4)^j X_{t+j}$ satisfies
$\eps_t-2.5\eps_{t-1}=X_t$.
\end{exercise}
\begin{solution}
\textbf{(a)} With $X_t=U_t-0.4U_{t-1}$, bilinearity gives
\[
\acvs_X(\tau)=\Cov(U_t-0.4U_{t-1},\,U_{t+\tau}-0.4U_{t+\tau-1})
=1.16\,\acvs_U(\tau)-0.4\,\acvs_U(\tau-1)-0.4\,\acvs_U(\tau+1),
\]
using $1+0.4^2=1.16$. In particular, with symmetry $\acvs_U(-1)=\acvs_U(1)$,
\[
\acvs_X(0)=1.16\,\acvs_U(0)-0.8\,\acvs_U(1)=\acvs_U(0)\{1.16-0.8\acf_U(1)\}.
\]
Dividing, and writing $\acvs_U(\tau)=\acf_U(\tau)\acvs_U(0)$,
\[
\acf_X(\tau)=\frac{1.16\,\acf_U(\tau)-0.4\,\acf_U(\tau-1)-0.4\,\acf_U(\tau+1)}{1.16-0.8\acf_U(1)} .
\]
Similarly, with $W_t=U_t-2.5U_{t-1}$ ($1+2.5^2=7.25$),
\[
\acf_W(\tau)=\frac{7.25\,\acf_U(\tau)-2.5\,\acf_U(\tau-1)-2.5\,\acf_U(\tau+1)}{7.25-5\,\acf_U(1)} .
\]
\textbf{(b)} Note $7.25=6.25\times1.16$, $2.5=6.25\times0.4$ and
$5=6.25\times0.8$. Factoring $6.25$ from numerator and denominator of
$\acf_W$,
\[
\acf_W(\tau)=\frac{6.25\{1.16\,\acf_U(\tau)-0.4\,\acf_U(\tau-1)-0.4\,\acf_U(\tau+1)\}}{6.25\{1.16-0.8\,\acf_U(1)\}}
=\acf_X(\tau).
\]
So the two have identical ACFs. (This is the MA non-uniqueness: the root
$z=2.5$ for $X$ and $z=0.4=1/2.5$ for $W$ give the same correlation structure;
only $X$, with root outside the circle, is invertible.)\\
\textbf{(c)} With $\eps_t=-\sum_{j\ge1}(0.4)^j X_{t+j}$,
\[
\begin{aligned}
\eps_t-2.5\,\eps_{t-1}
&=-\sum_{j\ge1}0.4^{j}X_{t+j}+2.5\sum_{j\ge1}0.4^{j}X_{t+j-1}\\
&=-\sum_{j\ge1}0.4^{j}X_{t+j}+2.5\sum_{l\ge0}0.4^{\,l+1}X_{t+l}
\qquad(l=j-1)\\
&=-\sum_{j\ge1}0.4^{j}X_{t+j}+\sum_{l\ge0}0.4^{\,l}X_{t+l}
\qquad(2.5\cdot0.4=1)\\
&=X_t+\sum_{l\ge1}0.4^{l}X_{t+l}-\sum_{j\ge1}0.4^{j}X_{t+j}=X_t .
\end{aligned}
\]
Thus the non-invertible $W$-type representation can still be inverted, but using
\emph{future} values $X_{t+j}$ --- exactly the ``future-noise'' branch of the
Laurent expansion (Prop.~\ref{prop:p3-laurent}).
\end{solution}

\begin{exercise}{Worked lecture example: ACVS of an $\mathrm{ARMA}(2,1)$ \quad\textnormal{[Notes]}}{p3-arma22ex}
For $\{I-\Bsh+\tfrac14\Bsh^2\}X_t=\{I+\Bsh\}\eps_t$ with $\eps_t$ white noise of
variance $\sigma^2$, determine the autocovariance sequence $\acvs_\tau$.
\end{exercise}
\begin{solution}
The model is $X_t-X_{t-1}+\tfrac14 X_{t-2}=\eps_t+\eps_{t-1}$, with
$\Phi(z)=1-z+\tfrac14 z^2=\tfrac14(z-2)^2$ (double root $z=2$ outside the unit
circle $\Rightarrow$ stationary) and $\Theta(z)=1+z$ (root $z=-1$ \emph{on} the
circle $\Rightarrow$ not invertible, but stationarity is all we need here).

\textbf{Recurrence for $\tau\ge2$.} Multiply by $X_{t-\tau}$ and take
expectations. For $\tau\ge2$ the right side is
$\E[(\eps_t+\eps_{t-1})X_{t-\tau}]=0$ (since $X_{t-\tau}$ involves only
$\eps_{t-\tau},\eps_{t-\tau-1},\dots$, all uncorrelated with $\eps_t,\eps_{t-1}$),
so
\[
\acvs_\tau-\acvs_{\tau-1}+\tfrac14\acvs_{\tau-2}=0,\qquad \tau\ge2 .
\]
The characteristic equation is $\Phi$ with a double root, so the general
solution is
\[
\acvs_\tau=(\beta_0+\beta_1\tau)\,2^{-\tau},\qquad \tau\ge0 .
\]
\textbf{Initial conditions ($\tau=0,1$).} First find the leading $\psi$-weights
from $\Phi(\Bsh)\psi(\Bsh)=\Theta(\Bsh)$: $\psi_0=1$ and
$\psi_1-\psi_0=1\Rightarrow\psi_1=2$. Then the cross terms are
$\E[\eps_t X_t]=\psi_0\sigma^2=\sigma^2$ and
$\E[\eps_{t-1}X_t]=\psi_1\sigma^2=2\sigma^2$. Hence
\[
\tau=0:\quad \acvs_0-\acvs_1+\tfrac14\acvs_2=\E[(\eps_t+\eps_{t-1})X_t]=\sigma^2+2\sigma^2=3\sigma^2,
\]
\[
\tau=1:\quad \acvs_1-\acvs_0+\tfrac14\acvs_1=\E[(\eps_t+\eps_{t-1})X_{t-1}]=\sigma^2 .
\]
(The second uses $\E[\eps_t X_{t-1}]=0$ and
$\E[\eps_{t-1}X_{t-1}]=\psi_0\sigma^2=\sigma^2$.) Substituting
$\acvs_\tau=(\beta_0+\beta_1\tau)2^{-\tau}$ gives the linear system
$\beta_0=\tfrac{32}{3}\sigma^2,\ \beta_1=8\sigma^2$. Therefore
\[
\boxed{\;\acvs_\tau=\Big(\tfrac{32}{3}+8\tau\Big)2^{-\tau}\,\sigma^2,\qquad \tau\ge0\;}
\]
(and $\acvs_{-\tau}=\acvs_\tau$). In particular
$\acvs_0=\tfrac{32}{3}\sigma^2,\ \acvs_1=\tfrac{28}{3}\sigma^2,\ \acvs_2=\tfrac{20}{3}\sigma^2$.\\
\textit{Check (via $\psi$-weights):} the recurrence
$\psi_j=\psi_{j-1}-\tfrac14\psi_{j-2}$ has closed form $\psi_j=(1+3j)2^{-j}$, and
$\sigma^2\sum_{j\ge0}\psi_j\psi_{j+\tau}$ reproduces
$\tfrac{32}{3},\tfrac{28}{3},\tfrac{20}{3},\dots$ exactly.
\end{solution}

\begin{exercise}{Prove the PACF cut-off for a causal $\mathrm{AR}(p)$ \quad\textnormal{[Sheet 2, Ex.~2.7]}}{p3-pacfex}
Prove that for a causal $\mathrm{AR}(p)$, $\pacf_\tau=0$ for all $\tau>p$.
Use the prediction equations
$\Cov\big(Y-P_{X_1,\dots,X_n}(Y),X_k\big)=0$ for all $k$, and the fact that for
causal processes $\Cov(X_t,\eps_s)=0$ when $t<s$.
\end{exercise}
\begin{solution}
Fix $\tau>p$. By definition $\pacf_\tau=\pacf_{\tau,\tau}$, the last coefficient
in the order-$\tau$ best linear predictor
\[
P_{X_0,\dots,X_{\tau-1}}(X_\tau)=\sum_{j=1}^{\tau}\pacf_{\tau,j}X_{\tau-j}.
\]
We show the AR coefficients themselves solve the prediction equations. Set
$\tilde\phi_j=\phi_j$ for $1\le j\le p$ and $\tilde\phi_j=0$ for $j>p$, so that
$X_\tau=\sum_{j=1}^{\tau}\tilde\phi_j X_{\tau-j}+\eps_\tau$. Then for every
$k\in\{0,\dots,\tau-1\}$,
\[
\Cov\Big(X_\tau-\sum_{j=1}^{\tau}\tilde\phi_j X_{\tau-j},\,X_k\Big)
=\Cov(\eps_\tau,X_k)=0,
\]
where the last equality is causality ($X_k$ depends only on
$\eps_k,\eps_{k-1},\dots$, all with index $<\tau$, so uncorrelated with
$\eps_\tau$). Hence $\{\tilde\phi_j\}$ satisfies the orthogonality (prediction)
equations, and by uniqueness of the best linear predictor
$\pacf_{\tau,j}=\tilde\phi_j$. In particular
$\pacf_\tau=\pacf_{\tau,\tau}=\tilde\phi_\tau=0$ because $\tau>p$. $\qed$
\end{solution}

\begin{exercise}{Causality \& invertibility of an $\mathrm{ARMA}(1,1)$, and its $\psi$-weights \quad\textnormal{[New]}}{p3-newarma11}
Consider $(1-0.5\Bsh)X_t=(1+0.3\Bsh)\eps_t$. (a) Is it stationary? Invertible?
(b) Find the $\mathrm{MA}(\infty)$ weights $\psi_j$. (c) Compute $\acvs_0$ and
$\acvs_1$.
\end{exercise}
\begin{solution}
\textbf{(a)} $\Phi(z)=1-0.5z$ has root $z=2$ ($\abs{2}>1$) $\Rightarrow$
\textbf{stationary}. $\Theta(z)=1+0.3z$ has root $z=-1/0.3\approx-3.33$
($\abs{z}>1$) $\Rightarrow$ \textbf{invertible}.\\
\textbf{(b)} Match coefficients in $(1-0.5z)\sum_{j\ge0}\psi_j z^j=1+0.3z$:
$\psi_0=1$; $\psi_1-0.5\psi_0=0.3\Rightarrow\psi_1=0.8$;
$\psi_j-0.5\psi_{j-1}=0\Rightarrow\psi_j=0.5\,\psi_{j-1}$ for $j\ge2$. Closed
form (general ARMA(1,1) with $X_t=\phi X_{t-1}+\eps_t+\theta\eps_{t-1}$):
$\psi_j=(\phi+\theta)\phi^{\,j-1}$ for $j\ge1$ with $\phi=0.5,\ \theta=0.3$, i.e.
\[
\psi_0=1,\qquad \psi_j=0.8\,(0.5)^{\,j-1}=1.6\,(0.5)^{j},\quad j\ge1 .
\]
\textbf{(c)} Using $\acvs_\tau=\sigma^2\sum_{j\ge0}\psi_j\psi_{j+\tau}$ with
$\sum_{j\ge1}\psi_j^2=0.8^2\sum_{j\ge1}0.25^{\,j-1}=0.64/0.75=\tfrac{64}{75}$,
\[
\acvs_0=\sigma^2\Big(1+\tfrac{64}{75}\Big)=\tfrac{139}{75}\sigma^2\approx1.853\,\sigma^2 .
\]
For $\acvs_1=\sigma^2\big(\psi_0\psi_1+\sum_{j\ge1}\psi_j\psi_{j+1}\big)$ with
$\sum_{j\ge1}\psi_j\psi_{j+1}=0.5\sum_{j\ge1}\psi_j^2=\tfrac12\cdot\tfrac{64}{75}=\tfrac{32}{75}$,
\[
\acvs_1=\sigma^2\Big(0.8+\tfrac{32}{75}\Big)=\tfrac{92}{75}\sigma^2\approx1.227\,\sigma^2 .
\]
(Quick check: $\acf_1=92/139\approx0.662$, sensibly large for a positively
correlated ARMA(1,1).)
\end{solution}

\begin{exercise}{Invertibility of an $\mathrm{MA}(2)$ and the unit-circle convention \quad\textnormal{[New]}}{p3-newma2}
Let $X_t=\eps_t-\theta_1\eps_{t-1}-\theta_2\eps_{t-2}$ with
$(\theta_1,\theta_2)=(0.2,\,0.48)$, i.e.\
$X_t=\eps_t-0.2\eps_{t-1}-0.48\eps_{t-2}$. Is it stationary? Invertible?
\end{exercise}
\begin{solution}
\textbf{Stationarity.} Every finite MA is stationary
(Prop.~\ref{prop:p3-alwaysinv}); no computation needed.\\
\textbf{Invertibility.} With $X_t=\Theta(\Bsh)\eps_t$ and
$\Theta(z)=1-0.2z-0.48z^2$, solve $0.48z^2+0.2z-1=0$:
\[
z=\frac{-0.2\pm\sqrt{0.04+1.92}}{0.96}=\frac{-0.2\pm1.4}{0.96}
=\Big\{\,\tfrac{5}{4},\ -\tfrac{5}{3}\,\Big\}.
\]
Both roots satisfy $\abs{z}>1$ ($1.25$ and $-1.667$), so the process is
\textbf{invertible}. Equivalently, $\Theta(z)=(1-0.8z)(1+0.6z)$ has reciprocal
roots $0.8$ and $-0.6$, both of modulus $<1$, the mirror statement of ``roots
outside''. (Had we instead picked coefficients whose root had $\abs{z}<1$, the
process would be non-invertible despite being perfectly stationary.)
\end{solution}

\begin{exercise}{Spectral density of an $\mathrm{ARMA}(1,1)$ \quad\textnormal{[New]}}{p3-newsdf}
Find the spectral density $\sdf(f)$ of the causal, invertible process
$(1-\phi\Bsh)X_t=(1-\theta\Bsh)\eps_t$ with $\abs\phi<1,\ \abs\theta<1$, and use
it to confirm $\acvs_0$ for the model
$(1-0.5\Bsh)X_t=(1+0.3\Bsh)\eps_t$ of Exercise~\ref{exo:p3-newarma11}.
\end{exercise}
\begin{solution}
By the ARMA spectral formula (Key Technique~\ref{kt:p3-sdf}),
\[
\sdf(f)=\sigma^2\,\frac{\abs{1-\theta e^{-2\pi i f}}^2}{\abs{1-\phi e^{-2\pi i f}}^2}
=\sigma^2\,\frac{1+\theta^2-2\theta\cos(2\pi f)}{1+\phi^2-2\phi\cos(2\pi f)},
\qquad -\tfrac12\le f\le\tfrac12 .
\]
Integrating over one period recovers the variance,
$\acvs_0=\int_{-1/2}^{1/2}\sdf(f)\,df$; doing the integral (or summing
$\sigma^2\sum_j\psi_j^2$ with $\psi_0=1,\ \psi_j=(\phi-\theta)\phi^{\,j-1}$) gives
the standard ARMA(1,1) variance
\[
\acvs_0=\sigma^2\,\frac{1+\theta^2-2\phi\theta}{1-\phi^2}.
\]
The model $(1-0.5\Bsh)X_t=(1+0.3\Bsh)\eps_t$ corresponds to $\phi=0.5$ and
$\theta=-0.3$ in the $(1-\theta\Bsh)$ convention. Then numerator
$=1+0.09-2(0.5)(-0.3)=1.39$, denominator $=1-0.25=0.75$, so
\[
\acvs_0=\frac{1.39}{0.75}\sigma^2=\frac{139}{75}\sigma^2\approx1.853\,\sigma^2,
\]
agreeing with Exercise~\ref{exo:p3-newarma11}. (Sanity check on the density:
since $\phi>0$ and the effective MA root is positive, $\sdf$ peaks at $f=0$, the
hallmark of a positively correlated, low-frequency-dominated process.)
\end{solution}

% ======================================================================
%  PART 4 -- Parametric Estimation: Yule--Walker and Least Squares
% ======================================================================
\section{Parametric Estimation: Yule--Walker and Least Squares}
\label{sec:estimation}

Having modelled $\{X_t\}$ as an ARMA process, we now \emph{fit} it: given data
$X_1,\dots,X_n$, estimate the autoregressive coefficients $\phi_1,\dots,\phi_p$,
the moving-average coefficients $\theta_1,\dots,\theta_q$ and the innovation
variance $\sigma_\eps^2$. This part develops the two workhorse estimators of the
course. The \textbf{Yule--Walker} (method-of-moments) estimator turns the AR
recursion into a linear system in the autocovariances and is solved by a single
matrix inversion (accelerated by the Levinson--Durbin algorithm). The
\textbf{least-squares} estimator minimises the sum of squared one-step
prediction errors; for pure AR models it is an ordinary linear regression
(forward, backward, or forward--backward), and for general ARMA models it
becomes a nonlinear minimisation in which the unobserved innovations
$\{\eps_t\}$ are reconstructed recursively. Throughout, $\eps_t\iid\mathcal
N(0,\sigma_\eps^2)$ is Gaussian white noise and the AR part is assumed causal
(roots of $\Phi(z)$ outside the unit circle), so that $\eps_t$ is uncorrelated
with the past $X_{t-k}$, $k>0$ --- the single fact that makes the whole
derivation work.

% ----------------------------------------------------------------------
\subsection{Definitions}

\begin{definition}{Yule--Walker estimators (AR$(p)$)}{p4-yw}
Let $\{X_t\}$ be a mean-zero causal AR$(p)$ process
$X_t=\sum_{k=1}^p\phi_kX_{t-k}+\eps_t$, observed at $X_1,\dots,X_n$. Form the
sample ACVS (mean known to be zero)
\[
\hat\acvs_\tau=\frac1n\sum_{t=1}^{n-\abs\tau}X_tX_{t+\abs\tau},
\]
collect $\hat{\boldsymbol\acvs}_p=(\hat\acvs_1,\dots,\hat\acvs_p)\Tr$ and the
symmetric Toeplitz matrix
\[
\hat\Gamma_p=
\begin{pmatrix}
\hat\acvs_0 & \hat\acvs_1 & \cdots & \hat\acvs_{p-1}\\
\hat\acvs_1 & \hat\acvs_0 & \cdots & \hat\acvs_{p-2}\\
\vdots & \vdots & \ddots & \vdots\\
\hat\acvs_{p-1} & \hat\acvs_{p-2} & \cdots & \hat\acvs_0
\end{pmatrix}.
\]
The \textbf{Yule--Walker estimators} are
\[
\hat{\boldsymbol\phi}_p=\hat\Gamma_p^{-1}\hat{\boldsymbol\acvs}_p,
\qquad
\hat\sigma_\eps^2=\hat\acvs_0-\sum_{j=1}^p\hat\phi_j\,\hat\acvs_j .
\]
They are the empirical analogue of the population identities
$\boldsymbol\acvs=\Gamma\boldsymbol\phi$ and
$\acvs_0=\sum_j\phi_j\acvs_j+\sigma_\eps^2$.
\end{definition}

\begin{definition}{Forward least-squares estimators (AR$(p)$)}{p4-fls}
Write the AR equations for the observable times $t=p+1,\dots,n$ as the linear
model $\mathbf X_F=F\boldsymbol\phi+\boldsymbol\eps$, where
\[
\mathbf X_F=\begin{pmatrix}X_{p+1}\\ \vdots\\ X_n\end{pmatrix},
\quad
F=\begin{pmatrix}
X_p & X_{p-1} & \cdots & X_1\\
X_{p+1} & X_p & \cdots & X_2\\
\vdots & \vdots & & \vdots\\
X_{n-1} & X_{n-2} & \cdots & X_{n-p}
\end{pmatrix},
\quad
\boldsymbol\phi=\begin{pmatrix}\phi_1\\ \vdots\\ \phi_p\end{pmatrix}.
\]
The \textbf{forward least-squares estimator} minimises
\[
\mathrm{SS}_F(\boldsymbol\phi)=\norm{\mathbf X_F-F\boldsymbol\phi}^2
=\sum_{t=p+1}^n\Big(X_t-\sum_{k=1}^p\phi_kX_{t-k}\Big)^2
=\sum_{t=p+1}^n\eps_t^2,
\]
and equals the ordinary least-squares solution
\[
\hat{\boldsymbol\phi}_F=(F\Tr F)^{-1}F\Tr\mathbf X_F,
\qquad
\hat\sigma_F^2=\frac{\norm{\mathbf X_F-F\hat{\boldsymbol\phi}_F}^2}{\,n-2p\,}.
\]
The first $p$ observations $X_1,\dots,X_p$ are discarded as conditioning values.
\end{definition}

\begin{definition}{Backward least-squares estimators (AR$(p)$)}{p4-bls}
Exploiting time reversibility of a Gaussian AR process, regress the
\emph{first} $n-p$ values on their \emph{future} lags: for $t=1,\dots,n-p$,
$X_t=\sum_{j=1}^p\phi_jX_{t+j}+\nu_t$. In matrix form
$\mathbf X_B=B\boldsymbol\phi+\boldsymbol\nu$ with
\[
\mathbf X_B=\begin{pmatrix}X_1\\ \vdots\\ X_{n-p}\end{pmatrix},
\qquad
B=\begin{pmatrix}
X_2 & X_3 & \cdots & X_{p+1}\\
X_3 & X_4 & \cdots & X_{p+2}\\
\vdots & \vdots & & \vdots\\
X_{n-p+1} & X_{n-p+2} & \cdots & X_n
\end{pmatrix}.
\]
The \textbf{backward least-squares estimator} minimises
$\mathrm{SS}_B(\boldsymbol\phi)=\norm{\mathbf X_B-B\boldsymbol\phi}^2
=\sum_{t=1}^{n-p}\big(X_t-\sum_{k=1}^p\phi_kX_{t+k}\big)^2$ and is
\[
\hat{\boldsymbol\phi}_B=(B\Tr B)^{-1}B\Tr\mathbf X_B,
\qquad
\hat\sigma_B^2=\frac{\norm{\mathbf X_B-B\hat{\boldsymbol\phi}_B}^2}{\,n-2p\,}.
\]
The same $\phi$-parameters appear because $Y_t=X_{-t}$ is an AR$(p)$ with
identical coefficients and an innovation $\tilde\nu_t\eqdist\eps_t$.
\end{definition}

\begin{definition}{Forward--backward least-squares estimator}{p4-fbls}
The \textbf{forward--backward} estimator $\hat{\boldsymbol\phi}_{FB}$ is the
minimiser of the \emph{combined} criterion
\[
\mathrm{SS}_F(\boldsymbol\phi)+\mathrm{SS}_B(\boldsymbol\phi).
\]
Equivalently it is the OLS solution of the stacked system obtained by piling
$\mathbf X_F,\mathbf X_B$ and $F,B$ on top of each other. Simulation studies
show it typically outperforms both the forward and the backward estimator.
\end{definition}

\begin{intuition}
The three least-squares variants only differ in \emph{which} $n-p$ one-step
residuals they square. Forward LS predicts each point from its past; backward
LS predicts it from its future (legitimate because a Gaussian AR run backwards
in time is the \emph{same} AR model); forward--backward LS uses both directions
at once, roughly doubling the effective data and stabilising the estimate near
the stationarity boundary.
\end{intuition}

\begin{definition}{Least squares for general ARMA via residual reconstruction}{p4-lsarma}
For an ARMA$(p,q)$ model $\Phi(\Bsh)X_t=\Theta(\Bsh)\eps_t$ the innovations
$\{\eps_t\}$ are unobserved. The least-squares estimator chooses the parameters
minimising $S(\boldsymbol\phi,\boldsymbol\theta)=\sum_t\hat\eps_t^{\,2}$, where
for each candidate parameter value the residuals are \textbf{reconstructed
recursively} from the model equation after fixing the $q$ pre-sample
innovations (and any pre-sample data) to zero, e.g.\
$\eps_0=\eps_{-1}=\cdots=0$. The minimisation is nonlinear (no closed form)
except in the pure-AR case, where it reduces to the linear forward least squares
above.
\end{definition}

% ----------------------------------------------------------------------
\subsection{Theorems \& Propositions}

\begin{proposition}{Yule--Walker equations}{p4-yweq}
Let $\{X_t\}$ be a mean-zero causal AR$(p)$ process with ACVS $\acvs_\tau$.
Then for every $k\ge1$,
\[
\acvs_k=\sum_{j=1}^p\phi_j\,\acvs_{k-j},
\]
and in particular, taking $k=1,\dots,p$ and using $\acvs_{-\tau}=\acvs_\tau$,
\[
\boldsymbol\acvs=\Gamma\boldsymbol\phi,
\qquad
\boldsymbol\acvs=(\acvs_1,\dots,\acvs_p)\Tr,\quad
\Gamma=(\acvs_{i-j})_{i,j=1}^p .
\]
Moreover $\acvs_0=\sum_{j=1}^p\phi_j\acvs_j+\sigma_\eps^2$.\\[2pt]
\textit{Proof idea:} multiply $X_t=\sum_j\phi_jX_{t-j}+\eps_t$ by $X_{t-k}$ and
take expectations. Because the process is \emph{causal}, $X_{t-k}$ depends only
on $\eps_{t-k},\eps_{t-k-1},\dots$, so for $k>0$ we have $\E[\eps_tX_{t-k}]=0$;
with mean zero $\E[X_{t-j}X_{t-k}]=\acvs_{k-j}$, giving the recursion. For the
variance, multiply by $X_t$ instead: the cross term becomes
$\E[\eps_tX_t]=\E[\eps_t(\sum_j\phi_jX_{t-j}+\eps_t)]=\sigma_\eps^2$ (all
$\E[\eps_tX_{t-j}]=0$, $j\ge1$).\\
\textit{Why:} these $p+1$ equations express the $p+1$ unknowns
$(\phi_1,\dots,\phi_p,\sigma_\eps^2)$ purely in terms of $\acvs_0,\dots,\acvs_p$;
replacing the population ACVS by its sample version yields the Yule--Walker
estimators. It is the canonical method-of-moments fit for AR models.
\end{proposition}

\begin{intuition}
The phrase ``causality $\Rightarrow\E[\eps_tX_{t-k}]=0$'' is the heart of the
derivation. Causality means $X_{t-k}$ is a function of innovations
\emph{strictly earlier} than $\eps_t$; uncorrelated white noise then kills the
cross term, leaving a clean linear system in the autocovariances. Without
causality the noise leaks into the past values and the Yule--Walker equations
no longer hold.
\end{intuition}

\begin{proposition}{Minimising the prediction MSE gives Yule--Walker}{p4-mse}
Let $\{Y_t\}$ be mean-zero stationary with ACVS $\acvs$. The coefficients
$\phi_1,\dots,\phi_k$ minimising the one-step mean-square prediction error
\[
\E\Big[\big(Y_t-\textstyle\sum_{i=1}^k\phi_iY_{t-i}\big)^2\Big]
\]
satisfy the Yule--Walker equations $\acvs_j=\sum_{i=1}^k\phi_i\,\acvs_{j-i}$ for
$j=1,\dots,k$.\\[2pt]
\textit{Proof idea:} since $\{Y_t\}$ is mean-zero, the objective equals the
variance
$\acvs_0-2\sum_i\phi_i\acvs_i+\sum_{i,j}\phi_i\phi_j\acvs_{i-j}$, a positive
quadratic in $\boldsymbol\phi$. Setting $\partial/\partial\phi_i=0$ gives
$-2\acvs_i+2\sum_j\phi_j\acvs_{i-j}=0$, i.e.\ the Yule--Walker equations.\\
\textit{Why:} it shows Yule--Walker is not merely a moment trick but the
\emph{best linear predictor}: the AR coefficients are exactly the optimal
linear-prediction weights, which is why YW and least squares agree
asymptotically for AR models.
\end{proposition}

\begin{proposition}{Forward / backward least squares are the OLS estimators}{p4-ols}
With the design matrices of Definitions~\ref{def:p4-fls}--\ref{def:p4-bls}, the
minimisers of $\mathrm{SS}_F$ and $\mathrm{SS}_B$ are
\[
\hat{\boldsymbol\phi}_F=(F\Tr F)^{-1}F\Tr\mathbf X_F,
\qquad
\hat{\boldsymbol\phi}_B=(B\Tr B)^{-1}B\Tr\mathbf X_B,
\]
with residual-variance estimators
$\hat\sigma^2=\norm{\text{residuals}}^2/(n-2p)$.\\[2pt]
\textit{Proof idea:} each is a standard linear-regression sum of squares
$\norm{\mathbf y-A\boldsymbol\phi}^2$; the normal equations
$A\Tr A\,\boldsymbol\phi=A\Tr\mathbf y$ give the stated solution, and the
mean-square residual divides by the degrees of freedom $(n-p)-p$.\\
\textit{Why:} it places AR estimation inside the familiar OLS framework, so all
regression machinery (existence, uniqueness when $A\Tr A$ is invertible,
$\hat\sigma^2$) applies directly.
\end{proposition}

\begin{proposition}{Time reversibility of a Gaussian AR process}{p4-timerev}
If $\{X_t\}$ is a stationary Gaussian AR$(p)$ process with coefficients
$\phi_1,\dots,\phi_p$ and innovation variance $\sigma_\eps^2$, then the
time-reversed process $Y_t=X_{-t}$ is an AR$(p)$ process with the
\emph{same} coefficients and innovation variance:
\[
Y_t=\sum_{j=1}^p\phi_jY_{t-j}+\tilde\nu_t,\qquad \tilde\nu_t\eqdist\eps_t,
\]
equivalently $X_t=\sum_{j=1}^p\phi_jX_{t+j}+\nu_t$ with $\nu_t=\tilde\nu_{-t}$.\\[2pt]
\textit{Proof idea:} a stationary Gaussian process is determined by its ACVS,
and the ACVS is symmetric, $\acvs_{-\tau}=\acvs_\tau$; hence the forward and
reversed processes have identical second-order structure and (being Gaussian)
identical distributions. The full argument is given later via frequency-domain
methods (the spectral density is even in $f$).\\
\textit{Why:} it justifies the backward least-squares regression --- predicting
the present from the future uses the very same parameters --- and underpins the
forward--backward estimator.
\end{proposition}

\begin{proposition}{Levinson--Durbin: solving Yule--Walker in $O(p^2)$}{p4-ld}
The Yule--Walker system
$\hat\Gamma_p\hat{\boldsymbol\phi}_p=\hat{\boldsymbol\acvs}_p$ has a symmetric
Toeplitz coefficient matrix. The \textbf{Levinson--Durbin algorithm} exploits
this structure to solve it recursively in $k=1,\dots,p$, costing $O(p^2)$
operations rather than the $O(p^3)$ of a naive matrix inversion. As a
by-product it produces the partial autocorrelations
$\hat\pacf_k=\hat\phi_{k,k}$ and the prediction-error variances at every order.\\[2pt]
\textit{Proof idea:} the order-$k$ solution is obtained from the order-$(k-1)$
solution by one rank-one update governed by the reflection coefficient
$\hat\phi_{k,k}=(\hat\acvs_k-\sum_{j=1}^{k-1}\hat\phi_{k-1,j}\hat\acvs_{k-j})/
v_{k-1}$, where $v_{k-1}$ is the running error variance; the remaining
coefficients update as
$\hat\phi_{k,j}=\hat\phi_{k-1,j}-\hat\phi_{k,k}\hat\phi_{k-1,k-j}$.\\
\textit{Why:} for high-order AR fits, spectral estimation, or repeatedly
solving nested systems, the quadratic cost (and the free PACF) makes
Levinson--Durbin the standard solver.
\end{proposition}

\begin{proposition}{Least-squares fits need not be stationary}{p4-nonstat}
The forward, backward and forward--backward least-squares estimators
$\hat{\boldsymbol\phi}_F,\hat{\boldsymbol\phi}_B,\hat{\boldsymbol\phi}_{FB}$ are
\emph{not} constrained to yield a stationary (causal) model: the implied
polynomial $\hat\Phi(z)$ may have roots inside the unit circle.\\[2pt]
\textit{Proof idea:} the OLS objective places no constraint on the root moduli;
it only minimises the residual sum of squares, which can be smaller for a
non-causal parameter vector.\\
\textit{Why:} non-stationary estimates are a problem for \emph{prediction}, but
acceptable for some uses such as \emph{spectral estimation}, where they still
produce a valid spectral density estimate. The Yule--Walker estimator, by
contrast, always yields a stationary fit.
\end{proposition}

\begin{intuition}[Yule--Walker vs.\ least squares]
For large $n$ the two approaches give very similar AR estimates: the
least-squares score equations reduce, to leading order, to
$\hat\mu\approx\bar Y$ and $\hat\phi\approx\hat\acf_1$, which is precisely what
method-of-moments delivers. The practical differences are at the margins:
Yule--Walker is a single linear solve and always stationary; least squares
extends cleanly to ARMA and to including a mean, but can stray outside the
stationary region.
\end{intuition}

% ----------------------------------------------------------------------
\subsection{Key Techniques}

\begin{keytech}{Yule--Walker from a sample ACVS / ACF}{p4-ktyw}
To fit AR$(p)$ by Yule--Walker:
\begin{enumerate}
  \item \textbf{Estimate the ACVS} $\hat\acvs_0,\dots,\hat\acvs_p$ (divide-by-$n$
        estimator on the mean-corrected series).
  \item \textbf{Build} the $p\times p$ Toeplitz matrix
        $\hat\Gamma_p=(\hat\acvs_{i-j})$ and the vector
        $\hat{\boldsymbol\acvs}_p=(\hat\acvs_1,\dots,\hat\acvs_p)\Tr$.
  \item \textbf{Solve} $\hat{\boldsymbol\phi}_p=\hat\Gamma_p^{-1}\hat{\boldsymbol\acvs}_p$
        (by hand for $p\le3$; otherwise Levinson--Durbin).
  \item \textbf{Variance:}
        $\hat\sigma_\eps^2=\hat\acvs_0-\sum_{j=1}^p\hat\phi_j\hat\acvs_j$.
\end{enumerate}
\textbf{ACF shortcut:} if only the autocorrelations $\hat\acf_\tau$ are given,
divide every entry by $\hat\acvs_0$. Writing $V=\hat\acvs_0^{-1}I$, the system
$\hat{\boldsymbol\acvs}_p=\hat\Gamma_p\hat{\boldsymbol\phi}_p$ becomes
$V\hat{\boldsymbol\acvs}_p=V\hat\Gamma_p\hat{\boldsymbol\phi}_p$, i.e.\
\[
\hat{\boldsymbol\phi}_p=\bar\Gamma_p^{-1}\hat{\boldsymbol\acf}_p,
\qquad
\bar\Gamma_p=(\hat\acf_{i-j}),\quad
\hat{\boldsymbol\acf}_p=(\hat\acf_1,\dots,\hat\acf_p)\Tr,
\]
so the \emph{same} solve works with the correlation matrix and the correlation
vector --- $\hat\acvs_0$ cancels. (It is needed only for $\hat\sigma_\eps^2$,
via $\hat\sigma_\eps^2=\hat\acvs_0(1-\sum_j\hat\phi_j\hat\acf_j)$.)
\end{keytech}

\begin{keytech}{Recursive innovation reconstruction for MA / ARMA least squares}{p4-ktrec}
To evaluate $S(\cdot)=\sum_t\hat\eps_t^2$ for an MA or ARMA model at a candidate
parameter value, rebuild the unobserved innovations one step at a time:
\begin{enumerate}
  \item \textbf{Solve the model equation for $\eps_t$.}
        MA(1) $X_t=\eps_t-\theta\eps_{t-1}\Rightarrow\eps_t=X_t+\theta\eps_{t-1}$;
        ARMA(1,1) $X_t=\phi X_{t-1}+\eps_t-\theta\eps_{t-1}
        \Rightarrow\eps_t=X_t-\phi X_{t-1}+\theta\eps_{t-1}$.
  \item \textbf{Initialise the $q$ pre-sample innovations (and pre-sample data)
        to zero:} $\eps_0=0$ (and $X_0=0$ for the ARMA case).
  \item \textbf{Recurse forward} $t=1,2,\dots,n$ to obtain
        $\hat\eps_1,\hat\eps_2,\dots,\hat\eps_n$, each a function of the unknown
        parameters.
  \item \textbf{Form $S$ and minimise} over the parameters (closed form only if
        $S$ is quadratic, as in the toy MA(1) of
        Exercise~\ref{exo:p4-e43}; otherwise numerically).
\end{enumerate}
\textbf{Why it is needed:} an MA/ARMA model is linear in the data but the
innovations are latent, so least squares is no longer a one-shot regression ---
the recursion supplies the residuals before squaring them. The method works for
any ARMA$(p,q)$; one always assumes the $q$ initial innovations vanish.
\end{keytech}

\begin{keytech}{Forward vs.\ backward design matrices (AR$(2)$ template)}{p4-ktfb}
For AR$(2)$ ($p=2$) the two regressions are
\[
\underbrace{\begin{pmatrix}X_3\\ X_4\\ \vdots\\ X_n\end{pmatrix}}_{\mathbf X_F}
=\underbrace{\begin{pmatrix}X_2 & X_1\\ X_3 & X_2\\ \vdots & \vdots\\ X_{n-1}&X_{n-2}\end{pmatrix}}_{F}
\begin{pmatrix}\phi_1\\ \phi_2\end{pmatrix}+\boldsymbol\eps,
\qquad
\underbrace{\begin{pmatrix}X_1\\ X_2\\ \vdots\\ X_{n-2}\end{pmatrix}}_{\mathbf X_B}
=\underbrace{\begin{pmatrix}X_2 & X_3\\ X_3 & X_4\\ \vdots & \vdots\\ X_{n-1}&X_n\end{pmatrix}}_{B}
\begin{pmatrix}\phi_1\\ \phi_2\end{pmatrix}+\boldsymbol\nu.
\]
Forward regresses the last $n-p$ points on their past lags; backward regresses
the first $n-p$ points on their future lags. Then
$\hat{\boldsymbol\phi}=(A\Tr A)^{-1}A\Tr\mathbf X$ with $A=F$ or $A=B$. Remove
the sample mean first if it is significantly non-zero.
\end{keytech}

% ----------------------------------------------------------------------
\subsection{Exercises}

\begin{exercise}{Yule--Walker on the sunspot series \quad\textnormal{[Sheet 4]}}{p4-e41}
For the Wolfer sunspot numbers the sample autocovariances are
$\hat\acvs_0=1382.185$, $\hat\acvs_1=1114.378$, $\hat\acvs_2=591.721$,
$\hat\acvs_3=96.215$. Fit the AR$(2)$ model
$Y_t=\phi_1Y_{t-1}+\phi_2Y_{t-2}+\eps_t$ ($Y_t$ mean-corrected) by Yule--Walker,
estimating $\phi_1,\phi_2$ and $\sigma_\eps^2$.
\end{exercise}
\begin{solution}
With $p=2$ the Yule--Walker system
$\hat\Gamma_2\hat{\boldsymbol\phi}=\hat{\boldsymbol\acvs}_2$ is
\[
\begin{pmatrix}\hat\acvs_0 & \hat\acvs_1\\ \hat\acvs_1 & \hat\acvs_0\end{pmatrix}
\begin{pmatrix}\hat\phi_1\\ \hat\phi_2\end{pmatrix}
=\begin{pmatrix}\hat\acvs_1\\ \hat\acvs_2\end{pmatrix}
\quad\Longrightarrow\quad
\begin{pmatrix}1382.185 & 1114.378\\ 1114.378 & 1382.185\end{pmatrix}
\begin{pmatrix}\hat\phi_1\\ \hat\phi_2\end{pmatrix}
=\begin{pmatrix}1114.378\\ 591.721\end{pmatrix}.
\]
The inverse of a $2\times2$ symmetric matrix
$\left(\begin{smallmatrix}a&b\\ b&a\end{smallmatrix}\right)$ is
$\tfrac{1}{a^2-b^2}\left(\begin{smallmatrix}a&-b\\ -b&a\end{smallmatrix}\right)$,
with here $a=1382.185,\ b=1114.378$ and determinant
$a^2-b^2=1.9104\times10^6-1.2418\times10^6=6.686\times10^5$. Hence
\[
\hat{\boldsymbol\phi}
=\frac{1}{6.686\times10^5}
\begin{pmatrix}1382.185 & -1114.378\\ -1114.378 & 1382.185\end{pmatrix}
\begin{pmatrix}1114.378\\ 591.721\end{pmatrix}
\approx\begin{pmatrix}1.3175\\ -0.6341\end{pmatrix}.
\]
The innovation variance is
\[
\hat\sigma_\eps^2=\hat\acvs_0-\hat\phi_1\hat\acvs_1-\hat\phi_2\hat\acvs_2
=1382.185-1.3175(1114.378)-(-0.6341)(591.721)\approx289.21 .
\]
So $\hat\phi_1\approx1.3175$, $\hat\phi_2\approx-0.6341$,
$\hat\sigma_\eps^2\approx289.21$. (One checks the fit is causal: the roots of
$1-1.3175z+0.6341z^2$ are complex with modulus
$1/\sqrt{0.6341}\approx1.256>1$.)
\end{solution}

\begin{exercise}{Deriving Yule--Walker by minimising the MSE \quad\textnormal{[Sheet 4]}}{p4-e42}
Let $\{Y_t\}$ be mean-zero stationary with ACVS $\acvs$. Show that the
$\phi_1,\dots,\phi_k$ minimising
$\E\big[(Y_t-\sum_{i=1}^k\phi_iY_{t-i})^2\big]$ satisfy
$\acvs_j=\sum_{i=1}^k\phi_i\acvs_{j-i}$, $j=1,\dots,k$.
\end{exercise}
\begin{solution}
Write $W=Y_t-\sum_{i=1}^k\phi_iY_{t-i}$. Decompose
$\E[W^2]=\Var(W)+(\E[W])^2$. Because $\{Y_t\}$ is stationary,
$\E[W]=(1-\sum_i\phi_i)\E[Y_t]$, and since the process is mean-zero this term
vanishes, so $\E[W^2]=\Var(W)$. Expanding the variance by bilinearity,
\[
\begin{aligned}
\E[W^2]
&=\Cov\Big(Y_t-\sum_{i=1}^k\phi_iY_{t-i},\ Y_t-\sum_{j=1}^k\phi_jY_{t-j}\Big)\\
&=\Cov(Y_t,Y_t)-2\sum_{i=1}^k\phi_i\Cov(Y_t,Y_{t-i})
   +\sum_{i=1}^k\sum_{j=1}^k\phi_i\phi_j\Cov(Y_{t-i},Y_{t-j})\\
&=\acvs_0-2\sum_{i=1}^k\phi_i\acvs_i
   +\sum_{i=1}^k\sum_{j=1}^k\phi_i\phi_j\acvs_{i-j}.
\end{aligned}
\]
This is a positive (semi-)definite quadratic in $\boldsymbol\phi$, so its unique
stationary point is the global minimum. Differentiating with respect to the
$i$-th coefficient (using $\acvs_{i-j}=\acvs_{j-i}$ to combine the two
double-sum terms) and equating to zero,
\[
\frac{\partial\E[W^2]}{\partial\phi_i}=-2\acvs_i+2\sum_{j=1}^k\phi_j\acvs_{i-j}=0
\quad\Longrightarrow\quad
\acvs_i=\sum_{j=1}^k\phi_j\acvs_{i-j},\qquad i=1,\dots,k,
\]
which are exactly the Yule--Walker equations.
\end{solution}

\begin{exercise}{Estimating $\theta$ of an MA(1) by recursive least squares \quad\textnormal{[Sheet 4]}}{p4-e43}
A mean-zero MA(1) $X_t=\eps_t-\theta\eps_{t-1}$ is observed at $X_1=0$, $X_2=1$,
$X_3=0.5$. Estimate $\theta$ by least squares.
\end{exercise}
\begin{solution}
The innovations are latent, so reconstruct them recursively. Solving the model
equation for $\eps_t$ gives $\eps_t=X_t+\theta\eps_{t-1}$. Take $t=0$ as the
start and set $\eps_0=0$:
\[
\eps_1=X_1+\theta\eps_0=0,\qquad
\eps_2=X_2+\theta\eps_1=1,\qquad
\eps_3=X_3+\theta\eps_2=0.5+\theta .
\]
The least-squares criterion is
\[
S(\theta)=\sum_{t=1}^3\eps_t^2=0^2+1^2+(0.5+\theta)^2=1+(0.5+\theta)^2 .
\]
This is minimised when $(0.5+\theta)^2$ is smallest, i.e.\ at $\theta=-0.5$.
Hence $\hat\theta=-0.5$. (Note $\abs{\hat\theta}<1$, so the fit is invertible.)
\end{solution}

\begin{exercise}{AR$(3)$ Yule--Walker from a given ACF \quad\textnormal{[Sheet 4]}}{p4-e44}
A series of length $100$ has sample ACF $\hat\acf_1=0.8$, $\hat\acf_2=0.5$,
$\hat\acf_3=0.4$. Assuming a zero-mean AR$(3)$ model, estimate
$\phi_1,\phi_2,\phi_3$.
\end{exercise}
\begin{solution}
Use the ACF shortcut: dividing the Yule--Walker system by $\acvs_0$ replaces
covariances by correlations, $\hat{\boldsymbol\phi}=\bar\Gamma_3^{-1}\hat{\boldsymbol\acf}_3$,
where $\bar\Gamma_3=(\hat\acf_{i-j})$ uses $\hat\acf_0=1$:
\[
\begin{pmatrix}1 & 0.8 & 0.5\\ 0.8 & 1 & 0.8\\ 0.5 & 0.8 & 1\end{pmatrix}
\begin{pmatrix}\hat\phi_1\\ \hat\phi_2\\ \hat\phi_3\end{pmatrix}
=\begin{pmatrix}0.8\\ 0.5\\ 0.4\end{pmatrix}.
\]
Solving this $3\times3$ system (e.g.\ by Gaussian elimination) gives
\[
\hat{\boldsymbol\phi}=(\,1.30909,\ -0.95455,\ 0.50909\,)\Tr .
\]
\textit{Check (row 1):} $1(1.30909)+0.8(-0.95455)+0.5(0.50909)
=1.30909-0.76364+0.25455=0.8$.\ \checkmark\;
\textit{Row 3:} $0.5(1.30909)+0.8(-0.95455)+1(0.50909)
=0.65455-0.76364+0.50909=0.4$.\ \checkmark\;
So $\hat\phi_1\approx1.309$, $\hat\phi_2\approx-0.955$,
$\hat\phi_3\approx0.509$. (The series length $n=100$ does not enter; only the
ACF values determine the point estimates.)
\end{solution}

\begin{exercise}{Forward and backward least squares (conceptual) \quad\textnormal{[Sheet 4]}}{p4-e45}
Describe how to compute the forward and backward least-squares AR$(2)$ fits to
the sunspot series, and interpret the reported output
$\hat{\boldsymbol\phi}_F\approx(1.4046,-0.7113)\Tr$,
$\hat\sigma_F^2\approx232.26$,
$\hat{\boldsymbol\phi}_B\approx(1.4061,-0.7076)\Tr$,
$\hat\sigma_B^2\approx231.03$.
\end{exercise}
\begin{solution}
\textbf{Procedure.} Centre the series, $x_t\mapsto x_t-\bar x$ (the sunspot mean
is clearly non-zero). For the forward fit build $\mathbf X_F=(x_3,\dots,x_n)\Tr$
and the $(n-2)\times2$ matrix $F$ with rows $(x_{t-1},x_{t-2})$, then
$\hat{\boldsymbol\phi}_F=(F\Tr F)^{-1}F\Tr\mathbf X_F$ and
$\hat\sigma_F^2=\norm{\mathbf X_F-F\hat{\boldsymbol\phi}_F}^2/(n-2p)$ with
$p=2$. For the backward fit build $\mathbf X_B=(x_1,\dots,x_{n-2})\Tr$ and $B$
with rows $(x_{t+1},x_{t+2})$, then
$\hat{\boldsymbol\phi}_B=(B\Tr B)^{-1}B\Tr\mathbf X_B$ with the analogous
$\hat\sigma_B^2$. In R these are two short functions building the lagged design
columns and calling \texttt{solve(t(A)\%*\%A)\%*\%t(A)\%*\%X}.

\textbf{Interpretation.} The forward and backward estimates are almost
identical ($\hat{\boldsymbol\phi}_F\approx\hat{\boldsymbol\phi}_B$), a direct
illustration of \emph{time reversibility}: a Gaussian AR run backwards has the
same coefficients, so both regressions estimate the same $\boldsymbol\phi$. The
small discrepancy is sampling noise from using different $n-p$ residuals. Both
$\hat\phi_1\approx1.40$, $\hat\phi_2\approx-0.71$ are close to (but not equal to)
the Yule--Walker values $(1.3175,-0.6341)$ of Exercise~\ref{exo:p4-e41}; the two
methods agree only asymptotically, and least squares does not constrain the fit
to be stationary. A forward--backward fit would minimise
$\mathrm{SS}_F+\mathrm{SS}_B$ jointly, typically landing between the two and with
slightly lower variance.
\end{solution}

\begin{exercise}{AR$(1)$ Yule--Walker, by hand \quad\textnormal{[New]}}{p4-e46}
For a mean-zero AR$(1)$ $X_t=\phi X_{t-1}+\eps_t$, a series of length $n=200$
gives $\hat\acvs_0=4.0$ and $\hat\acvs_1=2.4$. Find the Yule--Walker estimates
$\hat\phi$ and $\hat\sigma_\eps^2$, and state the implied stationarity check.
\end{exercise}
\begin{solution}
For $p=1$ the Yule--Walker system collapses to a scalar equation
$\hat\acvs_1=\hat\phi\,\hat\acvs_0$, so
\[
\hat\phi=\frac{\hat\acvs_1}{\hat\acvs_0}=\frac{2.4}{4.0}=0.6=\hat\acf_1 .
\]
(For AR(1) the YW estimate is just the lag-1 sample autocorrelation.) The
variance estimator is
\[
\hat\sigma_\eps^2=\hat\acvs_0-\hat\phi\,\hat\acvs_1
=4.0-0.6\times2.4=4.0-1.44=2.56 .
\]
This is consistent with the population identity
$\acvs_0=\sigma_\eps^2/(1-\phi^2)$: indeed
$\hat\acvs_0(1-\hat\phi^2)=4.0(1-0.36)=2.56=\hat\sigma_\eps^2$.\ \checkmark\;
\textbf{Stationarity:} $\Phi(z)=1-0.6z$ has root $z=1/0.6\approx1.667$, of
modulus $>1$, so the fit is causal/stationary.
\end{solution}

\begin{exercise}{ARMA(1,1) residual recursion \quad\textnormal{[New]}}{p4-e47}
For a mean-zero ARMA(1,1) $X_t=\phi X_{t-1}+\eps_t-\theta\eps_{t-1}$, write out
the recursion that reconstructs $\hat\eps_1,\dots,\hat\eps_4$ as functions of
$(\phi,\theta)$ from data $X_1,X_2,X_3,X_4$, stating the initialisation, and
exhibit the least-squares objective.
\end{exercise}
\begin{solution}
Solve the model equation for the current innovation:
$\eps_t=X_t-\phi X_{t-1}+\theta\eps_{t-1}$. Initialise the pre-sample data and
innovation to zero, $X_0=0$ and $\eps_0=0$. Then
\[
\begin{aligned}
\hat\eps_1&=X_1-\phi X_0+\theta\eps_0=X_1,\\
\hat\eps_2&=X_2-\phi X_1+\theta\hat\eps_1=X_2-\phi X_1+\theta X_1,\\
\hat\eps_3&=X_3-\phi X_2+\theta\hat\eps_2
            =X_3-\phi X_2+\theta\big(X_2-\phi X_1+\theta X_1\big),\\
\hat\eps_4&=X_4-\phi X_3+\theta\hat\eps_3 .
\end{aligned}
\]
Each $\hat\eps_t$ is an explicit (nonlinear in $\theta$) function of the
parameters. The least-squares estimator minimises
\[
S(\phi,\theta)=\sum_{t=1}^4\hat\eps_t^{\,2}
=\sum_{t=1}^4\big(X_t-\phi X_{t-1}+\theta\hat\eps_{t-1}\big)^2 ,
\]
which, because of the $\theta\hat\eps_{t-1}$ feedback, is \emph{not} quadratic
in $\theta$ and must be minimised numerically. This is the general
ARMA$(p,q)$ recipe with $q=1$ pre-sample innovation set to zero.
\end{solution}

\begin{exercise}{Why an MA(1) least-squares fit can be quadratic --- or not \quad\textnormal{[New]}}{p4-e48}
For the MA(1) $X_t=\eps_t-\theta\eps_{t-1}$ with $\eps_0=0$ and data
$X_1,\dots,X_n$, show that the reconstructed residual $\hat\eps_t$ is a
polynomial of degree $t-1$ in $\theta$, and explain why the toy three-point
problem of Exercise~\ref{exo:p4-e43} nevertheless had a clean quadratic
minimisation.
\end{exercise}
\begin{solution}
The recursion $\hat\eps_t=X_t+\theta\hat\eps_{t-1}$ with $\hat\eps_0=0$ unrolls
to the geometric-type sum
\[
\hat\eps_t=\sum_{k=0}^{t-1}\theta^{\,k}X_{t-k}
=X_t+\theta X_{t-1}+\theta^2X_{t-2}+\cdots+\theta^{\,t-1}X_1,
\]
a polynomial in $\theta$ of degree $t-1$. Hence
$S(\theta)=\sum_{t=1}^n\hat\eps_t^2$ is a polynomial of degree $2(n-1)$ in
$\theta$, generally with several stationary points --- so in general the MA(1)
least-squares fit requires numerical optimisation, not a single derivative.

In Exercise~\ref{exo:p4-e43} the data were special: $X_1=0$ forces
$\hat\eps_1=0$ and $\hat\eps_2=X_2=1$ (free of $\theta$), so only the last
residual $\hat\eps_3=X_3+\theta\hat\eps_2=0.5+\theta$ carries any
$\theta$-dependence, and it is \emph{linear}. Therefore
$S(\theta)=1+(0.5+\theta)^2$ is exactly quadratic with the unique minimiser
$\hat\theta=-0.5$. The clean answer was an artefact of $X_1=0$; with generic
data the degree-$2(n-1)$ surface would need a numerical solver.
\end{solution}

\begin{pitfall}
Three traps in this part. (1) \textbf{Indexing of the Toeplitz system:}
$\Gamma_p$ uses $\acvs_{i-j}$ with $\acvs_0$ on the diagonal, and the
\emph{right-hand side} is $(\acvs_1,\dots,\acvs_p)$ --- do not put $\acvs_0$ in
the vector. (2) \textbf{Degrees of freedom in $\hat\sigma^2$:} the least-squares
variance divides by $n-2p$ (there are $n-p$ residuals and $p$ fitted
coefficients), not by $n$. (3) \textbf{Forgetting to initialise the
recursion:} in MA/ARMA least squares you must set the $q$ pre-sample
innovations (and any pre-sample data) to zero before recursing, otherwise
$\hat\eps_1$ is undefined.
\end{pitfall}

% ======================================================================
%  PART 5 -- Trends, Seasonality and Differencing: SARIMA
% ======================================================================
\section{Trends, Seasonality and Differencing: SARIMA}
\label{sec:sarima}

Real signals rarely look stationary: they drift (a \emph{trend}) and they
repeat (a \emph{season}). The working model of this chapter is
$X_t=\mu_t+Y_t$ with $\mu_t$ a deterministic time-dependent mean and $\{Y_t\}$
stationary. There are two ways to peel off $\mu_t$: \emph{estimate and subtract}
(regression on $t$ and on seasonal dummies), or \emph{difference it away} with
the operator $\diff=I-B$. Differencing of order $d$ annihilates a polynomial
trend of degree $d-1$; a \emph{seasonal} difference $\diff_{(s)}=I-B^{s}$ kills
a period-$s$ season. Layering autoregressive and moving-average structure onto
\emph{both} the ordinary lags and the seasonal lags gives the SARMA family, and
folding the differencing in front gives the (S)ARIMA model
$\Phi(B)(1-B)^{d}X_t=\Theta(B)\eps_t$ --- the standard Box--Jenkins workhorse.
Throughout, $\{\eps_t\}$ is mean-zero white noise of variance $\sigma_\eps^2$,
AR/seasonal-AR coefficients are $\phi_j,\phi^{(s)}_j$ and MA/seasonal-MA
coefficients are $\theta_k,\theta^{(s)}_k$.

% ----------------------------------------------------------------------
\subsection{Definitions}

\begin{definition}{Mean-trend (signal-plus-noise) model}{p5-trendmodel}
A series with a slowly varying level is modelled as
\[
X_t=\mu_t+Y_t,
\]
where $\mu_t$ is a deterministic, time-dependent mean (the \textbf{trend}) and
$\{Y_t\}$ is a zero-mean second-order stationary process. The canonical linear
case is $\mu_t=a+bt$; more generally $\mu_t$ may be a degree-$(q-1)$ polynomial
in $t$. When the level instead repeats, $\mu_t=s_t$ with $s_{t+s}=s_t$ is a
\textbf{periodic} (seasonal) trend of period $s$; the two can coexist,
$X_t=a+bt+s_t+Y_t$.
\end{definition}

\begin{definition}{Difference operator $\diff$}{p5-diff}
The (first-)\textbf{difference operator} is
\[
\diff=I-B,\qquad \diff X_t=X_t-X_{t-1},
\]
where $B$ is the backshift operator $BX_t=X_{t-1}$. It is \emph{linear}. Its
$d$-th power is, by the binomial theorem,
\[
\diff^{d}=(I-B)^{d}=\sum_{k=0}^{d}\binom{d}{k}(-1)^{k}B^{k},
\qquad
\diff^{d}X_t=\sum_{k=0}^{d}\binom{d}{k}(-1)^{k}X_{t-k}.
\]
For example $\diff^{2}X_t=X_t-2X_{t-1}+X_{t-2}$.
\end{definition}

\begin{definition}{Seasonal difference operator $\diff_{(s)}$}{p5-seasdiff}
The \textbf{seasonal difference} of period $s$ is
\[
\diff_{(s)}=I-B^{s},\qquad \diff_{(s)}X_t=X_t-X_{t-s}.
\]
It removes a deterministic component of period $s$ (since
$\diff_{(s)}s_t=s_t-s_{t-s}=0$) and, applied to a linear trend, leaves a
constant: $\diff_{(s)}(bt)=bs$. The ordinary difference is the case $s=1$,
$\diff=\diff_{(1)}$. To kill a linear-plus-seasonal trend one composes them,
$\diff\diff_{(s)}=(I-B)(I-B^{s})$.
\end{definition}

\begin{definition}{Seasonal moving average $\mathrm{SMA}(Q)_s$}{p5-sma}
A \textbf{seasonal moving average} of order $Q$ at period $s$ is
\[
X_t=\eps_t-\sum_{j=1}^{Q}\theta^{(s)}_j\,\eps_{t-js}
=\Theta^{(s)}(B)\eps_t,
\qquad
\Theta^{(s)}(z)=1-\sum_{j=1}^{Q}\theta^{(s)}_j\,z^{sj}.
\]
Only lags that are \emph{multiples of $s$} appear. The basic seasonal MA(1) is
$X_t=\eps_t-\theta\eps_{t-s}$.
\end{definition}

\begin{definition}{Seasonal autoregression $\mathrm{SAR}(P)_s$}{p5-sar}
A \textbf{seasonal autoregression} of order $P$ at period $s$ is
\[
X_t=\eps_t+\sum_{j=1}^{P}\phi^{(s)}_j\,X_{t-js},
\qquad
\Phi^{(s)}(z)=1-\sum_{j=1}^{P}\phi^{(s)}_j\,z^{sj},
\qquad
\Phi^{(s)}(B)X_t=\eps_t .
\]
Again only the seasonal lags $s,2s,\dots,Ps$ enter.
\end{definition}

\begin{definition}{SARMA$(p,q)\times(P,Q)_s$}{p5-sarma}
Combining an ordinary $\mathrm{ARMA}(p,q)$ with a seasonal
$\mathrm{ARMA}(P,Q)_s$ gives the \textbf{multiplicative seasonal ARMA}
\[
\Phi(B)\,\Phi^{(s)}(B)\,X_t=\Theta(B)\,\Theta^{(s)}(B)\,\eps_t,
\]
with the four polynomials
\(\Phi(z)=1-\sum_{j=1}^p\phi_jz^j\),
\(\Theta(z)=1-\sum_{k=1}^q\theta_kz^k\),
\(\Phi^{(s)}(z)=1-\sum_{j=1}^P\phi^{(s)}_jz^{sj}\),
\(\Theta^{(s)}(z)=1-\sum_{k=1}^Q\theta^{(s)}_kz^{sk}\).
Special cases: $\mathrm{SMA}(Q)_s$ is $\mathrm{SARMA}(0,0)\times(0,Q)_s$ and
$\mathrm{SAR}(P)_s$ is $\mathrm{SARMA}(0,0)\times(P,0)_s$. Example, a
$\mathrm{SARMA}(0,1)\times(0,1)_{12}$:
\[
X_t=(1-\theta B)\bigl(1-\theta^{(12)}B^{12}\bigr)\eps_t .
\]
\end{definition}

\begin{definition}{ARIMA$(p,d,q)$}{p5-arima}
$\{X_t\}$ is an \textbf{ARIMA$(p,d,q)$} (``integrated ARMA'') process if its
$d$-th difference is a (causal, invertible) $\mathrm{ARMA}(p,q)$:
\[
\Phi(B)\,(1-B)^{d}X_t=\Theta(B)\,\eps_t .
\]
Equivalently, writing $Z_t=\diff^{d}X_t$,
\[
Z_t=\sum_{j=1}^{p}\phi_jZ_{t-j}+\eps_t-\sum_{j=1}^{q}\theta_j\eps_{t-j}.
\]
The ``I'' stands for \emph{integrated}: $X_t$ is recovered from $Z_t$ by summing
(the inverse of differencing) $d$ times. ARMA is the case $d=0$.
\end{definition}

\begin{definition}{SARIMA$(p,d,q)\times(P,D,Q)_s$}{p5-sarima}
The full \textbf{seasonal ARIMA} folds both an ordinary difference of order $d$
and a seasonal difference of order $D$ in front of a SARMA:
\[
\Phi(B)\,\Phi^{(s)}(B)\,(1-B)^{d}(1-B^{s})^{D}\,X_t
=\Theta(B)\,\Theta^{(s)}(B)\,\eps_t .
\]
Here $d$ removes a polynomial trend, $D$ removes a (possibly stochastic)
seasonal pattern, and the four polynomials capture the remaining short-range and
seasonal dependence.
\end{definition}

\begin{definition}{Random walk and the IMA$(1,1)$/ARI$(1,1)$ models}{p5-rw}
The \textbf{random walk} is the AR(1) at the boundary $\phi=1$,
\[
X_t=X_{t-1}+\eps_t ,
\]
which is \emph{not} stationary, but whose first difference $\diff X_t=\eps_t$ is
white noise: it is an ARIMA$(0,1,0)$. Two minimal one-step-integrated models:
\begin{itemize}
\item \textbf{IMA$(1,1)$} $=$ ARIMA$(0,1,1)$: $Z_t=\diff Y_t=\eps_t-\theta\eps_{t-1}$,
      i.e.\ $Y_t=Y_{t-1}+\eps_t-\theta\eps_{t-1}$;
\item \textbf{ARI$(1,1)$} $=$ ARIMA$(1,1,0)$: $Z_t=\diff Y_t=\phi Z_{t-1}+\eps_t$
      with $\abs\phi<1$, i.e.\ $Y_t-Y_{t-1}=\phi(Y_{t-1}-Y_{t-2})+\eps_t$.
\end{itemize}
Neither $\{Y_t\}$ is stationary.
\end{definition}

\begin{intuition}
``Regress-and-subtract'' versus ``difference'' is the central practical choice.
Regression on $a+bt$ (or on $\sin/\cos$ harmonics) gives an interpretable,
deterministic trend but assumes you know its functional form. Differencing is
model-free and robust --- $\diff$ knocks out any linear drift, $\diff^{q}$ any
degree-$(q-1)$ polynomial, $\diff_{(s)}$ any period-$s$ season --- at the price
of slightly distorting the noise ($\diff Y_t$ replaces $Y_t$) and risking
\emph{over}-differencing (injecting a spurious unit root into the MA part).
\end{intuition}

% ----------------------------------------------------------------------
\subsection{Theorems \& Propositions}

\begin{proposition}{$\diff$ removes a linear trend; $\diff Y$ stays stationary}{p5-lintrend}
Let $X_t=a+bt+Y_t$ with $\{Y_t\}$ stationary. Then
\[
\diff X_t=X_t-X_{t-1}=b+\diff Y_t ,
\]
so the trend $a+bt$ is reduced to the constant $b$, and the random part
$\diff Y_t=Y_t-Y_{t-1}$ is again stationary. Differencing once more removes the
constant too: $\diff^{2}X_t=\diff^{2}Y_t=Y_t-2Y_{t-1}+Y_{t-2}$.\\[2pt]
\textit{Proof idea:} substitute $a+bt+Y_t$ and cancel:
$a+bt-(a+b(t-1))=b$. Linearity of $\diff$ gives $\diff Y_t$; a finite linear
combination of a stationary process is stationary, hence so is $\diff Y_t$.\\
\textit{Why:} the formal justification for first-differencing a series that
``looks linear''; it converts a non-stationary level into a stationary
increment, up to a known constant that the mean estimate then absorbs.
\end{proposition}

\begin{proposition}{$\diff^{q}$ annihilates a degree-$(q-1)$ polynomial trend}{p5-polytrend}
Let $X_t=\mu_t+Y_t$ with $\mu_t$ a polynomial of degree $q-1$ in $t$. Then
$\diff^{q}\mu_t\equiv 0$, so
\[
\diff^{q}X_t=\diff^{q}Y_t=\sum_{k=0}^{q}\binom{q}{k}(-1)^{k}Y_{t-k}.
\]
\textit{Proof idea:} $\diff$ lowers the degree of a polynomial in $t$ by exactly
one (the leading $t^{m}$ becomes $t^{m}-(t-1)^{m}=m\,t^{m-1}+\cdots$); applying
$\diff$ a total of $q$ times to a degree-$(q-1)$ polynomial drives it to the
zero sequence. Linearity then leaves $\diff^{q}Y_t$, expanded by the binomial
theorem. (See Exercise~\ref{exo:p5-arimainvar} for the careful induction that
$\diff^{k}t^{k-1}=0$.)\\
\textit{Why:} tells you exactly how many differences $d=q$ are needed to flatten
a polynomial trend of a given degree, which is how one chooses $d$ in
ARIMA$(p,d,q)$.
\end{proposition}

\begin{theorem}{Combined $\diff\diff_{(12)}$ yields a stationary process}{p5-combined}
Let $X_t=a+bt+s_t+Y_t$, where $\{Y_t\}$ is zero-mean stationary with ACVS
$\acvs_\tau$, $a,b$ are constants and $s_{t+12}=s_t$ has period $12$. Then
\[
\diff\diff_{(12)}X_t=\diff\diff_{(12)}Y_t
=Y_t-Y_{t-1}-Y_{t-12}+Y_{t-13}
\]
is weakly stationary, with mean $0$ and
\[
\begin{aligned}
\Cov\bigl(\diff\diff_{(12)}X_t,\ \diff\diff_{(12)}X_{t-\tau}\bigr)
=\;&\acvs_\tau-\acvs_{\tau+1}-\acvs_{\tau+12}+\acvs_{\tau+13}\\
&-\acvs_{\tau-1}+\acvs_{\tau}+\acvs_{\tau+11}-\acvs_{\tau+12}\\
&-\acvs_{\tau-12}+\acvs_{\tau-11}+\acvs_{\tau}-\acvs_{\tau+1}\\
&+\acvs_{\tau-13}-\acvs_{\tau-12}-\acvs_{\tau-1}+\acvs_{\tau} .
\end{aligned}
\]
\textit{Proof idea:} first $\diff_{(12)}X_t=12b+\diff_{(12)}Y_t$ kills $s_t$ and
turns $bt$ into the constant $12b$; then $\diff$ kills the constant, leaving the
fixed linear combination $Y_t-Y_{t-1}-Y_{t-12}+Y_{t-13}$. Its mean is $0$ and
its covariance, expanded by bilinearity, is a function of $\tau$ \emph{only}
(the four-by-four table above).\\
\textit{Why:} the theoretical backing for the standard ``$\diff_{12}$ then
$\diff$'' pre-processing of monthly data (e.g.\ the Mauna Loa $\mathrm{CO}_2$
series); it shows no $t$-dependence survives.
\end{theorem}

\begin{proposition}{Multiplicative-trend variant: $\diff_{(12)}^{2}$}{p5-multtrend}
If instead $X_t=(a+bt)s_t+Y_t$ with $s_{t+12}=s_t$, then
\[
\diff_{(12)}X_t=12b\,s_{t-12}+Y_t-Y_{t-12},
\qquad
\diff_{(12)}^{2}X_t=Y_t-2Y_{t-12}+Y_{t-24},
\]
which has zero mean and a covariance depending only on $\abs\tau$, hence is
stationary.\\[2pt]
\textit{Proof idea:} $\diff_{(12)}$ once turns $(a+bt)s_t$ into $12b\,s_{t-12}$
(using $s_t=s_{t-12}$); a second $\diff_{(12)}$ removes that periodic term as
well, leaving a fixed combination of $Y$'s.\\
\textit{Why:} warns that when trend and season \emph{multiply} (amplitude grows
with level), a \emph{single seasonal} double-difference --- not
$\diff\diff_{(12)}$ --- is what flattens the series.
\end{proposition}

\begin{proposition}{ACVS of the seasonal MA(1)}{p5-smaacvs}
For $X_t=\eps_t-\theta\eps_{t-s}$ with white-noise variance $\sigma_\eps^2$,
\[
\acvs_\tau=\sigma_\eps^2\Bigl[(1+\theta^{2})\delta_\tau-\theta(\delta_{\tau-s}+\delta_{\tau+s})\Bigr]
=\begin{cases}
(1+\theta^{2})\sigma_\eps^2, & \tau=0,\\
-\theta\,\sigma_\eps^2, & \tau=\pm s,\\
0, &\text{otherwise,}
\end{cases}
\]
so the ACF is $\acf_{\pm s}=-\theta/(1+\theta^{2})$ and zero at every other lag.\\[2pt]
\textit{Proof idea:} apply the white-noise filter rule
$\acvs_\tau=\sigma_\eps^2\sum_j a_j a_{j+\abs\tau}$ to the taps
$a_0=1,\ a_s=-\theta$; only $\tau\in\{0,\pm s\}$ pair surviving noise indices.\\
\textit{Why:} a single nonzero ACF spike at the seasonal lag $s$ is the
fingerprint that flags an $\mathrm{SMA}(1)_s$ component in a correlogram.
\end{proposition}

\begin{proposition}{Random walk is non-stationary; its difference is white noise}{p5-rwprop}
For the random walk $X_t=\sum_{j=1}^{t}\eps_j$ (with $X_0=0$),
\[
\E[X_t]=0,\qquad \Var(X_t)=t\,\sigma_\eps^2,\qquad
\Cov(X_t,X_{t+\tau})=t\,\sigma_\eps^2\ \ (\tau\ge0),
\]
all of which depend on $t$, so $\{X_t\}$ is \textbf{not} stationary; but
$\diff X_t=\eps_t$ is white noise.\\[2pt]
\textit{Proof idea:} the variance of a sum of $t$ uncorrelated unit pieces grows
linearly in $t$; the overlap of $X_t$ and $X_{t+\tau}$ is the first $t$ shared
terms.\\
\textit{Why:} the prototype unit-root process --- exactly the ``$\phi=1$
boundary'' that ARIMA differencing is designed to handle.
\end{proposition}

\begin{theorem}{Second-order growth of IMA$(1,1)$; correlation $\to 1$}{p5-imagrowth}
Let $Y_t=Y_{t-1}+\eps_t-\theta\eps_{t-1}$ with $Y_0=0$. Then for $t\ge1$
\[
Y_t=\eps_t+(1-\theta)\sum_{j=1}^{t-1}\eps_j-\theta\eps_0 ,
\]
hence
\[
\Var(Y_t)=\sigma_\eps^2\bigl[\,1+(1-\theta)^{2}(t-1)+\theta^{2}\,\bigr],
\]
and for $\tau>0$
\[
\Cov(Y_t,Y_{t+\tau})=\sigma_\eps^2\bigl[\,1-\theta+(1-\theta)^{2}(t-1)+\theta^{2}\,\bigr].
\]
As $\tau/t\to0$ the correlation $\Corr(Y_t,Y_{t+\tau})\to1$.\\[2pt]
\textit{Proof idea:} unroll the recursion from $Y_0=0$ to get the explicit
linear combination of $\eps_0,\dots,\eps_t$; read off $\Var$ and $\Cov$ from the
shared coefficients (the bulk $(1-\theta)$ terms dominate). Both grow like
$(1-\theta)^2(t-1)\sigma_\eps^2$, so their ratio $\to1$.\\
\textit{Why:} shows quantitatively that an integrated series has
\emph{variance growing linearly in $t$} and \emph{nearby values almost perfectly
correlated} --- the diagnostic signature of needing a difference.
\end{theorem}

\begin{theorem}{Second-order growth of ARI$(1,1)$; correlation near unity}{p5-arigrowth}
Let $Y_t-Y_{t-1}=\phi(Y_{t-1}-Y_{t-2})+\eps_t$ with $\abs\phi<1$ and $Y_0=0$.
Then
\[
Y_t=\sum_{j=1}^{t}\ \sum_{i=0}^{t-j}\phi^{i}\,\eps_j
=\sum_{j=1}^{t}\frac{1-\phi^{\,t-j+1}}{1-\phi}\,\eps_j ,
\]
so that
\[
\begin{aligned}
\Var(Y_t)&=\sigma_\eps^2\sum_{j=1}^{t}\!\left(\frac{1-\phi^{\,t-j+1}}{1-\phi}\right)^{2},\\[2pt]
\Cov(Y_t,Y_{t+\tau})&=\sigma_\eps^2\sum_{j=1}^{t}
\frac{1-\phi^{\,t-j+1}}{1-\phi}\cdot\frac{1-\phi^{\,t+\tau-j+1}}{1-\phi}
\qquad (\tau>0).
\end{aligned}
\]
Again the correlation between nearby times is close to unity.\\[2pt]
\textit{Proof idea:} the increments $Z_t=\diff Y_t$ form a stationary AR(1),
$Z_t=\sum_{i\ge0}\phi^{i}\eps_{t-i}$; integrating ($Y_t=\sum_{m\le t}Z_m$ from
$0$) and summing the inner geometric series gives the closed coefficients
$(1-\phi^{t-j+1})/(1-\phi)$.\\
\textit{Why:} confirms that even when the increments are a well-behaved
stationary AR(1), \emph{integrating} them recreates the random-walk-like
variance growth and near-unit correlation, so the level $\{Y_t\}$ must be
differenced before fitting.
\end{theorem}

\begin{pitfall}
$\diff_{(s)}=I-B^{s}$ is \emph{not} $\diff^{s}=(I-B)^{s}$. The seasonal
difference subtracts the value $s$ steps back ($X_t-X_{t-s}$, two terms); the
$s$-th ordinary difference is an $(s+1)$-term binomial combination. They remove
different things: $\diff_{(s)}$ a period-$s$ season, $\diff^{s}$ a degree-$(s-1)$
polynomial. Mixing them up is a classic exam error.
\end{pitfall}

% ----------------------------------------------------------------------
\subsection{Key Techniques}

\begin{keytech}{Choosing the differencing orders $d$ and $D$}{p5-choosed}
\begin{enumerate}
  \item \textbf{Trend.} If the level looks like a degree-$(q-1)$ polynomial,
        take $d=q$ ordinary differences (linear drift $\Rightarrow d=1$;
        quadratic $\Rightarrow d=2$). Symptoms of a remaining unit root: a
        slowly decaying sample ACF and near-linear variance growth.
  \item \textbf{Season.} If a fixed pattern of period $s$ persists, take one
        seasonal difference $\diff_{(s)}$ ($D=1$); rarely $D=2$.
  \item \textbf{Order.} Apply $\diff_{(s)}$ first to remove the season, then
        $\diff^{d}$; for monthly data with linear drift this is
        $\diff\diff_{(12)}$.
  \item \textbf{Stop} when the differenced series has a stable mean, constant
        variance and a quickly decaying ACF. \textbf{Do not over-difference}:
        an extra $\diff$ injects a spurious MA unit root
        ($\acf_1\to-\tfrac12$) and inflates variance.
\end{enumerate}
\end{keytech}

\begin{keytech}{Expanding a SARMA into explicit lag form}{p5-expandsarma}
Multiply the polynomials out, treating $B$ as a formal variable:
\begin{enumerate}
  \item Write the four factors $\Phi(B),\Phi^{(s)}(B)$ (left) and
        $\Theta(B),\Theta^{(s)}(B)$ (right).
  \item Expand each product, e.g.
        $(1-\theta B)(1-\theta^{(12)}B^{12})
        =1-\theta B-\theta^{(12)}B^{12}+\theta\theta^{(12)}B^{13}$.
  \item Read off the surviving lags. The seasonal factor only contributes lags
        that are multiples of $s$; cross terms give lags like $s\pm1$.
  \item Translate $B^{k}X_t=X_{t-k}$, $B^{k}\eps_t=\eps_{t-k}$ to get the
        difference equation, then solve for $X_t$ (move all but $X_t$ to the
        right) for a one-step recursion.
\end{enumerate}
\end{keytech}

\begin{keytech}{Deriving the MA$(\infty)$ ($\psi$-weights) of a SARMA/ARMA}{p5-psiweights}
To get $X_t=\sum_{j\ge0}\psi_j\eps_{t-j}$ from $\Phi(B)X_t=\Theta(B)\eps_t$:
\begin{enumerate}
  \item Invert the AR factor as a geometric series:
        $\dfrac{1}{1-\phi B}=\sum_{l\ge0}\phi^{l}B^{l}$ (valid when
        $\abs\phi<1$).
  \item Multiply by the full MA numerator and collect the coefficient of
        $B^{j}$.
  \item \textbf{Match powers of $B$}: $\Phi(B)\sum_j\psi_jB^{j}=\Theta(B)$ gives
        the recursion
        $\psi_j=\theta^{\text{(num)}}_j+\sum_{i=1}^{p}\phi_i\psi_{j-i}$,
        $\psi_0=1$, with $\theta^{\text{(num)}}_j$ the coefficient of $B^{j}$ in
        the (expanded) MA numerator. Beyond the highest numerator lag the pure
        recursion $\psi_j=\sum_i\phi_i\psi_{j-i}$ takes over.
\end{enumerate}
Spot the recursion early (e.g.\ $\psi_3=\phi\psi_2$) --- it collapses the
algebra.
\end{keytech}

\begin{keytech}{ACVS/ACF of a seasonal MA$(\infty)$ via even/odd lags}{p5-acvsseasonal}
For a process with only even (or only seasonal) taps,
$X_t=\sum_j a_j\eps_{t-j}$:
\begin{enumerate}
  \item Use $\acvs_\tau=\sigma_\eps^2\sum_{j\ge0}a_j a_{j+\tau}$ and
        $\acf_\tau=\acvs_\tau/\acvs_0$.
  \item If the taps vanish at odd lags ($a_{\text{odd}}=0$), then so do
        $\acvs_\tau$ and $\acf_\tau$ at odd $\tau$ --- only even (seasonal) lags
        carry correlation.
  \item Sum the resulting geometric series in $\phi^{2}$ for the closed form;
        $\Var(X_t)=\acvs_0=\sigma_\eps^2\sum_j a_j^2$.
\end{enumerate}
\end{keytech}

% ----------------------------------------------------------------------
\subsection{Exercises}

\begin{exercise}{ARIMA invariance to an added polynomial \quad\textnormal{[Sheet 5]}}{p5-arimainvar}
Suppose $\{X_t\}$ is an ARIMA$(p,d,q)$ process satisfying
$\Phi(B)(I-B)^{d}X_t=\Theta(B)Z_t$ with $Z_t$ white noise. Show that the same
equations are satisfied by $W_t=X_t+A_0+A_1t+\cdots+A_{d-1}t^{d-1}$, where the
$A_j$ are independent random coefficients.
\end{exercise}
\begin{solution}
Because $\Phi(B)$ and $\Theta(B)$ act on whatever $(I-B)^{d}X_t$ produces, it
suffices to show that the added polynomial is annihilated by $(I-B)^{d}$, i.e.\
$(I-B)^{d}W_t=(I-B)^{d}X_t$.

First, $(I-B)\,\mathbf 0=\mathbf 0$ on the zero sequence. By linearity of
$(I-B)^{d}$,
\[
(I-B)^{d}W_t=(I-B)^{d}X_t
+\sum_{k=1}^{d}A_{k-1}\,(I-B)^{d}t^{\,k-1}.
\]
So we only need $(I-B)^{d}t^{\,k-1}=\mathbf 0$ for $k=1,\dots,d$; it is enough to
prove the stronger statement
\[
(I-B)^{k}\,t^{\,k-1}=\mathbf 0\qquad\text{for all }k\in\Nat ,
\]
by induction on $k$.

\emph{Base case $k=1$:} $(I-B)t^{0}=(I-B)\mathbf 1=\mathbf 1-\mathbf 1=\mathbf 0$
(here $\mathbf 1$ is the constant-one sequence).

\emph{Inductive step:} assume $(I-B)^{r}t^{\,r-1}=\mathbf 0$ for all $r\le k$.
Then
\[
(I-B)^{k+1}t^{\,k}=(I-B)^{k}(I-B)t^{\,k}
=(I-B)^{k}\bigl(t^{\,k}-(t-1)^{k}\bigr).
\]
By the binomial theorem the top power cancels,
\[
t^{\,k}-(t-1)^{k}
=-\sum_{r=0}^{k-1}\binom{k}{r}(-1)^{\,k-r}t^{\,r},
\]
a polynomial of degree $k-1$. Applying $(I-B)^{k}$ and using linearity together
with the induction hypothesis (each $(I-B)^{r+1}t^{\,r}=\mathbf 0$ for
$r\le k-1$, and one extra difference of $\mathbf 0$ is still $\mathbf 0$), every
term vanishes, so $(I-B)^{k+1}t^{\,k}=\mathbf 0$.

By induction the claim holds for all $k$, hence $(I-B)^{d}W_t=(I-B)^{d}X_t$ and
$\{W_t\}$ satisfies the identical ARIMA equations. \emph{Interpretation:} the
$d$-fold integration in an ARIMA$(p,d,q)$ leaves the degree-$(d-1)$ polynomial
``initial conditions'' completely undetermined by the model.
\end{solution}

\begin{exercise}{$\psi$-weights of a SARMA$(1,2)\times(0,1)_4$ \quad\textnormal{[Sheet 5]}}{p5-sarmapsi}
$U_t$ follows a $\mathrm{SARMA}(1,2)\times(0,1)_4$ model. Write the defining
equation and find the weights $\psi_j$ in the MA$(\infty)$ representation
$U_t=\sum_{j\ge0}\psi_j\eps_{t-j}$.
\end{exercise}
\begin{solution}
With one seasonal MA factor at $s=4$ the model is
\[
(1-\phi B)\,U_t=(1-\theta_1B-\theta_2B^{2})(1-\Theta B^{4})\,\eps_t
\;\Longleftrightarrow\;
U_t=\phi U_{t-1}+(1-\theta_1B-\theta_2B^{2})(1-\Theta B^{4})\eps_t .
\]
Invert the AR factor as a geometric series, $1/(1-\phi B)=\sum_{l\ge0}\phi^{l}B^{l}$:
\[
\begin{aligned}
U_t&=\sum_{l\ge0}(\phi B)^{l}(1-\theta_1B-\theta_2B^{2})(1-\Theta B^{4})\eps_t\\
&=\sum_{l\ge0}\bigl(\phi^{l}B^{l}-\theta_1\phi^{l}B^{l+1}-\theta_2\phi^{l}B^{l+2}
-\Theta\{\phi^{l}B^{l+4}-\theta_1\phi^{l}B^{l+5}-\theta_2\phi^{l}B^{l+6}\}\bigr)\eps_t .
\end{aligned}
\]
Collecting the coefficient of $B^{j}$ gives $\psi_j$. For the low lags
(before the seasonal term enters):
\[
\psi_0=1,\quad
\psi_1=-\theta_1+\phi,\quad
\psi_2=-\theta_2-\theta_1\phi+\phi^{2},\quad
\psi_3=-\theta_2\phi-\theta_1\phi^{2}+\phi^{3}=\phi\,\psi_2 .
\]
From $j=4$ the seasonal factor contributes:
\[
\begin{aligned}
\psi_4&=-\Theta-\theta_2\phi^{2}-\theta_1\phi^{3}+\phi^{4}=-\Theta+\phi\,\psi_3,\\
\psi_5&=-\Theta(\phi-\theta_1)-\theta_2\phi^{3}-\theta_1\phi^{4}+\phi^{5},\\
\psi_6&=-\Theta(\phi^{2}-\theta_1\phi-\theta_2)+(\phi^{6}-\theta_1\phi^{5}-\theta_2\phi^{4}),\\
\psi_7&=-\Theta(\phi^{3}-\theta_1\phi^{2}-\theta_2\phi)+(\phi^{7}-\theta_1\phi^{6}-\theta_2\phi^{5})=\phi\,\psi_6 .
\end{aligned}
\]
The pattern stabilises: for all $l\ge6$,
\[
\boxed{\;\psi_l=\bigl(\phi^{l}-\theta_1\phi^{\,l-1}-\theta_2\phi^{\,l-2}\bigr)
-\Theta\bigl(\phi^{\,l-4}-\theta_1\phi^{\,l-5}-\theta_2\phi^{\,l-6}\bigr).\;}
\]
(The recursion $\psi_l=\phi\psi_{l-1}$ holds once $l$ is past every fixed
numerator lag, here $l\ge7$, which is the labour-saving shortcut.)
\end{solution}

\begin{exercise}{Seasonal ARMA $(1-0.7B^{2})X=(1-0.3B^{2})\eps$ \quad\textnormal{[Sheet 5]}}{p5-seasarma}
Study $(1-0.7B^{2})X_t=(1-0.3B^{2})\eps_t$, $\eps_t$ white noise with unit
variance.
(1) Find $\{a_j\}$ in $X_t=\sum_{j\ge0}a_j\eps_{t-j}$.
(2) Find $\{b_j\}$ in $\eps_t=\sum_{j\ge0}b_jX_{t-j}$.
(3) Find and describe the autocorrelation $\acf_\tau$.
\end{exercise}
\begin{solution}
Write $\phi=0.7$, $\theta=0.3$, so $(1-\phi B^{2})X_t=(1-\theta B^{2})\eps_t$.
The AR polynomial $\Phi(z)=1-\phi z^{2}$ has roots $z=\pm\phi^{-1/2}$ with
modulus $>1$, so the process is stationary (and causal) and we may invert.

\textbf{(1) MA$(\infty)$.} Geometric expansion in $B^{2}$:
\[
X_t=\frac{1-\theta B^{2}}{1-\phi B^{2}}\eps_t
=\sum_{k\ge0}(\phi B^{2})^{k}(1-\theta B^{2})\eps_t .
\]
The coefficient of $B^{2k}$ is $\phi^{k}$ and of $B^{2k+2}$ is $-\phi^{k}\theta$,
so only even lags survive:
\[
a_0=1,\quad a_1=0,\quad a_2=\phi-\theta,\quad a_3=0,\quad a_4=\phi^{2}-\phi\theta,
\dots
\]
and in general $a_j=0$ for odd $j$, while for even $j\ge2$
\[
\boxed{\,a_j=\phi^{\,j/2}-\phi^{\,(j-2)/2}\theta\,}
=\phi^{\,(j-2)/2}(\phi-\theta).
\]

\textbf{(2) AR$(\infty)$ (invertibility).} The MA polynomial
$1-\theta z^{2}$ has roots $\pm\theta^{-1/2}$, modulus $>1$, so the process is
invertible. By symmetry, swap the roles of $\phi$ and $\theta$:
\[
\eps_t=\frac{1-\phi B^{2}}{1-\theta B^{2}}X_t=\sum_{k\ge0}(\theta B^{2})^{k}(1-\phi B^{2})X_t,
\]
hence $b_0=1$, $b_j=0$ for odd $j$, and for even $j\ge2$
\[
\boxed{\,b_j=\theta^{\,j/2}-\theta^{\,(j-2)/2}\phi\,.}
\]

\textbf{(3) Autocorrelation.} Using
$\acvs_\tau=\sigma_\eps^2\sum_{j\ge0}a_j a_{j+\tau}$ (white noise) and
$\acf_\tau=\acvs_\tau/\acvs_0$. Since the $a_j$ vanish at odd $j$, $\acf_\tau=0$
for all odd $\tau$, and $\acf_0=1$.

First the variance. With $a_0=1$ and $a_{2m}=\phi^{m-1}(\phi-\theta)$ for
$m\ge1$,
\[
\Var(X_t)=\acvs_0=\sigma_\eps^2\Bigl(1+(\phi-\theta)^{2}\sum_{m\ge1}\phi^{2(m-1)}\Bigr)
=\sigma_\eps^2\Bigl(1+\frac{(\phi-\theta)^{2}}{1-\phi^{2}}\Bigr).
\]
Writing $(\phi-\theta)^{2}=\phi^{2}(1-2\phi^{-1}\theta+\phi^{-2}\theta^{2})$
recovers the notes' form
\[
\Var(X_t)=\sigma_\eps^2\Bigl\{1+(1-2\phi^{-1}\theta+\phi^{-2}\theta^{2})\tfrac{\phi^{2}}{1-\phi^{2}}\Bigr\}.
\]
For even $\tau=2m\ (m\ge1)$,
\[
\acvs_\tau=\sigma_\eps^2\Bigl[\,a_0 a_\tau+\sum_{k\ge1}a_{2k}a_{2k+\tau}\Bigr]
=\sigma_\eps^2\Bigl[(\phi^{m}-\phi^{m-1}\theta)+(\phi-\theta)^{2}\phi^{m}\sum_{k\ge1}\phi^{2(k-1)}\Bigr],
\]
where the cross term used
$a_{2k}a_{2k+2m}=\phi^{k-1}(\phi-\theta)\cdot\phi^{k+m-1}(\phi-\theta)$.
Summing the geometric series and dividing by $\acvs_0$ (with $m=\tfrac\tau2$),
\[
\begin{aligned}
\acf_\tau
&=\frac{\phi^{m}-\phi^{m-1}\theta+(\phi-\theta)^{2}\dfrac{\phi^{m}}{1-\phi^{2}}}
{\,1+\dfrac{(\phi-\theta)^{2}}{1-\phi^{2}}\,}\\[4pt]
&=\phi^{\,m}\cdot
\frac{1-\phi^{-1}\theta+(1-2\phi^{-1}\theta+\phi^{-2}\theta^{2})\dfrac{\phi^{2}}{1-\phi^{2}}}
{\,1+(1-2\phi^{-1}\theta+\phi^{-2}\theta^{2})\dfrac{\phi^{2}}{1-\phi^{2}}\,}.
\end{aligned}
\]
With $\phi=0.7,\ \theta=0.3$: $\Var(X_t)=1+\tfrac{0.16}{0.51}\approx1.3137$, and
$\acf_2\approx0.472$, $\acf_4\approx0.330$, $\acf_6\approx0.231$,
$\acf_8\approx0.162$ --- a clean \emph{geometric decay at rate $\phi$ along the
even lags}, with all odd-lag correlations zero. (Numerically verified.) The
correlogram therefore shows spikes only at $\tau=2,4,6,\dots$ decaying like
$0.7^{\tau/2}$.
\end{solution}

\begin{exercise}{Order of $\diff\diff_{(s)}$ on trend-plus-season \quad\textnormal{[New]}}{p5-newdiff}
Let $X_t=a+bt+s_t+Y_t$ with $s_{t+4}=s_t$ (quarterly season) and $\{Y_t\}$
zero-mean stationary. Compute $\diff\diff_{(4)}X_t$ explicitly and verify it is a
fixed linear combination of $\{Y_t\}$ with zero mean.
\end{exercise}
\begin{solution}
Apply the seasonal difference first:
\[
\diff_{(4)}X_t=X_t-X_{t-4}
=\bigl(a+bt+s_t+Y_t\bigr)-\bigl(a+b(t-4)+s_{t-4}+Y_{t-4}\bigr)
=4b+(Y_t-Y_{t-4}),
\]
using $s_t=s_{t-4}$. The trend has collapsed to the constant $4b$ and the season
is gone. Now apply $\diff=I-B$:
\[
\diff\diff_{(4)}X_t=\bigl(4b+Y_t-Y_{t-4}\bigr)-\bigl(4b+Y_{t-1}-Y_{t-5}\bigr)
=Y_t-Y_{t-1}-Y_{t-4}+Y_{t-5}.
\]
This is a fixed linear combination of the stationary $\{Y_t\}$, so its mean is
$\E[Y_t-Y_{t-1}-Y_{t-4}+Y_{t-5}]=0$ and, by bilinearity, its covariance depends
on the lag only; hence $\diff\diff_{(4)}X_t$ is weakly stationary. (Same shape as
the $\diff\diff_{(12)}$ result with $12\to4$.)
\end{solution}

\begin{exercise}{Expand and recurse a SARMA$(0,1)\times(1,0)_{12}$ \quad\textnormal{[New]}}{p5-newsarma}
Write the explicit one-step recursion for a $\mathrm{SARMA}(0,1)\times(1,0)_{12}$
model $(1-\phi^{(12)}B^{12})X_t=(1-\theta B)\eps_t$, give its MA$(\infty)$
weights, and find its ACVS at lags $0,1,11,12,13$.
\end{exercise}
\begin{solution}
Expanding the seasonal AR factor and solving for $X_t$:
\[
X_t=\phi^{(12)}X_{t-12}+\eps_t-\theta\eps_{t-1}.
\]
This is an AR(1) acting only on the seasonal lag $12$, driven by an MA(1) input.
Its MA$(\infty)$ comes from inverting
$1/(1-\phi^{(12)}B^{12})=\sum_{l\ge0}(\phi^{(12)})^{l}B^{12l}$:
\[
X_t=\sum_{l\ge0}(\phi^{(12)})^{l}\bigl(\eps_{t-12l}-\theta\eps_{t-12l-1}\bigr),
\qquad
\psi_{12l}=(\phi^{(12)})^{l},\quad \psi_{12l+1}=-\theta(\phi^{(12)})^{l},
\]
all other $\psi_j=0$. By the white-noise filter rule
$\acvs_\tau=\sigma_\eps^2\sum_{j\ge0}\psi_j\psi_{j+\tau}$:
\[
\begin{aligned}
\acvs_0&=\sigma_\eps^2\,\frac{1+\theta^{2}}{1-(\phi^{(12)})^{2}},
&\qquad
\acvs_1&=\sigma_\eps^2\sum_{l\ge0}\psi_{12l}\psi_{12l+1}
=\frac{-\theta\,\sigma_\eps^2}{1-(\phi^{(12)})^{2}},\\
\acvs_{11}&=\sigma_\eps^2\sum_{l\ge0}\psi_{12l+1}\psi_{12(l+1)}
=\frac{-\theta\,\phi^{(12)}\sigma_\eps^2}{1-(\phi^{(12)})^{2}},
&\qquad
\acvs_{12}&=\sigma_\eps^2\sum_{l\ge0}\psi_{12l}\psi_{12l+12}
=\phi^{(12)}\acvs_0,\\
\acvs_{13}&=\sigma_\eps^2\sum_{l\ge0}\psi_{12l}\psi_{12(l+1)+1}
=\frac{-\theta\,\phi^{(12)}\sigma_\eps^2}{1-(\phi^{(12)})^{2}}=\acvs_{11}.
\end{aligned}
\]
So correlation appears in clusters around the seasonal lags $0,12,24,\dots$
(each carrying MA(1) side-lobes at $\pm1$), decaying geometrically in
$\phi^{(12)}$ --- the textbook shape of a seasonal AR with a non-seasonal MA.
\end{solution}

\begin{exercise}{IMA$(1,1)$ variance growth and the over-differencing trap \quad\textnormal{[New]}}{p5-newima}
For $Y_t=Y_{t-1}+\eps_t-\theta\eps_{t-1}$ with $Y_0=0$: (a) confirm
$\Var(Y_t)$ grows linearly in $t$ for $\theta\ne1$; (b) what happens at the
boundary $\theta=1$, and why is that the over-differencing warning?
\end{exercise}
\begin{solution}
\textbf{(a)} From $Y_t=\eps_t+(1-\theta)\sum_{j=1}^{t-1}\eps_j-\theta\eps_0$ and
uncorrelated noise,
\[
\Var(Y_t)=\sigma_\eps^2\bigl[1+(1-\theta)^{2}(t-1)+\theta^{2}\bigr],
\]
which is \emph{linear} in $t$ with slope $(1-\theta)^{2}\sigma_\eps^2>0$ whenever
$\theta\ne1$: the series wanders like a (damped) random walk, the hallmark of an
integrated process. The covariance
$\Cov(Y_t,Y_{t+\tau})=\sigma_\eps^2[1-\theta+(1-\theta)^{2}(t-1)+\theta^{2}]$
shares the same leading $(1-\theta)^{2}(t-1)$ term, so the correlation of nearby
points $\to1$.

\textbf{(b)} At $\theta=1$ the slope vanishes and
$\Var(Y_t)=\sigma_\eps^2[1+1]=2\sigma_\eps^2$ is \emph{bounded}. In fact
$Y_t=Y_{t-1}+\eps_t-\eps_{t-1}$ telescopes to $Y_t=\eps_t-\eps_0$, so
$\diff Y_t=\eps_t-\eps_{t-1}$ is a (non-invertible) MA(1) with an MA unit root
and lag-1 ACF $\acf_1=-\theta/(1+\theta^{2})=-1/2$. This is exactly the
signature of \emph{over-differencing}: if you difference a series that was
already stationary, you create a $(1-B)$ factor in the MA polynomial, i.e.\ a
non-invertible MA(1) with $\acf_1=-\tfrac12$. Seeing $\acf_1\approx-\tfrac12$
after differencing is the cue to difference one time fewer.
\end{solution}

% ======================================================================
%  PART 6 -- Fourier Theory, Spectral Density and the Periodogram
% ======================================================================
\section{Fourier Theory, Spectral Density and the Periodogram}
\label{sec:spectral}

The frequency domain is the second face of a stationary process. Every
second-order property carried by the autocovariance sequence (ACVS)
$\{\acvs_\tau\}$ in the time domain is mirrored, via the Fourier transform, by
the \textbf{spectral density function} (SDF) $\sdf(f)$ in the frequency domain.
This part builds that bridge in three stages. First, pure Fourier theory:
continuous-time, discrete-time and finite-data transforms, the convolution
theorem, aliasing and the Nyquist frequency. Second, the spectral
representation of a stationary process: $\sdf(f)=\sum_\tau\acvs_\tau
e^{-2\pi i f\tau}$, Herglotz's theorem, the integrated spectrum and its
Lebesgue decomposition, and the spectral representation theorem with its
orthogonal-increment process $\{Z(f)\}$. Third, estimation: the periodogram
$\hat\sdf^{(p)}(f)=\abs{J(f)}^2$, why it is (asymptotically) unbiased but
\emph{inconsistent}, and how tapering cures bias while multitapering cures
variance. Throughout we take the sampling interval $\Delta=1$ once we reach
the spectral chapters, and we always work with ordinary frequency $f$ (not
angular frequency $\omega=2\pi f$).

% ======================================================================
\subsection{Definitions}
% ======================================================================

\begin{definition}{Continuous-time Fourier transform and its inverse}{p6-ctft}
For $g:\Reals\to\Comp$ with $g\in L^1$, the \textbf{Fourier transform}
$G$ (upper case denotes the transform; some authors write $\hat g$) is
\[
G(f)=\int_{-\infty}^{\infty} g(t)\,e^{-2\pi i f t}\,dt,\qquad f\in\Reals .
\]
If in addition $G\in L^1$, the \textbf{inverse transform} recovers $g$ for
almost all $t$:
\[
g(t)=\int_{-\infty}^{\infty} G(f)\,e^{2\pi i f t}\,df .
\]
We write $g\longleftrightarrow G$ and call $(g,G)$ a \textbf{Fourier pair}.
\end{definition}

\begin{definition}{Properties of the continuous-time transform}{p6-ctprops}
For $g,h\in L^1$ with transforms $G,H$ and scalar $\alpha$, the basic pairs are
\[
\begin{aligned}
\conj{g(t)} &\longleftrightarrow \conj{G(-f)}, &\qquad
g(\alpha t)\;(\alpha\ne0) &\longleftrightarrow \tfrac{1}{\abs\alpha}G(f/\alpha),\\
g(t+\tau) &\longleftrightarrow G(f)\,e^{2\pi i \tau f}, &\qquad
\alpha g(t)+h(t) &\longleftrightarrow \alpha G(f)+H(f).
\end{aligned}
\]
(Conjugation reflects, scaling in time inversely scales frequency, a time
shift multiplies by a phase, and the transform is linear.)
\end{definition}

\begin{definition}{Continuous-time convolution}{p6-ctconv}
For $g,h\in L^1$, the \textbf{convolution} $u=g*h$ is
\[
u(t)=\int_{-\infty}^{\infty} g(t-u)\,h(u)\,du,\qquad t\in\Reals .
\]
Convolution at a point is a weighted average of one function by the other,
centred at that point.
\end{definition}

\begin{definition}{Discrete-time Fourier transform and its inverse}{p6-dtft}
For a sequence $\{g_t\}_{t\in\Delta\Ints}$ in $\ell^1$, the
\textbf{discrete-time Fourier transform} is
\[
G(f)=\Delta\sum_{t\in\Delta\Ints} g_t\,e^{-2\pi i t f},\qquad f\in\Reals,
\]
with inverse (a Fourier series with the roles of domains swapped)
\[
g_t=\int_{-1/(2\Delta)}^{1/(2\Delta)} G(f)\,e^{2\pi i t f}\,df,\qquad t\in\Delta\Ints .
\]
A defining feature is \textbf{periodicity}: $G(f)=G(f+k/\Delta)$ for every
$k\in\Ints$, so $G$ is determined by its values on one period of length
$1/\Delta$.
\end{definition}

\begin{definition}{Discrete-time and periodic convolution}{p6-dtconv}
For $g,h\in\ell^1$ the \textbf{discrete convolution} $u=g*h$ is
\[
u_t=\Delta\sum_{u\in\Delta\Ints} g_{t-u}\,h_u,\qquad t\in\Delta\Ints,
\]
(with $\Delta=1$ the prefactor disappears). The dual object, the
\textbf{periodic convolution} $U=G*H$ of two periodic spectra, is
\[
U(f)=\int_{-1/(2\Delta)}^{1/(2\Delta)} G(f-f')\,H(f')\,df' .
\]
\end{definition}

\begin{definition}{Finite-data Fourier transform and the Dirichlet kernel}{p6-fdft}
Given observations $g_0,\dots,g_{(n-1)\Delta}$ at times
$T=\{0,\Delta,\dots,(n-1)\Delta\}$, the \textbf{finite Fourier transform} is
\[
G_n(f)=\sum_{t\in T} g_t\,e^{-2\pi i t f},\qquad f\in\Reals .
\]
It equals the full transform of $g_t$ multiplied by the indicator
$d_t=\tfrac1\Delta\ind_T(t)$, whose own transform is the \textbf{Dirichlet
kernel}
\[
\Dir_T(f)=\sum_{t\in T} e^{-2\pi i f t}
=\frac{\sin(n\pi f\Delta)}{\sin(\pi f\Delta)}\,e^{-\pi i (n-1) f\Delta}.
\]
Thus $G_n=\Dir_T * G$: observing a finite window \emph{blurs} the true
transform.
\end{definition}

\begin{definition}{Fourier frequencies}{p6-fourierfreq}
The FFT evaluates $G_n$ at the \textbf{Fourier frequencies}
\[
f_k=\frac{k}{\Delta n},\qquad k\in\Ints,\quad -\frac{n}{2}\le k\le \frac{n-1}{2}.
\]
With $\Delta=1$ these are $f_k=k/n$. They are equally spaced over one period
and are the frequencies at which distinct periodogram ordinates are
(approximately) uncorrelated.
\end{definition}

\begin{definition}{Spectral density function (discrete and continuous)}{p6-sdf}
For a discrete-time stationary process $\{X_t\}_{t\in\Ints}$ with ACVS
$\{\acvs_\tau\}\in\ell^1$, the \textbf{spectral density function} is the
discrete-time Fourier transform of the ACVS (with $\Delta=1$):
\[
\sdf(f)=\sum_{\tau\in\Ints}\acvs_\tau\,e^{-2\pi i f\tau},\qquad f\in\Reals .
\]
For a continuous-time process with $\acvs\in L^1$,
$\sdf(f)=\int_{-\infty}^{\infty}\acvs(\tau)e^{-2\pi i f\tau}\,d\tau$. By Fourier
inversion the ACVS is recovered as
$\acvs_\tau=\int_{-1/2}^{1/2}\sdf(f)e^{2\pi i f\tau}\,df$; in particular
$\acvs_0=\Var(X_t)=\int_{-1/2}^{1/2}\sdf(f)\,df$, so the SDF distributes the
total ``power'' $\Var(X_t)$ across frequencies. A valid SDF is real, even
($\sdf(-f)=\sdf(f)$, since $\acvs_{-\tau}=\acvs_\tau$), non-negative, and
periodic with period $1$.
\end{definition}

\begin{definition}{Harmonic (sinusoidal) process}{p6-harmonic}
Let $A>0$ and $\Theta\sim\mathrm{Unif}(-\pi,\pi)$ be independent. The
\textbf{harmonic process} is
\[
X_t=A\cos(\nu t+\Theta),\qquad t\in\Ints,
\]
with $\nu\in\Reals$ the oscillation frequency. Then $\E[X_t]=0$ and
\[
\Cov(X_{t+\tau},X_t)=\tfrac12\,\E[A^2]\cos(\nu\tau),
\]
so it is stationary but its ACVS does \emph{not} decay; in particular
$\{\acvs_\tau\}\notin\ell^1$ and no ordinary SDF exists. Its spectrum is a
pair of lines (jumps) at $\pm\nu/(2\pi)$.
\end{definition}

\begin{definition}{Integrated spectrum (spectral distribution function)}{p6-integspec}
By Herglotz's theorem, every ACVS admits the representation
\[
\acvs_\tau=\int_{-1/2}^{1/2} e^{2\pi i \tau f}\,dS^{(I)}(f),
\]
where $S^{(I)}$, the \textbf{integrated spectrum}, is right-continuous,
bounded, non-decreasing, with $S^{(I)}(-\tfrac12)=0$ and
$S^{(I)}(\tfrac12)=\sigma^2=\acvs_0$. It behaves like a (scaled) cumulative
distribution function over frequency. When $\{\acvs_\tau\}\in\ell^1$ there is a
density $\sdf$ with $S^{(I)}(f)=\int_{-1/2}^{f}\sdf(\lambda)\,d\lambda$, i.e.\
$\sdf=dS^{(I)}/df$, recovering Definition~\ref{def:p6-sdf}.
\end{definition}

\begin{definition}{Periodogram}{p6-periodogram}
For data $X_1,\dots,X_n$ ($T=\{1,\dots,n\}$, $\Delta=1$), the
\textbf{periodogram} is the plug-in estimator obtained by Fourier transforming
the biased sample ACVS:
\[
\hat\sdf^{(p)}(f)=\sum_{\tau\in\Ints}\hat\acvs_\tau\,e^{-2\pi i f\tau},
\qquad
\hat\acvs_\tau=\frac1n\sum_{t=1}^{n-\abs\tau}(X_t-\bar X)(X_{t+\abs\tau}-\bar X)
\]
for $\abs\tau\le n-1$ ($0$ otherwise). Equivalently (Proposition
\ref{prop:p6-periogramJ}) $\hat\sdf^{(p)}(f)=\abs{J(f)}^2$ with
$J(f)=\sqrt{1/n}\sum_{t\in T}(X_t-\bar X)e^{-2\pi i t f}$.
\end{definition}

\begin{definition}{Data taper and tapered periodogram}{p6-taper}
A \textbf{data taper} is a sequence $\{h_t\}_{t\in T}$ with
$\norm{h}_2^2=\sum_t h_t^2=1$. The \textbf{tapered Fourier transform} and
\textbf{tapered periodogram} are
\[
J_h(f)=\sum_{t\in T} h_t\,(X_t-\bar X)\,e^{-2\pi i t f},
\qquad
\hat\sdf^{(p)}_h(f)=\abs{J_h(f)}^2 .
\]
Choosing $h_t=\sqrt{1/n}$ (the constant taper) recovers the ordinary
periodogram. The normalisation $\norm{h}_2^2=1$ keeps the estimator unbiased
for white noise (Exercise~\ref{exo:p6-tunbias}).
\end{definition}

\begin{definition}{Multitaper spectral estimator}{p6-multitaper}
Given $K$ \textbf{orthonormal} tapers $h_k=\{h_{t,k}\}_{t\in T}$ with
$\sum_{t\in T} h_{t,k}h_{t,k'}=\delta_{k,k'}$, the \textbf{multitaper
estimator} averages the individual tapered periodograms,
\[
\hat\sdf^{(mt)}(f)=\frac1K\sum_{k=1}^{K}\hat\sdf^{(p)}_{h_k}(f).
\]
Its expectation is a convolution with the \textbf{average kernel}
$\overline H(f)=\tfrac1K\sum_{k=1}^K\abs{H_k(f)}^2$, and its variance is
$\approx\sdf(f)^2/K$.
\end{definition}

% ======================================================================
\subsection{Theorems \& Propositions}
% ======================================================================

\begin{theorem}{Fourier inversion (continuous and discrete time)}{p6-inversion}
\emph{(Continuous time.)} If $g,G\in L^1$ then the inverse transform recovers
$g$ for almost all $t$:
$g(t)=\int_{-\infty}^{\infty} G(f)\,e^{2\pi i f t}\,df$.
\emph{(Discrete time.)} If $\{g_t\}\in\ell^1$ then $G$ is recovered on one
period as a Fourier series with the domains swapped:
\[
g_t=\int_{-1/(2\Delta)}^{1/(2\Delta)} G(f)\,e^{2\pi i t f}\,df,\qquad
t\in\Delta\Ints .
\]
\textit{Proof idea:} substitute the forward transform into the right-hand side
and use the Fourier orthogonality relation
($\int e^{2\pi i f(t-s)}\,df$ acts as a delta), the integrability hypotheses
justifying the interchange of integration/summation.\\
\textit{Why:} inversion is what makes $(g,G)$ a genuine \emph{Fourier pair} ---
no information is lost. It is the engine behind recovering the ACVS from the SDF
($\acvs_\tau=\int_{-1/2}^{1/2}\sdf(f)e^{2\pi i f\tau}\,df$) and behind every
``time $\leftrightarrow$ frequency'' duality used below.
\end{theorem}

\begin{theorem}{Convolution theorem (continuous time)}{p6-convthm}
For $g,h\in L^1$, if $u=g*h$ then $U=G\cdot H$: the transform of a convolution
is the product of the transforms. If also $G,H\in L^1$ and $v=g\cdot h$, then
dually $V=G*H$.\\[2pt]
\textit{Proof idea:} substitute the convolution integral into the definition of
$U$, swap the order of integration (Fubini, justified by
$\iint\abs{g(x-y)h(y)}\,dx\,dy=\norm g_1\norm h_1<\infty$), and recognise the
shift factor $e^{-2\pi i s f}$ as producing $G(f)$.\\
\textit{Why:} the central computational identity of Fourier analysis ---
``multiplication in one domain is convolution in the other''. It underlies
filtering, the SDF of a filtered process, and the periodogram bias formula.
\end{theorem}

\begin{theorem}{Convolution theorem (discrete time)}{p6-dconvthm}
For $g,h\in\ell^1$ both pairs hold:
\[
h\cdot g \longleftrightarrow H*G,\qquad
h*g \longleftrightarrow H\cdot G .
\]
\textit{Proof idea:} for $u=h*g$, expand $U(f)=\sum_t u_t e^{-2\pi i t f}$,
insert $e^{-2\pi i (t-s)f}e^{-2\pi i s f}$, factor the double sum into
$\big(\sum_s h_s e^{-2\pi i s f}\big)\big(\sum_x g_x e^{-2\pi i x f}\big)$;
interchange is legal because $g,h\in\ell^1$. For $v=h\cdot g$, replace one
factor by its inverse transform and integrate.\\
\textit{Why:} the discrete analogue used to derive the finite-data transform
($G_n=\Dir_T*G$) and the periodogram expectation (a frequency convolution with
the Fej\'er kernel).
\end{theorem}

\begin{theorem}{Aliasing theorem and the Nyquist frequency}{p6-aliasing}
Let $g\in L^1$ be continuous with transform $G$, and let $G$ denote the
transform of the sampled sequence $\{g(t)\}_{t\in\Delta\Ints}$. Under the
summability hypotheses, for every $f$
\[
G(f)=\sum_{k\in\Ints} G(f+k/\Delta).
\]
Sampling \emph{folds} all the continuous-frequency content at
$\{f+k/\Delta\}_{k}$ onto the single frequency $f$. For an SDF this reads
$\sdf(f)=\sum_{k\in\Ints}\sdf(f+k)$ (with $\Delta=1$). The
\textbf{Nyquist frequency} is $1/(2\Delta)$ (i.e.\ $\tfrac12$ when
$\Delta=1$): content above it cannot be recovered.\\[2pt]
\textit{Proof idea:} write $g(t)$ at integer times by both inversion formulas;
split $\int_{-\infty}^{\infty}G(f)e^{2\pi i f t}df$ into period-length blocks,
shift each block by $k/\Delta$ (the phase $e^{2\pi i t k/\Delta}=1$ for
$t\in\Delta\Ints$), swap sum and integral (summability), and match the two
representations of $g(t)$.\\
\textit{Why:} explains why we sample fast enough; if $\sdf$ is negligible
beyond $1/(2\Delta)$ the discrete spectrum faithfully matches the continuous
one, otherwise it is corrupted (poor agreement).
\end{theorem}

\begin{intuition}[Why aliasing matters]
A high-frequency sinusoid sampled too slowly looks identical to a
low-frequency one --- the wagon-wheel effect. Frequencies $f$ and $f+k/\Delta$
are indistinguishable from the samples; they are ``aliases''. The cure is to
sample finely (small $\Delta$, large Nyquist) or to low-pass filter before
sampling.
\end{intuition}

\begin{proposition}{FFT complexity}{p6-fft}
The finite Fourier transform at the $n$ Fourier frequencies $\{k/(\Delta n)\}$
can be computed in $O(n\log n)$ time via the \textbf{Fast Fourier Transform},
versus $O(n^2)$ for the naive sum.\\[2pt]
\textit{Proof idea:} divide-and-conquer --- a transform of length $n$ splits
into two of length $n/2$ (even/odd indices), recursing $\log_2 n$ levels with
$O(n)$ work each.\\
\textit{Why:} makes periodogram/spectral computation feasible at scale; one
can also invert the periodogram by FFT to get the sample ACVS cheaply.
\end{proposition}

\begin{theorem}{Herglotz's theorem}{p6-herglotz}
A complex sequence $\{g_t\}_{t\in\Ints}$ is \textbf{non-negative definite} if
and only if there exists a right-continuous, bounded, non-decreasing function
$G^{(I)}$ with $G^{(I)}(-\tfrac12)=0$ such that
\[
g_t=\int_{-1/2}^{1/2} e^{2\pi i t f}\,dG^{(I)}(f),\qquad t\in\Ints .
\]
\textit{Proof idea:} ($\Leftarrow$) any such integral representation yields a
non-negative-definite sequence because $G^{(I)}$ is non-decreasing.
($\Rightarrow$) construct $G^{(I)}$ as the (weak) limit of the non-negative
Ces\`aro/Fej\'er averages of the partial sums $\sum_{\abs t\le N} g_t
e^{-2\pi i t f}$, which form a tight family of measures.\\
\textit{Why:} it is the existence theorem for the spectrum. Since every ACVS is
non-negative definite (a foundational property of any stationary process),
\emph{every} stationary process has an integrated spectrum $S^{(I)}$ --- even
those, like harmonic processes, with no ordinary SDF.
\end{theorem}

\begin{corollary}{Integrated spectrum exists for every stationary process}{p6-integexists}
Applying Herglotz's theorem to the non-negative-definite ACVS gives
$\acvs_\tau=\int_{-1/2}^{1/2}e^{2\pi i \tau f}\,dS^{(I)}(f)$ with
$S^{(I)}(-\tfrac12)=0$, $S^{(I)}(\tfrac12)=\sigma^2$, and $S^{(I)}$
non-decreasing.\\[2pt]
\textit{Why:} guarantees a frequency-domain description always exists; the SDF
is merely the special (absolutely continuous) case $dS^{(I)}=\sdf\,df$.
\end{corollary}

\begin{theorem}{Lebesgue decomposition of the integrated spectrum}{p6-lebesgue}
Any integrated spectrum decomposes uniquely as
\[
S^{(I)}(\cdot)=S^{(I)}_1(\cdot)+S^{(I)}_2(\cdot)+S^{(I)}_3(\cdot),
\]
where each $S^{(I)}_j$ is non-negative, non-decreasing with
$S^{(I)}_j(-\tfrac12)=0$, and
\begin{enumerate}
  \item $S^{(I)}_1$ is \textbf{absolutely continuous}:
        $S^{(I)}_1(f)=\int_{-1/2}^{f}\sdf(\lambda)\,d\lambda$ (the ordinary SDF
        $\sdf$);
  \item $S^{(I)}_2$ is a \textbf{step function} with jumps $\{p_\ell\}$ at
        frequencies $\{f_\ell\}$ of pure sinusoids (a \emph{line spectrum});
  \item $S^{(I)}_3$ is a \textbf{continuous singular} function (pathological,
        of no practical use).
\end{enumerate}
\textit{Proof idea:} the Lebesgue decomposition theorem for monotone functions
/ measures, specialised to the finite interval $[-\tfrac12,\tfrac12]$.\\
\textit{Why:} classifies the possible spectral behaviours: purely continuous
($S^{(I)}_1$ only), purely discrete/line ($S^{(I)}_2$ only), or mixed.
\end{theorem}

\begin{intuition}[Reading the three components]
\textbf{Absolutely continuous} ($S^{(I)}_1$): the ACVS damps to zero,
$\acvs_\tau\to0$ as $\tau\to\infty$ (since for integrable $\sdf$,
$\int\cos(2\pi f\tau)\sdf\,df\to0$). Example: any ARMA process.
\textbf{Discrete/line} ($S^{(I)}_2$): the ACVS does \emph{not} damp to zero
(a sinusoid keeps oscillating). Example: a sinusoid with random phase and
amplitude (the harmonic process). \textbf{Mixed}: a signal-plus-noise sum of
the two, e.g.\ a periodic component buried in coloured noise.
\end{intuition}

\begin{example}{Spectral characterisation of common processes}{p6-speccases}
\begin{itemize}
  \item \textbf{Purely continuous spectrum} ($S^{(I)}_1\ge0$,
        $S^{(I)}_2\equiv0$): ARMA processes; the SDF $\sdf$ exists and
        $\acvs_\tau\to0$.
  \item \textbf{Purely discrete (line) spectrum} ($S^{(I)}_1\equiv0$,
        $S^{(I)}_2\ge0$): random-phase/amplitude sinusoid; the ACVS persists.
  \item \textbf{Mixed spectrum} ($S^{(I)}_1\ge0$, $S^{(I)}_2\ge0$): pointwise
        aggregation, e.g.\ harmonic signal plus an ARMA background.
\end{itemize}
\end{example}

\begin{theorem}{Spectral representation theorem}{p6-specrep}
Let $\{X_t\}$ be a real-valued discrete-time stationary process with mean
$\mu$. There exists an \textbf{orthogonal-increment process} $\{Z(f)\}$ on
$[-\tfrac12,\tfrac12]$ with
\[
X_t=\mu+\int_{-1/2}^{1/2} e^{2\pi i f t}\,dZ(f)
\quad\text{(equality in mean square),}
\]
where for $f\ne f'$,
\[
\E[dZ(f)]=0,\qquad
\Var\big(dZ(f)\big)=dS^{(I)}(f),\qquad
\Cov\big(dZ(f),dZ(f')\big)=0 .
\]
\textit{Proof idea:} build $Z$ as the $L^2$ (mean-square) limit of finite
Fourier sums of the $X_t$'s; orthogonality of increments mirrors the
non-decrease of $S^{(I)}$, and the variance of an increment is exactly the mass
$S^{(I)}$ places there.\\
\textit{Why:} the deep structural result: \emph{every} stationary process is a
superposition of uncorrelated random sinusoids $e^{2\pi i f t}\,dZ(f)$, with
``amount of variance'' $dS^{(I)}(f)$ at each frequency. It makes precise the
slogan ``$\sdf(f)\,df$ is the power in $[f,f+df]$''.
\end{theorem}

\begin{recallbox}[Recall: covariance of complex random variables]
For complex $Z_1,Z_2$ the covariance is
$\Cov(Z_1,Z_2)=\E[(Z_1-\E Z_1)\conj{(Z_2-\E Z_2)}]$ (note the conjugate on the
second factor). The \emph{complementary covariance} (relation) is
$\mathrm{Rel}\{Z_1,Z_2\}=\E[(Z_1-\E Z_1)(Z_2-\E Z_2)]$ (no conjugate).
\end{recallbox}

\begin{proposition}{Periodogram as squared modulus of $J$}{p6-periogramJ}
With $T=\{1,\dots,n\}$ and $Y_t=X_t-\bar X$,
\[
\hat\sdf^{(p)}(f)=\abs{J(f)}^2,\qquad
J(f)=\sqrt{\tfrac1n}\sum_{t\in T} Y_t\,e^{-2\pi i t f}.
\]
\textit{Proof idea:} starting from
$\hat\sdf^{(p)}(f)=\sum_{\tau}\frac1n\sum_t Y_tY_{t+\abs\tau}e^{-2\pi i\tau f}$,
re-index the double lag sum as a free double sum over $s,t\in T$,
\[
\hat\sdf^{(p)}(f)=\frac1n\sum_{s=1}^{n}\sum_{t=1}^{n}
Y_tY_s\,e^{-2\pi i s f}e^{2\pi i t f}
=\Big(\sqrt{\tfrac1n}\sum_t Y_t e^{-2\pi i t f}\Big)
\conj{\Big(\sqrt{\tfrac1n}\sum_s Y_s e^{-2\pi i s f}\Big)}
=\abs{J(f)}^2 .
\]
\textit{Why:} lets the periodogram be computed by a single FFT in
$O(n\log n)$, and makes manifest that $\hat\sdf^{(p)}\ge0$.
\end{proposition}

\begin{theorem}{Expectation of the periodogram: the Fej\'er kernel}{p6-periogexp}
For a zero-mean process,
\[
\E\big[\hat\sdf^{(p)}(f)\big]
=\sum_{\tau\in\Ints} w_\tau\,\acvs_\tau\,e^{-2\pi i\tau f},\qquad
w_\tau=\Big(1-\frac{\abs\tau}{n}\Big)\ind_{[-n,n]}(\tau),
\]
equivalently the frequency convolution
\[
\E\big[\hat\sdf^{(p)}(f)\big]=\int_{-1/2}^{1/2}\sdf(f')\,\Fej_n(f-f')\,df',
\qquad
\Fej_n(f)=
\begin{cases}
\dfrac1n\dfrac{\sin^2(n\pi f)}{\sin^2(\pi f)}, & f\ne0,\\[2.2ex]
n, & f=0,
\end{cases}
\]
the \textbf{Fej\'er kernel}.\\[2pt]
\textit{Proof idea:} $\E[\hat\acvs_\tau]=\big(1-\tfrac{\abs\tau}{n}\big)\acvs_\tau$
for $\abs\tau<n$; transform term by term. The triangular weight
$w_\tau=\tfrac1n(d*\tilde d)_\tau$ is the autocorrelation of the rectangular
window, whose transform is $\abs{\Dir_T(f)}^2/n=\Fej_n(f)$, so multiplication by
$w_\tau$ in time becomes convolution with $\Fej_n$ in frequency.\\
\textit{Why:} quantifies the \textbf{bias}. Because $\Fej_n$ has wide side
lobes that decay slowly, $\E[\hat\sdf^{(p)}]$ is a \emph{blurred} (leaked)
version of $\sdf$; peaks are smeared and power leaks into neighbouring
frequencies. As $n\to\infty$, $\Fej_n$ concentrates at $0$, so the periodogram
is asymptotically unbiased.
\end{theorem}

\begin{theorem}{Variance of the periodogram: inconsistency}{p6-periogvar}
For $f\ne0\ \mathrm{mod}\ \tfrac12$,
\[
\Var\big(\hat\sdf^{(p)}(f)\big)\xrightarrow[n\to\infty]{}\sdf(f)^2 \;>\;0,
\]
and at distinct Fourier frequencies the ordinates are asymptotically
uncorrelated.\\[2pt]
\textit{Proof idea:} for a Gaussian process $J(f)$ is asymptotically complex
Gaussian, so $\hat\sdf^{(p)}(f)=\abs{J(f)}^2$ is asymptotically
$\tfrac12\sdf(f)\chi^2_2=\sdf(f)\cdot\mathrm{Exp}(1)$, whose variance is
$\sdf(f)^2$.\\
\textit{Why:} the periodogram is \textbf{not consistent} --- its variance does
\emph{not} vanish as $n$ grows; more data merely gives more (equally noisy)
ordinates. This is the motivation for smoothing/averaging (multitapering,
or averaging neighbouring ordinates).
\end{theorem}

\begin{pitfall}
The periodogram is \emph{not} a consistent estimator of the SDF. Doubling $n$
does \textbf{not} halve the variance at a fixed $f$; it stays
$\approx\sdf(f)^2$. ``It got noisier as I added data'' is expected behaviour,
not a bug --- you see finer frequency resolution but the same per-ordinate
scatter. To reduce variance you must \emph{average} (multitaper or smooth).
\end{pitfall}

\begin{theorem}{Expectation of the tapered periodogram}{p6-tapexp}
With a taper $\{h_t\}$ (set $h_t=0$ off $T$) and mean replaced by $\mu$,
\[
\E\big[\hat\sdf^{(p)}_h(f)\big]
=\int_{-1/2}^{1/2}\sdf(f')\,\abs{H(f-f')}^2\,df',
\]
where $H$ is the discrete Fourier transform of $h$.\\[2pt]
\textit{Proof idea:} expand $\hat\sdf^{(p)}_h(f)=\sum_{t,s}h_th_sY_tY_s
e^{-2\pi i f(t-s)}$, take expectations to get $\sum_\tau w_\tau\acvs_\tau
e^{-2\pi i f\tau}$ with $w_\tau=\sum_t h_th_{t+\tau}=(\tilde h*\bar h)_\tau$
(autocorrelation of the taper). By the convolution theorem its transform is
$\tilde H(f)\bar H(f)=\abs{H(f)}^2$, giving the frequency convolution.\\
\textit{Why:} the kernel is now $\abs{H(f)}^2$ instead of the Fej\'er kernel
$\Fej_n$. A well-chosen taper (e.g.\ a dpss / Slepian taper) makes
$\abs{H(f)}^2$ far more concentrated near $0$ with much lower side lobes, so
\textbf{tapering reduces bias / leakage}.
\end{theorem}

\begin{corollary}{Unbiasedness of the tapered periodogram for white noise}{p6-tapunbias}
If $\{X_t\}$ is white noise with variance $\sigma^2$ (so $\sdf\equiv\sigma^2$)
and $\norm{h}_2^2=1$, then $\E[\hat\sdf^{(p)}_h(f)]=\sigma^2=\sdf(f)$ for all
$n$ and all $f$.\\[2pt]
\textit{Proof idea:} the kernel integrates out,
$\E[\hat\sdf^{(p)}_h(f)]=\sigma^2\int_{-1/2}^{1/2}\abs{H(f')}^2\,df'
=\sigma^2\sum_{t}h_t^2=\sigma^2\norm{h}_2^2$ (Parseval), which equals
$\sigma^2$ exactly when $\norm h_2^2=1$.\\
\textit{Why:} pins down the normalisation $\norm h_2^2=1$ in the definition of
a taper, and shows tapering costs nothing in bias for flat spectra.
\end{corollary}

\begin{theorem}{Variance reduction by multitapering}{p6-mtvar}
For $K$ orthonormal tapers, asymptotically
\[
\Cov\big(\hat\sdf^{(p)}_{h_k}(f),\hat\sdf^{(p)}_{h_{k'}}(f)\big)
\to\sdf(f)^2\,\delta_{k,k'},
\quad\text{hence}\quad
\Var\big(\hat\sdf^{(mt)}(f)\big)\approx\frac{\sdf(f)^2}{K}.
\]
\textit{Proof idea:} orthogonality of the tapers makes the $K$ individual
tapered periodograms asymptotically uncorrelated, each with variance
$\sdf(f)^2$; averaging $K$ uncorrelated terms divides the variance by $K$.\\
\textit{Why:} \textbf{multitapering restores consistency}: as $n\to\infty$ one
may let $K\to\infty$ (slowly) so $\Var\to0$. It is the standard fix for the
periodogram's non-vanishing variance, with only a modest bias increase.
\end{theorem}

\begin{intuition}[Bias vs.\ variance, the whole story in one line]
The raw periodogram has tolerable (vanishing) \emph{bias} but stubborn
\emph{variance}. \textbf{Tapering} attacks bias (sharper kernel
$\abs{H}^2$ vs.\ $\Fej_n$); \textbf{multitapering} attacks variance (average
$K$ near-independent estimates, $\Var\approx\sdf^2/K$). Together they give a
usable spectral estimate.
\end{intuition}

% ======================================================================
\subsection{Key Techniques}
% ======================================================================

\begin{keytech}{Compute an SDF from an ACVS}{p6-sdffromacvs}
To get $\sdf(f)=\sum_\tau\acvs_\tau e^{-2\pi i f\tau}$:
\begin{enumerate}
  \item Write down the ACVS (only finitely many lags for MA; geometric for AR).
  \item Pair $\tau$ with $-\tau$ and use
        $e^{-2\pi i f\tau}+e^{2\pi i f\tau}=2\cos(2\pi f\tau)$, giving
        $\sdf(f)=\acvs_0+2\sum_{\tau\ge1}\acvs_\tau\cos(2\pi f\tau)$ (real,
        even).
  \item For a finite linear filter, prefer the factored form: if
        $X_t=\sum_j\psi_j\eps_{t-j}$ with $\eps$ white noise of variance
        $\sigma^2$, then
        $\sdf(f)=\sigma^2\big|\sum_j\psi_j e^{-2\pi i j f}\big|^2$.
\end{enumerate}
For an AR/ARMA $\Phi(B)X_t=\Theta(B)\eps_t$,
$\sdf(f)=\sigma^2\,\dfrac{\abs{\Theta(e^{-2\pi i f})}^2}{\abs{\Phi(e^{-2\pi i f})}^2}$.
\textbf{Sanity checks:} $\sdf$ must be real, even, non-negative, and integrate
to $\acvs_0$ over $[-\tfrac12,\tfrac12]$.
\end{keytech}

\begin{keytech}{Test whether a candidate $\acvs$ is a valid ACVS}{p6-validacvs}
A sequence $\{\acvs_\tau\}$ is a valid ACVS \emph{iff} it is non-negative
definite, which (in the summable case) is equivalent to its Fourier transform
being non-negative:
\[
\sdf(f)=\sum_\tau\acvs_\tau e^{-2\pi i f\tau}\ \ge 0\quad\text{for all }f .
\]
Procedure: (i) check $\acvs_{-\tau}=\acvs_\tau$ and $\acvs_0\ge\abs{\acvs_\tau}$
(necessary); (ii) compute $\sdf(f)$; (iii) hunt for a frequency where
$\sdf(f)<0$ --- often the endpoints $f=0$ (gives $\sum_\tau\acvs_\tau$) or
$f=\tfrac12$ (gives $\sum_\tau(-1)^\tau\acvs_\tau$) are the culprits. A single
negative value \textbf{disqualifies} the candidate.
\end{keytech}

\begin{keytech}{Recover an ACVS from a given SDF}{p6-acvsfromsdf}
Use $\acvs_\tau=\int_{-1/2}^{1/2}\sdf(f)e^{2\pi i f\tau}\,df$. Practical moves:
\begin{itemize}
  \item Drop the imaginary part: since $\acvs_\tau$ is real and $\sdf$ even,
        $\acvs_\tau=\int_{-1/2}^{1/2}\sdf(f)\cos(2\pi f\tau)\,df
        =2\int_0^{1/2}\sdf(f)\cos(2\pi f\tau)\,df$.
  \item For polynomial $\sdf$ integrate by parts; the building blocks are
        $\int_0^{1/2}\cos(2\pi f\tau)\,df=\frac{\sin(\pi\tau)}{2\pi\tau}$ and
        $\int_0^{1/2}f\cos(2\pi f\tau)\,df
        =\frac{\sin(\pi\tau)}{4\pi\tau}+\frac{\cos(\pi\tau)-1}{4\pi^2\tau^2}$.
  \item For a rational $\sdf$ (ARMA), match it to
        $\sigma^2\abs{\Theta}^2/\abs{\Phi}^2$ and read off the recursion, or do
        a partial-fraction / contour integral.
  \item Treat $\tau=0$ separately (the $\sin(\pi\tau)/\tau$ terms need a limit).
\end{itemize}
\end{keytech}

\begin{keytech}{Read off periodogram bias and variance}{p6-readbias}
Given a periodogram problem, recall the three facts:
\begin{itemize}
  \item \textbf{Bias / blurring:} $\E[\hat\sdf^{(p)}]=\sdf*\Fej_n$; the Fej\'er
        kernel's slow-decaying side lobes leak power between frequencies.
        Sharp spectral peaks are the worst affected.
  \item \textbf{Variance:} $\Var(\hat\sdf^{(p)}(f))\to\sdf(f)^2\ne0$ --- not
        consistent.
  \item \textbf{Fixes:} tapering swaps $\Fej_n$ for the sharper $\abs{H}^2$
        (less bias); multitapering averages $K$ tapers
        ($\Var\approx\sdf^2/K$, less variance).
\end{itemize}
\end{keytech}

% ======================================================================
\subsection{Exercises}
% ======================================================================

\begin{exercise}{Fourier transform property table \quad\textnormal{[Sheet 6]}}{p6-ft-props}
For $g\in L^1$ with transform $G$, prove: (1) if $h(t)=\conj{g(t)}$ then
$H(f)=\conj{G(-f)}$; (2) if $h(t)=g(\alpha t)$, $\alpha\ne0$, then
$H(f)=G(f/\alpha)/\abs\alpha$; (3) if $h(t)=g(t+\tau)$ then
$H(f)=G(f)e^{2\pi i f\tau}$; (4) if $h(t)=\alpha g(t)+p(t)$ then
$H(f)=\alpha G(f)+P(f)$.
\end{exercise}
\begin{solution}
\textbf{(1)} Conjugation: pull the conjugate outside (it flips the sign of the
exponent),
\[
H(f)=\int_{-\infty}^{\infty}\conj{g(t)}e^{-2\pi i t f}\,dt
=\conj{\int_{-\infty}^{\infty} g(t)e^{2\pi i t f}\,dt}
=\conj{G(-f)} .
\]
\textbf{(2)} Scaling: substitute $t'=\alpha t$, $dt=dt'/\abs\alpha$ (the
$\abs\alpha$ comes from reversing limits when $\alpha<0$):
\[
H(f)=\int_{-\infty}^{\infty} g(\alpha t)e^{-2\pi i t f}\,dt
=\frac{1}{\abs\alpha}\int_{-\infty}^{\infty} g(t')e^{-2\pi i (f/\alpha)t'}\,dt'
=\frac{1}{\abs\alpha}\,G(f/\alpha).
\]
\textbf{(3)} Shift: insert $1=e^{-2\pi i f\tau}e^{2\pi i f\tau}$,
\[
H(f)=\int_{-\infty}^{\infty} g(t+\tau)e^{-2\pi i f t}\,dt
=\int_{-\infty}^{\infty} g(t+\tau)e^{-2\pi i f(t+\tau)}\,dt\;e^{2\pi i f\tau}
=G(f)e^{2\pi i f\tau}.
\]
\textbf{(4)} Linearity of the integral gives $H=\alpha G+P$ immediately.
\end{solution}

\begin{exercise}{Continuous-time convolution theorem \quad\textnormal{[Sheet 6]}}{p6-convproof}
Let $g,h\in L^1$ and $u=g*h$. Prove $U=G\cdot H$.
\end{exercise}
\begin{solution}
\[
\begin{aligned}
U(f)&=\int_{-\infty}^{\infty} u(x)e^{-2\pi i x f}\,dx
     =\int_{-\infty}^{\infty}\!\!\int_{-\infty}^{\infty} g(x-y)h(y)\,dy\;e^{-2\pi i x f}\,dx\\
    &=\int_{-\infty}^{\infty}\Big(\int_{-\infty}^{\infty} g(x-y)e^{-2\pi i x f}\,dx\Big)h(y)\,dy
     \qquad\text{(Fubini)}\\
    &=\int_{-\infty}^{\infty}\Big(\int_{-\infty}^{\infty} g(s)e^{-2\pi i s f}\,ds\Big)
       e^{-2\pi i y f}h(y)\,dy
     \qquad (s=x-y)\\
    &=G(f)\int_{-\infty}^{\infty} h(y)e^{-2\pi i y f}\,dy=G(f)H(f).
\end{aligned}
\]
The interchange is justified because
$\iint\abs{g(x-y)h(y)}\,dx\,dy=\norm g_1\norm h_1<\infty$, so Fubini applies.
\end{solution}

\begin{exercise}{Aliasing via Poisson summation \quad\textnormal{[Sheet 6]}}{p6-aliasproof}
For continuous $g\in L^1$ with transform $G$, and $G$ the transform of the
sampled sequence $\{g(t)\}_{t\in\Delta\Ints}$, show
$G(f)=\sum_{k\in\Ints} G(f+k/\Delta)$ (under the stated summability
hypotheses).
\end{exercise}
\begin{solution}
For $t\in\Delta\Ints$ both inversion formulas hold:
\[
g(t)=\int_{-\infty}^{\infty} G(f)e^{2\pi i f t}\,df,
\qquad
g(t)=\int_{-1/(2\Delta)}^{1/(2\Delta)} G(f)e^{2\pi i f t}\,df .
\]
Split the real line into period blocks centred at $k/\Delta$ and shift each:
\[
\begin{aligned}
g(t)&=\sum_{k\in\Ints}\int_{k/\Delta-1/(2\Delta)}^{k/\Delta+1/(2\Delta)} G(f)e^{2\pi i f t}\,df
     =\sum_{k\in\Ints}\int_{-1/(2\Delta)}^{1/(2\Delta)} G(f+k/\Delta)e^{2\pi i (f+k/\Delta)t}\,df\\
    &=\sum_{k\in\Ints}\int_{-1/(2\Delta)}^{1/(2\Delta)} G(f+k/\Delta)e^{2\pi i f t}\,df,
\end{aligned}
\]
because $t=z\Delta$ for some $z\in\Ints$ makes $e^{2\pi i (k/\Delta)t}
=e^{2\pi i kz}=1$. By the summability hypothesis
$\sum_k\int\abs{G(f+k/\Delta)}\,df<\infty$ we may swap sum and integral:
\[
\int_{-1/(2\Delta)}^{1/(2\Delta)}\Big(\sum_{k\in\Ints} G(f+k/\Delta)\Big)e^{2\pi i f t}\,df
=g(t)=\int_{-1/(2\Delta)}^{1/(2\Delta)} G(f)e^{2\pi i f t}\,df .
\]
Equality of these Fourier-series coefficients (for all $t\in\Delta\Ints$) forces
$\sum_k G(f+k/\Delta)=G(f)$ for a.e.\ $f$; both sides are continuous (by the
$\ell^1$/summability hypotheses), so equality holds for all $f$.
\end{solution}

\begin{exercise}{SDF of white noise \quad\textnormal{[Sheet 7]}}{p6-wnsdf}
Let $\{\eps_t\}$ be white noise with variance $\sigma_\eps^2$. Find its SDF.
\end{exercise}
\begin{solution}
The ACVS is $\acvs_\tau=\sigma_\eps^2$ for $\tau=0$ and $0$ otherwise, so only
the $\tau=0$ term survives:
\[
\sdf(f)=\sum_{\tau\in\Ints}\acvs_\tau e^{-2\pi i f\tau}=\sigma_\eps^2,
\qquad\text{for all }f .
\]
The spectrum is \textbf{flat} --- equal power at every frequency, which is why
it is called ``white''.
\end{solution}

\begin{exercise}{Is a truncated ACVS a valid ACVS? \quad\textnormal{[Sheet 7]}}{p6-truncacvs}
Let $\acvs_\tau=1$ for $\abs\tau\le K$ and $0$ otherwise, $K$ a positive
integer. Is $\{\acvs_\tau\}$ a valid ACVS? What would the SDF be?
\end{exercise}
\begin{solution}
It is summable and symmetric, so we may form
\[
\sdf(f)=\sum_{\tau=-K}^{K} e^{-2\pi i\tau f}
=1+2\sum_{\tau=1}^{K}\cos(2\pi\tau f)
=\frac{\sin\big((2K+1)\pi f\big)}{\sin(\pi f)}
\]
(the Dirichlet kernel). This is \emph{not} non-negative everywhere. Evaluate at
the Nyquist endpoint $f=\tfrac12$: using $\cos(\pi\tau)=(-1)^\tau$,
\[
\sdf(\tfrac12)=1+2\sum_{\tau=1}^{K}(-1)^\tau=(-1)^K .
\]
So $\sdf(\tfrac12)=-1<0$ whenever $K$ is \textbf{odd} (write $K=2q+1$). Since the
SDF goes negative, $\{\acvs_\tau\}$ is \textbf{not} a valid ACVS for any
stationary process. (It fails to be non-negative definite.)
\end{solution}

\begin{intuition}
Truncating a genuine ACVS to a finite window in general breaks
non-negative-definiteness --- exactly why one cannot just chop the ACVS and
hope. The clean way to ``truncate'' is to taper the spectrum, not the ACVS.
\end{intuition}

\begin{exercise}{ACVS from a triangular SDF \quad\textnormal{[Sheet 7]}}{p6-triangsdf}
Find the ACVS of $\{X_t\}$ with
$\sdf(f)=\dfrac{\sigma^2}{\pi}\,(1-2\abs f)\,\ind_{\{\abs f\le 1/2\}}$.
\end{exercise}
\begin{solution}
For $\tau\ne0$, since $\sdf$ is even the $\sin$ part drops out:
\[
\acvs_\tau=\int_{-1/2}^{1/2} e^{2\pi i f\tau}\sdf(f)\,df
=\frac{2\sigma^2}{\pi}\int_{0}^{1/2}\cos(2\pi f\tau)(1-2f)\,df
=\frac{2\sigma^2}{\pi}\big(I_1-2I_2\big),
\]
with $I_1=\int_0^{1/2}\cos(2\pi f\tau)\,df$ and
$I_2=\int_0^{1/2} f\cos(2\pi f\tau)\,df$. Computing,
\[
I_1=\frac{\sin(\pi\tau)}{2\pi\tau},\qquad
I_2=\frac{\sin(\pi\tau)}{4\pi\tau}+\frac{\cos(\pi\tau)-1}{4\pi^2\tau^2}.
\]
Then $I_1-2I_2=-\dfrac{\cos(\pi\tau)-1}{2\pi^2\tau^2}
=\dfrac{1-\cos(\pi\tau)}{2\pi^2\tau^2}$, so
\[
\boxed{\;\acvs_\tau=\frac{\sigma^2}{\pi^3}\,\frac{1-\cos(\pi\tau)}{\tau^2}\;}
\qquad(\tau\ne0).
\]
For $\tau=0$, $\acvs_0=\dfrac{2\sigma^2}{\pi}\int_0^{1/2}(1-2f)\,df
=\dfrac{2\sigma^2}{\pi}\cdot\tfrac14=\dfrac{\sigma^2}{2\pi}$. (Note
$1-\cos(\pi\tau)=0$ for even $\tau$, $=2$ for odd $\tau$, so the ACVS is
supported on odd lags plus $\tau=0$. Numerically: $\acvs_0=\sigma^2/(2\pi)$,
$\acvs_1=2\sigma^2/\pi^3$, $\acvs_2=0$, $\acvs_3=2\sigma^2/(9\pi^3)$, all
verified.)
\end{solution}

\begin{exercise}{SDF of an MA(q) process \quad\textnormal{[Sheet 7]}}{p6-maqsdf}
What is the SDF of an MA($q$) process
$X_t=\sum_{j=0}^{q}\theta_j\eps_{t-j}$ (here writing $\theta_0=1$ for
concreteness, $\eps$ white noise of variance $\sigma_\eps^2$)?
\end{exercise}
\begin{solution}
Its ACVS is $\acvs_\tau=\sigma_\eps^2\sum_{j=0}^{q}\theta_j\theta_{j+\abs\tau}$
for $\abs\tau\le q$ (zero otherwise), which we write as the double sum
$\acvs_\tau=\sigma_\eps^2\sum_{j=0}^{q}\sum_{k=0}^{q}
\theta_j\theta_k\ind_{\{\tau=k-j\}}$. Since $\acvs\in\ell^1$ the SDF exists:
\[
\begin{aligned}
\sdf(f)&=\sum_{\tau\in\Ints}\acvs_\tau e^{-2\pi i\tau f}
       =\sigma_\eps^2\sum_{j=0}^{q}\sum_{k=0}^{q}\theta_j\theta_k
        \sum_{\tau}\ind_{\{\tau=k-j\}}e^{-2\pi i\tau f}\\
       &=\sigma_\eps^2\sum_{j=0}^{q}\sum_{k=0}^{q}\theta_j\theta_k
        e^{-2\pi i(k-j)f}
       =\sigma_\eps^2\Big(\sum_{j=0}^{q}\theta_j e^{-2\pi i j f}\Big)
        \conj{\Big(\sum_{k=0}^{q}\theta_k e^{-2\pi i k f}\Big)}\\
       &=\sigma_\eps^2\,\Big|\,\sum_{j=0}^{q}\theta_j e^{-2\pi i j f}\,\Big|^2
        =\sigma_\eps^2\,\abs{\Theta(e^{-2\pi i f})}^2 .
\end{aligned}
\]
The SDF is the squared modulus of the MA transfer function --- automatically
non-negative, as any SDF must be.
\end{solution}

\begin{exercise}{Time-varying-amplitude harmonic process \quad\textnormal{[Sheet 7]}}{p6-tvharm}
Let $X_t=\eps_t\cos(\nu t+\Theta)$ with $\nu$ fixed, $\eps_t$ mean-zero white
noise of variance $\sigma^2$ (not constrained positive), independent of the
random phase $\Theta$ whose pdf is
$f_\Theta(\theta)=\tfrac1{2\pi}(1+\cos\theta)$ on $\abs\theta\le\pi$. Find the
first two moments of $X_t$; are they $t$-free?
\end{exercise}
\begin{solution}
\textbf{Mean.} By independence,
$\E[X_t]=\E[\eps_t]\,\E[\cos(\nu t+\Theta)]=0\cdot(\cdot)=0$, constant in $t$
(the second factor is finite since $\abs{\cos}\le1$).\\
\textbf{Autocovariance.} By independence,
$\E[X_tX_{t+\tau}]
=\E[\eps_t\eps_{t+\tau}]\,\E[\cos(\nu t+\Theta)\cos(\nu(t+\tau)+\Theta)]$.
White noise gives $\E[\eps_t\eps_{t+\tau}]=\sigma^2\ind_{\{\tau=0\}}$. Using
$\cos A\cos B=\tfrac12[\cos(A-B)+\cos(A+B)]$,
\[
\E\big[\cos(\nu t+\Theta)\cos(\nu(t+\tau)+\Theta)\big]
=\tfrac12\big[\cos(\nu\tau)+\E[\cos(\nu(2t+\tau)+2\Theta)]\big].
\]
The remaining expectation vanishes for this $\Theta$:
\[
\begin{aligned}
\E[\cos(\nu(2t+\tau)+2\Theta)]
&=\frac1{2\pi}\int_{-\pi}^{\pi}\cos(\nu(2t+\tau)+2\theta)\,d\theta\\
&\quad+\frac1{2\pi}\int_{-\pi}^{\pi}\cos(\nu(2t+\tau)+2\theta)\cos\theta\,d\theta .
\end{aligned}
\]
The first integral is
$\tfrac1{4\pi}\big[\sin(\nu(2t+\tau)+2\theta)\big]_{-\pi}^{\pi}=0$.
For the second, write
$\cos(\cdots+2\theta)\cos\theta=\tfrac12[\cos(\cdots+\theta)+\cos(\cdots+3\theta)]$;
both pieces integrate to $0$ over the full period. Hence
$\E[\cos(\nu(2t+\tau)+2\Theta)]=0$, and
\[
\E[X_tX_{t+\tau}]=\sigma^2\ind_{\{\tau=0\}}\cdot\tfrac12\cos(\nu\tau)
=\begin{cases}\sigma^2/2,&\tau=0,\\[2pt]0,&\tau\ne0.\end{cases}
\]
Both moments are independent of $t$, so $\{X_t\}$ is stationary; in fact it is
white noise with variance $\sigma^2/2$. (The non-uniform phase did not spoil
stationarity precisely because its $2\theta$-harmonics integrate out.)
\end{solution}

\begin{exercise}{Periodogram as $\abs{J(f)}^2$ \quad\textnormal{[Sheet 8]}}{p6-perioJ}
With $T=\{1,\dots,n\}$, show $\hat\sdf^{(p)}(f)=\abs{J(f)}^2$ for
$J(f)=\sqrt{1/n}\sum_{t\in T}(X_t-\bar X)e^{-2\pi i t f}$.
\end{exercise}
\begin{solution}
Put $Y_t=X_t-\bar X$. Starting from the definition and re-indexing the lag sum
as a free double sum over $s,t\in\{1,\dots,n\}$ (the substitution $\tau=t-s$
sweeps all $\abs\tau\le n-1$),
\[
\begin{aligned}
\hat\sdf^{(p)}(f)
&=\sum_{\tau=-n+1}^{n-1}\frac1n\sum_{t=1}^{n-\abs\tau} Y_tY_{t+\abs\tau}e^{-2\pi i\tau f}
 =\frac1n\sum_{s=1}^{n}\sum_{t=1}^{n} Y_tY_s\,e^{-2\pi i s f}e^{2\pi i t f}\\
&=\Big(\sqrt{\tfrac1n}\sum_{t=1}^{n} Y_t e^{-2\pi i t f}\Big)
  \conj{\Big(\sqrt{\tfrac1n}\sum_{s=1}^{n} Y_s e^{-2\pi i s f}\Big)}
 =J(f)\,\conj{J(f)}=\abs{J(f)}^2 .
\end{aligned}
\]
(In the second equality we relabel $s=t+|\tau|$ so the lag $\tau=t-s$ runs
freely, and $e^{-2\pi i\tau f}=e^{2\pi i t f}e^{-2\pi i s f}$ separates the two
conjugate sums.)
\end{solution}

\begin{exercise}{Expectation of the tapered periodogram \quad\textnormal{[Sheet 8]}}{p6-tapexp}
Replacing $\bar X$ by $\mu$, show
$\E[\hat\sdf^{(p)}_h(f)]=\int_{-1/2}^{1/2}\sdf(f')\abs{H(f-f')}^2\,df'$, with
$H$ the transform of $h$.
\end{exercise}
\begin{solution}
With $Y_t=X_t-\mu$ and $g_t=h_t$ ($g_t=0$ off $T$),
\[
\hat\sdf^{(p)}_h(f)=\Big|\sum_{t=1}^{n} g_t Y_t e^{-2\pi i f t}\Big|^2
=\sum_{\tau=-n+1}^{n-1}\sum_{t} g_t g_{t+\abs\tau}\,Y_t Y_{t+\tau}\,e^{-2\pi i f\tau}.
\]
Taking expectations ($\E[Y_tY_{t+\tau}]=\acvs_\tau$),
\[
\E\big[\hat\sdf^{(p)}_h(f)\big]
=\sum_{\tau\in\Ints} w_\tau\,\acvs_\tau\,e^{-2\pi i f\tau},
\qquad w_\tau=\sum_{t\in\Ints} h_t h_{t+\tau}.
\]
By the convolution theorem $\E[\hat\sdf^{(p)}_h(f)]=\int_{-1/2}^{1/2}
\sdf(f')W(f-f')\,df'$ where $W$ is the transform of $w$. Setting
$\tilde h_t=h_{-t}$ and $\bar h_t=h_t$, we recognise $w=\tilde h*\bar h$, so by
the convolution theorem again $W(f)=\tilde H(f)\bar H(f)=\abs{H(f)}^2$.
Therefore
$\E[\hat\sdf^{(p)}_h(f)]=\int_{-1/2}^{1/2}\sdf(f')\abs{H(f-f')}^2\,df'$.
\end{solution}

\begin{exercise}{Unbiased tapered periodogram for white noise \quad\textnormal{[Sheet 8]}}{p6-tunbias}
If $\{X_t\}$ is white noise with variance $\sigma^2$, prove
$\hat\sdf^{(p)}_h(f)$ is unbiased for all $n$ provided $\norm h_2^2=1$.
\end{exercise}
\begin{solution}
White noise has $\sdf\equiv\sigma^2$. Using the tapered-expectation formula and
relabelling (periodicity),
\[
\begin{aligned}
\E\big[\hat\sdf^{(p)}_h(f)\big]
&=\int_{-1/2}^{1/2}\sdf(f-f')\abs{H(f')}^2\,df'
 =\sigma^2\int_{-1/2}^{1/2}\abs{H(f')}^2\,df'\\
&=\sigma^2\sum_{t\in\Ints}\sum_{s\in\Ints} h_t h_s
   \int_{-1/2}^{1/2} e^{-2\pi i(t-s)f'}\,df'
 =\sigma^2\sum_{t} h_t^2=\sigma^2\norm h_2^2,
\end{aligned}
\]
where $\int_{-1/2}^{1/2}e^{-2\pi i\tau f}\,df=\ind_{\{\tau=0\}}$ collapses the
double sum (Parseval). Hence if $\norm h_2^2=1$ then
$\E[\hat\sdf^{(p)}_h(f)]=\sigma^2=\sdf(f)$, i.e.\ exactly unbiased.
\end{solution}

% ---------------------- NEW EXERCISES ----------------------------------

\begin{exercise}{SDF of an MA(1) process \quad\textnormal{[New]}}{p6-ma1sdf}
Let $X_t=\eps_t-\theta\eps_{t-1}$ with $\eps$ white noise of variance
$\sigma^2$. Find $\sdf(f)$ two ways (from the ACVS and from the transfer
function), and locate its maximum and minimum.
\end{exercise}
\begin{solution}
\textbf{From the ACVS.} $\acvs_0=\sigma^2(1+\theta^2)$,
$\acvs_{\pm1}=-\sigma^2\theta$, $0$ otherwise. Then
\[
\sdf(f)=\acvs_0+2\acvs_1\cos(2\pi f)
=\sigma^2\big(1+\theta^2-2\theta\cos(2\pi f)\big).
\]
\textbf{From the transfer function.} $\Theta(z)=1-\theta z$, so
\[
\sdf(f)=\sigma^2\abs{1-\theta e^{-2\pi i f}}^2
=\sigma^2\big(1-\theta e^{-2\pi i f}\big)\big(1-\theta e^{2\pi i f}\big)
=\sigma^2(1+\theta^2-2\theta\cos(2\pi f)),
\]
the same answer. Since $\cos(2\pi f)$ ranges over $[-1,1]$: for $\theta>0$ the
\textbf{minimum} is at $f=0$, $\sdf(0)=\sigma^2(1-\theta)^2$, and the
\textbf{maximum} at $f=\tfrac12$, $\sdf(\tfrac12)=\sigma^2(1+\theta)^2$
(high-pass shape); for $\theta<0$ these swap (low-pass). The SDF is
non-negative everywhere, as required (it touches $0$ only on the unit circle
when $\abs\theta=1$).
\end{solution}

\begin{exercise}{SDF of an AR(1) process \quad\textnormal{[New]}}{p6-ar1sdf}
Let $X_t=\phi X_{t-1}+\eps_t$, $\abs\phi<1$, $\eps$ white noise of variance
$\sigma^2$. Find $\sdf(f)$ from the ACVS and confirm it via the transfer
function. Describe its shape for $\phi>0$ and $\phi<0$.
\end{exercise}
\begin{solution}
The ACVS is $\acvs_\tau=\dfrac{\sigma^2}{1-\phi^2}\,\phi^{\abs\tau}$. Summing the
two geometric tails,
\[
\begin{aligned}
\sdf(f)&=\frac{\sigma^2}{1-\phi^2}\sum_{\tau\in\Ints}\phi^{\abs\tau}e^{-2\pi i f\tau}
        =\frac{\sigma^2}{1-\phi^2}\Big(1+\sum_{\tau\ge1}\phi^\tau
         (e^{-2\pi i f\tau}+e^{2\pi i f\tau})\Big)\\
       &=\frac{\sigma^2}{1-\phi^2}\Big(1+\frac{\phi e^{-2\pi i f}}{1-\phi e^{-2\pi i f}}
         +\frac{\phi e^{2\pi i f}}{1-\phi e^{2\pi i f}}\Big)
        =\frac{\sigma^2}{1-\phi^2}\cdot\frac{1-\phi^2}
         {1-2\phi\cos(2\pi f)+\phi^2}\\
       &=\frac{\sigma^2}{1-2\phi\cos(2\pi f)+\phi^2}.
\end{aligned}
\]
\textbf{Transfer-function check.} $\Phi(z)=1-\phi z$, so
$\sdf(f)=\sigma^2/\abs{1-\phi e^{-2\pi i f}}^2
=\sigma^2/(1-2\phi\cos(2\pi f)+\phi^2)$ --- identical. For $\phi>0$ the
denominator is smallest at $f=0$, so $\sdf$ peaks at low frequency (slowly
varying, ``red'' spectrum); for $\phi<0$ it peaks at $f=\tfrac12$ (rapidly
alternating, ``blue'' spectrum). As $\phi\to1$ the peak at $0$ blows up
(near-unit-root, very long-range dependence).
\end{solution}

\begin{exercise}{Endpoint test for a candidate ACVS \quad\textnormal{[New]}}{p6-endpoint}
Decide whether $\acvs_0=1,\ \acvs_{\pm1}=\rho,\ \acvs_\tau=0\ (\abs\tau\ge2)$ is
a valid ACVS, as a function of $\rho\in\Reals$.
\end{exercise}
\begin{solution}
This is the ACVS shape of an MA(1). Compute the SDF:
\[
\sdf(f)=1+2\rho\cos(2\pi f).
\]
Since $\cos(2\pi f)\in[-1,1]$, the minimum of $\sdf$ is $1-2\abs\rho$, attained
at $f=0$ (if $\rho<0$) or $f=\tfrac12$ (if $\rho>0$). Non-negativity
$\sdf(f)\ge0$ for all $f$ holds \emph{iff} $1-2\abs\rho\ge0$, i.e.
\[
\boxed{\;\abs\rho\le\tfrac12\;}
\]
For $\abs\rho>\tfrac12$ the SDF dips below zero and the sequence is \textbf{not}
a valid ACVS. (Consistent with the fact that an MA(1) cannot have
$\abs{\acf_1}>\tfrac12$.)
\end{solution}

\begin{exercise}{Variance from the SDF; Parseval \quad\textnormal{[New]}}{p6-parseval}
For a stationary process with SDF $\sdf$, show
$\Var(X_t)=\int_{-1/2}^{1/2}\sdf(f)\,df$, and use it to read off $\acvs_0$ for
the AR(1) SDF $\sdf(f)=\sigma^2/(1-2\phi\cos(2\pi f)+\phi^2)$.
\end{exercise}
\begin{solution}
By the inversion formula at lag $\tau=0$,
\[
\acvs_0=\int_{-1/2}^{1/2}\sdf(f)e^{2\pi i f\cdot 0}\,df
=\int_{-1/2}^{1/2}\sdf(f)\,df=\Var(X_t),
\]
the spectral statement of Parseval: total variance equals total spectral mass.
For the AR(1), rather than integrate directly we use the known
$\acvs_0=\sigma^2/(1-\phi^2)$; consistency is the standard integral
\[
\int_{-1/2}^{1/2}\frac{\sigma^2\,df}{1-2\phi\cos(2\pi f)+\phi^2}
=\frac{\sigma^2}{1-\phi^2},
\]
which one can verify by residues (substitute $z=e^{2\pi i f}$; the integrand has
poles at $z=\phi$ and $z=1/\phi$, only $z=\phi$ lying inside the unit circle).
Thus $\Var(X_t)=\sigma^2/(1-\phi^2)$, matching the time-domain computation.
\end{solution}

% ======================================================================
%  PART 7 -- Multivariate Time Series and Vector Autoregression (VAR)
% ======================================================================
\section{Multivariate Time Series and Vector Autoregression (VAR)}
\label{sec:multivariate}

So far each $X_t$ has been a scalar. Many real systems are
$d$-dimensional: a basket of normalised stock prices (Apple, Nasdaq-100, S\&P
500), or the longitude/latitude/temperature of an ocean drifter. We now let
$\{X_t\}$ take values in $\Reals^d$ and ask two questions: (i) how do we
\emph{describe} the second-order structure --- now a matrix-valued
autocovariance and a matrix-valued spectrum, with genuine \textbf{cross}-terms
between components, and (ii) how do we \emph{model} it. The canonical model is
the vector autoregression VAR($p$), and the key trick is that
\emph{every VAR($p$) is a VAR(1)} in disguise (companion form), so its theory
reduces to that of a single matrix $\Phi_1$ (or $F$).

% ----------------------------------------------------------------------
\subsection{Definitions}

\begin{definition}{Multivariate time series}{p7-mvts}
A \textbf{multivariate time series} is a real $d$-vector-valued discrete-time
stochastic process
\[
X_t=\big(X_t^{(1)},X_t^{(2)},\dots,X_t^{(d)}\big)\Tr,\qquad t\in\Ints,
\]
whose marginal processes $\{X_t^{(1)}\},\dots,\{X_t^{(d)}\}$ are each univariate
time series. When $d=2$ we call it \textbf{bivariate}.
\end{definition}

\begin{definition}{Multivariate second-order stationarity}{p7-mvstat}
$\{X_t\}$ is \textbf{multivariate second-order (weakly) stationary} if for all
$t,s,\tau\in\Ints$,
\[
\E[X_t]=\E[X_s],\qquad
\Cov(X_{t+\tau},X_t)=\Cov(X_{s+\tau},X_s),\qquad
\trace\big(\Var(X_t)\big)<\infty .
\]
Equivalently: each marginal $\{X_t^{(j)}\}$ is second-order stationary
\textbf{and} every pair has a shift-invariant cross-covariance,
$\Cov\!\big(X_{t+\tau}^{(j)},X_t^{(k)}\big)=\Cov\!\big(X_{s+\tau}^{(j)},X_s^{(k)}\big)$
for all $1\le j,k\le d$. (Here $\Cov(U,V)$ of two vectors is the matrix with
$(j,k)$ entry $\Cov(U^{(j)},V^{(k)})$.)
\end{definition}

\begin{definition}{Matrix autocovariance sequence (ACVS)}{p7-acvs}
For a multivariate second-order stationary $\{X_t\}$, the
\textbf{(matrix) autocovariance sequence} is the $d\times d$ matrix-valued
sequence
\[
\Gamma_\tau=\Cov(X_{t+\tau},X_t),\qquad \tau\in\Ints ,
\]
with $(j,k)$ entry $\acvs_\tau^{(j,k)}=\Cov\!\big(X_{t+\tau}^{(j)},X_t^{(k)}\big)$.
\begin{itemize}
  \item \textbf{The lag $\tau$ sits in the FIRST argument.} In the scalar case
        this is immaterial, but here it determines the convention for
        $\Gamma_{-\tau}=\Gamma_\tau\Tr$ below.
  \item When $j=k$, $\acvs_\tau^{(j,j)}$ is the ordinary ACVS of the marginal
        $\{X_t^{(j)}\}$.
  \item When $j\ne k$, $\acvs_\tau^{(j,k)}$ is called the
        \textbf{cross-covariance sequence} (CCS) between components $j$ and $k$.
\end{itemize}
\end{definition}

\begin{definition}{Cross-correlation sequence (CCS)}{p7-ccs}
The \textbf{cross-correlation sequence} between components $j$ and $k$ is
\[
\acf_\tau^{(j,k)}
=\frac{\acvs_\tau^{(j,k)}}{\sqrt{\acvs_0^{(j,j)}\,\acvs_0^{(k,k)}}} .
\]
For $j=k$ this is the ordinary ACF (so $\acf_0^{(j,j)}=1$). For $j\ne k$ it need
\emph{not} equal $1$ at $\tau=0$ and need \emph{not} be symmetric in $\tau$; it
satisfies $\acf_\tau^{(j,k)}=\acf_{-\tau}^{(k,j)}$.
\end{definition}

\begin{definition}{Multivariate white noise $\mathrm{WN}(0,\Sigma)$}{p7-wn}
A $d$-variate series $\{\eps_t\}$ is \textbf{white noise with mean $0$ and
covariance matrix $\Sigma$}, written $\{\eps_t\}\sim\mathrm{WN}(0,\Sigma)$, iff
it is stationary with mean vector $0$ and
\[
\Gamma_\tau^{(\eps)}=
\begin{cases}\Sigma & \tau=0,\\[2pt] 0 & \tau\ne0.\end{cases}
\]
\begin{itemize}
  \item There \emph{can} be correlation between different components at lag $0$
        (off-diagonal entries of $\Sigma$): the components are uncorrelated
        \textbf{across time} but possibly correlated \textbf{contemporaneously}.
  \item Often one further assumes $\Sigma$ \emph{diagonal} (componentwise
        independent noise), but this is an extra assumption, not part of the
        definition.
\end{itemize}
\end{definition}

\begin{definition}{Uncorrelated processes}{p7-uncorr}
Two jointly second-order stationary processes $\{X_t^{(1)}\}$ and
$\{X_t^{(2)}\}$ are \textbf{uncorrelated} if their cross-covariance vanishes at
\emph{every} lag:
\[
\acvs_\tau^{(1,2)}=0\qquad\text{for all }\tau\in\Ints .
\]
\end{definition}

\begin{definition}{Spectral density matrix function}{p7-sdm}
Assume $\{\acvs_\tau^{(j,k)}\}_{\tau\in\Ints}\in\ell^1$ for all $1\le j,k\le d$
(absolutely summable cross-covariances). The \textbf{spectral density matrix}
is the $d\times d$ matrix-valued function
\[
S(f)=\sum_{\tau\in\Ints}\Gamma_\tau\,e^{-2\pi i\tau f},
\qquad f\in[-\tfrac12,\tfrac12],
\]
with inverse $\displaystyle \Gamma_\tau=\int_{-1/2}^{1/2}S(f)\,e^{2\pi if\tau}\,df$.
\begin{itemize}
  \item The $(j,k)$ entry is $S_{j,k}(f)$. For $j=k$ it is the ordinary
        (real-valued) spectral density $\sdf_{j}(f)$.
  \item For $j\ne k$, $S_{j,k}(f)$ is the \textbf{cross-spectral density}, which
        is \emph{complex-valued} in general.
\end{itemize}
\end{definition}

\begin{definition}{Coherence}{p7-coh}
The \textbf{coherence} between components $j$ and $k$ is
\[
r_{j,k}(f)=\frac{S_{j,k}(f)}{\sqrt{S_{j,j}(f)\,S_{k,k}(f)}}\in\Comp ,
\]
the frequency-domain analogue of a correlation between the increments $dZ_j$ and
$dZ_k$. One reports:
\begin{itemize}
  \item its \textbf{magnitude} $\abs{r_{j,k}(f)}$ (or magnitude-squared
        $\abs{r_{j,k}(f)}^2\in[0,1]$), the strength of frequency-by-frequency
        association;
  \item its \textbf{phase} $\arg r_{j,k}(f)$, the ``sign'' / lead--lag angle in
        the complex plane.
\end{itemize}
\end{definition}

\begin{definition}{Vector autoregression VAR($p$)}{p7-var}
Let $\{\eps_t\}\sim\mathrm{WN}(0,\Sigma)$ be $d$-variate. A process $\{X_t\}$ is
a \textbf{vector autoregressive process of order $p$}, VAR($p$), if
\[
X_t=\Phi_1 X_{t-1}+\Phi_2 X_{t-2}+\cdots+\Phi_p X_{t-p}+\eps_t ,
\qquad \Phi_j\in\Reals^{d\times d},\ \Phi_p\ne0 .
\]
Defining the matrix polynomial $\Phi(z)=I_d-\sum_{j=1}^p\Phi_j z^j$, this is
written compactly as
\[
\Phi(\Bsh)X_t=\eps_t ,
\]
where $\Bsh$ is the backshift operator. The same regressors (past values of
the \emph{whole} vector $X_t$) appear in every component equation; the
off-diagonal entries of the $\Phi_j$ are exactly the cross-component dynamics.
\end{definition}

\begin{definition}{Companion form of a VAR($p$)}{p7-companion}
Stack the last $p$ vectors into $Y_t=(X_t\Tr,X_{t-1}\Tr,\dots,X_{t-p+1}\Tr)\Tr
\in\Reals^{dp}$. Then a VAR($p$) is equivalently the VAR(1)
\[
\underbrace{Y_t}_{dp\times1}
=\underbrace{F}_{dp\times dp}\,\underbrace{Y_{t-1}}_{dp\times1}
+\underbrace{U_t}_{dp\times1},
\qquad
F=\begin{pmatrix}
\Phi_1 & \Phi_2 & \cdots & \Phi_{p-1} & \Phi_p\\
I_d & 0 & \cdots & 0 & 0\\
0 & I_d & & \vdots & \vdots\\
\vdots & & \ddots & 0 & 0\\
0 & \cdots & 0 & I_d & 0
\end{pmatrix},
\quad
U_t=\begin{pmatrix}\eps_t\\0\\\vdots\\0\end{pmatrix}.
\]
The top block row reproduces the VAR($p$) equation; the lower rows are
identities $X_{t-i}=X_{t-i}$ that shift the stack. Recover the original process
via $X_t=G Y_t$ with $G=(I_d\ \ 0\ \cdots\ 0)\in\Reals^{d\times dp}$.
\end{definition}

\begin{definition}{Stability of a VAR($p$)}{p7-stable}
A VAR($p$) is \textbf{stable} if the roots of the \textbf{reverse
characteristic polynomial}
\[
\det\!\big(I_d-\Phi_1 z-\Phi_2 z^2-\cdots-\Phi_p z^p\big)=0
\]
all lie \emph{strictly outside} the complex unit circle ($\abs{z}>1$).
Equivalently, the companion matrix $F$ has all eigenvalues of modulus
\emph{strictly less than one} ($\abs{\lambda}<1$). Stability $\Rightarrow$
stationarity (but not conversely in general).
\end{definition}

\begin{intuition}[Two equivalent stability tests --- and why they are reciprocal]
The reverse polynomial $\det(I_d-\sum_j\Phi_j z^j)$ uses $z=1/\lambda$, so
``roots $z$ outside the unit circle'' is exactly ``eigenvalues $\lambda$ of $F$
inside the unit circle''. For a VAR(1), $\det(I_d-\Phi_1 z)=0$ means $1/z$ is an
eigenvalue of $\Phi_1$, i.e.\ $z=1/\lambda_i(\Phi_1)$; demanding $\abs{z}>1$ is
demanding $\abs{\lambda_i(\Phi_1)}<1$. This mirrors the scalar AR causality
condition ``roots of $\phi(z)$ outside the unit circle''.
\end{intuition}

% ----------------------------------------------------------------------
\subsection{Theorems \& Propositions}

\begin{proposition}{Symmetry of the matrix ACVS}{p7-gammasym}
For a multivariate second-order stationary process,
\[
\acvs_\tau^{(j,k)}=\acvs_{-\tau}^{(k,j)}
\qquad\Longleftrightarrow\qquad
\Gamma_{-\tau}=\Gamma_\tau\Tr .
\]
\textit{Proof idea:} $\acvs_\tau^{(j,k)}=\Cov(X_{t+\tau}^{(j)},X_t^{(k)})
=\Cov(X_t^{(k)},X_{t+\tau}^{(j)})$ (covariance is symmetric in its arguments)
$=\Cov(X_{s-\tau}^{(k)},X_s^{(j)})$ (stationarity, set $s=t+\tau$)
$=\acvs_{-\tau}^{(k,j)}$.\\
\textit{Why:} the cross-covariance need \emph{not} be even in $\tau$
($\acvs_\tau^{(j,k)}\ne\acvs_{-\tau}^{(j,k)}$ in general); only the
\emph{transpose} symmetry $\Gamma_{-\tau}=\Gamma_\tau\Tr$ holds. This asymmetry
is precisely what encodes lead--lag relationships between channels.
\end{proposition}

\begin{pitfall}
In the univariate case $\acvs_{-\tau}=\acvs_\tau$ (even). Do \emph{not} carry
this over: in the multivariate case the diagonal entries are still even, but the
off-diagonal cross-covariances generally satisfy only
$\acvs_\tau^{(j,k)}=\acvs_{-\tau}^{(k,j)}$. Forgetting to swap the indices
$(j,k)\to(k,j)$ when flipping the sign of $\tau$ is the classic error.
\end{pitfall}

\begin{proposition}{Positive semi-definiteness of the matrix ACVS}{p7-psd}
The matrix sequence $\{\Gamma_\tau\}$ is positive semi-definite: for every
$n\in\Nat$ and every $a_1,\dots,a_n\in\Reals^d$,
\[
\sum_{j=1}^n\sum_{k=1}^n a_j\Tr\,\Gamma_{k-j}\,a_k\ge0 .
\]
\textit{Proof idea:} the left side equals $\Var\!\big(\sum_{j}a_j\Tr X_{t_j}\big)\ge0$
for suitable times (the scalar random variable $\sum_j a_j\Tr X_{t_j}$ has
non-negative variance).\\
\textit{Why:} this is the defining algebraic constraint a valid matrix ACVS must
satisfy; a candidate matrix sequence violating it cannot be the autocovariance
of any process.
\end{proposition}

\begin{proposition}{Contemporaneous correlation induced by shared noise}{p7-contemp}
Let $X_t=(X_t^{(1)},X_t^{(2)})\Tr$ be second-order stationary with ACVS
$\Gamma_\tau^{(X)}$, and let $\{\phi_t\}$ be a univariate mean-zero white noise
(variance $\sigma_\phi^2$) independent of $\{X_t\}$. Define
$Y_t=X_t+\phi_t\mathbf{1}$ (the \emph{same} scalar noise added to both
components), where $\mathbf{1}=(1,1)\Tr$. Then
\[
\Gamma_\tau^{(Y)}=\Gamma_\tau^{(X)}+\sigma_\phi^2\,\delta_{\tau,0}\,\mathbf 1\mathbf 1\Tr ,
\qquad
\mathbf 1\mathbf 1\Tr=\begin{pmatrix}1&1\\1&1\end{pmatrix}.
\]
\textit{Proof idea:} $\Gamma_\tau^{(Y)}=\Cov(X_{t+\tau}+\phi_{t+\tau}\mathbf1,
X_t+\phi_t\mathbf1)$; independence kills the $X$--$\phi$ cross terms, and
$\Cov(\phi_{t+\tau},\phi_t)=\sigma_\phi^2\delta_{\tau,0}$ contributes
$\sigma_\phi^2\delta_{\tau,0}$ to \emph{every} entry.\\
\textit{Why:} adding a common shock creates lag-$0$ correlation between
$Y^{(1)}$ and $Y^{(2)}$ \emph{regardless} of the correlation in $X$. This is the
prototypical mechanism behind ``everything moves together'' contemporaneous
correlation.
\end{proposition}

\begin{proposition}{Hermitian symmetry of the spectral matrix}{p7-herm}
For all $f\in[-\tfrac12,\tfrac12]$ and all $1\le j,k\le d$,
\[
S_{j,k}(f)=\conj{S_{k,j}(f)},\qquad
S_{j,k}(f)=\conj{S_{j,k}(-f)},
\]
equivalently in matrix form
\[
S(f)=S(f)\Herm,\qquad S(f)=S(-f)\Tr ,
\]
and $S(f)$ is positive semi-definite for every $f$.
\textit{Proof idea:} substitute $\Gamma_{-\tau}=\Gamma_\tau\Tr$ into
$S(f)=\sum_\tau\Gamma_\tau e^{-2\pi i\tau f}$ and relabel $\tau\to-\tau$; the
conjugate of $e^{-2\pi i\tau f}$ flips the sign in the exponent.\\
\textit{Why:} $S(f)\Herm=S(f)$ (Hermitian) guarantees real diagonal spectra and
makes coherence well defined; positive semi-definiteness is the frequency-domain
counterpart of Proposition~\ref{prop:p7-psd}.
\end{proposition}

\begin{theorem}{Multivariate spectral representation theorem}{p7-spectral}
Let $\{X_t\}$ be a vector-valued discrete-time stationary process with mean
$\mu$. There exists a vector-valued orthogonal-increment process $\{Z(f)\}$ on
$[-\tfrac12,\tfrac12]$ such that
\[
X_t=\mu+\int_{-1/2}^{1/2} e^{2\pi i f t}\,dZ(f).
\]
If moreover $\{\acvs_\tau^{(j,k)}\}\in\ell^1$, then for $f,f'\in[-\tfrac12,\tfrac12]$,
\[
\E[dZ(f)]=0,\qquad
\Var\big(dZ(f)\big)=S(f)\,df,\qquad
\Cov\big(dZ(f),dZ(f')\big)=0\ \ (f\ne f'),
\]
where $S(f)$ is the spectral density \emph{matrix} and the $0$ in the last
identity is the zero matrix.\\
\textit{Proof idea:} the multivariate analogue of the scalar Cram\'er/spectral
representation; the increment process $Z(f)$ carries the random ``Fourier
content'' of each component, with covariance structure given by the spectral
matrix.\\
\textit{Why:} it justifies thinking of a multivariate stationary process as a
superposition of uncorrelated random sinusoids whose (matrix) variance is
$S(f)\,df$; coherence is then literally the correlation of the increments
$dZ_j(f),dZ_k(f)$.
\end{theorem}

\begin{proposition}{Cross-spectrum of a pure delay process}{p7-delayspec}
Let $\{Y_t\}$ be stationary and $X_t=(Y_t,\,Y_{t-\nu})\Tr$ for fixed
$\nu\in\Ints$. Then
\[
S_{1,1}(f)=S_{2,2}(f)=\sdf_Y(f),\qquad
S_{1,2}(f)=\sdf_Y(f)\,e^{2\pi i f\nu},
\]
so the coherence is $r_{1,2}(f)=e^{2\pi i f\nu}$, of magnitude $1$ and phase
$2\pi\nu f$.
\textit{Proof idea:} $\acvs_\tau^{(1,2)}=\Cov(Y_{t+\tau},Y_{t-\nu})
=\acvs_{\tau+\nu}^{(Y)}$; Fourier-transforming and shifting the summation index
pulls out $e^{2\pi i f\nu}$.\\
\textit{Why:} a pure time shift between two channels produces \emph{perfect}
coherence (magnitude $1$) with a phase that is \emph{linear in $f$} with slope
$2\pi\nu$ --- the textbook signature of a lead/lag of $\nu$ samples.
\end{proposition}

\begin{theorem}{Everything is a VAR(1) (companion reduction)}{p7-everyvar1}
Any VAR($p$) process $\{X_t\}$ admits the VAR(1) representation
$Y_t=FY_{t-1}+U_t$ of Definition~\ref{def:p7-companion}, with
$\{U_t\}$ a $dp$-variate white noise whose only nonzero covariance block is the
top-left $\Sigma$. Hence all VAR($p$) properties (stability, MA($\infty$)
expansion, covariances) follow from the VAR(1) theory applied to $F$ and then
projected back by $X_t=GY_t$.\\
\textit{Proof idea:} write out the stacked vector $Y_t$; its top block equals
the VAR($p$) recursion and its remaining blocks are trivial identities, giving
exactly $F$ and $U_t$.\\
\textit{Why:} it lets us prove one clean VAR(1) result and reuse it for every
order $p$, e.g.\ the stability test becomes ``eigenvalues of $F$ inside the unit
circle''.
\end{theorem}

\begin{proposition}{MA($\infty$) representation of a stable VAR(1)}{p7-mainf}
If the VAR(1) $X_t=\Phi_1 X_{t-1}+\eps_t$ is stable, then
\[
X_t=\sum_{j=0}^{\infty}\Phi_1^{\,j}\,\eps_{t-j}.
\]
\textit{Proof idea:} back-substitute repeatedly,
$X_t=\Phi_1(\Phi_1 X_{t-2}+\eps_{t-1})+\eps_t=\Phi_1^2X_{t-2}+\Phi_1\eps_{t-1}+\eps_t=\cdots$;
stability ($\Phi_1^j\to0$) makes the remainder term vanish and the series
converge (in mean square).\\
\textit{Why:} the matrix geometric series is the multivariate causal linear
filter of the noise --- it underlies the covariance formula below and all
forecasting.
\end{proposition}

\begin{proposition}{Covariance of a stable VAR(1)}{p7-var1cov}
For a stable VAR(1) with $\{\eps_t\}\sim\mathrm{WN}(0,\Sigma)$ and $\tau\ge0$,
\[
\Gamma_\tau=\Cov(X_{t+\tau},X_t)
=\sum_{k=0}^{\infty}\Phi_1^{\,k+\tau}\,\Sigma\,\big(\Phi_1^{\,k}\big)\Tr .
\]
\textit{Proof idea:} substitute the MA($\infty$) form into
$\Cov\big(\sum_j\Phi_1^{j}\eps_{t+\tau-j},\sum_k\Phi_1^{k}\eps_{t-k}\big)$;
white noise makes only the terms with $t+\tau-j=t-k$, i.e.\ $j=k+\tau$, survive,
each contributing $\Phi_1^{k+\tau}\Sigma(\Phi_1^{k})\Tr$.\\
\textit{Why:} closed-form second-order structure of a VAR(1); for $\tau=0$ it
gives $\Gamma_0=\sum_k\Phi_1^k\Sigma(\Phi_1^k)\Tr$, the stationary covariance,
which also solves the discrete Lyapunov equation
$\Gamma_0=\Phi_1\Gamma_0\Phi_1\Tr+\Sigma$.
\end{proposition}

\begin{corollary}{Covariance of a VAR($p$) via the companion form}{p7-varpcov}
For a stable VAR($p$) with companion matrix $F$ and
$\Sigma_U=\Var(U_t)$ (top-left block $\Sigma$, all else $0$),
\[
\Gamma_\tau^{(X)}=\Cov(X_{t+\tau},X_t)
=G\Big(\sum_{k=0}^{\infty}F^{\,k+\tau}\,\Sigma_U\,\big(F^{\,k}\big)\Tr\Big)G\Tr,
\qquad G=(I_d\ 0\ \cdots\ 0).
\]
\textit{Proof idea:} apply Proposition~\ref{prop:p7-var1cov} to the VAR(1)
$Y_t=FY_{t-1}+U_t$ to get $\Gamma_\tau^{(Y)}$, then use $X_t=GY_t$ so
$\Gamma_\tau^{(X)}=G\Gamma_\tau^{(Y)}G\Tr$.\\
\textit{Why:} reduces the order-$p$ covariance to one VAR(1) computation plus a
projection.
\end{corollary}

\begin{theorem}{Yule--Walker equations for a VAR($p$)}{p7-yw}
For a stable VAR($p$) and $\tau\ge0$,
\[
\Gamma_\tau=\Phi_1\Gamma_{\tau-1}+\Phi_2\Gamma_{\tau-2}
+\cdots+\Phi_p\Gamma_{\tau-p}+\delta_{\tau,0}\,\Sigma .
\]
In particular for VAR(1): $\Gamma_\tau=\Phi_1\Gamma_{\tau-1}+\delta_{\tau,0}\Sigma$.
\textit{Proof idea:} multiply the VAR equation on the right by $X_t\Tr$ and take
expectations: $\Cov(X_{t+\tau},X_t)=\Cov(\sum_i\Phi_i X_{t+\tau-i}+\eps_{t+\tau},X_t)$;
the noise term $\Cov(\eps_{t+\tau},X_t)$ contributes $\Sigma$ only when
$\tau=0$ (since $X_t$ depends only on $\eps_s$ for $s\le t$).\\
\textit{Why:} given $\Sigma$ and the $\Phi_j$, solving the first $p+1$ equations
yields $\Gamma_0,\dots,\Gamma_p$, and the recursion then generates all
$\Gamma_\tau$ for $\tau>p$. Run in reverse, it estimates the $\Phi_j$ from sample
covariances (VAR fitting).
\end{theorem}

\begin{intuition}[Why the noise term only shows up at $\tau=0$]
$X_t$ is built from $\eps_t,\eps_{t-1},\dots$ (the causal MA($\infty$)). For
$\tau>0$, $\eps_{t+\tau}$ is in the \emph{future} of $X_t$, hence uncorrelated
with it, so $\Cov(\eps_{t+\tau},X_t)=0$. Only at $\tau=0$ does
$\Cov(\eps_t,X_t)=\Cov(\eps_t,\eps_t)=\Sigma$ appear. This is the exact
multivariate analogue of the scalar Yule--Walker derivation.
\end{intuition}

% ----------------------------------------------------------------------
\subsection{Key Techniques}

\begin{keytech}{Test VAR stability via the reverse characteristic polynomial}{p7-ktstable}
Given $\Phi_1,\dots,\Phi_p$:
\begin{enumerate}
  \item Form $M(z)=I_d-\Phi_1 z-\cdots-\Phi_p z^p$ (a $d\times d$ matrix whose
        entries are polynomials in $z$).
  \item Compute $\det M(z)$ --- a scalar polynomial of degree at most $dp$.
        \emph{Factor it} if possible (block/triangular structure often factors
        it cleanly).
  \item Find all roots $z_i$. The VAR is \textbf{stable} iff $\abs{z_i}>1$ for
        \emph{every} root.
\end{enumerate}
\textbf{Shortcut for VAR(1):} $\det(I_d-\Phi_1 z)=0\iff z=1/\lambda_i(\Phi_1)$,
so just check $\abs{\lambda_i(\Phi_1)}<1$ (spectral radius of $\Phi_1$ below
one). For a triangular $\Phi_1$, the eigenvalues are the diagonal entries.
\emph{Sanity check:} $\det M(0)=\det I_d=1$ always (the constant term is $1$).
\end{keytech}

\begin{keytech}{Compute a cross-spectrum from a cross-covariance}{p7-ktcrossspec}
To get $S_{j,k}(f)$:
\begin{enumerate}
  \item Find the cross-covariance $\acvs_\tau^{(j,k)}=\Cov(X_{t+\tau}^{(j)},X_t^{(k)})$
        in closed form, usually a \emph{shift} of a known (auto)covariance, e.g.\
        $\acvs_\tau^{(j,k)}=\alpha\,\acvs_{\tau-\nu}^{(k,k)}$.
  \item Fourier transform: $S_{j,k}(f)=\sum_\tau\acvs_\tau^{(j,k)}e^{-2\pi if\tau}$,
        then \emph{re-index} the sum to extract any $e^{\pm2\pi if\nu}$ factor.
        A lag-shift $\tau\to\tau-\nu$ produces a phase factor $e^{-2\pi if\nu}$.
  \item Read off amplitude $\abs{S_{j,k}(f)}$ and phase $\arg S_{j,k}(f)$; divide
        by $\sqrt{S_{j,j}S_{k,k}}$ for coherence.
\end{enumerate}
Mnemonic: \emph{a delay of $\nu$ samples in the time domain $=$ a linear phase
$\mp2\pi\nu f$ in the frequency domain}.
\end{keytech}

\begin{keytech}{Derive the VAR(1) MA($\infty$) and its covariance}{p7-ktmainf}
For a stable VAR(1) $X_t=\Phi_1 X_{t-1}+\eps_t$:
\begin{enumerate}
  \item Write $X_t=\sum_{j\ge0}\Phi_1^j\eps_{t-j}$ (matrix geometric filter).
  \item Plug into $\Cov(X_{t+\tau},X_t)$; keep only index pairs with matching
        noise times ($j=k+\tau$). Result:
        $\Gamma_\tau=\sum_{k\ge0}\Phi_1^{k+\tau}\Sigma(\Phi_1^k)\Tr$.
  \item For a \emph{numerical} $\Gamma_0$, solve the Lyapunov equation
        $\Gamma_0=\Phi_1\Gamma_0\Phi_1\Tr+\Sigma$ (vectorise:
        $\mathrm{vec}\,\Gamma_0=(I-\Phi_1\otimes\Phi_1)^{-1}\mathrm{vec}\,\Sigma$),
        or sum the series until $\Phi_1^k$ is negligible.
\end{enumerate}
If $\Phi_1$ and $\Sigma$ are diagonal, everything decouples into $d$ scalar
AR(1)'s: $\Gamma_0^{(j,j)}=\sigma_j^2/(1-\phi_j^2)$.
\end{keytech}

\begin{keytech}{Run the VAR Yule--Walker recursion}{p7-ktyw}
Given $\Phi_1,\dots,\Phi_p,\Sigma$ and known $\Gamma_0,\dots,\Gamma_{p-1}$:
\begin{itemize}
  \item For $\tau>0$: $\Gamma_\tau=\sum_{i=1}^p\Phi_i\Gamma_{\tau-i}$, using
        $\Gamma_{-s}=\Gamma_s\Tr$ whenever a negative-lag matrix appears.
  \item To \emph{initialise}, write the $\tau=0,\dots,p$ equations as a linear
        system in $\Gamma_0,\dots,\Gamma_p$ (the $\tau=0$ equation carries the
        $+\Sigma$ term) and solve once.
\end{itemize}
Watch the transpose: a term like $\Phi_2\Gamma_{-1}$ means $\Phi_2\Gamma_1\Tr$.
\end{keytech}

% ----------------------------------------------------------------------
\subsection{Exercises}

\begin{exercise}{Common-noise cross-correlation \quad\textnormal{[Sheet 9]}}{p7-e1}
Let $Y_t=(Y_t^{(1)},Y_t^{(2)})\Tr$ with $Y_t^{(j)}=X_t^{(j)}+\eps_t$, $j=1,2$,
where $X^{(1)},X^{(2)},\eps$ are mutually independent and $\{\eps_t\}$ is white
noise. Write $\sigma_j^2=\Var(X_t^{(j)})$, $\sigma_\eps^2=\Var(\eps_t)$ and
$\tilde\acvs_\tau^{(1,2)}=\Cov(X_{t+\tau}^{(1)},X_t^{(2)})$. Find the
cross-correlation $\acf_\tau^{(1,2)}$ between $\{Y^{(1)}\}$ and $\{Y^{(2)}\}$ at
all lags.
\end{exercise}
\begin{solution}
The \emph{same} scalar noise $\eps_t$ is added to both channels, so expand
\[
\acvs_\tau^{(1,2)}=\Cov(Y_{t+\tau}^{(1)},Y_t^{(2)})
=\Cov(X_{t+\tau}^{(1)},X_t^{(2)})+\underbrace{\Cov(X_{t+\tau}^{(1)},\eps_t)}_{0}
+\underbrace{\Cov(\eps_{t+\tau},X_t^{(2)})}_{0}+\Cov(\eps_{t+\tau},\eps_t).
\]
The two middle terms vanish by independence of $X$ and $\eps$, and
$\Cov(\eps_{t+\tau},\eps_t)=\sigma_\eps^2\delta_{\tau,0}$. Hence
\[
\acvs_\tau^{(1,2)}=
\begin{cases}\tilde\acvs_0^{(1,2)}+\sigma_\eps^2 & \tau=0,\\[2pt]
\tilde\acvs_\tau^{(1,2)} & \tau\ne0.\end{cases}
\]
For the marginal variances,
$\acvs_0^{(1,1)}=\Var(X_t^{(1)})+\Var(\eps_t)=\sigma_1^2+\sigma_\eps^2$ and
likewise $\acvs_0^{(2,2)}=\sigma_2^2+\sigma_\eps^2$. Therefore
\[
\acf_\tau^{(1,2)}=\frac{\acvs_\tau^{(1,2)}}{\sqrt{(\sigma_1^2+\sigma_\eps^2)(\sigma_2^2+\sigma_\eps^2)}}
=
\begin{cases}
\dfrac{\tilde\acvs_0^{(1,2)}+\sigma_\eps^2}{\sqrt{(\sigma_1^2+\sigma_\eps^2)(\sigma_2^2+\sigma_\eps^2)}} & \tau=0,\\[10pt]
\dfrac{\tilde\acvs_\tau^{(1,2)}}{\sqrt{(\sigma_1^2+\sigma_\eps^2)(\sigma_2^2+\sigma_\eps^2)}} & \tau\ne0.
\end{cases}
\]
The shared noise injects extra correlation \emph{only} at lag $0$ (the
$+\sigma_\eps^2$), exactly Proposition~\ref{prop:p7-contemp} in disguise.
\end{solution}

\begin{exercise}{Delay process: cross-spectrum amplitude \& phase \quad\textnormal{[Sheet 9]}}{p7-e2}
For the delay process with $\alpha>0$, $\nu\in\Ints$,
\[
X_t^{(1)}=\alpha X_{t-\nu}^{(2)}+\eps_t,\qquad t\in\Ints,
\]
where $\{\eps_t\}$ is white noise (variance $\sigma_\eps^2$) and
$\{X_t^{(2)}\}$ is stationary and uncorrelated with $\{\eps_t\}$ at all lags.
\textbf{(a)} Find the cross-correlation between $\{X^{(1)}\}$ and $\{X^{(2)}\}$
at all lags. \textbf{(b)} Find the cross-spectrum $S_{1,2}(f)$ and its amplitude
and phase.
\end{exercise}
\begin{solution}
\textbf{(a)} Using independence of $X^{(2)}$ and $\eps$,
\[
\acvs_\tau^{(1,2)}=\Cov(X_{t+\tau}^{(1)},X_t^{(2)})
=\alpha\,\Cov(X_{t+\tau-\nu}^{(2)},X_t^{(2)})+\underbrace{\Cov(\eps_{t+\tau},X_t^{(2)})}_{0}
=\alpha\,\acvs_{\tau-\nu}^{(2,2)} .
\]
The marginal variances are
$\acvs_0^{(1,1)}=\Var(\alpha X_{t-\nu}^{(2)}+\eps_t)=\alpha^2\sigma_2^2+\sigma_\eps^2$
(with $\sigma_2^2=\acvs_0^{(2,2)}$) and $\acvs_0^{(2,2)}=\sigma_2^2$. Hence
\[
\acf_\tau^{(1,2)}=\frac{\alpha\,\acvs_{\tau-\nu}^{(2,2)}}
{\sqrt{\alpha^2\sigma_2^2+\sigma_\eps^2}\,\sqrt{\sigma_2^2}}
=\frac{\acf_{\tau-\nu}^{(2,2)}}{\sqrt{1+\sigma_\eps^2/(\alpha\sigma_2)^2}} ,
\]
i.e.\ the autocorrelation of $X^{(2)}$ shifted by $\nu$ and shrunk by the noise
factor.\\[4pt]
\textbf{(b)} Fourier-transform the cross-covariance and shift the index
$\tau\to\tau+\nu$:
\[
\begin{aligned}
S_{1,2}(f)
&=\sum_{\tau=-\infty}^{\infty}\acvs_\tau^{(1,2)}e^{-2\pi if\tau}
=\alpha\sum_{\tau}\acvs_{\tau-\nu}^{(2,2)}e^{-2\pi if\tau}\\
&=\alpha\,e^{-2\pi if\nu}\sum_{\tau}\acvs_{\tau}^{(2,2)}e^{-2\pi if\tau}
=\alpha\,S_{2,2}(f)\,e^{-2\pi if\nu}.
\end{aligned}
\]
Since $S_{2,2}(f)\ge0$ is real (the spectral density of $X^{(2)}$), the
\textbf{amplitude} is
$\abs{S_{1,2}(f)}=\alpha\,S_{2,2}(f)$ (the spectrum of $X^{(2)}$ scaled by
$\alpha$) and the \textbf{phase} is $\arg S_{1,2}(f)=-2\pi\nu f$, linear in $f$
with slope $-2\pi\nu$ --- the delay signature.
\end{solution}

\begin{exercise}{Expectation \& variance of the sample cross-covariance \quad\textnormal{[Sheet 9]}}{p7-e3}
Let $Y_t=(Y_t^{(1)},Y_t^{(2)})\Tr\overset{\text{ind}}{\sim}\mathcal N(0,\Sigma)$
in $t$, with $\Sigma=\begin{psmallmatrix}\sigma_{11}&\sigma_{12}\\
\sigma_{21}&\sigma_{22}\end{psmallmatrix}$. With
$\hat\Gamma_\tau=\frac1N\sum_{t=1}^{N-\abs\tau}Y_{t+\tau}Y_t\Tr$, determine
$\E[\hat\acvs_\tau^{(1,2)}]$ and $\Var(\hat\acvs_\tau^{(1,2)})$.
\end{exercise}
\begin{solution}
\textbf{Expectation.} By linearity,
\[
\E[\hat\acvs_\tau^{(1,2)}]=\frac1N\sum_{t=1}^{N-\abs\tau}\E[Y_{t+\tau}^{(1)}Y_t^{(2)}].
\]
Since the $Y_t$ are independent across $t$, $\E[Y_{t+\tau}^{(1)}Y_t^{(2)}]=0$ for
$\tau\ne0$, and $=\sigma_{12}$ for $\tau=0$. With $N-\abs\tau=N$ terms at
$\tau=0$,
\[
\E[\hat\acvs_\tau^{(1,2)}]=\sigma_{12}\,\delta_{\tau,0}.
\]
\textbf{Variance.} Write
$\Var(\hat\acvs_\tau^{(1,2)})=\E\big[(\hat\acvs_\tau^{(1,2)})^2\big]-\E[\hat\acvs_\tau^{(1,2)}]^2$
and expand the square:
\[
\E\big[(\hat\acvs_\tau^{(1,2)})^2\big]
=\frac{1}{N^2}\Big(\sum_{t=1}^{N-\abs\tau}\E\big[(Y_{t+\tau}^{(1)})^2(Y_t^{(2)})^2\big]
+\sum_{t\ne s}\E\big[Y_{t+\tau}^{(1)}Y_t^{(2)}Y_{s+\tau}^{(1)}Y_s^{(2)}\big]\Big).
\]
For jointly Gaussian zero-mean variables, Isserlis' (Wick's) theorem gives the
key fourth moment
\[
\E\big[(Y_t^{(1)})^2(Y_t^{(2)})^2\big]=\sigma_{11}\sigma_{22}+2\sigma_{12}^2 .
\]
\emph{Case $\tau=0$.} The diagonal sum has $N$ identical terms
$\sigma_{11}\sigma_{22}+2\sigma_{12}^2$. For the off-diagonal ($t\ne s$),
independence across $t$ factorises the expectation into
$\E[Y_t^{(1)}Y_t^{(2)}]\,\E[Y_s^{(1)}Y_s^{(2)}]=\sigma_{12}^2$, and there are
$N(N-1)$ such pairs. Thus
\[
\E\big[(\hat\acvs_0^{(1,2)})^2\big]
=\frac1{N^2}\big[N(\sigma_{11}\sigma_{22}+2\sigma_{12}^2)+N(N-1)\sigma_{12}^2\big]
=\frac{\sigma_{11}\sigma_{22}+2\sigma_{12}^2}{N}+\frac{N-1}{N}\sigma_{12}^2 .
\]
Subtracting $\E[\hat\acvs_0^{(1,2)}]^2=\sigma_{12}^2$,
\[
\Var(\hat\acvs_0^{(1,2)})
=\frac{\sigma_{11}\sigma_{22}+2\sigma_{12}^2}{N}+\frac{N-1}{N}\sigma_{12}^2-\sigma_{12}^2
=\frac{1}{N}\big(\sigma_{11}\sigma_{22}+\sigma_{12}^2\big).
\]
\emph{Case $\tau\ne0$.} Now $\E[\hat\acvs_\tau^{(1,2)}]=0$. The diagonal sum has
$N-\abs\tau$ terms, but because $Y_{t+\tau}$ and $Y_t$ are \emph{independent}
(different time indices), the Gaussian fourth moment splits as
$\E[(Y_{t+\tau}^{(1)})^2]\,\E[(Y_t^{(2)})^2]=\sigma_{11}\sigma_{22}$. Every
off-diagonal term contains a factor with a unique time index appearing once and
so vanishes. Hence
\[
\Var(\hat\acvs_\tau^{(1,2)})
=\frac{1}{N^2}\big[(N-\abs\tau)\sigma_{11}\sigma_{22}+0\big]
=\frac{N-\abs\tau}{N^2}\,\sigma_{11}\sigma_{22}.
\]
\textbf{Summary.}
\[
\E[\hat\acvs_\tau^{(1,2)}]=\sigma_{12}\,\delta_{\tau,0},\qquad
\Var(\hat\acvs_\tau^{(1,2)})=\frac{N-\abs\tau}{N^2}\big(\sigma_{11}\sigma_{22}+\sigma_{12}^2\,\delta_{\tau,0}\big).
\]
(At $\tau=0$ this reads $\tfrac1N(\sigma_{11}\sigma_{22}+\sigma_{12}^2)$,
matching the case above.)
\end{solution}

\begin{exercise}{Stability of a bivariate VAR(2) \quad\textnormal{[New]}}{p7-e4}
Consider $X_t=\Phi_1 X_{t-1}+\Phi_2 X_{t-2}+\eps_t$ with
$\Phi_1=\begin{psmallmatrix}0.5&0.1\\0.4&0.5\end{psmallmatrix}$,
$\Phi_2=\begin{psmallmatrix}0&0\\0.25&0\end{psmallmatrix}$. Form the reverse
characteristic polynomial, find its roots, and decide whether the VAR(2) is
stable.
\end{exercise}
\begin{solution}
We need $\det\!\big(I_2-\Phi_1 z-\Phi_2 z^2\big)$. The inner matrix is
\[
M(z)=\begin{pmatrix}1-0.5z & -0.1z\\ -0.4z-0.25z^2 & 1-0.5z\end{pmatrix}.
\]
Its determinant is
\[
\det M(z)=(1-0.5z)^2-(-0.1z)(-0.4z-0.25z^2)
=(1-z+0.25z^2)-(0.04z^2+0.025z^3),
\]
\[
\det M(z)=1-z+0.21z^2-0.025z^3 .
\]
Finding the roots of $1-z+0.21z^2-0.025z^3=0$ gives
\[
z_1=1.3,\qquad z_2=3.55+4.26i,\qquad z_3=3.55-4.26i .
\]
The moduli are $\abs{z_1}=1.3>1$ and
$\abs{z_2}=\abs{z_3}=\sqrt{3.55^2+4.26^2}=\sqrt{12.60+18.15}=\sqrt{30.75}\approx5.545>1$.
All three roots lie strictly outside the unit circle, so the VAR(2) is
\textbf{stable} (hence stationary). \emph{Check:} $\det M(0)=1$, as it must be
(the constant term equals $\det I_2=1$).
\end{solution}

\begin{exercise}{Stability of a triangular trivariate VAR(1) \quad\textnormal{[New]}}{p7-e5}
Let $X_t=\Phi_1 X_{t-1}+\eps_t$ with
$\Phi_1=\begin{psmallmatrix}0.5&0&0\\0.1&0.1&0.3\\0&0.2&0.3\end{psmallmatrix}$.
Factor the reverse characteristic polynomial and decide stability.
\end{exercise}
\begin{solution}
Because the first row of $\Phi_1$ has only its $(1,1)$ entry nonzero, expand
$\det(I_3-\Phi_1 z)$ along the first row of
\[
I_3-\Phi_1 z=\begin{pmatrix}1-0.5z & 0 & 0\\ -0.1z & 1-0.1z & -0.3z\\ 0 & -0.2z & 1-0.3z\end{pmatrix}.
\]
This gives
\[
\det(I_3-\Phi_1 z)=(1-0.5z)\det\begin{pmatrix}1-0.1z & -0.3z\\ -0.2z & 1-0.3z\end{pmatrix}
=(1-0.5z)\big[(1-0.1z)(1-0.3z)-0.06z^2\big].
\]
The $2\times2$ determinant is $1-0.4z+0.03z^2-0.06z^2=1-0.4z-0.03z^2$, so
\[
\det(I_3-\Phi_1 z)=(1-0.5z)\big(1-0.4z-0.03z^2\big).
\]
The first factor gives $z_1=2$. The quadratic $1-0.4z-0.03z^2=0$, i.e.\
$0.03z^2+0.4z-1=0$, has roots
$z=\dfrac{-0.4\pm\sqrt{0.16+0.12}}{0.06}=\dfrac{-0.4\pm\sqrt{0.28}}{0.06}$,
giving $z_2\approx2.153$ and $z_3\approx-15.486$. All three roots satisfy
$\abs{z}>1$, so the VAR(1) is \textbf{stable}. (Equivalently: $\Phi_1$ is block
lower-triangular with eigenvalues $0.5$ and those of
$\begin{psmallmatrix}0.1&0.3\\0.2&0.3\end{psmallmatrix}$, namely
$\approx0.465$ and $\approx-0.065$; all have modulus $<1$.)
\end{solution}

\begin{exercise}{MA($\infty$) covariance of a diagonal VAR(1) \quad\textnormal{[New]}}{p7-e6}
Let $X_t=\Phi_1 X_{t-1}+\eps_t$ with
$\Phi_1=\begin{psmallmatrix}0.5&0\\0&0.3\end{psmallmatrix}$ and
$\{\eps_t\}\sim\mathrm{WN}(0,I_2)$. (a) Write the MA($\infty$) form. (b) Compute
the stationary covariance $\Gamma_0$ and the lag-$1$ matrix $\Gamma_1$. (c) What
is the coherence between the two channels?
\end{exercise}
\begin{solution}
\textbf{(a)} Since $\Phi_1=\diag(0.5,0.3)$ has spectral radius $0.5<1$, the VAR
is stable and
\[
X_t=\sum_{j=0}^\infty\Phi_1^j\eps_{t-j},\qquad
\Phi_1^j=\diag(0.5^j,\,0.3^j).
\]
\textbf{(b)} By the VAR(1) covariance formula with $\Sigma=I_2$,
\[
\Gamma_0=\sum_{k=0}^\infty\Phi_1^k(\Phi_1^k)\Tr
=\sum_{k=0}^\infty\diag(0.25^k,\,0.09^k)
=\diag\!\Big(\tfrac{1}{1-0.25},\,\tfrac{1}{1-0.09}\Big)
=\diag\!\Big(\tfrac43,\,\tfrac{100}{91}\Big),
\]
i.e.\ $\Gamma_0\approx\diag(1.333,\,1.099)$. (Equivalently solve
$\Gamma_0=\Phi_1\Gamma_0\Phi_1\Tr+I_2$ componentwise:
$g_{jj}=\phi_j^2 g_{jj}+1$.) The lag-$1$ matrix is, by Yule--Walker
$\Gamma_1=\Phi_1\Gamma_0$,
\[
\Gamma_1=\diag(0.5,0.3)\,\diag(\tfrac43,\tfrac{100}{91})
=\diag\!\Big(\tfrac23,\,\tfrac{30}{91}\Big)\approx\diag(0.667,\,0.330).
\]
\textbf{(c)} All off-diagonal entries of $\Gamma_\tau$ are $0$ for every $\tau$
(the channels are independent AR(1)'s), so $S_{1,2}(f)\equiv0$ and the coherence
$r_{1,2}(f)=0$ at all frequencies: the two series are completely uncoupled. The
off-diagonal entries of $\Phi_1$ (here zero) are exactly the cross-channel
dynamics, confirming the decoupling.
\end{solution}

\begin{exercise}{Reading off a VAR(1) from its Yule--Walker recursion \quad\textnormal{[New]}}{p7-e7}
A stable bivariate VAR(1) $X_t=\Phi_1 X_{t-1}+\eps_t$ has
$\Phi_1=\begin{psmallmatrix}0.6&0.2\\0.1&0.5\end{psmallmatrix}$ and
$\{\eps_t\}\sim\mathrm{WN}(0,I_2)$. (a) Verify stability via the reverse
characteristic polynomial. (b) Express $\Gamma_2$ in terms of $\Gamma_0$ using
Yule--Walker. (c) Give $\Gamma_{-1}$ in terms of $\Gamma_1$.
\end{exercise}
\begin{solution}
\textbf{(a)} $\det(I_2-\Phi_1 z)=(1-0.6z)(1-0.5z)-(0.2z)(0.1z)
=1-1.1z+0.30z^2-0.02z^2=1-1.1z+0.28z^2$. Solving $0.28z^2-1.1z+1=0$,
\[
z=\frac{1.1\pm\sqrt{1.21-1.12}}{0.56}=\frac{1.1\pm0.3}{0.56}
\;\Rightarrow\; z_1=\frac{1.4}{0.56}=2.5,\quad z_2=\frac{0.8}{0.56}=\frac{10}{7}\approx1.43 .
\]
Both have modulus $>1$, so the VAR(1) is stable. (Equivalently $\Phi_1$ has
eigenvalues $0.7$ and $0.4$, both inside the unit circle.)\\
\textbf{(b)} Yule--Walker for VAR(1) gives $\Gamma_\tau=\Phi_1\Gamma_{\tau-1}$
for $\tau\ge1$. Iterating,
\[
\Gamma_1=\Phi_1\Gamma_0,\qquad
\Gamma_2=\Phi_1\Gamma_1=\Phi_1^2\,\Gamma_0,
\qquad \Phi_1^2=\begin{pmatrix}0.38&0.22\\0.11&0.27\end{pmatrix}.
\]
\textbf{(c)} By the transpose symmetry $\Gamma_{-\tau}=\Gamma_\tau\Tr$,
\[
\Gamma_{-1}=\Gamma_1\Tr=(\Phi_1\Gamma_0)\Tr=\Gamma_0\Phi_1\Tr
\]
(using $\Gamma_0=\Gamma_0\Tr$). Note $\Gamma_{-1}\ne\Gamma_1$ in general, since
$\Phi_1$ is not symmetric --- the multivariate ACVS is \emph{not} even in
$\tau$.
\end{solution}

% ======================================================================
%  PART 8 -- Forecasting
% ======================================================================
\section{Forecasting}
\label{sec:forecasting}

The purpose of time series analysis is often to \emph{forecast} future values
from the data observed up to now. The recipe has three steps: (1) assume the
form of the model; (2) estimate its parameters; (3) use the assumed structure
and the estimated parameters to form predictions. In this part we take the
model structure and the parameters as \textbf{known} and predict $h$ lags
ahead, i.e.\ we want $X_{n+h}$ from $X_1,\dots,X_n$. The integer $h>0$ is the
\textbf{lead} (or \textbf{horizon}) and $n$ is the \textbf{origin}. The central
theme is that, for the second-order theory we use, the optimal predictor depends
\emph{only} on the covariance structure $\{\acvs_\tau\}$, and for causal ARMA
models it is computed cleanly by setting future innovations to zero.

% ----------------------------------------------------------------------
\subsection{Definitions}

\begin{definition}{Prediction mean square error}{p8-pmse}
Let $g(X_1,\dots,X_n)$ be any estimator (predictor) of $X_{n+h}$ based on the
observations $X_1,\dots,X_n$. Its performance is measured by the
\textbf{prediction mean square error}
\[
P^{\,n}_{n+h}\;:=\;\E\!\left[\big(X_{n+h}-g(X_1,\dots,X_n)\big)^2\right].
\]
Smaller $P^{\,n}_{n+h}$ means a better forecast; minimising it over a class of
predictors defines the ``best'' predictor in that class.
\end{definition}

\begin{definition}{Best linear predictor}{p8-blp}
The \textbf{best linear predictor} of $X_{n+h}$ based on $X_1,\dots,X_n$ is the
linear function
\[
X^{\,n}_{n+h}\;=\;P_{X_1,\dots,X_n}(X_{n+h})\;=\;\sum_{j=1}^{n}\beta_j X_j
\]
that minimises the prediction mean square error among all linear combinations of
$X_1,\dots,X_n$. We write $P_{X_1,\dots,X_n}(\cdot)$ for the projection onto the
linear span of the data.
\begin{itemize}
  \item For \textbf{Gaussian} time series the conditional expectation
        $\E[X_{n+h}\mid X_1,\dots,X_n]$ is itself linear in the data, so the best
        linear predictor coincides with the (overall) best predictor.
  \item In general the conditional expectation is a \emph{nonlinear} function of
        the data; for tractability one restricts attention to linear predictors,
        which need only second-order information.
\end{itemize}
\end{definition}

\begin{definition}{$h$-step forecast error and its variance}{p8-ferr}
For the best linear predictor $X^{\,n}_{n+h}$ define the \textbf{$h$-step ahead
forecast error}
\[
e_n(h)\;:=\;X_{n+h}-X^{\,n}_{n+h}.
\]
It is unbiased, $\E[e_n(h)]=0$, and its variance equals the prediction MSE,
$\Var\big(e_n(h)\big)=P^{\,n}_{n+h}$. For a causal ARMA process with
$MA(\infty)$ weights $\{\psi_j\}$ one has
$e_n(h)=\sum_{j=0}^{h-1}\psi_j\eps_{n+h-j}$ and
$\Var\big(e_n(h)\big)=\sigma^2\sum_{j=0}^{h-1}\psi_j^2$.
\end{definition}

\begin{definition}{Sources of prediction uncertainty}{p8-srcunc}
There are \textbf{three} components of uncertainty in a forecast:
\begin{enumerate}
  \item \textbf{model uncertainty} --- the assumed model form may be wrong;
  \item \textbf{estimation uncertainty} --- the parameters are estimated, not known;
  \item \textbf{innovation uncertainty} --- the inherent randomness of the future
        innovations $\eps_{n+1},\dots,\eps_{n+h}$.
\end{enumerate}
Throughout we assume the model is \emph{known} and the parameters are estimated
\emph{without error}, which removes the first two sources and leaves only
innovation uncertainty. (Bootstrap methods can incorporate the first two, but
are not covered here.)
\end{definition}

\begin{definition}{Gaussian prediction interval}{p8-pint}
For a \textbf{Gaussian} process the $h$-step error $e_n(h)$ is normal with mean
$0$ and variance $\Var(e_n(h))$, so a $(1-\alpha)$ \textbf{prediction interval}
for $X_{n+h}$ based on $X_1,\dots,X_n$ is
\[
X^{\,n}_{n+h}\;\pm\;z_{1-\alpha/2}\,\sqrt{\Var\big(e_n(h)\big)},
\]
where $z_{1-\alpha/2}$ is the $(1-\alpha/2)$ quantile of the standard normal
distribution (e.g.\ $z_{0.975}\approx1.96$ for a $95\%$ interval).
\end{definition}

\begin{intuition}
A point forecast alone is almost useless without a measure of its spread. The
interval widens with $h$ because more unknown innovations
$\eps_{n+1},\dots,\eps_{n+h}$ accumulate in the error. For a stationary ARMA
model the width saturates: as $h\to\infty$, $\Var(e_n(h))\to\acvs_0$, i.e.\ the
interval converges to ``mean $\pm$ a multiple of the marginal standard
deviation'' --- forecasting far ahead is no better than quoting the
unconditional distribution. For integrated (ARIMA) processes the interval grows
\emph{without bound}.
\end{intuition}

% ----------------------------------------------------------------------
\subsection{Theorems \& Propositions}

\begin{lemma}{Conditional expectation minimises the prediction MSE}{p8-condexp}
The conditional expectation
\[
g(X_1,\dots,X_n)=\E\big[X_{n+h}\mid X_1,\dots,X_n\big]
\]
minimises the prediction mean square error for $X_{n+h}$ based on
$X_1,\dots,X_n$.\\[3pt]
\textit{Proof idea:} For any real random variable $Y$ with $\mu=\E[Y]$ and any
constant $c$,
\[
\MSE(c)=\E\big[(Y-c)^2\big]=\E\big[(Y-\mu+\mu-c)^2\big]
=\E\big[(Y-\mu)^2\big]+(\mu-c)^2=\Var(Y)+(\mu-c)^2,
\]
which is minimised at $c=\mu=\E[Y]$ (the cross term vanishes because
$\E[Y-\mu]=0$). Applying this \emph{conditionally on} $X_1,\dots,X_n$ and using
the tower property
$\E[(X_{n+h}-g)^2]=\E\big[\E[(X_{n+h}-g)^2\mid X_1,\dots,X_n]\big]$, the inner
MSE is minimised pointwise by taking $g=\E[X_{n+h}\mid X_1,\dots,X_n]$.\\[2pt]
\textit{Why:} this is the theoretical gold standard for forecasting. For
Gaussian processes the conditional expectation is linear, so the best predictor
\emph{is} the best linear predictor; in general it motivates restricting to
linear predictors, which only need the ACVS.
\end{lemma}

\begin{example}{Gaussian AR(1) predictor --- mean reversion}{p8-ar1}
Let $X_t=\phi X_{t-1}+\eps_t$ with $\eps_t\iid\mathcal N(0,\sigma^2)$ and
$\abs\phi<1$ (stationary). Because the process is Gaussian, the best linear
predictor equals the conditional expectation, and iterating the recursion,
\[
X^{\,n}_{n+h}=\E[X_{n+h}\mid X_1,\dots,X_n]
=\E[\phi X_{n+h-1}+\eps_{n+h}\mid X_1,\dots,X_n]
=\phi^{\,h}X_n .
\]
Since $\abs\phi<1$, $X^{\,n}_{n+h}=\phi^{\,h}X_n\to0$ (the process mean) as
$h\to\infty$: the forecast \textbf{reverts to the mean}. This is a general
property of stationary ARMA models. (A worked figure in the notes uses
$\phi=0.8$, $\sigma^2=1$.)
\end{example}

\begin{example}{Gaussian MA(1) predictor}{p8-ma1}
Let $X_t=\mu+\eps_t-\theta\eps_{t-1}$ with $\eps_t\iid\mathcal N(0,\sigma^2)$.
For $h>0$,
\[
X^{\,n}_{n+h}=\mu+\E[\eps_{n+h}\mid X_1,\dots,X_n]-\theta\,\E[\eps_{n+h-1}\mid X_1,\dots,X_n]
=\mu-\theta\,\E[\eps_{n+h-1}\mid X_1,\dots,X_n],
\]
because $\eps_t$ for $t>n$ is uncorrelated with $X_1,\dots,X_n$ (so its
conditional expectation is $0$). Hence:
\begin{itemize}
  \item $h=1$: $\;X^{\,n}_{n+1}=\mu-\theta\,\hat\eps_n$ (the one observable
        innovation $\eps_n$ remains);
  \item $h>1$: $\;X^{\,n}_{n+h}=\mu$ --- the forecast hits the mean immediately,
        consistent with an MA(1) having memory only of length $1$.
\end{itemize}
\end{example}

\begin{theorem}{Prediction equations (zero-mean stationary process)}{p8-predeq}
For a zero-mean stationary process $\{X_t\}$, the best linear predictor
$X^{\,n}_{n+h}=\sum_{j=1}^n\beta_j X_j$ is found by solving for
$\beta_1,\dots,\beta_n$ the \textbf{prediction equations}
\[
\E\!\big[(X_{n+h}-X^{\,n}_{n+h})\,X_k\big]=0,\qquad k=1,\dots,n,
\]
i.e.\ the forecast error is orthogonal to every observation. Equivalently, in
matrix form $\Gamma_n\beta=\acvs_{[h]}$:
\[
\begin{pmatrix}
\acvs_0 & \acvs_1 & \cdots & \acvs_{n-1}\\
\acvs_1 & \acvs_0 & \cdots & \acvs_{n-2}\\
\vdots & \vdots & \ddots & \vdots\\
\acvs_{n-1} & \acvs_{n-2} & \cdots & \acvs_0
\end{pmatrix}
\begin{pmatrix}\beta_1\\ \beta_2\\ \vdots\\ \beta_n\end{pmatrix}
=\begin{pmatrix}\acvs_{n+h-1}\\ \acvs_{n+h-2}\\ \vdots\\ \acvs_{h}\end{pmatrix}.
\]
\textit{Proof idea:} the best linear predictor is the orthogonal projection of
$X_{n+h}$ onto $\mathrm{span}\{X_1,\dots,X_n\}$; the residual
$X_{n+h}-X^{\,n}_{n+h}$ must be orthogonal (uncorrelated) to each $X_k$.
Expanding $\E[(X_{n+h}-\sum_j\beta_jX_j)X_k]=0$ gives
$\acvs_{n+h-k}=\sum_{j=1}^n\beta_j\acvs_{j-k}$, which is row $k$ of
$\Gamma_n\beta=\acvs_{[h]}$.\\[2pt]
\textit{Why:} this is the master equation for linear prediction of \emph{any}
stationary process --- it shows the optimal forecast uses only the ACVS.
\end{theorem}

\begin{corollary}{Matrix solution and prediction MSE}{p8-matsol}
If $\Gamma_n$ is invertible, then with $X=(X_1,\dots,X_n)\Tr$ the best linear
predictor and its MSE are
\[
X^{\,n}_{n+h}=\acvs_{[h]}\Tr\,\Gamma_n^{-1}\,X,
\qquad
P^{\,n}_{n+h}=\acvs_0-\acvs_{[h]}\Tr\,\Gamma_n^{-1}\,\acvs_{[h]} .
\]
\textit{Proof idea:} solve $\Gamma_n\beta=\acvs_{[h]}$ for
$\beta=\Gamma_n^{-1}\acvs_{[h]}$ and substitute into
$X^{\,n}_{n+h}=\beta\Tr X$ and into
$P^{\,n}_{n+h}=\acvs_0-2\beta\Tr\acvs_{[h]}+\beta\Tr\Gamma_n\beta$; the last
two terms collapse using $\Gamma_n\beta=\acvs_{[h]}$.\\[2pt]
\textit{Why:} gives the closed-form forecast for short series.
\begin{itemize}
  \item $\Gamma_n$ is invertible for \emph{invertible} ARMA models (and not in
        all other cases), but the best prediction $X^{\,n}_{n+h}$ is
        \textbf{always unique} even when $\Gamma_n$ is singular.
  \item Computing $\acvs_{[h]}\Tr\Gamma_n^{-1}X$ is inefficient for large $n$
        because the $n\times n$ matrix must be inverted; the Durbin--Levinson
        and innovations algorithms avoid this (Shumway \& Stoffer, \S3.5).
\end{itemize}
\end{corollary}

\begin{theorem}{Prediction equations with non-zero mean (and intercept $\beta_0$)}{p8-predeqmean}
For a stationary process with $\mu=\E[X_t]$, write
$X^{\,n}_{n+h}=\beta_0+\sum_{j=1}^n\beta_j X_j$. The coefficients solve the
prediction equations $\E[(X_{n+h}-X^{\,n}_{n+h})X_k]=0$ for $k=0,1,\dots,n$
(with the convention $X_0\equiv1$), which give
\[
\Gamma_n\beta=\acvs_{[h]},
\qquad
\beta_0=\mu\big(1-\beta\Tr\mathbf 1_n\big),
\]
and, when $\Gamma_n$ is invertible,
\[
X^{\,n}_{n+h}=\mu+\acvs_{[h]}\Tr\,\Gamma_n^{-1}(X-\mu\mathbf 1_n),
\qquad
P^{\,n}_{n+h}=\acvs_0-\acvs_{[h]}\Tr\,\Gamma_n^{-1}\,\acvs_{[h]} .
\]
\textit{Proof idea:} the predictor has the form
$X^{\,n}_{n+h}=\mu+\beta\Tr(X-\mu\mathbf 1_n)$, so centring removes the mean and
reduces the problem to the zero-mean case --- there is no loss of generality in
assuming $\mu=0$, where $\beta_0=0$. The MSE is \emph{identical} to the
zero-mean case because the constant shift does not affect the error variance.\\[2pt]
\textit{Why:} real data are rarely centred; this is the version actually used in
practice. (Proved as Exercise~\ref{exo:p8-ex2}.)
\end{theorem}

\begin{theorem}{Best linear predictor of a causal ARMA via $\psi$-weights}{p8-armapred}
For a causal ARMA process with $MA(\infty)$ representation
$X_t=\sum_{j=0}^\infty\psi_j\eps_{t-j}$, the best linear predictor of $X_{n+h}$
based on $X_1,\dots,X_n$ is
\[
X^{\,n}_{n+h}=\sum_{j=h}^{\infty}\psi_j\,\eps_{n+h-j}
=\psi_h\eps_n+\psi_{h+1}\eps_{n-1}+\cdots,
\]
with prediction mean square error
\[
P^{\,n}_{n+h}=\sigma^2\sum_{j=0}^{h-1}\psi_j^2 .
\]
\textit{Proof idea:} write the predictor as a linear filter of the past
innovations, $X^{\,n}_{n+h}=\sum_{j=0}^\infty c_j\eps_{n-j}$. Then
\[
\E\big[(X_{n+h}-X^{\,n}_{n+h})^2\big]
=\sigma^2\Big(\sum_{j=0}^{h-1}\psi_j^2+\sum_{j=h}^\infty(\psi_j-c_{j-h})^2\Big),
\]
which is minimised by matching $c_{j}=\psi_{j+h}$, killing the second sum and
leaving $\sigma^2\sum_{j=0}^{h-1}\psi_j^2$. Practically:
\textbf{set the unrealised future innovations $\eps_{n+1},\dots,\eps_{n+h}$ to
zero} in $X_{n+h}=\sum_{j=0}^\infty\psi_j\eps_{n+h-j}$.\\[2pt]
\textit{Why:} this is the engine of all ARMA forecasting. It also shows that for
a causal ARMA, $\psi_j\to0$ exponentially so $X^{\,n}_{n+h}\to\mu$, and
$P^{\,n}_{n+h}=\sigma^2\sum_{j=0}^{h-1}\psi_j^2\uparrow\sigma^2\sum_{j=0}^\infty\psi_j^2=\acvs_0$
as $h\to\infty$. (Proved as Exercise~\ref{exo:p8-ex3}.)
\end{theorem}

\begin{intuition}
The two ARMA forecasting theorems (prediction equations vs.\ $\psi$-weights) are
the same projection viewed differently. The prediction equations project onto
$\mathrm{span}\{X_1,\dots,X_n\}$ using the ACVS; the $\psi$-weight formula
projects onto $\mathrm{span}\{\eps_s:s\le n\}$, the \emph{innovation} space.
Because for an invertible ARMA these two spans coincide, the answers agree --- but
the innovation view is far easier to compute: ``write the future, erase the
future innovations.''
\end{intuition}

\begin{proposition}{Limiting behaviour of ARMA forecasts}{p8-armalimit}
For a causal ARMA model with mean $\mu$,
$X^{\,n}_{n+h}=\mu+\sum_{j=h}^\infty\psi_j\eps_{n+h-j}$ with $\psi_j\to0$
exponentially fast, hence
\[
X^{\,n}_{n+h}\longrightarrow\mu,
\qquad
\Var\big(e_n(h)\big)=\sigma^2\sum_{j=0}^{h-1}\psi_j^2\longrightarrow\acvs_0,
\qquad h\to\infty .
\]
\textit{Proof idea:} immediate from the $\psi$-weight theorem and exponential
decay of $\psi_j$ (a consequence of causality, roots of $\Phi$ outside the unit
circle).\\[2pt]
\textit{Why:} quantifies ``the distant future is the unconditional
distribution'' --- the point forecast converges to the mean and the error
variance converges to the marginal variance $\acvs_0$.
\end{proposition}

\begin{proposition}{Forecasts for integrated (ARIMA) processes}{p8-arimapred}
Let $\{X_t\}$ be ARIMA$(p,d,q)$, $\Phi(B)(1-B)^dX_t=\Theta(B)\eps_t$. Then
$Z_t=\diff^{\,d}X_t$ is the ARMA$(p,q)$ process
$Z_t=\sum_{j=1}^p\phi_jZ_{t-j}+\eps_t-\sum_{j=1}^q\theta_j\eps_{t-j}$, and one
forecasts $X$ by undoing the differencing,
\[
\diff^{\,d}X^{\,n}_{n+h}=Z^{\,n}_{n+h}.
\]
For $d=1$ (so $X_{n+h}-X_{n+h-1}=Z_{n+h}$),
\[
X^{\,n}_{n+h}=X_n+\sum_{j=0}^{h-1}Z^{\,n}_{n+h-j},\qquad h\ge1 .
\]
If $\{Z_t\}$ has mean $\alpha\ne0$ then $Z^{\,n}_{n+h}\to\alpha$ and
$X^{\,n}_{n+h}\approx X_n+\alpha h$ for large $h$ (a deterministic drift line).\\[2pt]
\textit{Proof idea:} apply the linear projection operator $P_{X_1,\dots,X_n}$,
which is linear, to $\diff^{\,d}X_{n+h}=Z_{n+h}$; since differencing is a fixed
linear filter it commutes with the projection, giving
$\diff^{\,d}X^{\,n}_{n+h}=Z^{\,n}_{n+h}$. Telescoping the $d=1$ identity
$X_{n+h}=X_{n+h-1}+Z_{n+h}$ down to the last observation $X_n$ yields
$X^{\,n}_{n+h}=X_n+\sum_{j=0}^{h-1}Z^{\,n}_{n+h-j}$.\\[2pt]
\textit{Why:} reduces ARIMA forecasting to ARMA forecasting on the differenced
series, then ``re-integrates''.
\begin{itemize}
  \item For ARIMA models the prediction MSE \textbf{diverges}:
        $P^{\,n}_{n+h}\to\infty$ as $h\to\infty$ (no mean reversion --- the
        process is non-stationary).
  \item For SARIMA models with $d=D=1$ and $\alpha=0$ the forecast is linear and
        seasonal in $h$.
\end{itemize}
\end{proposition}

\begin{pitfall}
After expanding an ARIMA$(p,1,q)$ into the form
$X_t=(1+\phi_1)X_{t-1}-(\phi_1-\phi_2)X_{t-2}-\cdots$, the equation \emph{looks}
like an ARMA$(p{+}1,0,q)$, but it is \textbf{not} a stationary ARMA: its AR
coefficients do not satisfy the causality/stationarity conditions (there is a
unit root). Treat it only as a convenient recursion for forecasting, never as a
stationary model. This is exactly why $P^{\,n}_{n+h}\to\infty$.
\end{pitfall}

% ----------------------------------------------------------------------
\subsection{Key Techniques}

\begin{keytech}{Solving the prediction equations}{p8-ktpredeq}
To get the best linear predictor of a zero-mean stationary process:
\begin{enumerate}
  \item Build $\Gamma_n=[\acvs_{i-j}]_{i,j=1}^n$ (Toeplitz, from the ACVS).
  \item Build the right-hand side
        $\acvs_{[h]}=(\acvs_{n+h-1},\acvs_{n+h-2},\dots,\acvs_h)\Tr$.
  \item Solve $\Gamma_n\beta=\acvs_{[h]}$ for $\beta$, set
        $X^{\,n}_{n+h}=\sum_j\beta_jX_j=\acvs_{[h]}\Tr\Gamma_n^{-1}X$.
  \item Report the MSE
        $P^{\,n}_{n+h}=\acvs_0-\acvs_{[h]}\Tr\Gamma_n^{-1}\acvs_{[h]}$.
\end{enumerate}
For \textbf{non-zero mean} $\mu$, centre first: use $X-\mu\mathbf1_n$ in place of
$X$ and add $\mu$ back,
$X^{\,n}_{n+h}=\mu+\acvs_{[h]}\Tr\Gamma_n^{-1}(X-\mu\mathbf1_n)$; the MSE is
unchanged. \emph{Shortcut for AR(p):} for $h=1$ and $n\ge p$ the answer is just
$X^{\,n}_{n+1}=\sum_{i=1}^p\phi_iX_{n+1-i}$ --- no inversion needed (the next
technique explains why).
\end{keytech}

\begin{keytech}{Causal-ARMA forecast by zeroing future innovations}{p8-ktpsi}
The fast route for any causal ARMA:
\begin{enumerate}
  \item Write $X_{n+h}=\sum_{j=0}^\infty\psi_j\eps_{n+h-j}$ (the $MA(\infty)$
        form), or just the ARMA equation rearranged for $X_{n+h}$.
  \item Take conditional expectation given the data: every future innovation
        $\eps_{n+1},\dots,\eps_{n+h}$ becomes $0$; every past/present
        innovation $\eps_s$ ($s\le n$) stays (replace by its residual
        $\hat\eps_s$); every past observation $X_s$ ($s\le n$) stays; every
        future observation $X_s$ ($s>n$) is replaced by its forecast
        $X^{\,n}_s$.
  \item The MSE is $P^{\,n}_{n+h}=\sigma^2\sum_{j=0}^{h-1}\psi_j^2$ --- obtain the
        first $h$ $\psi$-weights from $\Theta(B)/\Phi(B)=\sum_j\psi_jB^j$ by
        matching powers of $B$.
\end{enumerate}
\end{keytech}

\begin{keytech}{Recursive truncated forecasting recipe (ARMA \& ARIMA)}{p8-ktrec}
The general-purpose hand-computation method. \textbf{Three steps:}
\begin{enumerate}
  \item \textbf{Isolate $y_t$:} expand the (ARI)MA equation so that $y_t$ is
        alone on the left and \emph{everything else} (past $y$'s, present/past
        $\eps$'s) is on the right.
  \item \textbf{Shift to $n+h$:} replace $t$ by $n+h$ throughout.
  \item \textbf{Substitute:} on the right-hand side replace
        \begin{itemize}
          \item future observations $y_{s}$ ($s>n$) by their forecasts
                $\hat y_{s\mid n}=X^{\,n}_s$;
          \item future errors $\eps_s$ ($s>n$) by $0$;
          \item past errors $\eps_s$ ($s\le n$) by the corresponding residuals
                $\hat\eps_s$ (and any $\eps_s$ with $s\le0$ by $0$).
        \end{itemize}
\end{enumerate}
Run $h=1,2,3,\dots$ in order; each forecast feeds the next. The residuals are
themselves built recursively (for $t=1,\dots,n$, and $0$ otherwise) via
\[
\hat\eps_t=\Big(\tilde X^{\,n}_t-\sum_{i=1}^p\phi_i\tilde X^{\,n}_{t-i}\Big)
+\sum_{k=1}^q\theta_k\hat\eps_{t-k}.
\]
The truncation (setting pre-sample terms to $0$) makes $\tilde X^{\,n}_t$ an
\emph{approximation}; the forecast error is estimated by
$\sigma^2\sum_{j=0}^{h-1}\psi_j^2$.
\end{keytech}

\begin{keytech}{Re-integrating ARIMA forecasts}{p8-ktint}
For ARIMA$(p,d,q)$: difference $d$ times to get the ARMA series
$Z_t=\diff^{\,d}X_t$, forecast $Z^{\,n}_{n+h}$ by the ARMA technique, then undo
the differencing. For $d=1$,
\[
X^{\,n}_{n+1}=X_n+Z^{\,n}_{n+1},\qquad
X^{\,n}_{n+h}=X_n+\sum_{j=0}^{h-1}Z^{\,n}_{n+h-j}\ (h\ge1).
\]
Equivalently, expand $\Phi(B)(1-B)^d$ as a single (unit-root) AR polynomial in
$X$ and apply the recursive recipe directly (often easier in practice). Remember
the MSE \emph{grows without bound}.
\end{keytech}

\begin{keytech}{Building Gaussian prediction intervals}{p8-ktpi}
\begin{enumerate}
  \item Get the point forecast $X^{\,n}_{n+h}$.
  \item Get the $\psi$-weights $\psi_0=1,\psi_1,\dots,\psi_{h-1}$ from
        $\Theta(B)/\Phi(B)$.
  \item Forecast-error variance: $\Var(e_n(h))=\sigma^2\sum_{j=0}^{h-1}\psi_j^2$.
  \item Interval: $X^{\,n}_{n+h}\pm z_{1-\alpha/2}\sqrt{\Var(e_n(h))}$ (use
        $1.96$ for $95\%$).
\end{enumerate}
Sanity check: for stationary ARMA the half-width increases in $h$ toward
$z_{1-\alpha/2}\sqrt{\acvs_0}$; for ARIMA it keeps growing.
\end{keytech}

% ----------------------------------------------------------------------
\subsection{Exercises}

\begin{exercise}{One-step AR(2) predictors from 1 and 2 observations \quad\textnormal{[Sheet 10]}}{p8-ex1}
Let $\eps_t$ be white noise. For the causal AR(2) process
$X_t=\phi_1X_{t-1}+\phi_2X_{t-2}+\eps_t$, show that the best one-step-ahead
linear predictors based on one observation $X_1$ and on two observations
$X_1,X_2$ are
\[
X^{\,1}_2=\frac{\acvs_1}{\acvs_0}X_1=\acf_1X_1,
\qquad
X^{\,2}_3=\phi_1X_2+\phi_2X_1 .
\]
\end{exercise}
\begin{solution}
\textbf{Based on $X_1$ (so $n=1$, $h=1$).} The prediction equation
$\Gamma_1\beta=\acvs_{[1]}$ is the scalar equation $\acvs_0\beta_1=\acvs_1$, so
$\beta_1=\acvs_1/\acvs_0=\acf_1$ and
\[
X^{\,1}_2=\beta_1X_1=\frac{\acvs_1}{\acvs_0}X_1=\acf_1X_1 .
\]
\textbf{Based on $X_1,X_2$ (so $n=2$, $h=1$).} We could build $\Gamma_2$ and
solve, but it is quicker to use the model directly. Since
$X_3=\phi_1X_2+\phi_2X_1+\eps_3$ we have
$X_3-(\phi_1X_2+\phi_2X_1)=\eps_3$, and $\eps_3$ (white noise, future) is
uncorrelated with $X_1$ and $X_2$:
\[
\E\big[(X_3-\phi_1X_2-\phi_2X_1)X_1\big]=\E[\eps_3X_1]=0,
\qquad
\E\big[(X_3-\phi_1X_2-\phi_2X_1)X_2\big]=\E[\eps_3X_2]=0 .
\]
Thus the candidate $\phi_1X_2+\phi_2X_1$ has error orthogonal to both
observations, so by the prediction-equations theorem it \emph{is} the best
linear predictor: $X^{\,2}_3=\phi_1X_2+\phi_2X_1$. (This is the general fact: for
an AR$(p)$ with $n\ge p$ observations, the one-step forecast is simply
$\sum_{i=1}^p\phi_iX_{n+1-i}$.)
\end{solution}

\begin{exercise}{Prediction equations with non-zero mean \quad\textnormal{[Sheet 10]}}{p8-ex2}
For a stationary process with $\mu=\E[X_t]$ and predictor
$X^{\,n}_{n+h}=\beta_0+\sum_{j=1}^n\beta_jX_j$ (set $X_0\equiv1$), prove that the
prediction equations $\E[(X_{n+h}-X^{\,n}_{n+h})X_k]=0$, $k=0,\dots,n$, give
$\Gamma_n\beta=\acvs_{[h]}$ and $\beta_0=\mu(1-\beta\Tr\mathbf1_n)$, and that
when $\Gamma_n$ is invertible
$X^{\,n}_{n+h}=\mu+\acvs_{[h]}\Tr\Gamma_n^{-1}(X-\mu\mathbf1_n)$ with MSE
$\acvs_0-\acvs_{[h]}\Tr\Gamma_n^{-1}\acvs_{[h]}$.
\end{exercise}
\begin{solution}
\textbf{(a) MSE in matrix form.} With $X=(X_1,\dots,X_n)\Tr$ and
$\mathbf1_n$ the all-ones vector,
\[
\begin{aligned}
S(\beta_0,\dots,\beta_n)
&=\E\Big[\big(X_{n+h}-\beta_0-\beta\Tr X\big)^2\Big]\\
&=\Var\big(X_{n+h}-\beta\Tr X\big)+\big(\E[X_{n+h}-\beta_0-\beta\Tr X]\big)^2\\
&=\acvs_0-2\beta\Tr\Cov(X_{n+h},X)+\beta\Tr\Var(X)\beta+\big(\mu-\beta_0-\mu\beta\Tr\mathbf1_n\big)^2,
\end{aligned}
\]
a quadratic bounded below by $0$, hence minimised where all
$\partial S/\partial\beta_r=0$.\\
\textbf{(b) Solve for $\beta_0$.} Only the last term depends on $\beta_0$:
\[
\frac{\partial S}{\partial\beta_0}=2\big(\mu-\beta_0-\mu\beta\Tr\mathbf1_n\big)(-1)
=0
\;\Longrightarrow\;
\beta_0=\mu\big(1-\beta\Tr\mathbf1_n\big).
\]
\textbf{(c) Reduce to zero mean.} Substituting $\beta_0$ back,
\[
X^{\,n}_{n+h}=\beta_0+\beta\Tr X=\mu+\beta\Tr(X-\mu\mathbf1_n),
\]
so the predictor depends on the data only through the \emph{centred} vector
$X-\mu\mathbf1_n$. Hence there is no loss of generality in taking $\mu=0$, in
which case $\beta_0=0$ and $X^{\,n}_{n+h}=\beta\Tr X$.\\
\textbf{(d) Normal equations.} For $\mu=0$,
$S(\beta)=\acvs_0-2\beta\Tr\Cov(X_{n+h},X)+\beta\Tr\Var(X)\beta$; setting
$\nabla_\beta S=0$ gives
\[
\Var(X)\,\beta=\Cov(X_{n+h},X),\quad\text{i.e.}\quad \Gamma_n\beta=\acvs_{[h]}.
\]
\textbf{(e) Equivalence to the orthogonality form.} Row $k$ of
$\Gamma_n\beta=\acvs_{[h]}$ reads $\sum_{j=1}^n\beta_j\acvs_{j-k}=\acvs_{n+h-k}$,
which is exactly $\E[(X_{n+h}-\sum_j\beta_jX_j)X_k]=0$ for $k=1,\dots,n$. The
$k=0$ equation (with $X_0=1$) gives $\mu-\beta_0-\sum_j\beta_j\mu=0$, recovering
$\beta_0=\mu(1-\beta\Tr\mathbf1_n)$ from part (b).\\
\textbf{(2a) Closed form.} If $\Gamma_n$ is invertible,
$\beta=\Gamma_n^{-1}\acvs_{[h]}$, so
\[
X^{\,n}_{n+h}=\mu+\acvs_{[h]}\Tr\Gamma_n^{-1}(X-\mu\mathbf1_n),
\]
differing from the lecture's zero-mean formula only by the centring/shift by
$\mu$.\\
\textbf{(2b) MSE.} Since $\E[X_{n+h}-X^{\,n}_{n+h}]=0$,
\[
P^{\,n}_{n+h}=\Var\big(X_{n+h}-\beta\Tr X\big)
=\acvs_0-2\beta\Tr\acvs_{[h]}+\beta\Tr\Gamma_n\beta
=\acvs_0-\acvs_{[h]}\Tr\Gamma_n^{-1}\acvs_{[h]},
\]
identical to the zero-mean case (the mean does not affect forecast accuracy).
\end{solution}

\begin{exercise}{Causal-ARMA predictor via $\psi$-weights \quad\textnormal{[Sheet 10]}}{p8-ex3}
Prove that the best linear predictor $X^{\,n}_{n+h}$ for $X_{n+h}$ in a causal
ARMA process with linear representation
$X_t=\sum_{j=0}^\infty\psi_j\eps_{t-j}$ is
$X^{\,n}_{n+h}=\sum_{j=h}^\infty\psi_j\eps_{n+h-j}$, with MSE
$\sigma^2\sum_{j=0}^{h-1}\psi_j^2$.
\end{exercise}
\begin{solution}
Consider a linear predictor $X^{\,n}_{n+h}=\sum_{i=1}^n\beta_iX_i$. Substituting
the linear representation $X_i=\sum_{j=0}^\infty\psi_j\eps_{i-j}$,
\[
X^{\,n}_{n+h}=\sum_{i=1}^n\beta_i\sum_{j=0}^\infty\psi_j\eps_{i-j}
=\sum_{j=0}^\infty c_j\,\eps_{n-j}
\]
for some coefficients $c_j$; i.e.\ \emph{any} linear predictor based on
$X_1,\dots,X_n$ is a linear combination of innovations up to time $n$ (it
``starts at $\eps_n$''). Now compute the MSE, using
$X_{n+h}=\sum_{j=0}^\infty\psi_j\eps_{n+h-j}$ and splitting the future
innovations $\eps_{n+1},\dots,\eps_{n+h}$ (indices $j=0,\dots,h-1$) from the
rest:
\[
\begin{aligned}
\E\Big[\big(X_{n+h}-X^{\,n}_{n+h}\big)^2\Big]
&=\E\Big[\Big(\sum_{j=0}^\infty\psi_j\eps_{n+h-j}-\sum_{j=0}^\infty c_j\eps_{n-j}\Big)^2\Big]\\
&=\sigma^2\Big(\sum_{j=0}^{h-1}\psi_j^2+\sum_{j=h}^\infty(\psi_j-c_{j-h})^2\Big),
\end{aligned}
\]
where the first sum collects the future innovations (which no past-based
predictor can touch) and the second matches $\psi_j$ against $c_{j-h}$ for the
observable innovations (cross terms vanish because the $\eps$'s are
uncorrelated). The first sum is fixed; the second is minimised --- to zero --- by
choosing $c_{j-h}=\psi_j$, i.e.\ $c_j=\psi_{j+h}$ for $j=0,1,\dots$. The
predictor is then
\[
X^{\,n}_{n+h}=\sum_{j=0}^\infty c_j\eps_{n-j}=\sum_{j=0}^\infty\psi_{j+h}\eps_{n-j}
=\sum_{j=h}^\infty\psi_j\eps_{n+h-j},
\]
and the residual is $X_{n+h}-X^{\,n}_{n+h}=\sum_{j=0}^{h-1}\psi_j\eps_{n+h-j}$,
giving
\[
P^{\,n}_{n+h}=\E\Big[\big(X_{n+h}-X^{\,n}_{n+h}\big)^2\Big]=\sigma^2\sum_{j=0}^{h-1}\psi_j^2 .
\]
\end{solution}

\begin{exercise}{ARIMA(3,1,1) worked truncated forecast \quad\textnormal{[Lecture]}}{p8-ex4}
For the fitted model $(1-\phi_1B-\phi_2B^2-\phi_3B^3)(1-B)y_t=(1+\theta_1B)\eps_t$
(the convention used in the lecture's worked example, with the MA polynomial
written $\Theta(B)=1+\theta_1B$), derive the one-, two- and general $h$-step
truncated forecasts. Then plug in the fitted values $\hat\phi_1=0.0044$,
$\hat\phi_2=0.0916$, $\hat\phi_3=0.3698$, $\hat\theta_1=-0.3921$ and write the
explicit $\hat y_{T+1\mid T}$ and $\hat y_{T+2\mid T}$.
\end{exercise}
\begin{solution}
\textbf{Step 1 --- isolate $y_t$.} Multiply out the AR$\times$difference factor
$(1-\phi_1B-\phi_2B^2-\phi_3B^3)(1-B)$:
\[
1-(1+\phi_1)B+(\phi_1-\phi_2)B^2+(\phi_2-\phi_3)B^3+\phi_3B^4,
\]
so the model is
\[
y_t-(1+\phi_1)y_{t-1}+(\phi_1-\phi_2)y_{t-2}+(\phi_2-\phi_3)y_{t-3}+\phi_3y_{t-4}
=\eps_t+\theta_1\eps_{t-1}.
\]
Moving everything except $y_t$ to the right,
\[
y_t=\underbrace{(1+\phi_1)y_{t-1}-(\phi_1-\phi_2)y_{t-2}-(\phi_2-\phi_3)y_{t-3}-\phi_3y_{t-4}}_{\text{AR part}}
+\underbrace{\eps_t+\theta_1\eps_{t-1}}_{\text{MA part}} .
\]
(We use the lecture's convention $\Theta(B)=1+\theta_1B$, so the MA term is
$\eps_t+\theta_1\eps_{t-1}$ and the forecast picks up $+\theta_1\hat\eps$.
Beware: in the earlier general theory the MA polynomial was written
$\Theta(B)=1-\theta_1B$; that sign convention would replace every $+\theta_1$
below by $-\theta_1$. We follow the worked example exactly so the numbers match.)
\\[3pt]
\textbf{Step 2--3 --- shift and substitute.} For $h=1$ (replace $t$ by $n+1$):
the only error terms are $\eps_{n+1}\to0$ and $\eps_n\to\hat\eps_n$ (last
residual). Hence
\[
y^{\,n}_{n+1}=(1+\phi_1)y_n-(\phi_1-\phi_2)y_{n-1}-(\phi_2-\phi_3)y_{n-2}-\phi_3y_{n-3}+\theta_1\hat\eps_n .
\]
For $h>1$ all error terms are future ($\to0$) and the future observations are
replaced by forecasts:
\[
y^{\,n}_{n+h}=(1+\phi_1)y^{\,n}_{n+h-1}-(\phi_1-\phi_2)y^{\,n}_{n+h-2}-(\phi_2-\phi_3)y^{\,n}_{n+h-3}-\phi_3y^{\,n}_{n+h-4},
\]
with the convention $y^{\,n}_s=y_s$ for $s\le n$.\\[3pt]
\textbf{Numbers.} Using
$1+\hat\phi_1=1.0044$, $\hat\phi_1-\hat\phi_2=0.0044-0.0916=-0.0872$,
$\hat\phi_2-\hat\phi_3=0.0916-0.3698=-0.2782$, $\hat\phi_3=0.3698$,
$\hat\theta_1=-0.3921$, the one-step forecast (writing $e_T$ for the last
residual) is
\[
\hat y_{T+1\mid T}=1.0044\,y_T+0.0872\,y_{T-1}+0.2782\,y_{T-2}-0.3698\,y_{T-3}+\hat\theta_1\,e_T,
\]
i.e.\ with $\hat\theta_1=-0.3921$,
\[
\hat y_{T+1\mid T}=1.0044\,y_T+0.0872\,y_{T-1}+0.2782\,y_{T-2}-0.3698\,y_{T-3}-0.3921\,e_T .
\]
(The two middle signs are positive because $-(\hat\phi_1-\hat\phi_2)=+0.0872$ and
$-(\hat\phi_2-\hat\phi_3)=+0.2782$.) The two-step forecast replaces $y_{T+1}$ by
$\hat y_{T+1\mid T}$ and sets all errors to $0$:
\[
\hat y_{T+2\mid T}=1.0044\,\hat y_{T+1\mid T}+0.0872\,y_T+0.2782\,y_{T-1}-0.3698\,y_{T-2}.
\]
Continue likewise for $h=3,4,\dots$; each step uses previously computed
forecasts. Because $d=1$, the prediction MSE grows without bound as $h$
increases.
\end{solution}

\begin{exercise}{One- and two-step forecast of an MA(2) with interval \quad\textnormal{[New]}}{p8-ex5}
Let $X_t=\eps_t-\theta_1\eps_{t-1}-\theta_2\eps_{t-2}$ with
$\eps_t\iid\mathcal N(0,\sigma^2)$ (invertible, zero mean). Give
$X^{\,n}_{n+1}$, $X^{\,n}_{n+2}$, $X^{\,n}_{n+3}$, their forecast-error
variances, and a $95\%$ interval for $X_{n+2}$.
\end{exercise}
\begin{solution}
Here $\psi_0=1,\ \psi_1=-\theta_1,\ \psi_2=-\theta_2$ and $\psi_j=0$ for
$j\ge3$. By the $\psi$-weight theorem, set future innovations to $0$ in
$X_{n+h}=\eps_{n+h}-\theta_1\eps_{n+h-1}-\theta_2\eps_{n+h-2}$:
\[
\begin{aligned}
X^{\,n}_{n+1}&=-\theta_1\hat\eps_n-\theta_2\hat\eps_{n-1},\\
X^{\,n}_{n+2}&=-\theta_2\hat\eps_n,\\
X^{\,n}_{n+3}&=0\quad(=\mu,\ \text{the mean; an MA(2) is forecast as the mean for }h>2).
\end{aligned}
\]
Forecast-error variances $\Var(e_n(h))=\sigma^2\sum_{j=0}^{h-1}\psi_j^2$:
\[
\Var(e_n(1))=\sigma^2,\quad
\Var(e_n(2))=\sigma^2(1+\theta_1^2),\quad
\Var(e_n(3))=\sigma^2(1+\theta_1^2+\theta_2^2)=\acvs_0 .
\]
For $h\ge3$ the variance saturates at $\acvs_0$, as it must. A $95\%$ interval
for $X_{n+2}$ is
\[
-\theta_2\hat\eps_n\;\pm\;1.96\,\sigma\sqrt{1+\theta_1^2}.
\]
\end{solution}

\begin{exercise}{Two-step AR(1) forecast variance and interval \quad\textnormal{[New]}}{p8-ex6}
For the Gaussian AR(1) $X_t=\phi X_{t-1}+\eps_t$ ($\abs\phi<1$,
$\eps_t\iid\mathcal N(0,\sigma^2)$), compute $X^{\,n}_{n+1},X^{\,n}_{n+2}$,
their error variances, and verify $\Var(e_n(h))\to\acvs_0$.
\end{exercise}
\begin{solution}
The $MA(\infty)$ weights of a causal AR(1) are $\psi_j=\phi^j$. Point forecasts
(Example~\ref{exmp:p8-ar1}): $X^{\,n}_{n+1}=\phi X_n$,
$X^{\,n}_{n+2}=\phi^2X_n$, and generally $X^{\,n}_{n+h}=\phi^hX_n$.
Forecast-error variances:
\[
\Var(e_n(1))=\sigma^2\psi_0^2=\sigma^2,
\qquad
\Var(e_n(2))=\sigma^2(\psi_0^2+\psi_1^2)=\sigma^2(1+\phi^2),
\]
and in general
$\Var(e_n(h))=\sigma^2\sum_{j=0}^{h-1}\phi^{2j}=\sigma^2\dfrac{1-\phi^{2h}}{1-\phi^2}$.
As $h\to\infty$ this tends to $\dfrac{\sigma^2}{1-\phi^2}=\acvs_0$, the marginal
variance of an AR(1) --- confirming the saturation. A $95\%$ interval for
$X_{n+2}$ is $\phi^2X_n\pm1.96\,\sigma\sqrt{1+\phi^2}$.
\end{solution}

\begin{exercise}{Re-integrating a random walk with drift \quad\textnormal{[New]}}{p8-ex7}
Let $X_t$ be ARIMA$(0,1,0)$ with drift: $Z_t=X_t-X_{t-1}=\alpha+\eps_t$,
$\eps_t\iid\mathcal N(0,\sigma^2)$. Find $X^{\,n}_{n+h}$ and $\Var(e_n(h))$, and
show the MSE diverges.
\end{exercise}
\begin{solution}
The differenced series $Z_t$ is white noise plus a constant mean
$\E[Z_t]=\alpha$, so $Z^{\,n}_{n+h}=\alpha$ for all $h$. Re-integrating
(Proposition on integrated processes),
\[
X^{\,n}_{n+h}=X_n+\sum_{j=0}^{h-1}Z^{\,n}_{n+h-j}=X_n+\alpha h,
\]
a straight drift line from the last observation. The forecast error is
$e_n(h)=X_{n+h}-X^{\,n}_{n+h}=\sum_{j=1}^h\eps_{n+j}$ (the accumulated future
innovations), so
\[
\Var(e_n(h))=h\,\sigma^2\;\xrightarrow[h\to\infty]{}\;\infty .
\]
Thus $P^{\,n}_{n+h}=h\sigma^2$ grows linearly and the prediction interval
$X_n+\alpha h\pm1.96\,\sigma\sqrt h$ fans out without bound --- the hallmark of an
integrated process (no mean reversion).
\end{solution}

% ======================================================================
%  PART 9 -- Diagnostics, Model Selection and Long Memory
% ======================================================================
\section{Diagnostics, Model Selection and Long Memory}
\label{sec:diagnostics}

Once a candidate ARIMA model has been identified and fitted, two questions
remain: \emph{does it fit?} (diagnostics, based on the residuals) and \emph{is
it the right model among several?} (model selection, based on information
criteria). This part collects the standard residual checks (mean, variance,
Gaussianity, autocorrelation), the Box--Pierce and Ljung--Box portmanteau
tests, the AIC/AICc/BIC criteria, and the Box--Jenkins methodology that ties
them together. It then steps outside the ARMA world entirely: ARMA covariances
decay exponentially fast, but many real series have covariances that decay only
\emph{polynomially}. This is \textbf{long memory}, modelled by fractional
differencing $(1-\Bsh)^d$ and the FARIMA class, and intimately related to
self-similarity and fractional Brownian motion.

% ----------------------------------------------------------------------
\subsection{Definitions}

\begin{definition}{Residuals and standardized residuals}{p9-resid}
For a fitted model with fitted value $\hat X_t$ at time $t$, the
\textbf{residuals} are
\[
e_t = X_t - \hat X_t ,
\]
and the \textbf{standardized residuals} are
\[
\tilde e_t = \frac{e_t}{\sqrt{\Var(e_t)}} .
\]
Model checking is built almost entirely on the $\{e_t\}$ (or $\{\tilde e_t\}$):
if the model is correct, the residuals should behave like the driving white
noise. The four questions to ask are:
\begin{enumerate}
  \item Do the residuals have constant zero mean?
  \item Is their variance constant in $t$ (\emph{homoscedasticity})?
  \item Are they uncorrelated in $t$?
  \item Are they Gaussian?
\end{enumerate}
\end{definition}

\begin{intuition}[Why residuals should be the noise]
A correctly specified model has captured all the predictable (linearly
dependent) structure of $\{X_t\}$. What is left over, $e_t=X_t-\hat X_t$, should
be the \emph{unpredictable} part --- the innovation $\eps_t$. So every property
we demanded of the white-noise input (zero mean, constant variance, no
correlation, optionally Gaussianity) is exactly what we now \emph{test for} in
the residuals. Failure of any of the four points the modeller back to a
deficiency: a trend left in (mean), volatility clustering (variance), an
unmodelled AR/MA term (correlation), or non-normal innovations (Gaussianity).
\end{intuition}

\begin{definition}{Checking mean, variance and Gaussianity}{p9-checks}
Three of the four checks are visual:
\begin{itemize}
  \item \textbf{Mean:} plot $\{(t,\tilde e_t)\}$. There should be no trend and
        the points should hover around zero.
  \item \textbf{Variance (homoscedasticity):} plot $\{(t,\tilde e_t)\}$ and
        look for constant spread; a fanning-out or fanning-in band signals
        non-constant variance.
  \item \textbf{Gaussianity:} a \textbf{Q--Q plot} of the residual quantiles
        against normal quantiles; Gaussian residuals lie on a straight line,
        whereas heavy tails (e.g.\ Student) curve away at the ends. Note this
        checks only \emph{marginal} Gaussianity.
\end{itemize}
\end{definition}

\begin{definition}{Sample residual autocorrelation \& the correlogram band}{p9-residacf}
The residual sample autocorrelation at lag $\tau$ is
\[
\hat r_\tau = \hat\acf^{(e)}_\tau
= \frac{\sum_{t=1}^{n-\tau}(e_{t+\tau}-\bar e)(e_t-\bar e)}
       {\sum_{t=1}^{n}(e_t-\bar e)^2 }.
\]
Plotting $\hat r_\tau$ against $\tau$ gives the residual \textbf{correlogram}.
Under the null hypothesis that the residuals are white noise, each $\hat r_\tau$
is approximately $\mathcal N(0,1/n)$, so individual values are compared against
the band
\[
\left(-\frac{1.96}{\sqrt n},\ \frac{1.96}{\sqrt n}\right),
\]
the (pointwise) $95\%$ acceptance band for ``$\hat r_\tau$ is zero''.
\end{definition}

\begin{definition}{$R^2$ statistic}{p9-r2}
With $s_e^2$ the residual variance and $s_y^2$ the sample variance of the data,
\[
R^2 = 1 - \frac{s_e^2}{s_y^2}.
\]
It measures the fraction of variance ``explained'' by the fit. In time-series
model selection it is a \textbf{dangerous} criterion (see
Proposition~\ref{prop:p9-r2danger}): adding parameters can only increase $R^2$,
and even the \emph{true} model has a bounded $R^2$ because a stochastic process
is not perfectly predictable.
\end{definition}

\begin{definition}{Information criteria: AIC, AICc, BIC}{p9-ic}
For a model with $k$ estimated parameters fitted to $n$ observations, with
maximized likelihood $L(\theta\mid y)$:
\[
\begin{aligned}
\mathrm{AIC}(\theta)  &= -2\log L(\theta\mid y) + 2k, \\[2pt]
\mathrm{AICc}(\theta) &= -2\log L(\theta\mid y) + 2k\,\frac{n}{n-k-1}
                       = \mathrm{AIC}(\theta) + \frac{2k(k+1)}{n-k-1}, \\[2pt]
\mathrm{BIC}(\theta)  &= -2\log L(\theta\mid y) + k\log(n).
\end{aligned}
\]
\textbf{Smaller is better} for all three. For an ARMA$(p,q)$ the parameter count
is
\[
k = p+q+1
\]
($p$ AR coefficients, $q$ MA coefficients, plus the noise variance $\sigma^2$).
AIC is known to \emph{overestimate} $p$; AICc (Hurvich \& Tsai, 1989) corrects
the small-sample bias and coincides with AIC as $n\to\infty$; BIC penalizes
parameters more heavily (for $n>e^2\approx 7.4$) and so favours more
parsimonious models.
\end{definition}

\begin{definition}{Long memory (covariance definition)}{p9-lmcov}
A second-order stationary process $\{X_t\}$ with ACVS $\{\acvs_\tau\}$ is said to
exhibit \textbf{long memory} (long-range dependence) if its covariances decay
only polynomially:
\[
\acvs_\tau \sim L_\gamma(\tau)\,\abs{\tau}^{-\alpha},\qquad \abs{\tau}\to\infty,
\]
where $L_\gamma$ is a \emph{slowly varying} function and the exponent satisfies
$0<\alpha<1$. Contrast this with ARMA models, whose covariances are either
finite-support ($\acvs_\tau=0$ for $\abs\tau>q$) or exponentially decaying
($\acvs_\tau\sim\abs\phi^{\,\tau}$). The full admissible range $0<\alpha<2$
yields three regimes:
\[
\underbrace{0<\alpha<1}_{\text{long memory}}, \qquad
\underbrace{\alpha=1}_{\text{short-range dependence}}, \qquad
\underbrace{1<\alpha<2}_{\text{anti-persistence}} .
\]
\end{definition}

\begin{definition}{Long memory (spectral definition)}{p9-lmspec}
Equivalently, long memory is characterized by a spectral density that blows up
(or vanishes) like a power of the frequency near the origin:
\[
\sdf(f) \sim L_S(f)\,\abs{f}^{\alpha-1},\qquad f\to0 ,
\]
with $L_S$ slowly varying. For $0<\alpha<1$ the exponent $\alpha-1\in(-1,0)$, so
$\sdf(f)\to\infty$ as $f\to0$ (a spectral pole at zero): the process concentrates
power at the lowest frequencies, the spectral signature of persistence. The
slowly varying parts of the two definitions are linked by
\[
L_S(f) = 2\,\acvs(0)\,\Gamma(1-\alpha)\sin\!\frac{\pi\alpha}{2},\qquad
L_\gamma(\tau) = 2\,\Gamma(\alpha)\sin\!\frac{\pi(1-\alpha)}{2}.
\]
The case $\alpha=1$ (short-range dependence) gives $\sdf(0)$ finite and nonzero,
exactly the ARMA situation $\sdf(0)=\sum_\tau\acvs_\tau<\infty$; the
anti-persistent case $1<\alpha<2$ forces $\sum_k\acvs_k=0$ (and $\sdf(0)=0$).
\end{definition}

\begin{definition}{Fractional differencing operator \& FARIMA}{p9-farima}
The \textbf{fractional differencing operator} is defined for real $d$ via the
binomial expansion of $(1-\Bsh)^d$ in the backshift operator $\Bsh$:
\[
(1-\Bsh)^d
= \sum_{k=0}^{\infty}\binom{d}{k}(-1)^k \Bsh^{\,k}
= \sum_{k=0}^{\infty}\frac{\Gamma(d+1)}{\Gamma(k+1)\,\Gamma(d-k+1)}(-1)^k \Bsh^{\,k}
= \sum_{j=0}^{\infty}\frac{\Gamma(j-d)}{\Gamma(-d)\,\Gamma(j+1)}\Bsh^{\,j}.
\]
The simplest long-memory process is the \textbf{fractionally differenced} (pure)
process $(1-\Bsh)^d X_t=\eps_t$. The general \textbf{FARIMA} (fractional ARIMA)
class is the stationary solution of
\[
(1-\Bsh)^d\,\phi(\Bsh)\,X_t = \psi(\Bsh)\,\eps_t,
\qquad d=\frac{1-\alpha}{2}\in(-0.5,0.5),
\]
where $\eps_t$ are i.i.d.\ zero-mean with variance $\sigma_\eps^2>0$, and
$\phi(z)$, $\psi(z)$ are degree-$p$ and degree-$q$ polynomials with all roots
outside the unit circle. Integer $d$ recovers ordinary ARIMA; fractional
$d\in(0,0.5)$ gives long memory.
\end{definition}

\begin{definition}{Self-similarity \& fractional Brownian motion}{p9-fbm}
A process $\{Y_t\}$ is \textbf{self-similar} with self-similarity parameter
(Hurst exponent) $H$ if, for every $c>0$,
\[
Y_{ct} \eqdist c^{H} Y_t
\]
(equality in distribution, jointly over $t$). Self-similarity with stationary
increments forces the covariance structure
\[
\Cov(Y_t,Y_s) = \frac{\sigma^2}{2}\Big(\abs{t}^{2H}+\abs{s}^{2H}-\abs{t-s}^{2H}\Big).
\tag{$\star$}
\]
A \emph{Gaussian} process with covariance $(\star)$ is \textbf{fractional
Brownian motion} $B_H(t)$; it is the \emph{only} self-similar Gaussian process.
Its increment process $X_t=B_H(t)-B_H(t-1)$ is \textbf{fractional Gaussian
noise}, with
\[
\acvs_X(\tau)=\frac{\sigma^2}{2}\Big(\abs{\tau+1}^{2H}-2\abs{\tau}^{2H}+\abs{\tau-1}^{2H}\Big)
\sim H(2H-1)\,\abs{\tau}^{2H-2},\qquad\abs{\tau}\to\infty .
\]
The graph $(t,B_H(t))$ has Hausdorff (fractal) dimension $D=2-H$ in one
dimension. Fractal dimension is a \emph{local} property; long-range dependence
is a \emph{global} one --- the two can be made completely independent.
\end{definition}

% ----------------------------------------------------------------------
\subsection{Theorems \& Propositions}

\begin{proposition}{Residuals of a correctly specified AR(1) are the noise}{p9-ar1res}
For an AR(1) model $X_t=\phi X_{t-1}+\eps_t$ with known true parameters, the
fitted one-step value is $\hat X_t=\phi X_{t-1}$, so the residual is
\[
e_t = X_t - \phi X_{t-1} = \eps_t .
\]
\textit{Proof idea:} substitute the model equation; the AR structure cancels
exactly, leaving the innovation.\\
\textit{Why:} this is the template for \emph{all} residual diagnostics. If the
model and parameters are right, residuals reproduce the white noise: constant
zero mean, constant variance, uncorrelated, and Gaussian whenever $\eps_t$ is.
Any departure from these is evidence \emph{against} the model.
\end{proposition}

\begin{proposition}{Box--Pierce portmanteau statistic}{p9-boxpierce}
To test $H_0:\acf_1=\acf_2=\cdots=\acf_m=0$ (residuals white up to lag $m$)
against the alternative that some $\acf_\tau\ne0$, the \textbf{Box--Pierce}
statistic is
\[
Q_m = n\sum_{\tau=1}^{m}\hat r_\tau^{\,2}.
\]
Under $H_0$, for an ARMA$(p,q)$ fit, $Q_m$ is approximately
$\chi^2_{m-p-q}$.\\
\textit{Proof idea:} under white noise each $\sqrt n\,\hat r_\tau\todist
\mathcal N(0,1)$ and the $\hat r_\tau$ are asymptotically independent, so
$\sum_\tau(\sqrt n\,\hat r_\tau)^2$ is a sum of squared standard normals; fitting
$p+q$ parameters removes $p+q$ degrees of freedom.\\
\textit{Why:} it tests \emph{many} lags at once (a ``portmanteau''), avoiding the
multiple-testing problem of eyeballing the correlogram lag by lag.
\end{proposition}

\begin{proposition}{Ljung--Box statistic (modified Box--Pierce)}{p9-ljungbox}
The small-sample-corrected (\textbf{Ljung--Box}) statistic is
\[
Q_m = n(n+2)\sum_{\tau=1}^{m}\frac{\hat r_\tau^{\,2}}{n-\tau},
\]
again approximately $\chi^2_{m-p-q}$ under $H_0$ for an ARMA$(p,q)$.\\
\textit{Proof idea:} the factor $\frac{n+2}{n-\tau}$ rescales each term so its
finite-$n$ variance matches the $\chi^2$ target more closely than the plain
$n\hat r_\tau^2$; as $n\to\infty$, $\frac{n(n+2)}{n-\tau}\to n$ and the two
statistics agree.\\
\textit{Why:} the Box--Pierce $\chi^2$ approximation is poor at realistic sample
sizes, so Ljung--Box is the version used in practice (e.g.\ R's
\texttt{Box.test(type="Ljung-Box")}). One usually computes $Q_m$ over a range of
$m$ and plots the resulting $p$-values.
\end{proposition}

\begin{pitfall}
The degrees of freedom are $m-p-q$, \textbf{not} $m$. You must subtract the
number of \emph{estimated} ARMA parameters ($p+q$; the noise variance is
\emph{not} subtracted). Using $\mathrm{df}=m$ inflates the critical value and
makes you reject far too rarely. Also require $m>p+q$ so the df is positive, and
$m$ should be large enough to catch the lags of interest but not so large that
the high-lag $\hat r_\tau$ (estimated from few pairs) add only noise.
\end{pitfall}

\begin{proposition}{The danger of $R^2$: an AR(1) caps at $\phi^2$}{p9-r2danger}
For an AR(1) with innovation variance $\sigma^2$, the process variance is
$s_y^2=\sigma^2/(1-\phi^2)$ and the one-step residual variance is
$s_e^2=\sigma^2$, so
\[
R^2 = 1-\frac{s_e^2}{s_y^2}
    = 1-\frac{\sigma^2(1-\phi^2)}{\sigma^2}
    = \phi^2 .
\]
\textit{Proof idea:} substitute the AR(1) variance formula (geometric sum of the
$\mathrm{MA}(\infty)$ coefficients $\phi^{2j}$) and the fact that the best
one-step prediction error \emph{is} $\eps_t$.\\
\textit{Why:} even the \emph{true} model has bounded $R^2$ ($\phi=0.4\Rightarrow
R^2=0.16$, which ``looks bad'' but is optimal). A stochastic process is not
perfectly predictable, so a low $R^2$ is no indictment, and a high one (from
piling on parameters) is no endorsement. \textbf{Use information criteria, not
$R^2$.}
\end{proposition}

\begin{proposition}{AIC vs.\ AICc gap}{p9-aicgap}
The corrected and uncorrected AIC differ by exactly the parameter-count penalty
\[
\mathrm{AICc}(\theta)-\mathrm{AIC}(\theta)
= 2k\left(\frac{n}{n-k-1}-1\right) = \frac{2k(k+1)}{n-k-1}\ \ (>0).
\]
\textit{Proof idea:} algebra --- put $\frac{n}{n-k-1}-1=\frac{k+1}{n-k-1}$ over a
common denominator.\\
\textit{Why:} the correction is always positive (it penalizes complexity more),
vanishes as $n\to\infty$ with $k$ fixed, and stays negligible whenever
$k^2/n\to0$. AICc is therefore the safe default in small samples or when $k$ is
comparable to $n$.
\end{proposition}

\begin{proposition}{BIC penalty exceeds AIC penalty for $n>e^2$}{p9-bicpen}
The per-parameter penalties are $2$ (AIC) and $\log n$ (BIC). Hence the BIC
penalty is the larger of the two iff
\[
\log n > 2 \iff n > e^2 \approx 7.39 .
\]
\textit{Proof idea:} compare the per-parameter coefficients $2$ and $\log n$
directly.\\
\textit{Why:} for any realistic $n\ (\ge 8)$ BIC charges more per parameter,
hence is more conservative and consistent (selects the true order with
probability $\to1$), whereas AIC is more permissive and tends to over-fit
$p$. Choose BIC when you want parsimony / consistency, AIC(c) for predictive
performance.
\end{proposition}

\begin{proposition}{Long-memory inflation of $\Var(\bar X)$}{p9-varxbar}
For a stationary process with absolutely continuous spectrum,
$\Var(\bar X)\sim \sdf(0)/n$ (the classical $1/n$ rate, since
$\sdf(0)=\sum_\tau\acvs_\tau<\infty$). Under long memory with
$\acvs_\tau\sim\abs\tau^{-\alpha}$ ($0<\alpha<1$), however,
\[
\Var(\bar X) \sim \frac{c}{n^{\alpha}},\qquad 0<\alpha<1 ,
\]
i.e.\ the variance decays \emph{slower} than $1/n$.\\
\textit{Proof idea:} $\Var(\bar X)=\frac1{n^2}\sum_{i,j}\acvs_{j-i}$; aggregating
$\sum_{j=-n}^{n}\acvs_j$ with $\acvs_j\sim\abs j^{-\alpha}$ recovers a power
$n^{1-\alpha}$, so dividing by $n$ gives $n^{-\alpha}$.\\
\textit{Why:} this is the empirical fingerprint of long memory (Smith, 1938):
sample means converge much more slowly than i.i.d.\ or short-range theory
predicts, so naive standard errors $\sqrt{\sdf(0)/n}$ are badly too small.
\end{proposition}

\begin{theorem}{Spectral density of a FARIMA / fractionally differenced process}{p9-farimaspec}
The spectral density of the FARIMA process
$(1-\Bsh)^d\,\phi(\Bsh)X_t=\psi(\Bsh)\eps_t$ is
\[
\sdf(f) = \sigma_\eps^2\,
\frac{\abs{\psi(e^{-2\pi i f})}^2}{\abs{\phi(e^{-2\pi i f})}^2}\,
\abs{1-e^{-2\pi i f}}^{-2d}
\;\overset{f\to0}{\sim}\; c_f\,\abs{f}^{-2d},
\qquad
c_f = \sigma_\eps^2\,\frac{\abs{\psi(1)}^2}{\abs{\phi(1)}^2}.
\]
\textit{Proof idea:} the spectral density of a filtered process multiplies the
input spectrum by the squared transfer function; the fractional difference
contributes $\abs{1-e^{-2\pi i f}}^{-2d}$, and near $f=0$,
$\abs{1-e^{-2\pi i f}}\sim 2\pi\abs{f}$, leaving the pole $\abs{f}^{-2d}$.\\
\textit{Why:} it shows the long-memory parameter $\alpha$ ($d=(1-\alpha)/2$, so
$-2d=\alpha-1$) is read off the slope of $\log\sdf$ vs.\ $\log\abs f$ near the
origin --- the basis of log-periodogram (GPH) estimation of $d$.
\end{theorem}

\begin{intuition}[Where the spectral pole comes from]
Differencing once, $(1-\Bsh)$, has transfer function $1-e^{-2\pi i f}$, which is
\emph{zero} at $f=0$: it kills the trend / mean (a high-pass filter). Taking a
\emph{fractional} negative power, $(1-\Bsh)^{-|d|}$, inverts a fraction of that,
producing a mild \emph{pole} at $f=0$ instead. So a long-memory series is, loosely,
``a little bit non-stationary'': it has more low-frequency power than any ARMA
but not so much that it fails to be stationary (as long as $d<1/2$).
\end{intuition}

\begin{proposition}{Estimating $d$ from the periodogram (log-periodogram idea)}{p9-gph}
Since $\E\,I(f)\approx \sdf(f)\sim L_S(f)\abs f^{\alpha-1}$ near $f=0$, taking
logarithms gives, for low Fourier frequencies $f_j$,
\[
\log I(f_j) \approx \mathrm{const} + (\alpha-1)\log\abs{f_j} + \text{error}
= \mathrm{const} - 2d\,\log\abs{f_j} + \text{error}.
\]
\textit{Proof idea:} the periodogram is an (asymptotically unbiased but
inconsistent) estimate of the spectral density; regressing $\log I(f_j)$ on
$\log\abs{f_j}$ over the lowest frequencies estimates the slope $\alpha-1=-2d$.\\
\textit{Why:} it turns long-memory estimation into a simple linear regression on
the log--log periodogram; one must take $\log$ \emph{before} taking expectations,
because $\E\log I\ne\log\E I$.
\end{proposition}

% ----------------------------------------------------------------------
\subsection{Key Techniques}

\begin{keytech}{Running a Ljung--Box test (df $=m-p-q$)}{p9-ktljung}
Given residuals from a fitted ARMA$(p,q)$ on $n$ observations:
\begin{enumerate}
  \item Compute the residual sample autocorrelations $\hat r_1,\dots,\hat r_m$
        (choose $m>p+q$, often $m\approx 10$ or $\approx\log n$).
  \item Form the Ljung--Box statistic
        $\displaystyle Q_m=n(n+2)\sum_{\tau=1}^{m}\frac{\hat r_\tau^2}{n-\tau}$.
  \item Degrees of freedom: $\mathrm{df}=m-p-q$.
  \item Critical value at level $a$: the upper-$a$ quantile of
        $\chi^2_{\mathrm{df}}$ (in R: \texttt{qchisq(1-a, df)}); reject $H_0$ if
        $Q_m$ exceeds it (equivalently if the $p$-value
        $=\,$\texttt{1-pchisq(Q,df)} is $<a$).
  \item \textbf{Reject} $\Rightarrow$ residual autocorrelation remains
        $\Rightarrow$ the model is inadequate. \textbf{Fail to reject}
        $\Rightarrow$ no evidence against whiteness $\Rightarrow$ model adequate.
\end{enumerate}
\end{keytech}

\begin{keytech}{Choosing between AIC, AICc and BIC}{p9-ktic}
\begin{enumerate}
  \item Fit every candidate model and record $-2\log L$ and $k=p+q+1$.
  \item Compute the desired criterion; \textbf{pick the smallest}.
  \item \textbf{Which one?}
    \begin{itemize}
      \item \textbf{AICc} when $n$ is small or $k$ is comparable to $n$ (it
            corrects AIC's tendency to over-fit $p$); always safe since
            AICc $\to$ AIC for large $n$.
      \item \textbf{BIC} when you want a \emph{parsimonious}, consistent choice
            (its $\log n$ penalty dominates for $n>e^2$).
      \item \textbf{AIC} for predictive/forecasting goals where slight
            over-fitting is tolerable.
    \end{itemize}
  \item For an extra parameter (going from $k-1$ to $k$ estimated parameters) to
        be worth it, the drop in $-2\log L$ must beat the increase in the penalty:
        $2$ (AIC), $\log n$ (BIC), or the (larger) AICc increment
        \[
        \Delta_{\mathrm{AICc}}
        = 2 + \frac{2k(k+1)}{n-k-1} - \frac{2(k-1)k}{n-k}\ \ (>2).
        \]
\end{enumerate}
\end{keytech}

\begin{keytech}{Box--Jenkins model-building loop}{p9-ktboxjenkins}
The Box--Jenkins methodology iterates four steps:
\begin{enumerate}
  \item \textbf{Identify} reasonable $p,d,q$:
        \begin{itemize}
          \item \emph{Choosing $d$:} plot the data; if non-stationary, plot
                differences; use the smallest order of differencing that
                stabilizes the series.
          \item \emph{Choosing $p,q$ (on the differenced data):} read the ACF
                and PACF --- a sharp drop in the ACF at lag $q$ suggests
                MA$(q)$; a sharp drop in the PACF at lag $p$ suggests AR$(p)$;
                no sharp drop suggests a mixed ARMA; very slow decay suggests
                non-stationarity (re-difference).
        \end{itemize}
  \item \textbf{Estimate} the parameters of the proposed ARIMA model.
  \item \textbf{Diagnose} via residuals (mean/variance/QQ/correlogram +
        Ljung--Box); if inadequate, return to step 1.
  \item \textbf{Use} the validated model (forecasting or other inference).
\end{enumerate}
\end{keytech}

\begin{keytech}{Identifying $p,q$ from ACF/PACF cut-offs}{p9-ktidentify}
Memorize the dual table (computed on the appropriately differenced series):
\begin{center}
\begin{tabular}{lll}
\toprule
Model & ACF & PACF \\
\midrule
AR$(p)$  & tails off (decays) & \textbf{cuts off after lag $p$} \\
MA$(q)$  & \textbf{cuts off after lag $q$} & tails off (decays) \\
ARMA$(p,q)$ & tails off & tails off \\
non-stationary & very slow decay & --- \\
\bottomrule
\end{tabular}
\end{center}
``Cuts off'' means the values fall inside the $\pm1.96/\sqrt n$ band beyond the
order; ``tails off'' means a geometric/sinusoidal decay that never abruptly
stops.
\end{keytech}

\begin{keytech}{Reading the long-memory regime from $\alpha$ (or $d$, or $H$)}{p9-ktregime}
The same phenomenon has three equivalent parameters. Convert and classify:
\[
d=\frac{1-\alpha}{2},\qquad H = d+\tfrac12 = 1-\tfrac{\alpha}{2}.
\]
\begin{center}
\begin{tabular}{llll}
\toprule
Regime & $\alpha$ & $d$ & $H$ \\
\midrule
Long memory (persistent)  & $0<\alpha<1$ & $0<d<\tfrac12$  & $\tfrac12<H<1$ \\
Short-range dependence    & $\alpha=1$   & $d=0$           & $H=\tfrac12$ \\
Anti-persistence          & $1<\alpha<2$ & $-\tfrac12<d<0$ & $0<H<\tfrac12$ \\
\bottomrule
\end{tabular}
\end{center}
Signatures: long memory $\Rightarrow$ $\sdf(0)=\infty$ (spectral pole),
$\sum_\tau\acvs_\tau=\infty$; anti-persistence $\Rightarrow$ $\sdf(0)=0$ and
$\sum_k\acvs_k=0$; short-range $\Rightarrow$ $\sdf(0)$ finite, nonzero (ARMA).
\end{keytech}

% ----------------------------------------------------------------------
\subsection{Exercises}

\begin{exercise}{When are AIC and AICc equivalent? \quad\textnormal{[Sheet 11]}}{p9-ex1}
The criteria are
$\mathrm{AIC}=-2\ln L+2k$ and
$\mathrm{AICc}=-2\ln L+2k\,\frac{n}{n-k-1}$.
Explain under what conditions on $n$ and $k$ the two become approximately
equivalent.
\end{exercise}
\begin{solution}
Subtract:
\[
\mathrm{AICc}-\mathrm{AIC}
= 2k\!\left(\frac{n}{n-k-1}-1\right)
= 2k\cdot\frac{n-(n-k-1)}{n-k-1}
= 2k\cdot\frac{k+1}{n-k-1}
= \frac{2k(k+1)}{n-k-1}.
\]
The two criteria are approximately equivalent exactly when this correction term
is small. In particular:
\begin{itemize}
  \item If $k$ is \textbf{fixed} and $n\to\infty$, then
        $\dfrac{2k(k+1)}{n-k-1}\to 0$, so AICc $\to$ AIC.
  \item If $k=k_n$ grows with $n$, the correction is still small provided $n$ is
        large relative to $k_n^2$, e.g.\ $k_n^2/n\to0$. If $k$ grows too fast
        ($k_n^2/n\not\to0$) the gap need not vanish and AICc stays meaningfully
        larger than AIC.
\end{itemize}
So: \emph{large samples relative to the number of parameters} make the two
agree; small samples or richly parameterized models make AICc strictly more
conservative.
\end{solution}

\begin{exercise}{AICc with $m$ coefficients fixed to zero \quad\textnormal{[Sheet 11]}}{p9-ex2}
Determine the form of AICc for a mean-zero ARMA$(p,q)$ model when exactly $m$ of
the AR and MA coefficients are known in advance and fixed to zero.
\end{exercise}
\begin{solution}
If $m$ of the $p+q$ AR/MA coefficients are fixed (not estimated), only $p+q-m$
of them are estimated. Adding the innovation variance $\sigma^2$, the
estimated-parameter count is
\[
k = p+q+1-m .
\]
Substituting into the AICc formula,
\[
\mathrm{AICc}(\theta)
= -2\ln L(\theta\mid y)
+ 2(p+q+1-m)\,\frac{n}{\,n-(p+q+1-m)-1\,}.
\]
(For the mean-zero model no separate mean parameter is counted.) Fixing
coefficients to zero \emph{lowers} $k$, which both reduces the penalty and
\emph{increases} the residual degrees of freedom in any subsequent Ljung--Box
test.
\end{solution}

\begin{exercise}{Compute $Q_4$ and test (Ljung--Box) \quad\textnormal{[Sheet 11]}}{p9-ex3}
For $n=150$ observations a mean-zero ARMA$(1,1)$ is fitted; the residual
autocorrelations at lags $1$--$4$ are
\[
\hat r_1=0.03,\quad \hat r_2=-0.09,\quad \hat r_3=0.11,\quad \hat r_4=-0.05 .
\]
(a) Compute the Ljung--Box statistic $Q_4$.
(b) Test $H_0:\acf_1=\cdots=\acf_4=0$ at the $5\%$ level (state the df, the
critical value, and the conclusion).
(c) Would the conclusion change for an AR$(1)$ fit?
\end{exercise}
\begin{solution}
\textbf{(a)} With $n=150$, $n(n+2)=150\cdot152=22800$, and $m=4$:
\[
Q_4 = 22800\left(\frac{0.03^2}{149}+\frac{(-0.09)^2}{148}
+\frac{0.11^2}{147}+\frac{(-0.05)^2}{146}\right).
\]
Term by term,
\[
\frac{0.0009}{149}=6.04\times10^{-6},\quad
\frac{0.0081}{148}=5.47\times10^{-5},\quad
\frac{0.0121}{147}=8.23\times10^{-5},\quad
\frac{0.0025}{146}=1.71\times10^{-5}.
\]
Summing: $6.04\times10^{-6}+5.47\times10^{-5}+8.23\times10^{-5}+1.71\times10^{-5}
=1.60\times10^{-4}$. Hence
\[
Q_4 \approx 22800\times 1.60\times10^{-4} \approx 3.65 .
\]
\textbf{(b)} For ARMA$(1,1)$, $p=q=1$, so $\mathrm{df}=m-p-q=4-1-1=2$. The $5\%$
critical value is the upper-$5\%$ point of $\chi^2_2$, i.e.\
\texttt{qchisq(0.95, df=2)} $=5.991$. Since $Q_4\approx3.65<5.991$, we
\textbf{do not reject} $H_0$: there is no significant evidence of residual
autocorrelation, and the ARMA$(1,1)$ appears adequate.\\
\textbf{(c)} For AR$(1)$, $p=1,q=0$, so $\mathrm{df}=4-1-0=3$ and the critical
value rises to \texttt{qchisq(0.95, df=3)} $=7.815$. Since $3.65<7.815$ the
conclusion is unchanged (do not reject). In general, fewer fitted parameters
leave \emph{more} degrees of freedom, widening the acceptance region.
\end{solution}

\begin{exercise}{AIC vs.\ BIC penalty comparison \quad\textnormal{[Sheet 11]}}{p9-ex4}
With $\mathrm{AIC}=-2\log L+2k$ and $\mathrm{BIC}=-2\log L+k\log n$:
(a) Show the BIC penalty exceeds the AIC penalty iff $n>e^2\approx7.4$.
(b) For fixed $\log L$, a $k=3$ model competes with a $k=2$ model at $n=50$.
Which criterion is more likely to select the larger model?
\end{exercise}
\begin{solution}
\textbf{(a)} Per parameter the penalties are $2$ (AIC) and $\log n$ (BIC), so
\[
\log n > 2 \iff n > e^2 \approx 7.39 .
\]
For any realistic $n\ (\ge 8)$ BIC charges the larger penalty per parameter and
hence favours more parsimonious models.\\
\textbf{(b)} Going from $k=2$ to $k=3$ adds one parameter. The extra penalty is
\[
\text{AIC: } 2\times1=2,\qquad
\text{BIC: } \log(50)\approx 3.91 .
\]
So AIC selects the larger model whenever $-2\log L$ drops by more than $2$ (i.e.\
the log-likelihood improves by $>1$), while BIC requires the drop to exceed
$3.91$ (log-likelihood improvement $>1.96$). \textbf{AIC is more willing to
select the larger model}; BIC is harder to satisfy and protects against
over-fitting.
\end{solution}

\begin{exercise}{Box--Pierce vs.\ Ljung--Box on the same data \quad\textnormal{[New]}}{p9-ex5}
Using the data of Exercise~\ref{exo:p9-ex3} ($n=150$, $\hat r_{1:4}=
0.03,-0.09,0.11,-0.05$, ARMA$(1,1)$ fit), compute the \emph{Box--Pierce}
statistic $Q_4^{\mathrm{BP}}$ and compare it with the Ljung--Box value
$Q_4\approx3.65$. Does the test decision change?
\end{exercise}
\begin{solution}
The Box--Pierce statistic is $Q_m^{\mathrm{BP}}=n\sum_{\tau=1}^m\hat r_\tau^2$:
\[
\sum_{\tau=1}^4 \hat r_\tau^2 = 0.03^2+0.09^2+0.11^2+0.05^2
= 0.0009+0.0081+0.0121+0.0025 = 0.0236 .
\]
Hence
\[
Q_4^{\mathrm{BP}} = 150\times 0.0236 = 3.54 .
\]
This is slightly smaller than the Ljung--Box $Q_4\approx3.65$ (Ljung--Box
inflates each term by $\frac{n+2}{n-\tau}>1$). Both use $\mathrm{df}=4-1-1=2$ and
the same critical value $5.991$; since $3.54<5.991$ and $3.65<5.991$, \textbf{both
fail to reject} $H_0$. With $n=150$ the two agree closely; the difference matters
only at small $n$, where the Box--Pierce $\chi^2$ approximation is unreliable and
Ljung--Box is preferred.
\end{solution}

\begin{exercise}{$R^2$ of a fitted AR(1) \quad\textnormal{[New]}}{p9-ex6}
An AR$(1)$ model is fitted to a long stationary series and the estimated
coefficient is $\hat\phi=0.7$. (a) What value of $R^2$ should you expect?
(b) A colleague adds three more lags (fitting AR$(4)$) and reports
$R^2=0.55$. Has the model genuinely improved? How would you decide?
\end{exercise}
\begin{solution}
\textbf{(a)} For an AR$(1)$ the (population) $R^2$ equals $\phi^2$:
\[
R^2 = 1-\frac{s_e^2}{s_y^2}
    = 1-\frac{\sigma^2}{\sigma^2/(1-\phi^2)} = \phi^2 = 0.7^2 = 0.49 .
\]
So roughly $R^2\approx 0.49$ even for the correct, optimal model --- a
``modest'' value that is in fact the best attainable.\\
\textbf{(b)} A rise from $0.49$ to $0.55$ is \emph{not} evidence of a better
model: adding parameters can only increase the in-sample $R^2$. To decide,
compare \emph{information criteria}: compute AICc (or BIC) for the AR$(1)$ and
AR$(4)$ and pick the smaller. Equivalently, test whether the extra coefficients
are significant and whether the AR$(1)$ residuals already pass a Ljung--Box test
--- if they do, the extra lags are spurious and the parsimonious AR$(1)$ wins.
\end{solution}

\begin{exercise}{Classify a process by its long-memory parameter \quad\textnormal{[New]}}{p9-ex7}
A stationary series has ACVS $\acvs_\tau\sim c\,\abs\tau^{-0.4}$ as
$\abs\tau\to\infty$. (a) Find $\alpha$, the fractional difference parameter $d$,
and the Hurst exponent $H$. (b) Classify the dependence regime. (c) Is
$\sum_\tau\acvs_\tau$ finite? What does $\sdf(f)$ do near $f=0$?
\end{exercise}
\begin{solution}
\textbf{(a)} Matching $\acvs_\tau\sim L_\gamma(\tau)\abs\tau^{-\alpha}$ gives
$\alpha=0.4$. Then
\[
d=\frac{1-\alpha}{2}=\frac{1-0.4}{2}=0.3,\qquad
H=d+\tfrac12=0.8\ \ \big(=1-\tfrac{\alpha}{2}=1-0.2\big).
\]
\textbf{(b)} Since $0<\alpha=0.4<1$ (equivalently $0<d=0.3<\tfrac12$, or
$\tfrac12<H=0.8<1$), the process exhibits \textbf{long memory / positive
persistence}.\\
\textbf{(c)} The tail $\abs\tau^{-0.4}$ is \emph{not} summable
($\sum\abs\tau^{-\alpha}$ diverges for $\alpha\le1$), so
$\sum_\tau\acvs_\tau=\infty$. Correspondingly $\sdf(f)\sim L_S(f)\abs
f^{\alpha-1}=L_S(f)\abs f^{-0.6}\to\infty$ as $f\to0$: a spectral pole at the
origin, the hallmark of long memory.
\end{solution}

\begin{exercise}{Variance of $B_H$ and the self-similarity scaling \quad\textnormal{[New]}}{p9-ex8}
For fractional Brownian motion $B_H(t)$ with covariance
$\Cov(B_H(t),B_H(s))=\frac{\sigma^2}{2}(\abs t^{2H}+\abs s^{2H}-\abs{t-s}^{2H})$:
(a) compute $\Var(B_H(t))$; (b) verify the self-similarity relation
$B_H(ct)\eqdist c^H B_H(t)$ at the level of variances; (c) for $H=\tfrac12$
identify the process.
\end{exercise}
\begin{solution}
\textbf{(a)} Set $s=t$ in the covariance:
\[
\Var(B_H(t))=\Cov(B_H(t),B_H(t))
=\frac{\sigma^2}{2}\big(\abs t^{2H}+\abs t^{2H}-0\big)
=\sigma^2\abs t^{2H}.
\]
\textbf{(b)} Then $\Var(B_H(ct))=\sigma^2\abs{ct}^{2H}=c^{2H}\sigma^2\abs t^{2H}
=\Var(c^H B_H(t))$, consistent with $B_H(ct)\eqdist c^H B_H(t)$. (For a centred
Gaussian process, matching the full covariance structure --- which scales by
$c^{2H}$ in every entry --- gives equality in distribution.)\\
\textbf{(c)} For $H=\tfrac12$, $\Var(B_{1/2}(t))=\sigma^2\abs t$ and
$\Cov(B_{1/2}(t),B_{1/2}(s))=\frac{\sigma^2}{2}(\abs t+\abs s-\abs{t-s})
=\sigma^2\min(t,s)$ for $t,s\ge0$ --- exactly \textbf{standard Brownian motion}.
Its increments are ordinary (uncorrelated) white noise: $H=\tfrac12$ is the
short-range / memoryless boundary ($d=0$, $\alpha=1$).
\end{solution}

% ======================================================================
%  PART 10 -- Non-Linear Time Series: ARCH and GARCH
% ======================================================================
\section{Non-Linear Time Series: ARCH and GARCH}
\label{sec:archgarch}

Linear models (AR, MA, ARMA) describe the conditional \emph{mean} of a series and
assume a constant innovation variance. Financial log-returns violate this: they
are nearly uncorrelated (so a linear conditional-mean model is almost useless),
yet their \emph{squares} are strongly dependent --- large moves cluster with
large moves and quiet periods with quiet periods. This \textbf{volatility
clustering}, together with heavy (\textbf{fat}) tails, motivates models for the
\emph{conditional variance}. \textbf{ARCH} and \textbf{GARCH} keep the series
mean-zero and uncorrelated (white noise) while letting the conditional variance
$\sigma_t^2$ evolve in time as a function of past squared returns (and, for
GARCH, past variances). Throughout, $\mathcal{F}_t$ denotes the information
available up to and including time $t$, and $\{\eps_t\}$ is white noise with
$\E[\eps_t]=0$, $\Var(\eps_t)=1$, independent of $\mathcal{F}_{t-1}$.

% ----------------------------------------------------------------------
\subsection{Definitions}

\begin{definition}{Log-returns \& volatility clustering}{p10-logret}
For a price series $\{X_t\}$ (e.g.\ a stock index), the \textbf{log-returns} are
\[
r_t=\log\!\Big(\frac{X_t}{X_{t-1}}\Big)=\log(X_t)-\log(X_{t-1}),
\]
measuring relative changes in price. Empirically log-returns show:
\begin{itemize}
  \item \textbf{little linear dependence:} the sample ACF/PACF of $r_t$ is
        essentially that of white noise;
  \item \textbf{dependent squares:} the sample ACF/PACF of $r_t^2$ decays slowly
        $\Rightarrow$ the conditional variance changes over time;
  \item \textbf{volatility clustering:} large changes tend to be followed by
        large changes, small by small;
  \item \textbf{fat tails:} kurtosis $>3$ (heavier tails than the normal),
        visible as deviations in a QQ-plot against $\mathcal{N}(0,1)$.
\end{itemize}
\end{definition}

\begin{definition}{ARCH(1) model}{p10-arch1}
$\{r_t\}$ is an \textbf{AutoRegressive Conditionally Heteroscedastic} model of
order $1$, $\mathrm{ARCH}(1)$, if for $t\in\Ints$
\[
r_t=\sigma_t\eps_t,\qquad \sigma_t^2=\alpha_0+\alpha_1 r_{t-1}^2,
\]
with $\{\eps_t\}$ white noise of mean $0$ and variance $1$, independent of
$\mathcal F_{t-1}$, and parameters $\alpha_0>0,\ \alpha_1\ge 0$. The volatility
$\sigma_t$ is $\mathcal F_{t-1}$-measurable (known one step ahead). Immediate
properties (Proposition~\ref{prop:p10-archmoments}):
\[
\E[r_t\mid\mathcal F_{t-1}]=0,\qquad
\Var(r_t\mid\mathcal F_{t-1})=\sigma_t^2=\alpha_0+\alpha_1 r_{t-1}^2 .
\]
\end{definition}

\begin{definition}{ARCH(p) model}{p10-archp}
$\{r_t\}$ is an $\mathrm{ARCH}(p)$ model if $r_t=\sigma_t\eps_t$ with
\[
\sigma_t^2=\alpha_0+\alpha_1 r_{t-1}^2+\cdots+\alpha_p r_{t-p}^2
=\alpha_0+\sum_{j=1}^{p}\alpha_j r_{t-j}^2,
\]
$\{\eps_t\}$ white noise $(0,1)$ independent of $\mathcal F_{t-1}$, and
$\alpha_0>0,\ \alpha_p>0,\ \alpha_j\ge0$ for $1\le j<p$. Volatility may now
depend on several past squared shocks, so as before
$\Var(r_t\mid\mathcal F_{t-1})=\sigma_t^2$.
\end{definition}

\begin{definition}{GARCH(p,q) model}{p10-garch}
$\{r_t\}$ is a \textbf{Generalised ARCH} model $\mathrm{GARCH}(p,q)$ if
$r_t=\sigma_t\eps_t$ with
\[
\sigma_t^2=\alpha_0+\sum_{j=1}^{p}\alpha_j r_{t-j}^2+\sum_{j=1}^{q}\beta_j\sigma_{t-j}^2,
\]
$\{\eps_t\}$ white noise $(0,1)$ independent of $\mathcal F_{t-1}$, and parameters
$\alpha_0>0$, $\alpha_p>0$, $\beta_q>0$, with $\alpha_j,\beta_j\ge0$ otherwise.
Feeding back past \emph{variances} $\sigma_{t-j}^2$ lets volatility persist over
time far more parsimoniously than an $\mathrm{ARCH}(p)$ with large $p$. A
$\mathrm{GARCH}(p,0)$ is exactly an $\mathrm{ARCH}(p)$. As before
$\Var(r_t\mid\mathcal F_{t-1})=\sigma_t^2$.
\end{definition}

\begin{intuition}[Why feed back $\sigma_{t-1}^2$?]
In $\mathrm{ARCH}(1)$ the conditional variance depends only on the \emph{single}
most recent squared return $r_{t-1}^2$. One quiet day ($r_{t-1}^2$ small)
instantly collapses $\sigma_t^2$ back toward $\alpha_0$, so the model cannot
sustain long high-volatility regimes. GARCH adds the smooth, slowly varying term
$\beta_1\sigma_{t-1}^2$: yesterday's \emph{variance estimate} is mostly carried
forward, so volatility is sticky. In practice fitted $\mathrm{GARCH}(1,1)$ models
have $\alpha_1+\beta_1$ close to $1$, giving long memory in volatility (still
geometric decay, but very slow).
\end{intuition}

\begin{definition}{Standardised residuals (diagnostics)}{p10-resid}
After fitting, the \textbf{standardised residuals} are
$\tilde e_t=r_t/\hat\sigma_t$. If the model is correct these should look like
i.i.d.\ $(0,1)$ noise: their ACF/PACF \emph{and} the ACF/PACF of their
\emph{squares} $\tilde e_t^2$ should resemble white noise, and a QQ-plot against
the assumed innovation law should be roughly straight. Under Gaussian
innovations one also overlays in-sample bands $\pm1.96\,\hat\sigma_t$ on $r_t$.
\end{definition}

\begin{definition}{ARMA--GARCH model}{p10-armagarch}
Because a GARCH process is itself white noise, it can serve as the innovation of
an ARMA model: with $r_t$ a $\mathrm{GARCH}(\tilde p,\tilde q)$ process,
\[
X_t=\sum_{j=1}^{p}\phi_j X_{t-j}-\sum_{k=0}^{q}\theta_k\,r_{t-k},
\qquad \theta_0=-1,
\]
where the $\phi_j,\theta_k$ satisfy the usual $\mathrm{ARMA}(p,q)$ stationarity
and invertibility conditions. This captures a conditional mean (ARMA) and a
time-varying conditional variance (GARCH) at once; one must only be careful with
the likelihood, as the errors are no longer i.i.d.\ Gaussian.
\end{definition}

% ----------------------------------------------------------------------
\subsection{Theorems \& Propositions}

\begin{proposition}{Conditional and unconditional moments of ARCH(1)}{p10-archmoments}
For the $\mathrm{ARCH}(1)$ model with $\alpha_1<1$:
\[
\E[r_t\mid\mathcal F_{t-1}]=0,\quad
\Var(r_t\mid\mathcal F_{t-1})=\sigma_t^2,\quad
\E[r_t]=0,\quad
\Var(r_t)=\frac{\alpha_0}{1-\alpha_1}.
\]
\textit{Proof idea:} $\sigma_t$ is $\mathcal F_{t-1}$-measurable and $\eps_t$ is
independent of $\mathcal F_{t-1}$, so
$\E[r_t\mid\mathcal F_{t-1}]=\sigma_t\E[\eps_t]=0$ and
$\E[r_t^2\mid\mathcal F_{t-1}]=\sigma_t^2\E[\eps_t^2]=\sigma_t^2$. Iterated
expectation gives $\E[r_t]=\E[\E[r_t\mid\mathcal F_{t-1}]]=0$ and
$\Var(r_t)=\E[\sigma_t^2]=\alpha_0+\alpha_1\Var(r_{t-1})$; under stationarity
$\Var(r_t)=\Var(r_{t-1})$, solve to get $\alpha_0/(1-\alpha_1)$.\\
\textit{Why:} establishes that ARCH is mean-zero with a positive finite
unconditional variance exactly when $\alpha_1<1$ (necessary \emph{and}, for
$\mathrm{ARCH}(1)$, sufficient for weak stationarity).
\end{proposition}

\begin{proposition}{ARCH/GARCH is white noise (uncorrelated but dependent)}{p10-whitenoise}
Any $\mathrm{ARCH}(p)$ or $\mathrm{GARCH}(p,q)$ process is mean-zero and serially
uncorrelated: for $\tau\ne0$,
\[
\E[r_t]=0,\qquad \Cov(r_{t+\tau},r_t)=0,\qquad\text{hence}\qquad \Corr(r_{t+\tau},r_t)=0 .
\]
\textit{Proof idea:} take $\tau>0$. Since $r_t$ is
$\mathcal F_{t+\tau-1}$-measurable and $\sigma_{t+\tau}$ is too, condition on
$\mathcal F_{t+\tau-1}$:
\[
\Cov(r_{t+\tau},r_t)=\E[r_{t+\tau}r_t]
=\E\!\big[\,r_t\,\sigma_{t+\tau}\,\E[\eps_{t+\tau}\mid\mathcal F_{t+\tau-1}]\big]
=\E[r_t\sigma_{t+\tau}\cdot 0]=0 .
\]
\textit{Why:} the central qualitative fact --- ARCH/GARCH leaves the
\emph{linear} correlation structure flat (white noise) while creating strong
\emph{nonlinear} dependence in the squares. The conditional-mean and
conditional-variance modelling problems decouple.
\end{proposition}

\begin{theorem}{Isserlis' theorem (Wick formula)}{p10-isserlis}
If $(X_1,\dots,X_N)$ is a zero-mean multivariate normal vector and $N$ is even,
\[
\E[X_1\cdots X_N]=\sum_{p\in P_N}\ \prod_{\{i,j\}\in p}\Cov(X_i,X_j),
\]
where $P_N$ is the set of all pairings (perfect matchings) of $\{1,\dots,N\}$.
For $N=4$,
\[
\E[X_1X_2X_3X_4]=\E[X_1X_2]\E[X_3X_4]+\E[X_1X_3]\E[X_2X_4]+\E[X_1X_4]\E[X_2X_3].
\]
\textit{Proof idea:} expand the Gaussian moment generating function / cumulant
expansion; only pair-cumulants (covariances) survive for a Gaussian.\\
\textit{Why:} gives $\E[\eps_t^4]=3(\E[\eps_t^2])^2=3$ for standard Gaussian
innovations, the key input to the ARCH fourth-moment and kurtosis computations
below.
\end{theorem}

\begin{proposition}{Fourth moment and kurtosis of ARCH(1)}{p10-kurtosis}
For $\mathrm{ARCH}(1)$ with Gaussian innovations, the unconditional fourth moment
is finite iff $0\le\alpha_1^2<\tfrac13$, in which case
\[
\E[r_t^4]=\frac{3\alpha_0^2(1+\alpha_1)}{(1-\alpha_1)(1-3\alpha_1^2)},
\qquad
\kappa=\frac{\E[r_t^4]}{\Var(r_t)^2}=\frac{3\,(1-\alpha_1^2)}{1-3\alpha_1^2}\ >\ 3
\quad(\alpha_1>0).
\]
\textit{Proof idea:} $\E[r_t^4\mid\mathcal F_{t-1}]
=\sigma_t^4\,\E[\eps_t^4\mid\mathcal F_{t-1}]=3\sigma_t^4=3(\alpha_0+\alpha_1 r_{t-1}^2)^2$
(Isserlis). Take expectations, substitute $\E[r_{t-1}^2]=\alpha_0/(1-\alpha_1)$
and $\E[r_{t-1}^4]=\E[r_t^4]$ by stationarity, and solve the linear equation;
divide by $\Var(r_t)^2=\alpha_0^2/(1-\alpha_1)^2$ for $\kappa$.\\
\textit{Why:} shows ARCH \emph{generates} excess kurtosis (fat tails) endogenously,
even with Gaussian $\eps_t$ --- matching the heavy tails seen in returns.
\end{proposition}

\begin{intuition}[Two nested constraints]
$\alpha_1<1$ buys a finite \emph{variance}; the strictly tighter
$\alpha_1^2<\tfrac13$ (i.e.\ $\alpha_1<1/\sqrt3\approx0.577$) buys a finite
\emph{fourth} moment. As $\alpha_1\uparrow 1/\sqrt3$ the kurtosis
$\to\infty$: more ARCH effect $\Rightarrow$ heavier tails. The fourth moment is
the relevant one whenever you want the ACF of the squares to be well-defined.
\end{intuition}

\begin{proposition}{ACF of the squares of an ARCH(1)}{p10-acfsq}
For $\mathrm{ARCH}(1)$ with $0<\alpha_0$ and $0<\alpha_1<1/\sqrt3$,
\[
\Corr(r_{t+\tau}^2,r_t^2)=\alpha_1^{\abs\tau},\qquad \tau\in\Ints,
\]
i.e.\ the autocorrelation of the squared series decays \emph{geometrically}.\\
\textit{Proof idea:} condition on $\mathcal F_{t+\tau-1}$ to peel one step:
$\E[r_{t+\tau}^2 r_t^2\mid\mathcal F_{t+\tau-1}]=r_t^2\sigma_{t+\tau}^2
=r_t^2(\alpha_0+\alpha_1 r_{t+\tau-1}^2)$. Taking expectations and subtracting
$(\E[r_t^2])^2$ yields the recursion
$\Cov(r_{t+\tau}^2,r_t^2)=\alpha_1\Cov(r_{t+\tau-1}^2,r_t^2)$ (the constant terms
cancel using $\E[r_t^2]=\alpha_0/(1-\alpha_1)$). Iterating gives
$\Cov(r_{t+\tau}^2,r_t^2)=\alpha_1^{\tau}\Var(r_t^2)$, hence the stated
correlation.\\
\textit{Why:} explains the slowly decaying ACF of squared returns and warns that
ARCH(1) imposes \emph{purely geometric} memory in volatility --- often too fast
for real data, motivating GARCH.
\end{proposition}

\begin{proposition}{Stationarity \& variance of GARCH(p,q)}{p10-garchstat}
A $\mathrm{GARCH}(p,q)$ process has a weakly stationary solution \textbf{iff}
\[
\sum_{j=1}^{p}\alpha_j+\sum_{j=1}^{q}\beta_j<1,
\]
and then its unconditional variance is
\[
\Var(r_t)=\frac{\alpha_0}{1-\sum_{j=1}^{p}\alpha_j-\sum_{j=1}^{q}\beta_j}.
\]
\textit{Proof idea:} take unconditional expectations of
$\sigma_t^2=\alpha_0+\sum_j\alpha_j r_{t-j}^2+\sum_j\beta_j\sigma_{t-j}^2$ using
$\E[r_{t-j}^2]=\E[\sigma_{t-j}^2]=\Var(r_t)$ at stationarity, then solve; positivity
of the denominator forces the stated bound.\\
\textit{Why:} the master condition for fitting GARCH --- it must hold for the
fitted parameters, and the closed-form variance gives the long-run volatility
level. (More: $\sum\alpha+\sum\beta<1$ is sufficient for strict stationarity; for
some noise laws including Gaussian, the boundary $\sum\alpha+\sum\beta=1$ still
admits a strictly --- but not weakly --- stationary solution: the IGARCH case.)
\end{proposition}

\begin{proposition}{Fourth-moment condition for GARCH(1,1)}{p10-garch4}
For a $\mathrm{GARCH}(1,1)$ with Gaussian noise, the unconditional fourth moment
$\E[r_t^4]$ exists \textbf{iff}
\[
3\alpha_1^2+2\alpha_1\beta_1+\beta_1^2<1 .
\]
\textit{Proof idea:} square the recursion for $\sigma_t^2$, take expectations
using $\E[\eps_t^4]=3$ (Isserlis) and stationarity of $\E[\sigma_t^4]$; the
coefficient of $\E[\sigma_{t}^4]$ on the right is $3\alpha_1^2+2\alpha_1\beta_1+\beta_1^2$,
which must be $<1$ for a finite solution.\\
\textit{Why:} the analogue of $\alpha_1^2<1/3$; needed before quoting any
formula for the kurtosis or the ACF of squared GARCH returns.
\end{proposition}

\begin{proposition}{Maximum likelihood estimation (conditional Gaussian)}{p10-mle}
Conditioning on the first observation $r_1$ to avoid specifying a marginal law
for the initial value, the conditional likelihood factorises as
\[
L_c(\theta\mid r_1,\dots,r_n)=\prod_{t=2}^{n}f(r_t\mid r_{t-1},\dots,r_1;\theta).
\]
Under Gaussian innovations the $\mathrm{ARCH}(1)$ model gives
$r_t\mid\mathcal F_{t-1}\sim\mathcal N(0,\ \alpha_0+\alpha_1 r_{t-1}^2)$, so
\[
L_c(\alpha_0,\alpha_1)=\prod_{t=2}^{n}\frac{1}{\sqrt{2\pi(\alpha_0+\alpha_1 r_{t-1}^2)}}
\exp\!\left(-\frac{r_t^2}{2(\alpha_0+\alpha_1 r_{t-1}^2)}\right),
\]
maximised over $\alpha_0>0,\ \alpha_1\ge0$ (numerically; GARCH is analogous with
$\sigma_t^2$ computed by its recursion).\\
\textit{Proof idea:} the chain rule for densities plus conditional normality of
each $r_t$ given the past.\\
\textit{Why:} the standard estimation route for ARCH/GARCH; the per-term variance
$\sigma_t^2$ is what is being fitted, not a constant $\sigma^2$.
\end{proposition}

\begin{intuition}[Reading the (G)ARCH likelihood]
Each factor is a normal density with its \emph{own} variance $\sigma_t^2$. The
log-likelihood is $-\tfrac12\sum_{t}\big[\log\sigma_t^2+r_t^2/\sigma_t^2\big]$
(up to constants): the fit trades off ``don't make $\sigma_t^2$ needlessly large''
(the $\log\sigma_t^2$ penalty) against ``explain big $r_t^2$ with big
$\sigma_t^2$'' (the $r_t^2/\sigma_t^2$ term). That is exactly volatility tracking.
\end{intuition}

\begin{intuition}[Summary: pros \& cons]
\textbf{Pros:} ARCH/GARCH capture volatility clusters and time-varying
conditional variance; they reproduce fat tails (large kurtosis) and large shocks;
GARCH gives persistent volatility parsimoniously; and crucially they do
\emph{not} disturb the (white-noise) correlation structure, so they slot in as
ARMA innovations. \textbf{Cons:} positive and negative shocks affect volatility
\emph{symmetrically} (no leverage effect); volatility dependence decays only
\emph{geometrically}, often too fast for real markets. Extensions (EGARCH,
GJR-GARCH, long-memory GARCH) address these but are out of scope here.
\end{intuition}

% ----------------------------------------------------------------------
\subsection{Key Techniques}

\begin{keytech}{Unconditional variance \& stationarity of any (G)ARCH}{p10-ktvar}
To get the long-run variance and the stationarity condition in one move:
\begin{enumerate}
  \item Write the $\sigma_t^2$ recursion and take \emph{unconditional}
        expectations.
  \item Replace, by stationarity, $\E[r_{t-j}^2]=\E[\sigma_{t-j}^2]=\Var(r_t)=:v$
        (this uses $\E[r_{t-j}^2\mid\mathcal F_{t-j-1}]=\sigma_{t-j}^2$ and iterated
        expectation, so $\E[r_{t-j}^2]=\E[\sigma_{t-j}^2]$).
  \item Solve the scalar equation $v=\alpha_0+\big(\sum_j\alpha_j+\sum_j\beta_j\big)v$:
        \[
        v=\frac{\alpha_0}{1-\sum_j\alpha_j-\sum_j\beta_j}.
        \]
  \item Positivity ($v>0$, finite) $\Leftrightarrow$
        $\sum_j\alpha_j+\sum_j\beta_j<1$ --- the stationarity condition.
\end{enumerate}
Special cases: $\mathrm{ARCH}(1)$ gives $\alpha_0/(1-\alpha_1)$;
$\mathrm{ARCH}(p)$ gives $\alpha_0/(1-\sum_{j=1}^p\alpha_j)$;
$\mathrm{GARCH}(1,1)$ gives $\alpha_0/(1-\alpha_1-\beta_1)$.
\end{keytech}

\begin{keytech}{Proving uncorrelatedness via $\mathcal F_{t-1}$-measurability}{p10-ktwn}
To show $\Cov(r_{t+\tau},r_t)=0$ for $\tau>0$ in any (G)ARCH:
\begin{enumerate}
  \item Note $\E[r_t]=0$, so $\Cov(r_{t+\tau},r_t)=\E[r_{t+\tau}r_t]$.
  \item Condition on $\mathcal F_{t+\tau-1}$ (which contains $r_t$ and $\sigma_{t+\tau}$):
        \[
        \E[r_{t+\tau}r_t]=\E\big[\,\E[\sigma_{t+\tau}\eps_{t+\tau}r_t\mid\mathcal F_{t+\tau-1}]\big].
        \]
  \item Pull out the $\mathcal F_{t+\tau-1}$-measurable factors $r_t,\sigma_{t+\tau}$:
        $=\E\big[r_t\sigma_{t+\tau}\,\E[\eps_{t+\tau}\mid\mathcal F_{t+\tau-1}]\big]
        =\E[r_t\sigma_{t+\tau}\cdot0]=0$.
\end{enumerate}
The exact same template, with $\eps_{t+\tau}^2$ in place of $\eps_{t+\tau}$ (and
$\E[\eps^2\mid\mathcal F]=1$), peels one lag off $\Cov(r_{t+\tau}^2,r_t^2)$.
\end{keytech}

\begin{keytech}{Geometric recursion for the ACF of squares}{p10-ktacfsq}
To get $\Corr(r_{t+\tau}^2,r_t^2)$ for $\mathrm{ARCH}(1)$:
\begin{enumerate}
  \item Condition on $\mathcal F_{t+\tau-1}$ and use $\E[\eps^2\mid\cdot]=1$,
        $\sigma_{t+\tau}^2=\alpha_0+\alpha_1 r_{t+\tau-1}^2$:
        \[
        \E[r_{t+\tau}^2 r_t^2\mid\mathcal F_{t+\tau-1}]
        =r_t^2(\alpha_0+\alpha_1 r_{t+\tau-1}^2).
        \]
  \item Take expectations, subtract $(\E[r_t^2])^2$, and insert
        $\E[r_t^2]=\alpha_0/(1-\alpha_1)$; the non-covariance terms cancel,
        leaving $\Cov(r_{t+\tau}^2,r_t^2)=\alpha_1\Cov(r_{t+\tau-1}^2,r_t^2)$.
  \item Iterate $\tau$ times: $\Cov(r_{t+\tau}^2,r_t^2)=\alpha_1^{\tau}\Var(r_t^2)$,
        hence $\Corr=\alpha_1^{\abs\tau}$.
\end{enumerate}
The cancellation step is the crux; verify it explicitly (see
Exercise~\ref{exo:p10-e3}).
\end{keytech}

\begin{keytech}{Writing down the conditional-Gaussian (G)ARCH likelihood}{p10-ktmle}
To set up MLE:
\begin{enumerate}
  \item Under Gaussian noise, $r_t\mid\mathcal F_{t-1}\sim\mathcal N(0,\sigma_t^2)$,
        with $\sigma_t^2$ given by the model's recursion.
  \item Condition on enough initial values to start the recursion (e.g.\ $r_1$
        for $\mathrm{ARCH}(1)$; for GARCH also seed $\sigma_t^2$, e.g.\ at the
        sample variance).
  \item Form $\log L_c=-\tfrac12\sum_{t}\big[\log(2\pi\sigma_t^2)+r_t^2/\sigma_t^2\big]$
        and maximise numerically subject to positivity and the stationarity
        constraint $\sum\alpha+\sum\beta<1$.
\end{enumerate}
\end{keytech}

\begin{pitfall}
\textbf{Do not} confuse the two thresholds. $\alpha_1<1$ is for a finite
\emph{variance}; the kurtosis / ACF-of-squares formulas require the
\emph{stronger} $\alpha_1^2<1/3$ (ARCH(1)) or
$3\alpha_1^2+2\alpha_1\beta_1+\beta_1^2<1$ (GARCH(1,1)). Also, ARCH/GARCH are
uncorrelated \textbf{but not independent}: $\Corr(r_{t+\tau},r_t)=0$ yet
$\Corr(r_{t+\tau}^2,r_t^2)\ne0$. ``White noise'' $\ne$ ``no structure''.
\end{pitfall}

% ----------------------------------------------------------------------
\subsection{Exercises}

\begin{exercise}{ARCH(2): mean, variance, stationarity \quad\textnormal{[Sheet 12]}}{p10-e1}
Let $r_t=\sigma_t\eps_t$ with $\sigma_t^2=\beta_0+\beta_1 r_{t-1}^2+\beta_2 r_{t-2}^2$,
the $\beta_j\ge0$, and $\eps_t$ a zero-mean unit-variance process independent of
$\mathcal F_{t-1}$.
(a) Compute $\E[r_t]$ via iterated expectation.
(b) Assuming stationarity, find the marginal variance $\Var(r_t)$.
(c) Give necessary conditions for stationarity in terms of the $\beta_j$.
\end{exercise}
\begin{solution}
\textbf{(a)} Since $\sigma_t$ is $\mathcal F_{t-1}$-measurable and $\eps_t$ is
independent of $\mathcal F_{t-1}$ with $\E[\eps_t]=0$,
\[
\E[r_t\mid\mathcal F_{t-1}]=\E[\sigma_t\eps_t\mid\mathcal F_{t-1}]
=\sigma_t\,\E[\eps_t]=0
\ \Longrightarrow\
\E[r_t]=\E\big[\E[r_t\mid\mathcal F_{t-1}]\big]=\E[0]=0 .
\]
\textbf{(b)} By the law of total variance, with $\Var(\eps_t)=1$,
\[
\Var(r_t\mid\mathcal F_{t-1})=\sigma_t^2\Var(\eps_t)=\beta_0+\beta_1 r_{t-1}^2+\beta_2 r_{t-2}^2 .
\]
Then, since $\E[r_t\mid\mathcal F_{t-1}]=0$ has zero variance,
\[
\Var(r_t)=\E[\Var(r_t\mid\mathcal F_{t-1})]+\Var(\E[r_t\mid\mathcal F_{t-1}])
=\E[\beta_0+\beta_1 r_{t-1}^2+\beta_2 r_{t-2}^2]+0 .
\]
Using $\E[r_{t-j}^2]=\Var(r_{t-j})$ (mean zero) and stationarity
$\Var(r_t)=\Var(r_{t-1})=\Var(r_{t-2})=:v$,
\[
v=\beta_0+\beta_1 v+\beta_2 v
\ \Longrightarrow\
(1-\beta_1-\beta_2)\,v=\beta_0
\ \Longrightarrow\
\Var(r_t)=\frac{\beta_0}{1-\beta_1-\beta_2}\quad(\beta_1+\beta_2<1).
\]
\textbf{(c)} For the variance to be positive and finite we need
$\beta_0>0$ and $\beta_1+\beta_2<1$ (with $\beta_1,\beta_2\ge0$). This is the
necessary stationarity condition (the $p=2,q=0$ instance of
$\sum\alpha+\sum\beta<1$).
\end{solution}

\begin{exercise}{GARCH is zero-mean and uncorrelated \quad\textnormal{[Sheet 12]}}{p10-e2}
For a $\mathrm{GARCH}(m,r)$ with Gaussian noise,
$\sigma_t^2=\alpha_0+\sum_{j=1}^{m}\alpha_j r_{t-j}^2+\sum_{j=1}^{r}\beta_j\sigma_{t-j}^2$,
$\alpha_j,\beta_j\ge0$, prove that $\E[r_t]=0$ and $\Corr(r_{t+\tau},r_t)=0$.
\end{exercise}
\begin{solution}
\textbf{Zero mean.} $\sigma_t$ is determined by the past, hence
$\mathcal F_{t-1}$-measurable, and $\eps_t$ is independent of $\mathcal F_{t-1}$
with $\E[\eps_t]=0$:
\[
\E[r_t]=\E\big[\E[r_t\mid\mathcal F_{t-1}]\big]
=\E\big[\E[\sigma_t\eps_t\mid\mathcal F_{t-1}]\big]
=\E\big[\sigma_t\,\E[\eps_t\mid\mathcal F_{t-1}]\big]
=\E[\sigma_t\cdot0]=0 .
\]
\textbf{Uncorrelated.} Fix $\tau>0$. As $\E[r_t]=0$,
$\Cov(r_{t+\tau},r_t)=\E[r_{t+\tau}r_t]$. Both $r_t$ and $\sigma_{t+\tau}$ are
$\mathcal F_{t+\tau-1}$-measurable, so by iterated expectation
\[
\E[r_{t+\tau}r_t]
=\E\big[\E[\sigma_{t+\tau}\eps_{t+\tau}r_t\mid\mathcal F_{t+\tau-1}]\big]
=\E\big[r_t\,\sigma_{t+\tau}\,\E[\eps_{t+\tau}\mid\mathcal F_{t+\tau-1}]\big]
=\E[r_t\sigma_{t+\tau}\cdot0]=0 .
\]
By symmetry the same holds for $\tau<0$, and at $\tau=0$ the covariance is the
variance. Therefore $\Cov(r_{t+\tau},r_t)=0$ and so $\Corr(r_{t+\tau},r_t)=0$ for
all $\tau\ne0$: GARCH is white noise.
\end{solution}

\begin{exercise}{ARCH(1): geometric ACF of the squares \quad\textnormal{[Sheet 12]}}{p10-e3}
For an $\mathrm{ARCH}(1)$ with $0<\alpha_0$ and $0<\alpha_1<1/\sqrt3$, show
$\Corr(r_{t+\tau}^2,r_t^2)=\alpha_1^{\abs\tau}$.
\end{exercise}
\begin{solution}
Take $\tau>0$ and start from the covariance:
\[
\Cov(r_{t+\tau}^2,r_t^2)=\E[r_{t+\tau}^2 r_t^2]-\E[r_{t+\tau}^2]\,\E[r_t^2]
=\E[r_{t+\tau}^2 r_t^2]-\big(\E[r_t^2]\big)^2,
\]
the last step by stationarity ($\E[r_{t+\tau}^2]=\E[r_t^2]$). Condition on
$\mathcal F_{t+\tau-1}$; both $r_t^2$ and $\sigma_{t+\tau}^2$ are measurable there
and $\E[\eps_{t+\tau}^2\mid\mathcal F_{t+\tau-1}]=1$:
\[
\E[r_{t+\tau}^2 r_t^2\mid\mathcal F_{t+\tau-1}]
=r_t^2\,\sigma_{t+\tau}^2\,\E[\eps_{t+\tau}^2\mid\mathcal F_{t+\tau-1}]
=r_t^2(\alpha_0+\alpha_1 r_{t+\tau-1}^2).
\]
Taking expectations,
\[
\E[r_{t+\tau}^2 r_t^2]
=\alpha_0\,\E[r_t^2]+\alpha_1\,\E[r_{t+\tau-1}^2 r_t^2].
\]
Hence, writing $C_\tau:=\Cov(r_{t+\tau}^2,r_t^2)$ and using
$\E[r_{t+\tau-1}^2 r_t^2]=C_{\tau-1}+\E[r_{t+\tau-1}^2]\E[r_t^2]
=C_{\tau-1}+(\E[r_t^2])^2$,
\[
C_\tau=\alpha_1 C_{\tau-1}
+\underbrace{\alpha_0\E[r_t^2]+\alpha_1(\E[r_t^2])^2-(\E[r_t^2])^2}_{=:R}.
\]
Now substitute $\E[r_t^2]=\Var(r_t)=\alpha_0/(1-\alpha_1)$:
\[
R=\alpha_0\frac{\alpha_0}{1-\alpha_1}-(1-\alpha_1)\Big(\frac{\alpha_0}{1-\alpha_1}\Big)^2
=\frac{\alpha_0^2}{1-\alpha_1}-\frac{\alpha_0^2}{1-\alpha_1}=0 .
\]
So $C_\tau=\alpha_1 C_{\tau-1}$, and iterating $\tau$ times from
$C_0=\Var(r_t^2)$ gives $C_\tau=\alpha_1^{\tau}\Var(r_t^2)$. Therefore
\[
\Corr(r_{t+\tau}^2,r_t^2)=\frac{C_\tau}{\Var(r_t^2)}=\alpha_1^{\tau},
\]
and by symmetry of the autocovariance $\Corr(r_{t+\tau}^2,r_t^2)=\alpha_1^{\abs\tau}$.
(The condition $\alpha_1^2<1/3$ guarantees $\Var(r_t^2)=\E[r_t^4]-(\E[r_t^2])^2<\infty$,
so the correlation is well-defined.)
\end{solution}

\begin{exercise}{GARCH(1,1): variance and fourth-moment check \quad\textnormal{[New]}}{p10-e4}
A $\mathrm{GARCH}(1,1)$ is fitted to log-returns with
$\alpha_0=3.686\times10^{-6}$, $\alpha_1=0.1713$, $\beta_1=0.7897$ (the lecture's
S\&P~500 fit). (a) Verify stationarity and compute the implied unconditional
variance and standard deviation (daily volatility). (b) Does a finite fourth
moment exist? (c) Comment on the persistence $\alpha_1+\beta_1$.
\end{exercise}
\begin{solution}
\textbf{(a)} $\alpha_1+\beta_1=0.1713+0.7897=0.9610<1$, so the process is weakly
stationary. The unconditional variance is
\[
\Var(r_t)=\frac{\alpha_0}{1-\alpha_1-\beta_1}
=\frac{3.686\times10^{-6}}{1-0.9610}
=\frac{3.686\times10^{-6}}{0.0390}\approx 9.45\times10^{-5}.
\]
The implied daily volatility is $\sqrt{9.45\times10^{-5}}\approx 9.72\times10^{-3}$,
i.e.\ about $0.97\%$ per day --- a sensible equity figure.\\
\textbf{(b)} Check $3\alpha_1^2+2\alpha_1\beta_1+\beta_1^2$:
\[
3(0.1713)^2+2(0.1713)(0.7897)+(0.7897)^2
=0.08803+0.27055+0.62363=0.9822<1 .
\]
Since $0.9822<1$, the fourth moment exists (only just): the fitted model has
finite kurtosis, but the closeness to $1$ signals very heavy tails.\\
\textbf{(c)} $\alpha_1+\beta_1=0.961$ is close to $1$, the hallmark of financial
GARCH fits: volatility shocks are highly persistent (slow geometric decay of the
ACF of squares), reproducing long quiet/turbulent regimes, while the process
remains stationary because the sum is still strictly below $1$.
\end{solution}

\begin{exercise}{ARCH(1) parameter from a given kurtosis \quad\textnormal{[New]}}{p10-e5}
An $\mathrm{ARCH}(1)$ with Gaussian innovations is observed to have kurtosis
$\kappa=9$. Find $\alpha_1$. What is the largest kurtosis an $\mathrm{ARCH}(1)$
can attain while keeping a finite fourth moment?
\end{exercise}
\begin{solution}
Use $\kappa=\dfrac{3(1-\alpha_1^2)}{1-3\alpha_1^2}$ and set it equal to $9$:
\[
9(1-3\alpha_1^2)=3(1-\alpha_1^2)
\;\Rightarrow\;9-27\alpha_1^2=3-3\alpha_1^2
\;\Rightarrow\;6=24\alpha_1^2
\;\Rightarrow\;\alpha_1^2=\tfrac14
\;\Rightarrow\;\alpha_1=\tfrac12 .
\]
(Take the positive root; check $\alpha_1^2=1/4<1/3$, so the fourth moment is
indeed finite.) As $\alpha_1^2\uparrow1/3$ (i.e.\ $\alpha_1\uparrow1/\sqrt3$) the
denominator $1-3\alpha_1^2\downarrow0^+$ while the numerator stays positive, so
$\kappa\to+\infty$: there is \textbf{no} finite upper bound on the attainable
kurtosis short of losing the fourth moment. Conversely $\kappa\to3$ as
$\alpha_1\to0$ (back to constant variance), confirming $\kappa>3$ for all
$\alpha_1>0$.
\end{solution}

\begin{exercise}{ARCH(1) one-step variance forecast \& predictive band \quad\textnormal{[New]}}{p10-e6}
For an $\mathrm{ARCH}(1)$ with $\alpha_0=0.0001$, $\alpha_1=0.4$, suppose the most
recent return is $r_t=0.03$. (a) Give the one-step-ahead conditional variance
$\sigma_{t+1}^2$ and a $95\%$ predictive interval for $r_{t+1}$ under Gaussian
innovations. (b) Compare with the unconditional ($\pm1.96\,\sigma$) band.
(c) Is this model fourth-moment stationary, and what is its kurtosis?
\end{exercise}
\begin{solution}
\textbf{(a)} $\sigma_{t+1}^2=\alpha_0+\alpha_1 r_t^2
=0.0001+0.4(0.03)^2=0.0001+0.4(0.0009)=0.0001+0.00036=0.00046$. Thus
$\sigma_{t+1}=\sqrt{0.00046}\approx0.02145$ and, since
$r_{t+1}\mid\mathcal F_t\sim\mathcal N(0,\sigma_{t+1}^2)$, the conditional $95\%$
interval is
\[
0\pm1.96\,\sigma_{t+1}=\pm1.96(0.02145)\approx\pm0.0421 .
\]
\textbf{(b)} The unconditional variance is
$\Var(r_t)=\alpha_0/(1-\alpha_1)=0.0001/0.6\approx1.667\times10^{-4}$, so
$\sigma\approx0.01291$ and the unconditional band is
$\pm1.96(0.01291)\approx\pm0.0253$. After the large move $r_t=0.03$ the
conditional band ($\pm0.042$) is markedly \emph{wider} than the unconditional one
($\pm0.025$): the model has raised its volatility forecast --- exactly the
clustering behaviour we want.\\
\textbf{(c)} Variance stationarity: $\alpha_1=0.4<1$ holds. Fourth moment:
$\alpha_1^2=0.16<1/3\approx0.333$ holds, so the fourth moment (hence kurtosis) is
finite, with
\[
\kappa=\frac{3(1-\alpha_1^2)}{1-3\alpha_1^2}=\frac{3(1-0.16)}{1-0.48}
=\frac{3(0.84)}{0.52}=\frac{2.52}{0.52}\approx4.85>3 .
\]
\end{solution}

\begin{exercise}{GARCH(1,1) as an ARCH($\infty$) \quad\textnormal{[New]}}{p10-e7}
Explain --- qualitatively, then with the recursion algebra --- why a
$\mathrm{GARCH}(1,1)$ can mimic an $\mathrm{ARCH}(\infty)$ with geometrically
decaying coefficients, and why this is more parsimonious than $\mathrm{ARCH}(p)$.
\end{exercise}
\begin{solution}
\textbf{Algebra.} Iterate the GARCH(1,1) variance recursion
$\sigma_t^2=\alpha_0+\alpha_1 r_{t-1}^2+\beta_1\sigma_{t-1}^2$ by repeatedly
substituting $\sigma_{t-1}^2$:
\[
\sigma_t^2=\alpha_0\sum_{k=0}^{\infty}\beta_1^{k}
+\alpha_1\sum_{k=0}^{\infty}\beta_1^{k}\,r_{t-1-k}^2
=\frac{\alpha_0}{1-\beta_1}+\alpha_1\sum_{k=0}^{\infty}\beta_1^{k}\,r_{t-1-k}^2 ,
\]
valid when $0\le\beta_1<1$ (so the geometric series converges). This is an
$\mathrm{ARCH}(\infty)$ with intercept $\alpha_0/(1-\beta_1)$ and coefficients
$\alpha_1\beta_1^{k}$ on $r_{t-1-k}^2$ that decay \emph{geometrically} in the lag.\\
\textbf{Parsimony.} Reproducing the same slowly decaying dependence with a finite
$\mathrm{ARCH}(p)$ would need many free $\alpha_j$'s (large $p$), whereas
GARCH(1,1) encodes the entire geometric tail with just \emph{three} parameters
$(\alpha_0,\alpha_1,\beta_1)$. This is exactly the ``feed past variance back in''
idea: $\beta_1$ sets the persistence/decay rate, $\alpha_1$ the immediate shock
impact. It also explains the geometric ACF-of-squares decay carried over from the
ARCH(1) case.
\end{solution}

% ======================================================================
%  PART 11 -- Linear Time-Invariant (LTI) Filters
% ======================================================================
\section{Linear Time-Invariant (LTI) Filters}
\label{sec:lti}

A \emph{(digital) filter} is a map that turns one sequence into another,
$\{Y_t\}=\filt[\{X_t\}]$, with index set $\Ints$. The filters that matter for
time series are the \textbf{linear time-invariant} (LTI) ones: they treat all
time points alike and act linearly. Two ``test inputs'' completely characterise
such a filter --- a single impulse (giving the \emph{impulse response} $h$) and
a pure complex wave (giving the \emph{transfer function} $H$). Their power is
twofold: every LTI filter is a \emph{convolution} with $h$, and in the frequency
domain it acts \emph{diagonally}, multiplying the spectral density by
$\abs{H(f)}^2$. This single fact, $\sdf_Y(f)=\abs{H(f)}^2\sdf_X(f)$, is the
engine behind the closed-form spectral density of any ARMA process. Throughout
we allow time series to be \emph{complex-valued} and assume the (possibly
infinite) sums involved converge, the standard sufficient condition being
$h\in\ell^1$.

% ----------------------------------------------------------------------
\subsection{Definitions}

\begin{definition}{Complex-valued time series \& its autocovariance}{p11-complex}
A \textbf{complex-valued} time series can be written $X_t=U_t+iV_t$ with
$\{U_t\},\{V_t\}$ real-valued, and has mean
$\E[X_t]=\E[U_t]+i\,\E[V_t]$. Its \textbf{autocovariance} uses a
\emph{conjugate} on the second factor:
\[
\acvs^{(X)}_\tau=\Cov(X_{t+\tau},X_t)=\E\!\big[(X_{t+\tau}-\mu)\conj{(X_t-\mu)}\big].
\]
This reduces to the usual real ACVS when $X_t\in\Reals$. (Stationarity is
required jointly of the real and imaginary parts.) Throughout we drop the index
set: $\{X_t\}=\{X_t\}_{t\in\Ints}$.
\end{definition}

\begin{definition}{Digital filter \& output notation}{p11-filter}
A \textbf{digital filter} $\filt$ maps an input sequence to an output sequence,
$\{Y_t\}=\filt[\{X_t\}]$. To name a single output entry we put the time index
\emph{after} the bracket,
\[
Y_u=\filt[\{X_t\}]_u,\qquad u\in\Ints .
\]
The $t$ in $\{X_t\}$ is a \emph{dummy} variable; the $u\in\Ints$ after the
bracket has a definite meaning (the output time).
\end{definition}

\begin{recallbox}[Recall: the backshift operator]
The \textbf{backshift} $\Bsh$ satisfies $\Bsh[\{X_t\}]_u=X_{u-1}$ for every
$u\in\Ints$. Iterating, $\Bsh^u[\{X_t\}]_s=X_{s-u}$ for all $s\in\Ints$ (for
$u<0$ this applies $\Bsh^{-1}$ a total of $-u$ times). $\Bsh$ is itself the
prototypical LTI filter.
\end{recallbox}

\begin{definition}{Linear time-invariant (LTI) filter}{p11-lti}
A digital filter $\filt$ is \textbf{linear time-invariant} if for all sequences
$\{X_t\},\{Y_t\}$ and all $\alpha\in\Comp$:
\begin{enumerate}
  \item \textbf{Linearity:}
        $\displaystyle \filt[\{\alpha X_t+Y_t\}]=\alpha\,\filt[\{X_t\}]+\filt[\{Y_t\}]$;
  \item \textbf{Time invariance:}
        $\displaystyle \filt\big[\Bsh[\{X_t\}]\big]=\Bsh\big[\filt[\{X_t\}]\big]$,
        i.e.\ $\filt$ \emph{commutes with the backshift}, $\filt\Bsh=\Bsh\filt$.
\end{enumerate}
Time invariance for the single shift $\Bsh$ already implies it for every shift:
$\filt[\Bsh^u\{X_t\}]=\Bsh^u\filt[\{X_t\}]$ for all $u\in\Ints$. (The
continuous-time analogue is usually stated with $\Bsh^u$, but the weaker
one-step condition suffices.)
\end{definition}

\begin{example}{First LTI filters}{p11-examples}
\textbf{(i)} The backshift $\Bsh$ itself is LTI. \\
\textbf{(ii)} For fixed $\phi\in\Reals$, the one-tap filter
$\filt[\{X_t\}]_u=X_u+\phi X_{u-1}$ is LTI: linearity is immediate from
\[
\filt[\{\alpha X_t+Y_t\}]_u=\alpha(X_u+\phi X_{u-1})+(Y_u+\phi Y_{u-1})
=\alpha\,\filt[\{X_t\}]_u+\filt[\{Y_t\}]_u,
\]
and time invariance is checked the same way. This is exactly $\filt=I+\phi\Bsh$.
\end{example}

\begin{definition}{Impulse sequence \& impulse response}{p11-impulse}
For $m\in\Ints$ the \textbf{impulse sequence} $\{\delta_{t,m}\}$ is
$\delta_{t,m}=1$ if $t=m$ and $0$ otherwise. The \textbf{impulse response} of an
LTI filter $\filt$ is the sequence $\{h_m\}$ defined by
\[
h_m=\filt[\{\delta_{t,-m}\}]_0,\qquad m\in\Ints,
\]
where the subscript $0$ is the \emph{output time index}. Informally $h_m$ is the
response, read off at time $0$, to a unit shock placed at time $-m$. A common
sufficient regularity condition is $h\in\ell^1$
(then $\norm{h}_\infty\le\norm{h}_1<\infty$).
\end{definition}

\begin{definition}{Complex wave \& transfer function}{p11-transfer}
For $f\in\Reals$ the \textbf{complex wave} is the sequence
$\{\xi_{t,f}\}$ with $\xi_{t,f}=e^{2\pi i f t}$. The \textbf{transfer function}
of an LTI filter $\filt$ that can be applied to these waves is
\[
H(f)=\filt[\{\xi_{t,f}\}]_0,\qquad f\in\Reals .
\]
It encodes the frequency-domain behaviour of $\filt$: how each pure frequency is
re-weighted in amplitude and shifted in phase.
\end{definition}

\begin{recallbox}[Recall: mean-square equality]
For random variables $U,V\in L^2$ we write $U\ms V$ to mean
$\E[\,\abs{U-V}^2\,]=0$ (equality in mean square). This is the natural notion of
equality for the stochastic infinite sums (convolutions, spectral integrals)
appearing below.
\end{recallbox}

% ----------------------------------------------------------------------
\subsection{Theorems \& Propositions}

\begin{proposition}{Linear combination of LTI filters is LTI}{p11-lincomb}
If $\filt_1,\filt_2$ are LTI and $\alpha\in\Comp$, then
$\filt=\alpha\filt_1+\filt_2$ is LTI.\\[2pt]
\textit{Proof idea:} apply linearity of $\filt_1,\filt_2$ to
$\filt[\beta\{X_t\}+\{Y_t\}]$ and regroup; for time invariance use that each
$\filt_j$ commutes with $\Bsh$ and that $\Bsh$ is linear, so
$\filt[\Bsh\{X_t\}]=\Bsh[\filt\{X_t\}]$. (Full computation in
Exercise~\ref{exo:p11-e1}.)\\
\textit{Why:} it makes the LTI filters a \emph{vector space} (over $\Comp$); in
particular any polynomial in $\Bsh$, e.g.\ $\Phi(\Bsh)=\sum_j\phi_j\Bsh^j$, is
LTI.
\end{proposition}

\begin{proposition}{Cascade of LTI filters is LTI}{p11-cascade}
If $\filt_1,\filt_2$ are LTI, then the \textbf{cascade}
$\filt=\filt_1\filt_2$, defined by
$\filt[\{X_t\}]=\filt_1\big[\filt_2[\{X_t\}]\big]$, is LTI.\\[2pt]
\textit{Proof idea:} linearity is the composition of two linear maps; time
invariance follows by pushing $\Bsh$ through both filters,
$\filt_1[\filt_2[\Bsh\{X_t\}]]=\filt_1[\Bsh[\filt_2\{X_t\}]]
=\Bsh[\filt_1[\filt_2\{X_t\}]]$. (Exercise~\ref{exo:p11-e2}.)\\
\textit{Why:} filters can be \emph{chained} and the family is closed under
composition; combined with the previous proposition, all of $\{$polynomials and
power series in $\Bsh\}$ are LTI.
\end{proposition}

\begin{intuition}
LTI filters are exactly the maps that ``do the same thing at every time and
respect superposition''. Commuting with $\Bsh$ \emph{is} time invariance:
shifting the input then filtering gives the same as filtering then shifting. Both
closure properties (linear combinations, cascades) are inherited directly from
those two structural rules.
\end{intuition}

\begin{theorem}{LTI $\Leftrightarrow$ convolution with the impulse response}{p11-conv}
Under suitable convergence assumptions, an LTI filter with impulse response $h$
acts as a \textbf{convolution}:
\[
\filt[\{X_t\}]_u\ms\sum_{m\in\Ints}h_{u-m}X_m,\qquad u\in\Ints .
\]
Conversely, any filter defined by such a convolution (with, say, $h\in\ell^1$)
is LTI. The impulse response gives the convolution weights.\\[2pt]
\textit{Proof idea (deterministic $\ell^1$ version, Exercise~\ref{exo:p11-e3}):}
write $X^{(N)}=\sum_{\abs m\le N}X_m\delta_{\cdot,m}$; by continuity in
$\ell^1$, $\filt[\{X_t\}]_u=\lim_N\filt[X^{(N)}]_u$. Linearity turns each finite
truncation into a finite sum, and time invariance gives
$\filt[\{\delta_{t,m}\}]_u=\filt[\{\delta_{t,m-u}\}]_0=h_{u-m}$; absolute
convergence ($h\in\ell^1$, $X\in\ell^1$) lets the partial sums pass to the full
series. The converse checks linearity and time invariance directly on the
convolution (re-indexing $m'=m-1$ for the shift).\\
\textit{Why:} this is the master representation --- it reduces \emph{any} LTI
filter to a single object, the weights $h$, and underlies every spectral
computation below.
\end{theorem}

\begin{intuition}
Two complementary pictures of a time series, and an LTI filter acts simply on
both ``building blocks'':
\[
X_t=\sum_{m\in\Ints}X_m\,\delta_{m,t}\quad(\text{sum of impulses}),
\qquad
X_t\ms\mu+\int_{-1/2}^{1/2}e^{2\pi i f t}\,dZ(f)\quad(\text{sum of waves}).
\]
Applying $\filt$ gives, respectively,
$\filt[\{X_t\}]_u\ms\sum_m X_m\,\filt[\{\delta_{m,t}\}]_u$ (the impulse-response
view) and
$\filt[\{X_t\}]_u\ms\mu H(0)+\int_{-1/2}^{1/2}H(f)e^{2\pi i f u}\,dZ(f)$ (the
transfer-function view).
\end{intuition}

\begin{theorem}{Complex waves are eigensequences; $H(f)$ is the eigenvalue}{p11-eigen}
Let $\filt$ be an LTI filter applicable to the complex waves $\{\xi_{t,f}\}$.
Then each wave is an \textbf{eigensequence}: for every $f\in\Reals$,
\[
\filt[\{\xi_{t,f}\}]=H(f)\,\{\xi_{t,f}\},
\]
i.e.\ $\filt$ only rescales amplitude and shifts phase while preserving the
frequency.\\[2pt]
\textit{Proof idea:} since
$\Bsh^{-u}[\{\xi_{t,f}\}]_\tau=\xi_{\tau+u,f}=e^{2\pi i u f}\xi_{\tau,f}$, we have
$\Bsh^{-u}[\{\xi_{t,f}\}]=e^{2\pi i u f}\{\xi_{t,f}\}$. Then for any $u$,
\[
\filt[\{\xi_{t,f}\}]_u
=\Bsh^{u}\!\big[\filt[\Bsh^{-u}\{\xi_{t,f}\}]\big]_u
=\filt\big[e^{2\pi i f u}\{\xi_{t,f}\}\big]_0
=e^{2\pi i f u}\filt[\{\xi_{t,f}\}]_0=\xi_{u,f}\,H(f),
\]
using time invariance, the shift identity, linearity and the definition of $H$.\\
\textit{Why:} it diagonalises every LTI filter in the frequency basis ---
filtering becomes \emph{pointwise multiplication} by $H(f)$, the source of all
clean frequency-domain formulas.
\end{theorem}

\begin{theorem}{Transfer function $=$ Fourier transform of $h$}{p11-Hfourier}
If the impulse response satisfies $h\in\ell^1$, the transfer function is the
(discrete) Fourier transform of $h$:
\[
H(f)=\sum_{m\in\Ints}h_m\,e^{-2\pi i f m},\qquad f\in\Reals .
\]
\textit{Proof idea:} feed the wave $\xi_{t,f}$ into the convolution
representation: $H(f)=\filt[\{\xi_{t,f}\}]_0=\sum_m h_{-m}e^{2\pi i f m}
=\sum_m h_m e^{-2\pi i f m}$ (re-index $m\mapsto -m$); $h\in\ell^1$ makes the sum
absolutely convergent. (Exercise.)\\
\textit{Why:} the practical recipe --- read off $h$ from the filter, then sum the
series to get $H$ (and conversely, recover $h$ as Fourier coefficients of $H$).
\end{theorem}

\begin{theorem}{Filtering preserves stationarity}{p11-statpres}
Let $\{X_t\}$ be stationary and $\filt$ an LTI filter with $h\in\ell^1$. Then
$\{Y_t\}=\filt[\{X_t\}]$ is stationary.\\[2pt]
\textit{Proof idea:} by the convolution theorem $Y_t\ms\sum_m h_m X_{t-m}$.
By Fubini, $\E[Y_t]=\sum_m h_m\E[X_{t-m}]=\mu_X\sum_m h_m$ is $t$-free; the
covariance
$\acvs^{(Y)}_\tau=\sum_{m,s}h_m\conj{h_s}\,\acvs^{(X)}_{\tau-m+s}$
depends only on $\tau$; and $\acvs^{(Y)}_0\le\acvs^{(X)}_0\norm{h}_1^2<\infty$,
so all is finite.\\
\textit{Why:} it licenses speaking of $\sdf_Y$ at all --- only a stationary
output has a spectral density.
\end{theorem}

\begin{theorem}{Output spectrum: $\sdf_Y(f)=\abs{H(f)}^2\sdf_X(f)$}{p11-outspec}
Let $\{X_t\}$ be stationary with $\acvs^{(X)}\in\ell^1$, and $\filt$ an LTI
filter with $h\in\ell^1$ and transfer function $H$. If
$\{Y_t\}=\filt[\{X_t\}]$ then
\[
\sdf_Y(f)=\abs{H(f)}^2\,\sdf_X(f),\qquad f\in\Reals .
\]
\textit{Proof idea:} from the covariance formula above,
$\acvs^{(Y)}_\tau=\sum_u w_u\,\acvs^{(X)}_{\tau-u}$ where
$w_u=\sum_s h_{u+s}\conj{h_s}$ (set $u=m-s$). With $\tilde h_t=\conj{h_{-t}}$ one
has $w=h*\tilde h\in\ell^1$, so the convolution theorem gives
$\sdf_Y=W\,\sdf_X$ with
$W(f)=H(f)\,\widetilde H(f)=H(f)\conj{H(f)}=\abs{H(f)}^2$, since
$\widetilde H(f)=\conj{H(f)}$.\\
\textit{Why:} the workhorse identity. It converts \emph{any} filtering question
into multiplying the input spectrum by a gain $\abs{H(f)}^2$ --- and, applied to
$\Phi(\Bsh)$ and $\Theta(\Bsh)$, it yields the ARMA spectral density below.
\end{theorem}

\begin{lemma}{Transfer function of a polynomial filter $\Phi(\Bsh)$}{p11-polyfilter}
For a degree-$p$ polynomial $\Phi(z)=\sum_{j=0}^p\phi_j z^j$, the filter
$\Phi(\Bsh)$ is LTI with transfer function
\[
H(f)=\Phi\!\big(e^{-2\pi i f}\big),\qquad f\in\Reals .
\]
\textit{Proof idea:} $\Phi(\Bsh)$ is a linear combination of the LTI filters
$\Bsh^j$, hence LTI, with impulse response $h_m=\phi_m$ for $0\le m\le p$ and $0$
otherwise. Then $H(f)=\sum_m h_m e^{-2\pi i f m}
=\sum_{j=0}^p\phi_j e^{-2\pi i f j}=\Phi(e^{-2\pi i f})$.\\
\textit{Why:} it turns the operator $\Phi(\Bsh)$ into a concrete function of
$f$; the bridge from time-domain difference equations to the spectrum.
\end{lemma}

\begin{theorem}{ARMA spectral density}{p11-arma}
For a stationary $\mathrm{ARMA}(p,q)$ process $\Phi(\Bsh)X_t=\Theta(\Bsh)\eps_t$
with $\eps_t$ white noise of variance $\sigma^2$,
\[
\sdf_X(f)=\sigma^2\,\frac{\abs{\Theta(e^{-2\pi i f})}^2}{\abs{\Phi(e^{-2\pi i f})}^2}
=\sigma^2\,\left|\frac{\sum_{j=0}^q\theta_j e^{-2\pi i j f}}
{\sum_{j=0}^p\phi_j e^{-2\pi i j f}}\right|^2 .
\]
\textit{Proof idea:} write $\Phi(\Bsh)[\{X_t\}]=\Theta(\Bsh)[\{\eps_t\}]$.
Applying the output-spectrum theorem to each side and the polynomial-filter
lemma,
$\abs{\Phi(e^{-2\pi i f})}^2\sdf_X(f)=\abs{\Theta(e^{-2\pi i f})}^2\sdf_\eps(f)$
with $\sdf_\eps(f)=\sigma^2$; divide, legal because stationarity forbids zeros
of $\Phi$ on the unit circle, so $\Phi(e^{-2\pi i f})\ne0$.\\
\textit{Why:} ARMA processes have \emph{no} closed-form ACVS but a clean
closed-form spectral density --- the practical way to read off their
frequency content.
\end{theorem}

\begin{pitfall}
The conjugate sits on the \emph{second} factor in the complex ACVS,
$\Cov(X_{t+\tau},X_t)=\E[(X_{t+\tau}-\mu)\conj{(X_t-\mu)}]$. Forgetting it
breaks Hermitian symmetry. Also: the gain is $\abs{H(f)}^2$, \emph{not} $H(f)$
or $H(f)^2$ --- the $\sdf$ multiplier is always the squared modulus, which is
real and nonnegative as a spectral density must be.
\end{pitfall}

% ----------------------------------------------------------------------
\subsection{Key Techniques}

\begin{keytech}{Spectral density of a filtered process}{p11-ktspec}
To get $\sdf_Y$ when $\{Y_t\}=\filt[\{X_t\}]$ for an LTI filter:
\begin{enumerate}
  \item \textbf{Identify the impulse response} $h$ (the convolution weights), or
        write $\filt$ as a polynomial $\Phi(\Bsh)$.
  \item \textbf{Form the transfer function}
        $H(f)=\sum_m h_m e^{-2\pi i f m}$ (for $\Phi(\Bsh)$ this is
        $\Phi(e^{-2\pi i f})$).
  \item \textbf{Multiply by the squared gain:}
        $\sdf_Y(f)=\abs{H(f)}^2\sdf_X(f)$.
\end{enumerate}
Handy simplification: $\abs{1-e^{-2\pi i f}}^2=2-2\cos(2\pi f)=4\sin^2(\pi f)$.
\end{keytech}

\begin{keytech}{Transfer function from an impulse response (and back)}{p11-kth}
\textbf{Forward} ($h\to H$): sum the Fourier series
$H(f)=\sum_m h_m e^{-2\pi i f m}$; for a \emph{finite} filter this is a finite
trigonometric polynomial. \textbf{Backward} ($H\to h$): the weights are the
Fourier coefficients, $h_m=\int_{-1/2}^{1/2}H(f)e^{2\pi i f m}\,df$. To read $h$
straight off a convolution filter, recall $\filt[\{X_t\}]_u=\sum_m h_{u-m}X_m
=\sum_k h_k X_{u-k}$, so the coefficient of $X_{u-k}$ is $h_k$.
\end{keytech}

\begin{keytech}{Composing filters (cascades) in the frequency domain}{p11-ktcascade}
Filtering is multiplication by $H$ on each frequency, so a cascade
$\filt=\filt_1\filt_2$ has transfer function $H(f)=H_1(f)H_2(f)$ and squared
gain $\abs{H_1(f)}^2\abs{H_2(f)}^2$. Likewise convolution in time corresponds to
\emph{product} in frequency: $h^{(1)}*h^{(2)}\leftrightarrow H_1 H_2$. This is
why ARMA decomposes as ``MA part on top, AR part on the bottom''.
\end{keytech}

\begin{keytech}{Proving a filter is LTI}{p11-ktlti}
Check the two axioms on a generic input:
\begin{enumerate}
  \item \textbf{Linearity:} expand $\filt[\{\alpha X_t+Y_t\}]_u$ and match
        $\alpha\filt[\{X_t\}]_u+\filt[\{Y_t\}]_u$.
  \item \textbf{Time invariance:} show $\filt[\Bsh\{X_t\}]_u=\Bsh[\filt\{X_t\}]_u
        =\filt[\{X_t\}]_{u-1}$, typically by re-indexing the sum
        ($m'=m-1$); absolute convergence ($h\in\ell^1$) justifies the
        re-indexing.
\end{enumerate}
Shortcut: any linear combination or cascade of known LTI filters is automatically
LTI --- and any filter expressible as a convolution with $h\in\ell^1$ is LTI.
\end{keytech}

% ----------------------------------------------------------------------
\subsection{Exercises}

\begin{exercise}{Linear combination of LTI filters \quad\textnormal{[Sheet 13]}}{p11-e1}
Let $\filt_1,\filt_2$ be LTI and $\alpha\in\Comp$. Show
$\filt=\alpha\filt_1+\filt_2$ is LTI.
\end{exercise}
\begin{solution}
\textbf{Linearity.} For sequences $\{X_t\},\{Y_t\}$ and $\beta\in\Comp$, using
linearity of $\filt_1,\filt_2$,
\[
\begin{aligned}
\filt[\beta\{X_t\}+\{Y_t\}]
&=\alpha\filt_1[\beta\{X_t\}+\{Y_t\}]+\filt_2[\beta\{X_t\}+\{Y_t\}]\\
&=\alpha\big(\beta\filt_1[\{X_t\}]+\filt_1[\{Y_t\}]\big)
 +\big(\beta\filt_2[\{X_t\}]+\filt_2[\{Y_t\}]\big)\\
&=\beta\big(\alpha\filt_1[\{X_t\}]+\filt_2[\{X_t\}]\big)
 +\big(\alpha\filt_1[\{Y_t\}]+\filt_2[\{Y_t\}]\big)\\
&=\beta\,\filt[\{X_t\}]+\filt[\{Y_t\}].
\end{aligned}
\]
\textbf{Time invariance.} Using that each $\filt_j$ commutes with $\Bsh$ and that
$\Bsh$ is linear,
\[
\begin{aligned}
\filt\big[\Bsh[\{X_t\}]\big]
&=\alpha\filt_1\big[\Bsh[\{X_t\}]\big]+\filt_2\big[\Bsh[\{X_t\}]\big]
 =\alpha\,\Bsh\big[\filt_1[\{X_t\}]\big]+\Bsh\big[\filt_2[\{X_t\}]\big]\\
&=\Bsh\big[\alpha\filt_1[\{X_t\}]+\filt_2[\{X_t\}]\big]
 =\Bsh\big[\filt[\{X_t\}]\big].
\end{aligned}
\]
Both axioms hold, so $\filt$ is LTI. \hfill$\square$
\end{solution}

\begin{exercise}{Cascade of LTI filters \quad\textnormal{[Sheet 13]}}{p11-e2}
Let $\filt_1,\filt_2$ be LTI. Show the cascade
$\filt=\filt_1\filt_2$, $\filt[\{X_t\}]=\filt_1[\filt_2[\{X_t\}]]$, is LTI.
\end{exercise}
\begin{solution}
\textbf{Linearity.} For $\alpha\in\Comp$,
\[
\begin{aligned}
\filt[\alpha\{X_t\}+\{Y_t\}]
&=\filt_1\big[\filt_2[\alpha\{X_t\}+\{Y_t\}]\big]
 =\filt_1\big[\alpha\filt_2[\{X_t\}]+\filt_2[\{Y_t\}]\big]\\
&=\alpha\filt_1\big[\filt_2[\{X_t\}]\big]+\filt_1\big[\filt_2[\{Y_t\}]\big]
 =\alpha\,\filt[\{X_t\}]+\filt[\{Y_t\}].
\end{aligned}
\]
\textbf{Time invariance.} Push $\Bsh$ through each filter in turn:
\[
\filt\big[\Bsh[\{X_t\}]\big]
=\filt_1\big[\filt_2[\Bsh[\{X_t\}]]\big]
=\filt_1\big[\Bsh[\filt_2[\{X_t\}]]\big]
=\Bsh\big[\filt_1[\filt_2[\{X_t\}]]\big]
=\Bsh\big[\filt[\{X_t\}]\big].
\]
Hence $\filt$ is LTI. (Note the order is preserved; for general operators
$\filt_1\filt_2\ne\filt_2\filt_1$, though for polynomials in $\Bsh$ they do
commute.) \hfill$\square$
\end{solution}

\begin{exercise}{LTI on $\ell^1$ $\Leftrightarrow$ convolution \quad\textnormal{[Sheet 13]}}{p11-e3}
Let $\filt:\ell^1\to\Comp^{\Ints}$ be a digital filter such that for each fixed
$u$ the map $\{X_t\}\mapsto\filt[\{X_t\}]_u$ is continuous in the $\ell^1$ norm.
Define $h_m=\filt[\{\delta_{t,-m}\}]_0$ and assume $h\in\ell^1$. Prove $\filt$ is
LTI on $\ell^1$ iff $\filt[\{X_t\}]_u=\sum_{m\in\Ints}h_{u-m}X_m$ for all
$\{X_t\}\in\ell^1$ and $u\in\Ints$. Justify all limit/sum interchanges.
\end{exercise}
\begin{solution}
\textbf{($\Rightarrow$)} Assume $\filt$ is LTI on $\ell^1$. Truncate:
$X^{(N)}_t=\sum_{\abs m\le N}X_m\delta_{t,m}$. Since $\{X_t\}\in\ell^1$,
$\norm{X^{(N)}-X}_1\to0$, so by continuity of
$\{X_t\}\mapsto\filt[\{X_t\}]_u$,
\[
\filt[\{X_t\}]_u=\lim_{N\to\infty}\filt\big[\{X^{(N)}_t\}\big]_u .
\]
For each $N$, \emph{linearity} (finite sum) gives
$\filt[\{X^{(N)}_t\}]_u=\sum_{\abs m\le N}X_m\,\filt[\{\delta_{t,m}\}]_u$, and
\emph{time invariance} gives
$\filt[\{\delta_{t,m}\}]_u=\filt[\{\delta_{t,m-u}\}]_0=h_{u-m}$, so
\[
\filt\big[\{X^{(N)}_t\}\big]_u=\sum_{\abs m\le N}h_{u-m}X_m .
\]
Since $h\in\ell^1$ we have $\norm{h}_\infty\le\norm{h}_1<\infty$, hence
$\sum_{m\in\Ints}\abs{h_{u-m}X_m}\le\norm{h}_\infty\sum_m\abs{X_m}<\infty$:
the series converges absolutely, the partial sums converge to it, and
$\filt[\{X_t\}]_u=\sum_{m\in\Ints}h_{u-m}X_m$.

\textbf{($\Leftarrow$)} Suppose $\filt[\{X_t\}]_u=\sum_m h_{u-m}X_m$ with
$h\in\ell^1$. Absolute convergence holds because
$\sum_m\abs{h_{u-m}X_m}\le\norm{h}_\infty\norm{X}_1\le\norm{h}_1\norm{X}_1$.
\emph{Linearity:} for $\alpha\in\Comp$,
\[
\begin{aligned}
\filt[\alpha\{X_t\}+\{Y_t\}]_u
&=\sum_m h_{u-m}(\alpha X_m+Y_m)\\
&=\alpha\sum_m h_{u-m}X_m+\sum_m h_{u-m}Y_m
 =\alpha\filt[\{X_t\}]_u+\filt[\{Y_t\}]_u .
\end{aligned}
\]
\emph{Time invariance:} absolute convergence allows reindexing $m'=m-1$,
\[
\filt[\Bsh[\{X_t\}]]_u=\sum_m h_{u-m}X_{m-1}=\sum_{m'}h_{u-1-m'}X_{m'}
=\filt[\{X_t\}]_{u-1}=\Bsh[\filt[\{X_t\}]]_u .
\]
Hence $\filt$ is LTI on $\ell^1$. \hfill$\square$
\end{solution}

\begin{exercise}{Tapered DFT as convolution with $dZ$ \quad\textnormal{[Sheet 13]}}{p11-e4}
Let $\{X_t\}$ be stationary, mean zero, with spectral representation
$X_t=\int_{-1/2}^{1/2}e^{2\pi i f t}\,dZ(f)$. Observe $X_1,\dots,X_n$ and define
the \emph{tapered DFT}
$J_h(f)=\sum_{t=1}^n h_t X_t e^{-2\pi i f t}$, with $\norm{h_t}_2^2=1$. Show
$J_h(f)=\int_{-1/2}^{1/2}H(f-f')\,dZ(f')$, where
$H(f)=\sum_{t=1}^n h_t e^{-2\pi i f t}$.
\end{exercise}
\begin{solution}
Substitute the spectral representation and interchange the finite sum with the
integral:
\[
\begin{aligned}
J_h(f)&=\sum_{t=1}^n h_t X_t e^{-2\pi i f t}
       =\sum_{t=1}^n h_t\!\left(\int_{-1/2}^{1/2}e^{2\pi i f' t}\,dZ(f')\right)\!e^{-2\pi i f t}\\
      &=\int_{-1/2}^{1/2}\Big(\sum_{t=1}^n h_t\,e^{2\pi i (f'-f) t}\Big)dZ(f')
       =\int_{-1/2}^{1/2}H(f-f')\,dZ(f'),
\end{aligned}
\]
since $\sum_{t=1}^n h_t e^{2\pi i(f'-f)t}=\sum_{t=1}^n h_t e^{-2\pi i(f-f')t}
=H(f-f')$. Thus the tapered DFT is the spectral increment $dZ$ smeared by the
taper's transfer function $H$ centred at $f$. \hfill$\square$
\end{solution}

\begin{exercise}{Spectral density of the differenced process \quad\textnormal{[Sheet 13]}}{p11-e5}
Let $\{X_t\}$ be stationary and $\{Y_t\}=(I-\Bsh)[\{X_t\}]$, i.e.\
$Y_t=X_t-X_{t-1}$. Is $\{Y_t\}$ stationary? If so, express $\sdf_Y$ in terms of
$\sdf_X$.
\end{exercise}
\begin{solution}
$I-\Bsh$ is LTI with impulse response $h_0=1,\ h_1=-1$ (others $0$), so
$h\in\ell^1$ and $\{Y_t\}$ is stationary (filtering preserves stationarity). Its
transfer function is
\[
H(f)=h_0+h_1 e^{-2\pi i f}=1-e^{-2\pi i f}.
\]
Therefore
\[
\sdf_Y(f)=\abs{H(f)}^2\sdf_X(f)=\abs{1-e^{-2\pi i f}}^2\sdf_X(f).
\]
Simplify the gain: $\abs{1-e^{-2\pi i f}}^2=(1-\cos 2\pi f)^2+\sin^2 2\pi f
=2-2\cos(2\pi f)=4\sin^2(\pi f)$. Hence
\[
\boxed{\ \sdf_Y(f)=4\sin^2(\pi f)\,\sdf_X(f).\ }
\]
The differencing filter \emph{kills} $f=0$ (since $H(0)=0$) and amplifies high
frequencies, as expected of a high-pass operation. \hfill$\square$
\end{solution}

\begin{exercise}{Time reversibility of AR($p$) \quad\textnormal{[Sheet 13]}}{p11-e6}
In an AR($p$) process $X_t=\sum_{j=1}^p\phi_j X_{t-j}+\eps_t$ (Gaussian noise,
variance $\sigma_\eps^2$, mean zero), set $Y_t=X_{-t}$. Show $\{Y_t\}$ is an
AR($p$) with the \emph{same} parameters, driven by white noise with the same
distribution as $\eps$.
\end{exercise}
\begin{solution}
Put $\phi_0=-1$ and define $\tilde\nu_t=\sum_{j=0}^p\phi_j Y_{t-j}
=-Y_t+\sum_{j=1}^p\phi_j Y_{t-j}$, i.e.\ $Y_t=\sum_{j=1}^p\phi_j Y_{t-j}
-\tilde\nu_t$; we must show $\{\tilde\nu_t\}$ is white noise $\sim$ that of
$\eps$ (the sign is immaterial for noise). It is Gaussian (a linear
combination of Gaussians) with $\E[\tilde\nu_t]=0$ since $\E[X_t]=0$.

\emph{Spectral matching.} For every $\tau$,
\[
\acvs^{(Y)}_\tau=\E[Y_t Y_{t+\tau}]=\E[X_{-t}X_{-t-\tau}]
=\acvs^{(X)}_{-\tau}=\acvs^{(X)}_\tau,
\]
so $\sdf_Y(f)=\sdf_X(f)$ for all $f$. Now $\{\tilde\nu_t\}=\Phi(\Bsh)[\{Y_t\}]$
is an LTI filtering of $\{Y_t\}$, so by the polynomial-filter lemma and the
output-spectrum theorem,
\[
\sdf_{\tilde\nu}(f)=\abs{\Phi(e^{-2\pi i f})}^2\sdf_Y(f)
=\abs{\Phi(e^{-2\pi i f})}^2\sdf_X(f).
\]
But the AR spectral density (the $q=0$ ARMA case) gives
$\sdf_X(f)=\sigma_\eps^2/\abs{\Phi(e^{-2\pi i f})}^2$, hence
\[
\sdf_{\tilde\nu}(f)=\abs{\Phi(e^{-2\pi i f})}^2\cdot
\frac{\sigma_\eps^2}{\abs{\Phi(e^{-2\pi i f})}^2}=\sigma_\eps^2 .
\]
A constant spectral density means $\{\tilde\nu_t\}$ is white noise of variance
$\sigma_\eps^2$; being Gaussian and mean-zero it has the same distribution as
$\{\eps_t\}$. Therefore the reversed process is AR($p$) with identical
coefficients. \hfill$\square$
\end{solution}

\begin{exercise}{Spectral density of an MA(1) via a filter \quad\textnormal{[New]}}{p11-n1}
Let $\eps_t$ be white noise of variance $\sigma^2$ and $X_t=\eps_t-\theta\eps_{t-1}$.
Using the LTI machinery, find $\sdf_X(f)$ and check it against the direct ACVS.
\end{exercise}
\begin{solution}
Here $X_t=\Theta(\Bsh)\eps_t$ with $\Theta(z)=1-\theta z$, an LTI filter with
impulse response $h_0=1,\ h_1=-\theta$. Its transfer function is
$H(f)=1-\theta e^{-2\pi i f}$, so
\[
\sdf_X(f)=\abs{H(f)}^2\sdf_\eps(f)=\sigma^2\abs{1-\theta e^{-2\pi i f}}^2 .
\]
Expand:
\[
\abs{1-\theta e^{-2\pi i f}}^2=(1-\theta e^{-2\pi i f})(1-\theta e^{2\pi i f})
=1+\theta^2-2\theta\cos(2\pi f),
\]
giving $\sdf_X(f)=\sigma^2\big(1+\theta^2-2\theta\cos(2\pi f)\big)$.

\emph{Check via ACVS.} For this MA(1), $\acvs_0=(1+\theta^2)\sigma^2$,
$\acvs_{\pm1}=-\theta\sigma^2$, else $0$. Then
$\sdf_X(f)=\sum_\tau\acvs_\tau e^{-2\pi i f\tau}
=(1+\theta^2)\sigma^2-\theta\sigma^2(e^{-2\pi i f}+e^{2\pi i f})
=\sigma^2(1+\theta^2-2\theta\cos 2\pi f)$, identical. \hfill$\square$
\end{solution}

\begin{exercise}{Transfer function and gain of a two-sided moving average \quad\textnormal{[New]}}{p11-n2}
Let $\filt[\{X_t\}]_u=\tfrac13(X_{u-1}+X_u+X_{u+1})$ (a centred 3-point average).
Find the impulse response, the transfer function, and the squared gain
$\abs{H(f)}^2$. What does it do to $f=0$ versus $f=\tfrac13$?
\end{exercise}
\begin{solution}
Writing $\filt[\{X_t\}]_u=\sum_m h_{u-m}X_m=\sum_k h_k X_{u-k}$, the coefficients
of $X_{u-1},X_u,X_{u+1}$ are $h_1,h_0,h_{-1}$, so
$h_{-1}=h_0=h_1=\tfrac13$ (others $0$). Then
\[
H(f)=\sum_m h_m e^{-2\pi i f m}=\tfrac13\big(e^{2\pi i f}+1+e^{-2\pi i f}\big)
=\tfrac13\big(1+2\cos(2\pi f)\big).
\]
$H$ is real (the filter is symmetric, zero phase), so
\[
\abs{H(f)}^2=\tfrac19\big(1+2\cos 2\pi f\big)^2 .
\]
At $f=0$: $H(0)=1$, gain $1$ (the mean / DC component passes untouched). At
$f=\tfrac13$: $\cos(2\pi/3)=-\tfrac12$, so $H=\tfrac13(1-1)=0$ --- that frequency
is annihilated. The 3-point average is a low-pass smoother with a null at
$f=\tfrac13$. \hfill$\square$
\end{solution}

\begin{exercise}{Spectral density of an AR(1) \quad\textnormal{[New]}}{p11-n3}
Let $X_t=\phi X_{t-1}+\eps_t$ with $\abs\phi<1$ and $\eps$ white noise of
variance $\sigma^2$. Use the polynomial-filter lemma to obtain $\sdf_X(f)$ and
state where it peaks for $\phi>0$ and for $\phi<0$.
\end{exercise}
\begin{solution}
Rewrite as $\Phi(\Bsh)X_t=\eps_t$ with $\Phi(z)=1-\phi z$. By the
polynomial-filter lemma the AR filter has transfer function
$\Phi(e^{-2\pi i f})=1-\phi e^{-2\pi i f}$, so the output-spectrum theorem
applied to $\Phi(\Bsh)[\{X_t\}]=\{\eps_t\}$ gives
$\abs{1-\phi e^{-2\pi i f}}^2\sdf_X(f)=\sigma^2$, hence
\[
\sdf_X(f)=\frac{\sigma^2}{\abs{1-\phi e^{-2\pi i f}}^2}
=\frac{\sigma^2}{1+\phi^2-2\phi\cos(2\pi f)} .
\]
The denominator is smallest where $\phi\cos(2\pi f)$ is largest. For
$\phi>0$ this is $f=0$ (spectrum peaks at low frequency --- slowly varying,
persistent series); for $\phi<0$ it is $f=\tfrac12$ (peak at high frequency ---
rapidly oscillating series). Stationarity ($\abs\phi<1$) keeps the denominator
$\ge(1-\abs\phi)^2>0$, so $\sdf_X$ is finite everywhere. \hfill$\square$
\end{solution}

\begin{exercise}{Cascade gain of difference then smooth \quad\textnormal{[New]}}{p11-n4}
A series is first differenced, $\{V_t\}=(I-\Bsh)[\{X_t\}]$, then run through the
3-point average $\filt_2$ of Exercise~\ref{exo:p11-n2}, giving
$\{Y_t\}=\filt_2[\{V_t\}]$. Find $\sdf_Y$ in terms of $\sdf_X$.
\end{exercise}
\begin{solution}
The whole operation is the cascade $\filt=\filt_2(I-\Bsh)$, which is LTI with
transfer function the \emph{product} of the two transfer functions:
\[
H(f)=\underbrace{(1-e^{-2\pi i f})}_{\text{difference}}\cdot
\underbrace{\tfrac13(1+2\cos 2\pi f)}_{\text{3-point average}} .
\]
The squared gains multiply, so using
$\abs{1-e^{-2\pi i f}}^2=4\sin^2(\pi f)$,
\[
\sdf_Y(f)=\abs{H(f)}^2\sdf_X(f)
=4\sin^2(\pi f)\cdot\tfrac19\big(1+2\cos 2\pi f\big)^2\,\sdf_X(f).
\]
This both removes the $f=0$ component (the difference) and nulls $f=\tfrac13$
(the average) --- a band-suppressing combination. \hfill$\square$
\end{solution}

% ======================================================================
%  PART 12 -- Exam Strategy Checklist + Classic Pitfalls
% ======================================================================
\section{Exam Strategy Checklist}
\label{sec:strategy}

The exam (cf.\ the mock) is \textbf{closed-book, no calculator}, and made of a
few multi-part questions mixing \emph{stated definitions}, \emph{short proofs},
and \emph{calculations}. The recurring tasks are: write down a definition;
compute an ACVS/ACF; find an MA$(\infty)$ representation; check
causality/invertibility; derive a spectral density; forecast and give the
error; prove (joint) second-order stationarity; and read cross-spectra. This
checklist maps a \emph{question cue} to the \emph{tool} you deploy.

\begin{tcolorbox}[enhanced,breakable,colback=orange!6!white,colframe=orange!75!black,
  arc=1.5mm,boxrule=0.7pt,title=\textbf{Cue $\rightarrow$ Tool: the decision table},
  coltitle=white,colbacktitle=orange!75!black,fonttitle=\bfseries]
\renewcommand{\arraystretch}{1.35}
\begin{tabular}{@{}p{0.46\textwidth}@{\hspace{0.5em}}p{0.5\textwidth}@{}}
\toprule
\textbf{When the question says\dots} & \textbf{Deploy\dots}\\
\midrule
``State the definition of (weak/strict) stationarity'' &
The three moment conditions (mean const, finite const variance, covariance a
function of $\tau$ only) / equality of all finite-dim.\ distributions under
shift.\\
``Is $\{X_t\}$ (weakly) stationary?'' &
Check the 3 conditions in order; stop at first failure. Watch for undefined
moments and a residual $t$-dependence in the covariance.\\
``Compute the ACVS/ACF of an MA$(q)$'' &
$\acvs_\tau=\sigma^2\sum_j\theta_j\theta_{j+\abs\tau}$ for $\abs\tau\le q$, else
$0$. ACF \emph{cuts off} after lag $q$.\\
``Compute the ACVS of an AR/ARMA'' &
Multiply the defining equation by $X_{t-\tau}$, take expectations $\Rightarrow$
a \textbf{difference equation} $\acvs_\tau=\sum\phi_j\acvs_{\tau-j}$; solve with
initial conditions (handle small $\tau$ separately because of the MA/noise
terms).\\
``Find $\psi_j$ in $X_t=\sum\psi_j\eps_{t-j}$'' &
Write $X_t=\Theta(B)/\Phi(B)\,\eps_t$; expand by partial fractions /
$1/(1-gB)=\sum g^kB^k$, or match coefficients in $\Phi(B)\sum\psi_jB^j=\Theta(B)$.\\
``Is it causal/stationary? invertible?'' &
\textbf{Roots of $\Phi(z)$ outside the unit circle} $\Rightarrow$
causal/stationary; \textbf{roots of $\Theta(z)$ outside} $\Rightarrow$
invertible. (AR always invertible; MA always stationary.)\\
``Find the PACF'' &
Durbin--Levinson recursion on the ACF; for AR$(p)$, $\pacf_\tau=0$ for
$\tau>p$. PACF \emph{cuts off} after $p$.\\
``Identify a model from ACF/PACF plots'' &
ACF cuts at $q$ \& PACF tails $\Rightarrow$ MA$(q)$; PACF cuts at $p$ \& ACF
tails $\Rightarrow$ AR$(p)$; both tail $\Rightarrow$ ARMA; slow ACF decay
$\Rightarrow$ difference (non-stationary).\\
``Estimate AR parameters'' &
Yule--Walker $\hat\phi=\hat\Gamma^{-1}\hat\gamma$,
$\hat\sigma^2=\hat\gamma_0-\sum\hat\phi_j\hat\gamma_j$; or (forward/backward) LS.\\
``Estimate MA/ARMA by LS'' &
Reconstruct innovations recursively ($\eps_0=0$), minimise $\sum\hat\eps_t^2$.\\
``Remove a trend / make stationary'' &
Difference: $\diff^q$ kills a degree-$(q-1)$ polynomial trend; $\diff_{(s)}$
removes period-$s$ seasonality.\\
``Forecast $X_{n+h}$ (causal ARMA)'' &
$X^n_{n+h}=\sum_{j\ge h}\psi_j\eps_{n+h-j}$ (future $\eps\equiv0$); MSE
$=\sigma^2\sum_{j=0}^{h-1}\psi_j^2$. Or the recursive truncated recipe.\\
``Forecast an ARIMA'' &
Forecast the differenced ARMA $Z=\diff^dX$, then un-difference (cumulate).
$P^n_{n+h}\to\infty$.\\
``Give a prediction interval'' &
$X^n_{n+h}\pm z_{1-\alpha/2}\sqrt{\Var(e_n(h))}$ (Gaussian).\\
``Find the spectral density'' &
$\sdf(f)=\sum_\tau\acvs_\tau e^{-2\pi i f\tau}$; for ARMA use
$\sdf(f)=\sigma^2\abs{\Theta(e^{-2\pi i f})}^2/\abs{\Phi(e^{-2\pi i f})}^2$.\\
``Spectrum of a filtered series $Y=\filt X$'' &
$\sdf_Y(f)=\abs{H(f)}^2\sdf_X(f)$, $H(f)=\sum_m h_m e^{-2\pi i f m}$.\\
``Periodogram expectation/variance'' &
$\E[\hat\sdf^{(p)}]=\Fej_n*\sdf$ (biased, blurring); $\Var\to\sdf(f)^2$ (does
not vanish). Taper for bias, multitaper for variance.\\
``Is this $\acvs_\tau$ a valid ACVS?'' &
Compute its SDF; it is valid \emph{iff} $\sdf(f)\ge0$ for all $f$ (equivalently
$\acvs$ is positive semi-definite).\\
``Cross-spectrum / coherence of a delay'' &
$S_{1,2}(f)=c\,e^{-2\pi i f d}S_{2,2}(f)$: amplitude $\abs c\,S_{2,2}$, phase
$-2\pi f d$.\\
``Is the VAR stable?'' &
Roots of $\det(I-\Phi_1z-\cdots-\Phi_pz^p)$ outside unit circle (companion
eigenvalues inside).\\
``Test residual autocorrelation'' &
Ljung--Box $Q_m=n(n+2)\sum_{\tau=1}^m\hat r_\tau^2/(n-\tau)\sim\chi^2_{m-p-q}$.\\
``Compare models'' &
AIC/AICc/BIC ($k=p+q+1$); smaller is better; BIC penalises complexity more.\\
``Model changing variance / volatility'' &
ARCH/GARCH on the conditional variance; recall it is uncorrelated but
dependent; stationarity $\sum\alpha+\sum\beta<1$.\\
\bottomrule
\end{tabular}
\end{tcolorbox}

\begin{keytech}{The six moves you will almost certainly need}{p12-six}
\begin{enumerate}
  \item \textbf{Stationarity check} (3 conditions).
  \item \textbf{ACVS via the white-noise filter trick} (only equal noise
        indices survive) or via the \textbf{difference equation} (multiply by
        $X_{t-\tau}$, take expectations).
  \item \textbf{Root test} for causality/invertibility/stability.
  \item \textbf{$\Theta/\Phi$ expansion} for $\psi$-weights (partial fractions
        or coefficient matching).
  \item \textbf{Recursive forecasting} and its MSE.
  \item \textbf{Spectral density} from ACVS, from a filter, or from ARMA
        polynomials.
\end{enumerate}
\end{keytech}

\begin{keytech}{Exam hygiene}{p12-hygiene}
\begin{itemize}
  \item State which \emph{theorem/definition} you invoke (examiners reward
        naming the tool: ``by causality $\E[\eps_tX_{t-k}]=0$''; ``by Herglotz
        / the prediction equations / Isserlis'').
  \item No calculator: keep fractions exact; factor polynomials by hand
        ($z=\frac{-b\pm\sqrt{b^2-4ac}}{2a}$) and report $\abs{z}$ vs $1$.
  \item For ``$\abs z>1$?'' with complex roots, compare
        $\abs z^2=(\Re z)^2+(\Im z)^2$ to $1$ --- no need to take the square
        root.
  \item Always sanity-check: $\acvs_0\ge0$, $\abs{\acf_\tau}\le1$, MA(1) has
        $\abs{\acf_1}\le\tfrac12$, an SDF must be $\ge0$ and real, a spectral
        \emph{matrix} must be Hermitian.
\end{itemize}
\end{keytech}

% ======================================================================
\section{Classic Pitfalls to Avoid}
\label{sec:pitfalls}

\begin{pitfall}
\textbf{Roots inside vs outside the unit circle.} Stationarity/causality
requires the roots of $\Phi(z)$ to lie \emph{outside} $\abs z=1$ (equivalently
$\Phi(z)\ne0$ for $\abs z\le1$). Students routinely flip this. The
\emph{characteristic} root condition (``$\abs z>1$'') is the opposite of the
companion-eigenvalue condition (``$\abs\lambda<1$''): both say the same thing,
so do not mix the two phrasings in one argument.
\end{pitfall}

\begin{pitfall}
\textbf{Strict vs weak stationarity.} They do not imply each other. Strict but
not weak: i.i.d.\ Cauchy (no finite variance). Weak but not strict: a process
with matched first two moments but non-Gaussian higher structure. They coincide
\emph{only} for Gaussian processes.
\end{pitfall}

\begin{pitfall}
\textbf{Lag goes in the FIRST argument (multivariate).}
$\Gamma_\tau=\Cov(X_{t+\tau},X_t)$, so $\Gamma_{-\tau}=\Gamma_\tau\Tr$ (not
$\Gamma_\tau$). Cross-covariance sequences are \emph{not} symmetric in $\tau$:
$\gamma^{(j,k)}_\tau=\gamma^{(k,j)}_{-\tau}$. In the univariate real case this
distinction is invisible, which is exactly why it trips people in the
multivariate case.
\end{pitfall}

\begin{pitfall}
\textbf{ACF cut-off, PACF cut-off --- which is which.} MA$(q)$: \emph{ACF} cuts
off after $q$ (PACF tails). AR$(p)$: \emph{PACF} cuts off after $p$ (ACF tails).
ARMA: both tail off. Swapping these is the single most common identification
error.
\end{pitfall}

\begin{pitfall}
\textbf{Identifiability of MA / sign of $\theta_0$.} Fixing $\theta_0=-1$ is not
cosmetic: without it, scaling
$(\theta_j,\sigma^2)\mapsto(c\theta_j,c^{-2}\sigma^2)$ gives the \emph{same}
ACVS, so the model is unidentifiable. Likewise a non-invertible MA and its
invertible ``twin'' (root $g$ vs $1/g$) share the same ACF --- pick the
invertible one.
\end{pitfall}

\begin{pitfall}
\textbf{Biased is not bad.} The divide-by-$n$ ACVS estimator $\hat\acvs_\tau$ is
biased yet positive semi-definite and has \emph{smaller} MSE than the unbiased
$\tilde\acvs_\tau$. Unbiasedness is not the objective; positive
semi-definiteness and low MSE are. Do not ``correct'' the estimator on autopilot.
\end{pitfall}

\begin{pitfall}
\textbf{The periodogram is not consistent.} Its variance tends to $\sdf(f)^2$
and does \emph{not} shrink with $n$ --- more data makes it wigglier, not
smoother. Reducing bias needs \emph{tapering}; reducing variance needs
\emph{smoothing/multitapering} ($\Var\approx\sdf(f)^2/K$).
\end{pitfall}

\begin{pitfall}
\textbf{Forecasts of stationary vs integrated processes.} A stationary ARMA
forecast reverts to the mean and its error variance saturates at $\gamma_0$. An
\emph{integrated} (ARIMA, $d\ge1$) forecast does \emph{not} revert and its error
variance $\to\infty$. Quoting a finite long-horizon variance for a random walk
is wrong.
\end{pitfall}

\begin{pitfall}
\textbf{Setting future innovations to zero (only).} In the $\psi$-weight
forecast, set $\eps_{n+1},\dots,\eps_{n+h}=0$ but keep the \emph{observed past}
innovations as residuals; in the recursive recipe, replace future
\emph{observations} by their forecasts, future errors by $0$, and past errors by
residuals --- do not zero the past errors.
\end{pitfall}

\begin{pitfall}
\textbf{ARCH/GARCH are uncorrelated but dependent.} $\Corr(r_{t+\tau},r_t)=0$,
yet $\Corr(r^2_{t+\tau},r^2_t)\ne0$ (geometric for ARCH(1):
$\alpha_1^{\abs\tau}$). ``Uncorrelated'' $\ne$ ``independent'' here. Also check
\emph{moment} conditions ($\alpha_1<1$ for variance, $\alpha_1^2<1/3$ for finite
kurtosis under Gaussian noise) separately from stationarity.
\end{pitfall}

\begin{pitfall}
\textbf{Ljung--Box degrees of freedom.} Use $\chi^2_{m-p-q}$, subtracting the
number of fitted ARMA parameters --- not $\chi^2_m$. Forgetting the $-(p+q)$
shifts the critical value and can flip the conclusion.
\end{pitfall}

\begin{pitfall}
\textbf{Angular vs ordinary frequency.} This course uses ordinary frequency $f$
(period $1/f$), so the SDF integrates over $[-\tfrac12,\tfrac12]$ and
$\sdf(f)=\sum\acvs_\tau e^{-2\pi i f\tau}$. Switching to $\omega=2\pi f$ scatters
stray $2\pi$ factors --- keep one convention.
\end{pitfall}

\begin{pitfall}
\textbf{A valid ACVS must be positive semi-definite.} A plausible-looking
sequence (e.g.\ $\acvs_\tau=1$ for $\abs\tau\le K$, else $0$) need \emph{not} be
a real ACVS: its SDF can go negative. Always test via $\sdf(f)\ge0$ before
claiming a sequence is an autocovariance.
\end{pitfall}

\begin{center}
\begin{tcolorbox}[width=0.9\textwidth,colback=green!5!white,colframe=green!48!black,
  arc=2mm,boxrule=0.8pt,halign=center,fontupper=\large\bfseries]
You are ready when you can do every \textcolor{cyan!55!black}{Exercise} in this
document unaided, name the tool behind every \textcolor{red!62!black}{Theorem},
and avoid every \textcolor{red!70!black}{Pitfall} above. Good luck!
\end{tcolorbox}
\end{center}


\end{document}
